% Gerado automaticamente. Nao editar.
\providecommand{\fvMeshCW}[3]{
  \draw[->,loopcol,line width=1.1pt] (#1)
    arc[start angle=0,end angle=-305,radius=#2cm];
  \node[loopcol,fill=white,inner sep=1pt] at (#1){\footnotesize #3};}

\section{Eletrostática}

\subsection{Cargas elétricas}

\blocofixacao
\begin{questao}[1]{}

O princípio da conservação da carga elétrica estabelece que:
\begin{alternativas}
\item as cargas elétricas de mesmo sinal se repelem.
\item cargas elétricas de sinais opostos se atraem.
\item a soma algébrica das cargas elétricas é constante em um sistema eletricamente isolado.
\item a soma das cargas positivas e negativas é diferente de zero em um sistema neutro.
\item os elétrons livres se atraem.
\end{alternativas}
\resposta{(c)}


\end{questao}
\begin{questao}[1]{}

Dois corpos de materiais diferentes, inicialmente neutros, são atritados entre si,
de forma isolada do meio. É correto afirmar que:
\begin{alternativas}
\item um fica eletrizado positivamente e o outro, negativamente.
\item um fica eletrizado negativamente e o outro permanece neutro.
\item um fica eletrizado positivamente e o outro permanece neutro.
\item ambos ficam eletrizados negativamente.
\item ambos ficam eletrizados positivamente.
\end{alternativas}
\resposta{(a)}


\end{questao}
\begin{questao}[1]{}

Atrita-se um bastão de vidro com um pano de lã, ambos inicialmente neutros.
Pode-se afirmar que:
\begin{alternativas}
\item só a lã fica eletrizada.
\item só o bastão fica eletrizado.
\item o bastão e a lã se eletrizam com cargas de mesmo sinal.
\item o bastão e a lã se eletrizam com cargas de mesmo valor absoluto e sinais opostos.
\item nenhuma das anteriores.
\end{alternativas}
\resposta{(d)}


\end{questao}
\begin{questao}[1]{}

A tabela apresenta parte de uma \textbf{série triboelétrica}: ao serem atritados,
o material que aparece \emph{mais acima} tende a ceder elétrons ao que aparece
mais abaixo.
\begin{table}[H]
\centering
\small\sffamily
\setlength{\arrayrulewidth}{0.5pt}
\begin{tabular}{@{}c l@{}}
\toprule
\textbf{Ordem} & \textbf{Material}\\
\midrule
1 & Pele de coelho\\
2 & Vidro\\
3 & Lã\\
4 & Seda\\
5 & Plástico (PVC)\\
6 & Teflon\\
\bottomrule
\end{tabular}
\caption{Série triboelétrica simplificada.}
\end{table}
Atritando-se um bastão de \textbf{PVC} com um pedaço de \textbf{lã}, os sinais das
cargas adquiridas pelo PVC e pela lã são, respectivamente:
\begin{alternativas}[2]
\item positivo e positivo.
\item positivo e negativo.
\item negativo e positivo.
\item negativo e negativo.
\item nulo e nulo.
\end{alternativas}
\resposta{(c)}


\end{questao}
\begin{questao}[1]{}

Duas cargas elétricas $Q_1$ e $Q_2$ atraem-se quando colocadas próximas uma da outra.
\begin{itensq}
\item O que se pode afirmar sobre os sinais de $Q_1$ e $Q_2$?
\item Sabe-se que $Q_1$ é repelida por uma terceira carga $Q_3$, positiva.
      Qual é o sinal de $Q_2$?
\end{itensq}
\resposta{(a) sinais opostos; (b) negativo}


\end{questao}
\begin{questao}[1]{}

Com relação à eletrização de um corpo, assinale \textbf{todas} as afirmativas corretas:
\begin{alternativas}
\item Um corpo eletricamente neutro que perde elétrons fica eletrizado positivamente.
\item Um corpo eletricamente neutro não possui cargas elétricas.
\item Um dos processos de eletrização consiste em retirar prótons do corpo.
\item Um corpo eletricamente neutro não pode ser atraído por um corpo eletrizado.
\item Friccionando-se dois corpos constituídos do mesmo material, um se eletriza
      positivamente e o outro, negativamente.
\end{alternativas}
\resposta{Apenas (a)}


\end{questao}
\begin{questao}[1]{}

Em uma cabine de \textbf{pintura eletrostática}, as gotículas de tinta são
eletrizadas negativamente e a peça metálica a ser pintada é ligada ao terminal
positivo da fonte. Esse arranjo é usado porque:
\begin{alternativas}
\item as gotículas se repelem entre si e são atraídas pela peça, reduzindo o desperdício.
\item as gotículas se atraem entre si, formando gotas maiores.
\item a tinta se torna condutora e adere melhor.
\item o processo elimina a necessidade de secagem.
\item a peça precisa estar eletrizada com o mesmo sinal da tinta.
\end{alternativas}
\resposta{(a)}


\end{questao}
\begin{questao}[1]{}

Um corpo, inicialmente neutro, é eletrizado com carga $Q = \SI{48}{\micro\coulomb}$.
\dados{carga elementar $e = \SI{1.6e-19}{\coulomb}$.}
\begin{itensq}
\item O corpo ficou com falta ou com excesso de elétrons?
\item Qual é o número de elétrons retirado do corpo ou a ele fornecido?
\end{itensq}
\resposta{(a) falta de elétrons; (b) $n = \pot{3.0}{14}$ elétrons}


\end{questao}
\begin{questao}[1]{}

Um bastão carregado positivamente atrai um objeto isolado, suspenso por um fio.
Sobre esse objeto, é correto afirmar que:
\begin{alternativas}
\item necessariamente possui elétrons em excesso.
\item é obrigatoriamente condutor.
\item trata-se de um isolante.
\item está carregado positivamente.
\item pode estar neutro.
\end{alternativas}
\resposta{(e)}


\end{questao}

\blocoaplicacao
\begin{questao}[2]{}

Uma barra eletrizada negativamente é aproximada de um corpo metálico AB,
inicialmente neutro e apoiado sobre suportes isolantes, conforme a figura.

\begin{figure}[H]
\centering
\begin{tikzpicture}[scale=1,>=Stealth]
  
  \draw[fill=laranja!15,draw=laranja,thick,rounded corners=2pt]
        (-4.3,-0.3) rectangle (-2.6,0.3);
  \foreach \x in {-4.1,-3.7,-3.3,-2.9}
     \node[laranja,font=\bfseries\small] at (\x,0) {$-$};
  \node[below,font=\footnotesize\sffamily] at (-3.45,-0.45) {barra eletrizada};
  
  \shade[left color=cinza!35,right color=cinza!15,draw=cinza,thick,rounded corners=6pt]
        (-1.9,-0.45) rectangle (2.4,0.45);
  \node[font=\bfseries] at (-1.55,0) {A};
  \node[font=\bfseries] at ( 2.05,0) {B};
  
  \foreach \x in {-1.2,1.7}{
    \draw[fill=cinza!25,draw=cinza] (\x-0.18,-1.1) rectangle (\x+0.18,-0.45);
    \draw[fill=cinza!45,draw=cinza] (\x-0.45,-1.25) rectangle (\x+0.45,-1.1);}
  \node[font=\footnotesize\sffamily,cinza] at (0.25,-1.5) {suportes isolantes};
  
  \draw[cinza!60] (-2.6,-1.25) -- (2.9,-1.25);
\end{tikzpicture}
\caption{Barra negativa aproximada do condutor neutro AB.}
\label{fig:inducao-ab}
\end{figure}

Podemos afirmar que:
\begin{alternativas}
\item não haverá movimento de elétrons livres no corpo AB.
\item os elétrons livres do corpo AB deslocam-se para a extremidade A.
\item o sinal da carga que aparece em B é positivo.
\item ocorre no corpo metálico o fenômeno da indução eletrostática.
\item após a separação de cargas, a carga total do corpo é não nula.
\end{alternativas}
\resposta{(d)}


\end{questao}
\begin{questao}[2]{}

Um corpo eletrizado \textbf{positivamente} é aproximado de um corpo metálico
neutro, cujas extremidades são X (mais próxima) e Y (mais afastada).

\begin{figure}[H]
\centering
\begin{tikzpicture}[scale=1,>=Stealth]
  \draw[fill=Vcol!12,draw=Vcol,thick,rounded corners=2pt] (-4.2,-0.3) rectangle (-2.7,0.3);
  \foreach \x in {-4.0,-3.6,-3.2,-2.9} \node[Vcol,font=\bfseries\small] at (\x,0) {$+$};
  \node[below,font=\footnotesize\sffamily] at (-3.45,-0.45) {corpo eletrizado};
  \shade[left color=cinza!35,right color=cinza!15,draw=cinza,thick,rounded corners=6pt]
        (-2.0,-0.45) rectangle (2.3,0.45);
  \node[font=\bfseries] at (-1.65,0) {X};
  \node[font=\bfseries] at ( 1.95,0) {Y};
  \foreach \x in {-1.3,1.6}{
    \draw[fill=cinza!25,draw=cinza] (\x-0.18,-1.1) rectangle (\x+0.18,-0.45);
    \draw[fill=cinza!45,draw=cinza] (\x-0.45,-1.25) rectangle (\x+0.45,-1.1);}
  \draw[cinza!60] (-2.7,-1.25) -- (2.8,-1.25);
\end{tikzpicture}
\caption{Indução eletrostática em um condutor neutro.}
\label{fig:inducao-xy}
\end{figure}

Podemos afirmar que:
\begin{alternativas}
\item não haverá movimentação de cargas negativas no corpo neutro.
\item a carga que aparece em X é positiva.
\item a carga que aparece em Y é negativa.
\item haverá força de interação elétrica entre os dois corpos.
\item todas as afirmativas acima estão erradas.
\end{alternativas}
\resposta{(d)}


\end{questao}
\begin{questao}[2]{}

Duas esferas condutoras descarregadas, $x$ e $y$, colocadas sobre suportes
isolantes, estão inicialmente em contato. Um bastão carregado positivamente é
aproximado da esfera $x$, como mostra a figura.

\begin{figure}[H]
\centering
\begin{tikzpicture}[scale=1]
  \draw[fill=Vcol!12,draw=Vcol,thick,rounded corners=2pt] (-4.6,-0.25) rectangle (-3.2,0.25);
  \foreach \x in {-4.4,-4.0,-3.6,-3.35} \node[Vcol,font=\bfseries\small] at (\x,0) {$+$};
  
  \shade[ball color=cinza!30] (-1.9,0) circle (0.75);
  \shade[ball color=cinza!30] ( 0.1,0) circle (0.75);
  \draw[cinza,thick] (-1.9,0) circle (0.75);
  \draw[cinza,thick] ( 0.1,0) circle (0.75);
  \node at (-1.9,0) {\bfseries $x$};
  \node at ( 0.1,0) {\bfseries $y$};
  
  \draw[line width=2pt,cinza] (-1.15,0) -- (-0.65,0);
  
  \foreach \x in {-1.9,0.1}{
    \draw[fill=cinza!25,draw=cinza] (\x-0.12,-1.55) rectangle (\x+0.12,-0.75);
    \draw[fill=cinza!45,draw=cinza] (\x-0.42,-1.7) rectangle (\x+0.42,-1.55);}
  \draw[cinza!60] (-3.0,-1.7) -- (1.1,-1.7);
\end{tikzpicture}
\caption{Eletrização por indução com separação das esferas.}
\label{fig:inducao-xy2}
\end{figure}

Em seguida, a esfera $y$ é afastada de $x$, \textbf{mantendo-se o bastão em sua
posição}. Ao final, as cargas de $x$ e $y$ são, respectivamente:
\begin{alternativas}[2]
\item nula e positiva.
\item negativa e positiva.
\item nula e nula.
\item negativa e nula.
\item positiva e negativa.
\end{alternativas}
\resposta{(b)}


\end{questao}
\begin{questao}[2]{}

A figura mostra um eletroscópio de folhas, inicialmente neutro. Um bastão
eletrizado negativamente é aproximado do disco, \textbf{sem tocá-lo}.

\begin{figure}[H]
\centering
\begin{tikzpicture}[scale=1]
  
  \draw[fill=laranja!15,draw=laranja,thick,rounded corners=2pt] (-2.9,1.9) rectangle (-1.3,2.3);
  \foreach \x in {-2.7,-2.35,-2.0,-1.6} \node[laranja,font=\bfseries\small] at (\x,2.1) {$-$};
  
  \draw[fill=cinza!30,draw=cinza,thick] (-0.9,2.0) rectangle (0.9,2.2);
  
  \draw[line width=1.6pt,cinza] (0,2.0) -- (0,0.55);
  
  \draw[azulprof!35,thick,fill=azulclaro!35]
        (-1.05,-1.5) -- (-1.05,1.35) .. controls (-1.05,1.65) and (-0.6,1.7) ..
        (-0.45,1.7) -- (0.45,1.7) .. controls (0.6,1.7) and (1.05,1.65) ..
        (1.05,1.35) -- (1.05,-1.5) -- cycle;
  \draw[fill=cinza!45,draw=cinza] (-1.25,-1.65) rectangle (1.25,-1.5);
  
  \draw[line width=1.1pt,laranja] (0,0.55) -- (-0.42,-0.75);
  \draw[line width=1.1pt,laranja] (0,0.55) -- ( 0.42,-0.75);
  \node[font=\footnotesize\sffamily,cinza,right] at (1.35,-0.2) {folhas metálicas};
  \node[font=\footnotesize\sffamily,cinza,right] at (1.35,2.1) {disco};
\end{tikzpicture}
\caption{Eletroscópio de folhas sob indução.}
\label{fig:eletroscopio}
\end{figure}

Sobre a situação descrita, é correto afirmar que:
\begin{alternativas}
\item as folhas permanecem fechadas, pois o eletroscópio continua neutro.
\item as folhas se abrem e o disco fica com carga positiva.
\item as folhas se abrem e o disco fica com carga negativa.
\item as folhas se abrem porque o eletroscópio adquiriu carga negativa total.
\item as folhas se fecham ainda mais.
\end{alternativas}
\resposta{(b)}


\end{questao}
\begin{questao}[2]{}

Três bolas metálicas podem ser carregadas eletricamente. Observa-se que cada uma
das três bolas \textbf{atrai uma e repele outra}. Considere as hipóteses:

\smallskip
\begin{tabular}{@{}ll@{}}
I   & Apenas uma das bolas está carregada.\\
II  & Duas das bolas estão carregadas.\\
III & As três bolas estão carregadas.\\
\end{tabular}
\smallskip

O fenômeno pode ser explicado:
\begin{alternativas}
\item somente pelas hipóteses II ou III.
\item somente pela hipótese I.
\item somente pela hipótese III.
\item somente pela hipótese II.
\item somente pelas hipóteses I ou II.
\end{alternativas}
\resposta{(c)}


\end{questao}
\begin{questao}[2]{}

Três esferas metálicas podem ser carregadas eletricamente. Aproximando-as duas a
duas, observa-se que \textbf{em todos os casos ocorre atração}. Considere:

\smallskip
\begin{tabular}{@{}ll@{}}
I   & Somente uma das esferas está carregada.\\
II  & Duas esferas estão carregadas.\\
III & As três esferas estão carregadas.\\
\end{tabular}
\smallskip

Quais hipóteses explicam o fenômeno descrito?
\begin{alternativas}
\item Apenas a hipótese I.
\item Apenas a hipótese II.
\item Apenas a hipótese III.
\item Apenas as hipóteses I e II.
\item Nenhuma das três hipóteses.
\end{alternativas}
\resposta{(d)}


\end{questao}
\begin{questao}[2]{}

Uma estudante colocou sobre uma mesa horizontal três pêndulos eletrostáticos
idênticos, equidistantes entre si, como se cada um ocupasse o vértice de um
triângulo equilátero. As esferas metálicas dos pêndulos atraíram-se mutuamente.
A estudante pode concluir corretamente que:
\begin{alternativas}
\item as três esferas estavam eletrizadas com cargas de mesmo sinal.
\item duas esferas estavam eletrizadas com cargas de mesmo sinal e uma com sinal oposto.
\item duas esferas estavam eletrizadas com cargas de mesmo sinal e uma neutra.
\item duas esferas estavam eletrizadas com cargas de sinais opostos e uma neutra.
\item uma esfera estava eletrizada e duas, neutras.
\end{alternativas}
\resposta{(d)}


\end{questao}
\begin{questao}[2]{}

Eletriza-se por atrito um bastão de plástico com um pedaço de papel. Em seguida,
o bastão é aproximado de um pêndulo eletrostático já eletrizado e observa-se
\textbf{repulsão}. Em qual alternativa a carga de cada elemento corresponde a
essa descrição?

\begin{table}[H]
\centering
\small\sffamily
\setlength{\arraystretch}{1.25}
\begin{tabular}{@{}c c c c@{}}
\toprule
 & \textbf{Papel} & \textbf{Bastão} & \textbf{Pêndulo}\\
\midrule
a) & positiva & positiva & positiva\\
b) & negativa & positiva & negativa\\
c) & negativa & negativa & positiva\\
d) & positiva & positiva & negativa\\
e) & positiva & negativa & negativa\\
\bottomrule
\end{tabular}
\caption{Alternativas da questão.}
\end{table}
\resposta{(e)}


\end{questao}
\begin{questao}[2]{}

Duas esferas condutoras de dimensões iguais, A e B, estão eletrizadas com
$Q_A = \SI{+5}{\coulomb}$ e $Q_B = \SI{-1}{\coulomb}$, e são ligadas por um fio
condutor com uma chave C, inicialmente aberta.

\begin{figure}[H]
\centering
\begin{tikzpicture}[scale=1]
  \shade[ball color=Vcol!25] (-2.6,0) circle (0.8);
  \shade[ball color=Icol!25] ( 2.6,0) circle (0.8);
  \draw[Vcol,thick] (-2.6,0) circle (0.8);
  \draw[Icol,thick] ( 2.6,0) circle (0.8);
  \node[font=\bfseries] at (-2.6,0.18) {A};
  \node[font=\small] at (-2.6,-0.28) {$+\SI{5}{\coulomb}$};
  \node[font=\bfseries] at ( 2.6,0.18) {B};
  \node[font=\small] at ( 2.6,-0.28) {$-\SI{1}{\coulomb}$};
  
  \draw[line width=1.1pt,cinza] (-1.8,0) -- (-0.55,0);
  \draw[line width=1.1pt,cinza] ( 0.55,0) -- ( 1.8,0);
  \fill[cinza] (-0.55,0) circle (0.06);
  \fill[cinza] ( 0.55,0) circle (0.06);
  \draw[line width=1.1pt,cinza] (-0.55,0) -- (0.42,0.52);
  \node[font=\footnotesize\sffamily,cinza,above] at (0,0.62) {C (aberta)};
  
  \foreach \x in {-2.6,2.6}{
    \draw[fill=cinza!25,draw=cinza] (\x-0.12,-1.6) rectangle (\x+0.12,-0.8);
    \draw[fill=cinza!45,draw=cinza] (\x-0.42,-1.75) rectangle (\x+0.42,-1.6);}
  \draw[cinza!60] (-3.6,-1.75) -- (3.6,-1.75);
\end{tikzpicture}
\caption{Esferas idênticas ligadas por fio condutor com chave.}
\label{fig:esferas-chave-ab}
\end{figure}

Quando a chave é fechada, passam elétrons:
\begin{alternativas}
\item de A para B, e a nova carga de A é $\SI{+2}{\coulomb}$.
\item de A para B, e a nova carga de B é $\SI{-1}{\coulomb}$.
\item de B para A, e a nova carga de A é $\SI{+1}{\coulomb}$.
\item de B para A, e a nova carga de B é $\SI{-1}{\coulomb}$.
\item de B para A, e a nova carga de A é $\SI{+2}{\coulomb}$.
\end{alternativas}
\resposta{(e)}


\end{questao}
\begin{questao}[2]{}

Duas esferas condutoras idênticas, K e L, com cargas
$Q_K = \SI{+3}{\micro\coulomb}$ e $Q_L = \SI{+1}{\micro\coulomb}$, são ligadas por
um fio condutor com uma chave C, inicialmente aberta.

\begin{figure}[H]
\centering
\begin{tikzpicture}[scale=1]
  \shade[ball color=Vcol!25] (-2.6,0) circle (0.8);
  \shade[ball color=Vcol!12] ( 2.6,0) circle (0.8);
  \draw[Vcol,thick] (-2.6,0) circle (0.8);
  \draw[Vcol,thick] ( 2.6,0) circle (0.8);
  \node[font=\bfseries] at (-2.6,0.18) {K};
  \node[font=\small] at (-2.6,-0.30) {$+\SI{3}{\micro\coulomb}$};
  \node[font=\bfseries] at ( 2.6,0.18) {L};
  \node[font=\small] at ( 2.6,-0.30) {$+\SI{1}{\micro\coulomb}$};
  \draw[line width=1.1pt,cinza] (-1.8,0) -- (-0.55,0);
  \draw[line width=1.1pt,cinza] ( 0.55,0) -- ( 1.8,0);
  \fill[cinza] (-0.55,0) circle (0.06);
  \fill[cinza] ( 0.55,0) circle (0.06);
  \draw[line width=1.1pt,cinza] (-0.55,0) -- (0.42,0.52);
  \node[font=\footnotesize\sffamily,cinza,above] at (0,0.62) {C (aberta)};
  \foreach \x in {-2.6,2.6}{
    \draw[fill=cinza!25,draw=cinza] (\x-0.12,-1.6) rectangle (\x+0.12,-0.8);
    \draw[fill=cinza!45,draw=cinza] (\x-0.42,-1.75) rectangle (\x+0.42,-1.6);}
  \draw[cinza!60] (-3.6,-1.75) -- (3.6,-1.75);
\end{tikzpicture}
\caption{Esferas idênticas com cargas de mesmo sinal.}
\label{fig:esferas-chave-kl}
\end{figure}

Quando a chave é fechada:
\begin{alternativas}
\item não ocorre fluxo de elétrons, pois ambas estão com cargas positivas.
\item passam elétrons de L para K, transportando carga de $\SI{+2}{\micro\coulomb}$.
\item passam elétrons de L para K, transportando carga de $\SI{-1}{\micro\coulomb}$.
\item passam elétrons de K para L, transportando carga de $\SI{+1}{\micro\coulomb}$.
\item passam elétrons de K para L, transportando carga de $\SI{-1}{\micro\coulomb}$.
\end{alternativas}
\resposta{(c)}


\end{questao}
\begin{questao}[2]{}

Duas esferas metálicas idênticas estão carregadas com $5Q$ e $-3Q$.
Encostando-se uma na outra e separando-as em seguida, cada esfera ficará com
carga elétrica igual a:
\begin{alternativas}[3]
\item $8Q$
\item $5Q$
\item $3Q$
\item $2Q$
\item $Q$
\end{alternativas}
\resposta{(e)}


\end{questao}
\begin{questao}[2]{}

Uma pequena esfera metálica está eletrizada com carga $\pot{8.0}{-8}\,\si{\coulomb}$.
Colocando-a em contato com outra idêntica, porém neutra, o número de elétrons que
passa de uma esfera para a outra é:
\dados{$e = \SI{1.6e-19}{\coulomb}$.}
\begin{alternativas}[3]
\item $\pot{4.0}{12}$
\item $\pot{4.0}{11}$
\item $\pot{4.0}{10}$
\item $\pot{2.5}{12}$
\item $\pot{2.5}{11}$
\end{alternativas}
\resposta{(e)}


\end{questao}
\begin{questao}[2]{}

Um estudante afirma ter medido, em um pequeno corpo eletrizado, a carga
$Q = \pot{2.4}{-19}\,\si{\coulomb}$. Sabendo que $e = \SI{1.6e-19}{\coulomb}$,
esse resultado é possível? Justifique.
\resposta{Não: a carga não é múltipla inteira de $e$.}


\end{questao}

\blocoaprofundamento
\begin{questao}[3]{}

Três esferas metálicas A, B e C, idênticas e no vácuo, encontram-se inicialmente
com A carregada com $+Q$ e B e C neutras. A esfera A é colocada em contato com B
e, posteriormente, com C. Os valores finais das cargas de A, B e C são,
respectivamente:
\begin{alternativas}[2]
\item $\dfrac{Q}{3},\ \dfrac{Q}{3},\ \dfrac{Q}{3}$
\item $\dfrac{Q}{2},\ \dfrac{Q}{2},\ \dfrac{Q}{2}$
\item $\dfrac{Q}{4},\ \dfrac{Q}{4},\ \dfrac{Q}{4}$
\item $\dfrac{Q}{4},\ \dfrac{Q}{2},\ \dfrac{Q}{4}$
\item $\dfrac{Q}{2},\ \dfrac{Q}{2},\ \dfrac{Q}{4}$
\end{alternativas}
\resposta{(d)}


\end{questao}
\begin{questao}[3]{}

Quatro esferas metálicas idênticas estão isoladas entre si. As esferas A, B e C
estão neutras e a esfera D está eletrizada com carga $Q$. A esfera D é colocada
em contato sucessivamente com A, depois com B e, por fim, com C.
Ao final do processo, as cargas de C e D são, respectivamente:
\begin{alternativas}[2]
\item $\dfrac{Q}{8}$ e $\dfrac{Q}{8}$
\item $\dfrac{Q}{8}$ e $\dfrac{Q}{4}$
\item $\dfrac{Q}{4}$ e $\dfrac{Q}{8}$
\item $\dfrac{Q}{2}$ e $\dfrac{Q}{2}$
\item $Q$ e $-Q$
\end{alternativas}
\resposta{(a)}


\end{questao}
\begin{questao}[3]{}

Considere três esferas metálicas X, Y e Z, de diâmetros iguais. Y e Z estão fixas
e suficientemente afastadas para que efeitos de indução sejam desprezíveis. A
situação inicial é: X neutra, Y com carga $+Q$ e Z com carga $-Q$. As esferas não
trocam cargas com o ambiente. Fazendo-se X tocar \textbf{primeiro} Y e
\textbf{depois} Z, a carga final de X será:
\begin{alternativas}[2]
\item zero
\item $\dfrac{2Q}{3}$
\item $-\dfrac{Q}{2}$
\item $\dfrac{Q}{8}$
\item $-\dfrac{Q}{4}$
\end{alternativas}
\resposta{(e)}


\end{questao}
\begin{questao}[3]{}

Duas esferas metálicas idênticas, A e B, estão carregadas com
$\SI{16}{\micro\coulomb}$ e $\SI{4}{\micro\coulomb}$, respectivamente. Uma terceira
esfera C, idêntica às anteriores, está inicialmente descarregada. Coloca-se C em
contato com A; desfeito o contato, C é colocada em contato com B. Supondo que não
haja troca de cargas com o exterior, a carga final de C é:
\begin{alternativas}[3]
\item $\SI{8}{\micro\coulomb}$
\item $\SI{6}{\micro\coulomb}$
\item $\SI{4}{\micro\coulomb}$
\item $\SI{3}{\micro\coulomb}$
\item nula
\end{alternativas}
\resposta{(b)}


\end{questao}
\begin{questao}[3]{}

Três esferas condutoras A, B e C têm o mesmo diâmetro. A esfera A está
inicialmente neutra, e as outras duas estão carregadas com
$q_B = \SI{6}{\micro\coulomb}$ e $q_C = \SI{7}{\micro\coulomb}$. Com a esfera A,
toca-se primeiro B e depois C. As cargas de A, B e C após os contatos são,
respectivamente:
\begin{alternativas}
\item zero, zero e $\SI{13}{\micro\coulomb}$
\item $\SI{7}{\micro\coulomb}$, $\SI{3}{\micro\coulomb}$ e $\SI{5}{\micro\coulomb}$
\item $\SI{5}{\micro\coulomb}$, $\SI{3}{\micro\coulomb}$ e $\SI{5}{\micro\coulomb}$
\item $\SI{6}{\micro\coulomb}$, $\SI{7}{\micro\coulomb}$ e zero
\item todas iguais a $\SI{4.3}{\micro\coulomb}$
\end{alternativas}
\resposta{(c)}


\end{questao}
\begin{questao}[3]{}

Na eletrização por contato entre duas esferas condutoras de raios $R$ e $2R$,
sabe-se que, após o contato, a razão entre a carga adquirida por cada esfera e o
seu respectivo raio é constante. Se, inicialmente, a esfera de raio $R$ possui
carga $3Q$ e a de raio $2R$ possui $2Q$, calcule a carga final de cada esfera.
\resposta{$Q_1 = \dfrac{5Q}{3}$ (raio $R$) e $Q_2 = \dfrac{10Q}{3}$ (raio $2R$)}


\end{questao}
\begin{questao}[3]{}

Uma esfera metálica isolada possui carga inicial $Q_0$. Ela é colocada em contato,
uma única vez e sucessivamente, com $n$ esferas idênticas a ela e inicialmente
neutras.
\begin{itensq}
\item Obtenha a expressão da carga final da esfera original em função de $Q_0$ e $n$.
\item A partir de quantos contatos a carga da esfera original cai abaixo de
      \SI{1}{\percent} de $Q_0$?
\end{itensq}
\resposta{(a) $Q_n = Q_0/2^{\,n}$; \quad (b) $n = 7$}


\end{questao}
\begin{questao}[3]{}

Um gerador de Van de Graaff transporta carga para sua cúpula a uma taxa
constante de $\SI{1}{\micro\ampere}$.
\begin{itensq}
\item Que carga é acumulada em $\SI{10}{\second}$?
\item Quantos elétrons isso representa?
\item Se a cúpula é uma esfera de raio $\SI{25}{\centi\meter}$, qual é o potencial
      atingido? O ar rompe a cerca de $\SI{3}{\mega\volt\per\meter}$ --- haverá
      descarga?
\end{itensq}
\dados{$e=\SI{1.6e-19}{\coulomb}$; $\ke=\SI{9e9}{\newton\meter\squared\per\coulomb\squared}$.}
\resposta{(a) \SI{10}{\micro\coulomb}; (b) $\pot{6.25}{13}$;
(c) $V=\SI{360}{\kilo\volt}$, $E=\SI{1.44}{\mega\volt\per\meter}$ --- sem descarga}


\end{questao}
\begin{questao}[3]{}

Um fio de cobre tem massa $\SI{1}{\gram}$. Sabendo que cada átomo de cobre
contribui com \emph{um} elétron livre:
\begin{itensq}
\item Calcule o número de elétrons livres no fio.
\item Determine a carga total desses elétrons.
\item Compare com a carga típica obtida por atrito ($\sim\SI{1}{\micro\coulomb}$)
      e interprete o resultado.
\end{itensq}
\dados{$N_A=\SI{6.02e23}{\per\mole}$; $M_{\text{Cu}}=\SI{63.5}{\gram\per\mole}$;
$e=\SI{1.6e-19}{\coulomb}$.}
\resposta{(a) $\approx\pot{9.5}{21}$; (b) $\approx\SI{1517}{\coulomb}$;
(c) razão $\sim10^9$}


\end{questao}
\begin{questao}[3]{}

Três esferas condutoras isoladas têm raios $R$, $2R$ e $3R$. A primeira está
carregada com $Q$ e as demais estão neutras. Sabe-se que, ao se tocarem, duas
esferas de raios diferentes repartem a carga de modo que o \emph{potencial} de
ambas fique igual, o que implica $\dfrac{Q_1}{R_1}=\dfrac{Q_2}{R_2}$.

Encosta-se a esfera $A$ (raio $R$) na esfera $B$ (raio $2R$); separa-se; em
seguida, encosta-se $A$ na esfera $C$ (raio $3R$).
\begin{itensq}
\item Determine a carga final de cada esfera.
\item Verifique a conservação da carga total.
\item Generalize: qual é a carga final de $A$ após tocar sucessivamente esferas de
      raios $2R, 3R, \dots, nR$?
\end{itensq}
\resposta{(a) $Q_A=\dfrac{Q}{12}$, $Q_B=\dfrac{2Q}{3}$, $Q_C=\dfrac{Q}{4}$;
(b) soma $=Q$; (c) $Q_A^{(n)}=Q\displaystyle\prod_{j=2}^{n}\frac{1}{j+1}$}


\end{questao}

\blocoaprofundamento
\begin{questao}[3]{}

Um sistema isolado contém $5$ prótons e $3$ elétrons livres. Após uma
reorganização interna, dois elétrons passam a orbitar dois dos prótons,
formando átomos neutros.
\begin{itensq}
\item Determine a carga total antes e depois.
\item Justifique o resultado.
\end{itensq}
\dados{$e=\pot{1.6}{-19}\,\si{\coulomb}$}
\resposta{$Q=+2e=\pot{3.2}{-19}\,\si{\coulomb}$ nos dois instantes}


\end{questao}
\begin{questao}[3]{}

Um bastão isolante adquire, por atrito, carga de $\SI{-5}{\nano\coulomb}$.
\begin{itensq}
\item Determine o número de elétrons transferidos.
\item Determine a massa correspondente e compare com a massa do bastão
      ($\SI{20}{\gram}$).
\end{itensq}
\dados{$e=\pot{1.6}{-19}\,\si{\coulomb}$;
$m_e=\pot{9.11}{-31}\,\si{\kilogram}$}
\resposta{$n=\pot{3.1}{10}$ elétrons; $\Delta m=\pot{2.8}{-20}\,\si{\kilogram}$}


\end{questao}

\blocodesafio
\begin{questao}[4]{}

\textbf{(Série geométrica dos contatos.)} Uma esfera $A$ com carga $Q$ é
tocada sucessivamente por esferas idênticas \emph{neutras}
$B_1,B_2,\dots,B_n$, uma de cada vez, sendo cada uma retirada após o contato.
\begin{itensq}
\item Determine a carga de $A$ após $n$ contatos.
\item Determine a carga de cada esfera $B_k$.
\item Verifique a conservação da carga total.
\end{itensq}
\resposta{$Q_A=Q/2^{n}$; $Q_{B_k}=Q/2^{k}$; a soma reconstitui $Q$}


\end{questao}
\begin{questao}[4]{}

\textbf{(Limite de eletrização.)} Uma esfera condutora de raio
$R=\SI{10}{\centi\meter}$ é carregada no ar, cuja rigidez dielétrica é
$E_{rup}=\SI{3e6}{\volt\per\meter}$.
\begin{itensq}
\item Determine a carga máxima que ela pode reter.
\item Determine o potencial correspondente.
\item Explique o que ocorre se tentarmos ultrapassar esse valor.
\end{itensq}
\dados{$\ke=\SI{9e9}{\newton\meter\squared\per\coulomb\squared}$}
\resposta{$Q_{\max}=\SI{3.33}{\micro\coulomb}$; $V=\SI{300}{\kilo\volt}$}


\end{questao}
\begin{questao}[4]{}

\textbf{(Poder das pontas.)} Duas esferas condutoras de raios
$R=\SI{10}{\centi\meter}$ e $r=\SI{1}{\centi\meter}$ são ligadas por um fio
longo e carregadas até $V=\SI{30}{\kilo\volt}$.
\begin{itensq}
\item Determine a razão entre as cargas e entre as densidades superficiais.
\item Determine o campo na superfície de cada uma.
\item Explique o efeito e cite duas aplicações.
\end{itensq}
\resposta{$Q_R/Q_r=10$; $\sigma_r/\sigma_R=10$;
$E_R=\SI{300}{\kilo\volt\per\meter}$ e $E_r=\SI{3}{\mega\volt\per\meter}$}


\end{questao}
\begin{questao}[4]{}

\textbf{(Ordem de grandeza do atrito.)} Um bastão de vidro de $\SI{20}{\gram}$
adquire $\SI{5}{\nano\coulomb}$ por atrito. A massa molar média do vidro é
cerca de $\SI{60}{\gram\per\mole}$.
\begin{itensq}
\item Estime o número de átomos do bastão.
\item Determine a fração de átomos que perdeu um elétron.
\item Interprete o resultado.
\end{itensq}
\dados{$N_A=\pot{6.02}{23}\,\si{\per\mole}$}
\resposta{$\approx\pot{2}{23}$ átomos; fração de $\pot{1.6}{-13}$}


\end{questao}
\begin{questao}[4]{}

\textbf{(Força elétrica contra peso.)} Duas esferas idênticas de
$\SI{1}{\gram}$ estão a $\SI{5}{\centi\meter}$ uma da outra, cada uma com
carga $q$.
\begin{itensq}
\item Determine a força para $q=\SI{5}{\nano\coulomb}$ e para
      $q=\SI{50}{\nano\coulomb}$.
\item Determine a carga necessária para que a força iguale o peso.
\item Comente a viabilidade.
\end{itensq}
\dados{$g=\SI{9.8}{\meter\per\second\squared}$}
\resposta{$\SI{90}{\micro\newton}$ e $\SI{9}{\milli\newton}$;
$q\approx\SI{52}{\nano\coulomb}$}


\end{questao}
\begin{questao}[4]{}

\textbf{(Conservação em processos nucleares.)} Verifique a conservação da carga
nos processos abaixo.
\begin{itensq}
\item Decaimento beta: $n\to p+e^{-}+\bar\nu_e$.
\item Aniquilação: $e^{-}+e^{+}\to2\gamma$.
\item Discuta o alcance do princípio.
\end{itensq}
\resposta{Conservada em ambos: $0=+e-e+0$ e $-e+e=0$}


\end{questao}
\begin{questao}[4]{}

\textbf{(Copo de Faraday.)} Um corpo carregado é introduzido, sem tocar as
paredes, no interior de um recipiente metálico fechado ligado a um eletrômetro.
\begin{itensq}
\item Explique o que o instrumento indica e por quê.
\item Explique por que a leitura não depende da posição do corpo dentro do
      recipiente.
\item Descreva o que ocorre se o corpo tocar a parede interna.
\end{itensq}
\resposta{O eletrômetro indica exatamente a carga do corpo, independentemente da
posição}


\end{questao}
\begin{questao}[4]{}

\textbf{(Blindagem eletrostática.)} Uma malha metálica aterrada envolve um
instrumento sensível.
\begin{itensq}
\item Explique por que o campo externo não penetra.
\item Determine se a blindagem funciona também no sentido inverso.
\item Discuta a diferença entre malha aterrada e não aterrada.
\end{itensq}
\resposta{Blinda o interior contra campos externos; o aterramento é necessário
para blindar o exterior contra fontes internas}


\end{questao}

\gabarito[1.1 Cargas elétricas]

\subsection{Lei de Coulomb}

\blocofixacao
\begin{questao}[1]{}

Duas cargas $q_1 = \SI{+2}{\micro\coulomb}$ e $q_2 = \SI{-3}{\micro\coulomb}$
estão separadas por $\SI{4}{\centi\meter}$ no vácuo. Determine o módulo e o tipo
(atração ou repulsão) da força entre elas.
\dados{$\ke = \SI{9e9}{\newton\meter\squared\per\coulomb\squared}$.}
\resposta{$F = \SI{33.75}{\newton}$, atrativa}


\end{questao}
\begin{questao}[1]{}

Duas cargas $q_1 = q_2 = \SI{+5}{\micro\coulomb}$ estão a
$\SI{2}{\centi\meter}$ no vácuo. Determine o módulo e o tipo da força.
\dados{$\ke = \SI{9e9}{\newton\meter\squared\per\coulomb\squared}$.}
\resposta{$F = \SI{562.5}{\newton}$, repulsiva}


\end{questao}
\begin{questao}[1]{}

Duas cargas puntiformes se repelem com força de intensidade
$\pot{42}{-5}\,\si{\newton}$. Reduzindo-se a distância entre elas à
\textbf{metade}, a nova força de repulsão será:
\begin{alternativas}[2]
\item $\pot{10.5}{-5}\,\si{\newton}$
\item $\pot{21}{-5}\,\si{\newton}$
\item $\pot{84}{-5}\,\si{\newton}$
\item $\pot{168}{-5}\,\si{\newton}$
\item inalterada
\end{alternativas}
\resposta{(d)}


\end{questao}

\blocoaplicacao
\begin{questao}[2]{}

A força elétrica entre duas partículas de cargas $q$ e $q/2$, separadas por uma
distância $d$ no vácuo, é $F$. A força entre duas partículas de cargas $q$ e $2q$,
separadas por $d/2$, também no vácuo, é:
\begin{alternativas}[3]
\item $F$
\item $2F$
\item $4F$
\item $8F$
\item $16F$
\end{alternativas}
\resposta{(e)}


\end{questao}
\begin{questao}[2]{}

Se a força entre dois objetos carregados se mantém \emph{constante} mesmo quando a
carga de cada um é reduzida à metade, então a distância entre eles:
\begin{alternativas}
\item foi quadruplicada.
\item foi duplicada.
\item foi reduzida à quarta parte.
\item foi reduzida à metade.
\item permaneceu constante.
\end{alternativas}
\resposta{(d)}


\end{questao}
\begin{questao}[2]{}

Três cargas iguais de $\SI{+1}{\micro\coulomb}$ ocupam os vértices de um triângulo
equilátero de lado $\SI{5}{\centi\meter}$. Determine o módulo da força resultante
sobre uma delas.

\circuitoq{
\begin{tikzpicture}[scale=1,line width=0.8pt,>=Stealth]
 \coordinate (A) at (0,0);
 \coordinate (B) at (3,0);
 \coordinate (C) at (1.5,2.6);
 \foreach \p/\n in {A/A,B/B,C/C}{
   \shade[ball color=Vcol!25] (\p) circle (0.28);
   \draw[Vcol] (\p) circle (0.28);
   \node at (\p) {\footnotesize$+q$};}
 \draw[dashed,cinza] (A)--(B)--(C)--cycle;
 \node[below] at (1.5,-0.35) {$\SI{5}{\centi\meter}$};
\end{tikzpicture}}{Três cargas iguais nos vértices de um triângulo equilátero.}

\resposta{$F_R \approx \SI{6.24}{\newton}$}


\end{questao}
\begin{questao}[2]{}

Quatro cargas positivas iguais de $\SI{1}{\micro\coulomb}$ ocupam os vértices de
um quadrado de lado $\SI{10}{\centi\meter}$. Determine o módulo da força
resultante sobre uma delas.
\dados{$\ke = \SI{9e9}{\newton\meter\squared\per\coulomb\squared}$; $\sqrt2\approx1{,}41$.}

\circuitoq{
\begin{tikzpicture}[scale=1,line width=0.8pt,>=Stealth]
 \foreach \p in {(0,0),(2.4,0),(2.4,2.4),(0,2.4)}{
   \shade[ball color=Vcol!25] \p circle (0.22);}
 \node at (0,0){\footnotesize$+q$};\node at (2.4,0){\footnotesize$+q$};
 \node at (2.4,2.4){\footnotesize$+q$};\node at (0,2.4){\footnotesize$+q$};
 \draw[dashed,cinza] (0,0)--(2.4,0)--(2.4,2.4)--(0,2.4)--cycle;
 \draw[dashed,cinza] (0,0)--(2.4,2.4);
 \node[below] at (1.2,-0.3){$\SI{10}{\centi\meter}$};
\end{tikzpicture}}{Quatro cargas iguais nos vértices de um quadrado.}

\resposta{$F_R \approx \SI{1.72}{\newton}$}


\end{questao}
\begin{questao}[2]{}

Uma pequena esfera de massa $\SI{2}{\gram}$, eletrizada com carga $q$, permanece
em equilíbrio, suspensa, sob a ação combinada do peso e de um campo elétrico
uniforme vertical de intensidade $\SI{e5}{\newton\per\coulomb}$, dirigido para
cima. Determine a carga $q$.
\dados{$g = \SI{10}{\meter\per\second\squared}$.}
\resposta{$q = \SI{0.2}{\micro\coulomb}$}


\end{questao}

\blocoaprofundamento
\begin{questao}[3]{}

Duas cargas $Q_1$ e $Q_2 = 4Q_1$ (ambas positivas) estão fixas nos pontos A e B,
distantes $\SI{30}{\centi\meter}$. Em que posição $x$, medida a partir de A sobre
a reta AB, deve-se colocar uma carga $Q_3$ para que ela fique em equilíbrio sob
ação apenas de forças elétricas?

\circuitoq{
\begin{tikzpicture}[scale=1,line width=0.8pt]
 \draw[cinza] (0,0)--(6,0);
 \shade[ball color=Vcol!25] (0,0) circle (0.25);
 \shade[ball color=Vcol!25] (6,0) circle (0.25);
 \node at (0,0) {\footnotesize$Q_1$}; \node at (6,0) {\footnotesize$Q_2$};
 \node[below] at (0,-0.35) {A}; \node[below] at (6,-0.35) {B};
 \fill[laranja] (2,0) circle (0.09);
 \node[above,laranja] at (2,0.15) {$Q_3$};
 \draw[<->,cinza] (0,-0.9)--(2,-0.9) node[midway,below]{$x$};
 \draw[<->,cinza] (0,-1.5)--(6,-1.5) node[midway,below]{$\SI{30}{\centi\meter}$};
\end{tikzpicture}}{Posição de equilíbrio entre duas cargas.}

\resposta{$x = \SI{10}{\centi\meter}$ a partir de A}


\end{questao}
\begin{questao}[3]{}

Quatro cargas ocupam os vértices de um quadrado: duas delas são $+Q$ e $-Q$ (mesmo
módulo), e as outras duas, $q_1$ e $q_2$, são desconhecidas. Coloca-se uma carga
de prova positiva no centro e observa-se que a força resultante sobre ela aponta
na direção da diagonal, \emph{do vértice $-Q$ para o vértice $+Q$}. Sobre $q_1$ e
$q_2$ (nos outros dois vértices), pode-se afirmar que, para que essa força
resultante fique exatamente sobre a diagonal $-Q\,$--$\,+Q$:
\begin{alternativas}
\item devem ser iguais em módulo e de mesmo sinal entre si.
\item devem ser iguais em módulo e de sinais opostos.
\item devem ser ambas nulas.
\item devem ser ambas positivas e diferentes.
\item nada se pode afirmar.
\end{alternativas}
\resposta{(a)}


\end{questao}
\begin{questao}[3]{}

Duas cargas fixas, $q_1 = \SI{+9}{\micro\coulomb}$ em $x=0$ e
$q_2 = \SI{+4}{\micro\coulomb}$ em $x = \SI{10}{\centi\meter}$, estão sobre o eixo
$x$. Uma terceira carga $q_3$ é colocada sobre esse eixo.
\begin{itensq}
\item Determine a posição em que $q_3$ fica em equilíbrio.
\item Mostre que existe uma \emph{segunda} raiz matemática e explique por que ela
      deve ser descartada.
\item O equilíbrio encontrado é \emph{estável} ou \emph{instável}? Analise para
      $q_3$ positiva e negativa.
\end{itensq}
\resposta{(a) $x = \SI{6}{\centi\meter}$; (b) $x=\SI{30}{\centi\meter}$ (fora do
segmento); (c) estável para $q_3<0$, instável para $q_3>0$}


\end{questao}
\begin{questao}[3]{}

Compare a força elétrica e a força gravitacional entre dois elétrons.
\begin{itensq}
\item Mostre que a razão $F_e/F_g$ \emph{independe} da distância.
\item Calcule seu valor numérico.
\item Que conclusão isso permite sobre a estrutura da matéria e sobre a
      astronomia?
\end{itensq}
\dados{$\ke=\SI{9e9}{}$; $G=\SI{6.67e-11}{}$; $e=\SI{1.6e-19}{\coulomb}$;
$m_e=\SI{9.11e-31}{\kilogram}$ (SI).}
\resposta{(a) demonstração; (b) $\approx\pot{4.2}{42}$; (c) ver comentário}


\end{questao}
\begin{questao}[3]{}

Duas pequenas esferas idênticas, de massa $\SI{0.5}{\gram}$ cada, são suspensas
do mesmo ponto por fios isolantes de $\SI{20}{\centi\meter}$. Ao receberem cargas
iguais $q$, elas se afastam até que cada fio forme $\SI{10}{\degree}$ com a
vertical.

\circuitoq{
\begin{tikzpicture}[scale=1,line width=0.8pt,>=Stealth]
 \draw[cinza!60,line width=2pt] (-1.6,3)--(1.6,3);
 \coordinate (O) at (0,3);
 \coordinate (P1) at ({-2.4*sin(10)},{3-2.4*cos(10)});
 \coordinate (P2) at ({ 2.4*sin(10)},{3-2.4*cos(10)});
 \draw[cinza,thick] (O)--(P1); \draw[cinza,thick] (O)--(P2);
 \draw[dashed,cinza!70] (O)--(0,0.3);
 \shade[ball color=Vcol!30] (P1) circle (0.2);
 \shade[ball color=Vcol!30] (P2) circle (0.2);
 \node[font=\footnotesize] at (0.34,2.55) {$10^\circ$};
 \node[left,font=\footnotesize] at ($(P1)+(-0.25,0)$) {$q$};
 \node[right,font=\footnotesize] at ($(P2)+(0.25,0)$) {$q$};
 \draw[<->,cinza] ($(P1)+(0,-0.45)$)--($(P2)+(0,-0.45)$)
       node[midway,below,font=\footnotesize]{$d$};
 \node[right,font=\footnotesize,cinza] at (0.9,1.9) {$L=\SI{20}{\centi\meter}$};
\end{tikzpicture}}{Pêndulos eletrostáticos idênticos.}

\begin{itensq}
\item Determine a separação $d$ entre as esferas.
\item Calcule a força elétrica entre elas.
\item Determine a carga $q$ de cada esfera.
\end{itensq}
\dados{$g=\SI{10}{\meter\per\second\squared}$; $\sin10^\circ=0{,}174$;
$\tan10^\circ=0{,}176$; $\ke=\SI{9e9}{}$.}
\resposta{(a) $d\approx\SI{6.9}{\centi\meter}$; (b) $F\approx\SI{0.88}{\milli\newton}$;
(c) $q\approx\SI{22}{\nano\coulomb}$}


\end{questao}
\begin{questao}[3]{}

Quatro cargas de mesmo módulo $q=\SI{1}{\micro\coulomb}$ ocupam os vértices de um
quadrado de lado $a=\SI{10}{\centi\meter}$, \textbf{alternando os sinais}
($+,-,+,-$ ao percorrer o perímetro). Determine o módulo e a direção da força
resultante sobre uma das cargas positivas.
\dados{$\ke=\SI{9e9}{}$; $\sqrt2\approx1{,}41$.}
\resposta{$F_R \approx \SI{0.82}{\newton}$, dirigida ao centro do quadrado}


\end{questao}
\begin{questao}[3]{}

Duas cargas fixas, $q_1=+q$ e $q_2=+4q$, estão separadas por $d$. Deseja-se
colocar uma terceira carga $q_3$ de modo que \textbf{as três} fiquem em
equilíbrio simultaneamente.
\begin{itensq}
\item Determine a posição de $q_3$.
\item Determine o valor e o sinal de $q_3$.
\item Esse equilíbrio é estável?
\end{itensq}
\resposta{(a) a $d/3$ de $q_1$; (b) $q_3=-\dfrac{4q}{9}$; (c) instável}


\end{questao}
\begin{questao}[3]{}

Uma pequena esfera de massa $m = \SI{10}{\gram}$, eletrizada com carga $q$, é
suspensa por um fio isolante de massa desprezível em uma região de campo elétrico
\emph{horizontal} e uniforme, de intensidade $E = \SI{e4}{\newton\per\coulomb}$.
No equilíbrio, o fio forma $\SI{30}{\degree}$ com a vertical.

\circuitoq{
\begin{tikzpicture}[scale=1,line width=0.8pt,>=Stealth]
 \draw[cinza!60,line width=2pt] (-2,2.6)--(2,2.6);
 \coordinate (O) at (0,2.6);
 \coordinate (P) at ({2.2*sin(30)},{2.6-2.2*cos(30)});
 \draw[cinza,thick] (O)--(P);
 \draw[dashed,cinza!70] (O)--(0,0.2);
 \shade[ball color=Vcol!30] (P) circle (0.22);
 \node[right,font=\footnotesize] at ($(P)+(0.28,0)$) {$m,\,q$};
 \draw[->,Vcol,thick] (P)--($(P)+(1.0,0)$) node[right,font=\footnotesize]{$\vec F_e$};
 \draw[->,cinza,thick] (P)--($(P)+(0,-0.9)$) node[below,font=\footnotesize]{$\vec P$};
 \node[font=\footnotesize] at (0.42,1.9) {$30^\circ$};
 \foreach \y in {0.4,1.0,1.6} \draw[->,Vcol!45] (-1.9,\y)--(-1.2,\y);
 \node[Vcol,font=\footnotesize] at (-1.55,2.05) {$\vec E$};
\end{tikzpicture}}{Pêndulo eletrostático em campo horizontal.}

\begin{itensq}
\item Determine a carga $q$ da esfera.
\item Calcule a tração no fio.
\item Se o campo dobrasse, o ângulo dobraria? Justifique quantitativamente.
\end{itensq}
\dados{$g = \SI{10}{\meter\per\second\squared}$; $\tan30^\circ\approx0{,}577$.}
\resposta{(a) $q \approx \SI{5.8}{\micro\coulomb}$; (b) $T\approx\SI{0.115}{\newton}$;
(c) não --- o ângulo cresce menos que proporcionalmente}


\end{questao}
\begin{questao}[3]{}

Três cargas puntiformes estão sobre uma reta: $+Q$ na origem, $-2Q$ em $x=a$ e
$+Q$ em $x=2a$ (um \emph{quadrupolo linear}).
\begin{itensq}
\item Determine a força resultante sobre a carga central.
\item Determine a força resultante sobre a carga da extremidade em $x=0$.
\item Explique por que o sistema, embora tenha carga total nula, ainda exerce
      força sobre uma carga externa distante.
\end{itensq}
\resposta{(a) nula, por simetria; (b) $F=\dfrac{7\ke Q^{2}}{4a^{2}}$ dirigida para
$+x$; (c) ver comentário}


\end{questao}

\blocodesafio
\begin{questao}[4]{}

Uma carga $Q$ é dividida em duas partes, $q$ e $(Q-q)$, que são então colocadas a
uma distância fixa $d$ uma da outra.
\begin{itensq}
\item Obtenha a expressão da força entre elas em função de $q$.
\item Determine o valor de $q$ que \textbf{maximiza} essa força.
\item Qual é a força máxima?
\end{itensq}
\resposta{(a) $F=\dfrac{\ke\,q(Q-q)}{d^{2}}$; (b) $q=Q/2$;
(c) $F_{\max}=\dfrac{\ke Q^{2}}{4d^{2}}$}


\end{questao}

\blocofixacao
\begin{questao}[1]{}

Duas cargas puntiformes $q_1=\SI{20}{\nano\coulomb}$ e
$q_2=\SI{30}{\nano\coulomb}$ estão separadas por $\SI{3}{\centi\meter}$ no
vácuo. Determine o módulo da força elétrica e classifique-a.
\dados{$\ke=\SI{9e9}{\newton\meter\squared\per\coulomb\squared}$}
\resposta{$F=\SI{6}{\milli\newton}$, repulsiva}


\end{questao}
\begin{questao}[1]{}

Duas esferas idênticas recebem cargas iguais $q$ e se repelem com força de
$\SI{0.90}{\newton}$ quando separadas por $\SI{20}{\centi\meter}$.
\begin{itensq}
\item Determine $q$.
\item Quantos elétrons foram \textbf{retirados} de cada esfera?
\end{itensq}
\dados{$e=\pot{1.6}{-19}\,\si{\coulomb}$}
\resposta{(a) $q=\SI{2}{\micro\coulomb}$; (b) $n=\pot{1.25}{13}$ elétrons}


\end{questao}
\begin{questao}[1]{}

Duas cargas de $\SI{5}{\micro\coulomb}$ cada devem exercer uma sobre a outra
força de $\SI{2.5}{\newton}$ no vácuo. A que distância devem ser colocadas?
\resposta{$r=\SI{0.30}{\meter}$}


\end{questao}
\begin{questao}[1]{}

Duas cargas puntiformes se atraem com força $F$. Mantendo as cargas e
\textbf{dobrando} a distância entre elas, a nova força será:
\begin{alternativas}[3]
\item $F/4$
\item $F/2$
\item $F$
\item $2F$
\item $4F$
\end{alternativas}
\resposta{(a)}


\end{questao}
\begin{questao}[1]{}

Duas cargas separadas por uma distância fixa exercem força de
$\SI{0.80}{\newton}$ no vácuo. Todo o conjunto é então imerso em um líquido
isolante de permissividade relativa $\varepsilon_r=4$. Determine a nova força.
\resposta{$F'=\SI{0.20}{\newton}$}


\end{questao}
\begin{questao}[1]{}

Um bastão carregado positivamente é aproximado --- sem tocar --- de uma esfera
metálica \textbf{neutra} e isolada. A força entre eles é:
\begin{alternativas}[2]
\item nula, pois a esfera é neutra
\item atrativa
\item repulsiva
\item atrativa ou repulsiva, dependendo da distância
\end{alternativas}
\resposta{(b)}


\end{questao}
\begin{questao}[1]{}

Estime a força entre duas cargas de $\SI{1}{\coulomb}$ separadas por
$\SI{1}{\kilo\meter}$ e comente o resultado.
\resposta{$F=\SI{9}{\kilo\newton}$ (o peso de cerca de \SI{900}{\kilogram})}


\end{questao}

\blocoaplicacao
\begin{questao}[2]{}

Três cargas estão fixas sobre o eixo $x$: $q_1=\SI{+3}{\micro\coulomb}$ em
$x=0$, $q_2=\SI{+2}{\micro\coulomb}$ em $x=\SI{20}{\centi\meter}$ e
$q_3=\SI{-5}{\micro\coulomb}$ em $x=\SI{50}{\centi\meter}$. Determine a força
resultante sobre $q_2$ (módulo e sentido).
\resposta{$F=\SI{2.35}{\newton}$, no sentido $+x$}


\end{questao}
\begin{questao}[2]{}

Duas esferas metálicas idênticas com $q_1=\SI{+8}{\micro\coulomb}$ e
$q_2=\SI{-2}{\micro\coulomb}$ estão a $\SI{30}{\centi\meter}$. Elas são postas
em contato e recolocadas na mesma posição.
\begin{itensq}
\item Calcule a força antes do contato.
\item Calcule a força depois.
\end{itensq}
\resposta{(a) $\SI{1.6}{\newton}$, atrativa; (b) $\SI{0.9}{\newton}$, repulsiva}


\end{questao}
\begin{questao}[2]{}

Duas cargas fixas exercem entre si $\SI{0.60}{\newton}$ no vácuo. O conjunto é
mergulhado em óleo com $\varepsilon_r=2{,}5$.
\begin{itensq}
\item Qual é a nova força, mantida a distância?
\item Que fração da distância original restauraria a força de
      $\SI{0.60}{\newton}$?
\end{itensq}
\resposta{(a) $\SI{0.24}{\newton}$; (b) $r'=r/\sqrt{2{,}5}\approx 0{,}63\,r$}


\end{questao}
\begin{questao}[2]{}

A força entre duas cargas $q_1$ e $q_2$ separadas por $d$ vale $F$. Se $q_1$ for
\textbf{triplicada} e a distância for reduzida à \textbf{metade}, a nova força
será:
\begin{alternativas}[3]
\item $1{,}5F$
\item $3F$
\item $6F$
\item $12F$
\item $24F$
\end{alternativas}
\resposta{(d)}


\end{questao}
\begin{questao}[2]{}

Em uma impressora a jato eletrostático, cada gota de tinta tem massa
$\pot{1.0}{-11}\,\si{\kilogram}$ e é carregada pela remoção de
$\pot{2.0}{5}$ elétrons. Duas gotas vizinhas ficam a $\SI{1.0}{\milli\meter}$.
Compare a força elétrica entre elas com o peso de uma gota.
\dados{$e=\pot{1.6}{-19}\,\si{\coulomb}$; $g=\SI{9.8}{\meter\per\second\squared}$}
\resposta{$F_e=\pot{9.2}{-12}\,\si{\newton}$; $P=\pot{9.8}{-11}\,\si{\newton}$;
$F_e/P\approx 0{,}094$}


\end{questao}
\begin{questao}[2]{}

Duas cargas positivas, $q_1=\SI{4}{\micro\coulomb}$ em $x=0$ e
$q_2=\SI{9}{\micro\coulomb}$ em $x=\SI{50}{\centi\meter}$, estão fixas.
Determine o ponto onde uma terceira carga ficaria em equilíbrio.
\resposta{$x=\SI{20}{\centi\meter}$}


\end{questao}
\begin{questao}[2]{}

Uma carga $q$ e outra $100q$ interagem. Sobre os módulos das forças que uma
exerce sobre a outra, é correto afirmar:
\begin{alternativas}[2]
\item a força sobre $q$ é 100 vezes maior
\item a força sobre $100q$ é 100 vezes maior
\item as forças têm o mesmo módulo
\item depende das massas
\end{alternativas}
\resposta{(c)}


\end{questao}
\begin{questao}[2]{}

No modelo de Bohr para o hidrogênio, o elétron orbita o próton a
$\pot{5.3}{-11}\,\si{\meter}$. Determine a força elétrica e a aceleração
centrípeta do elétron.
\dados{$e=\pot{1.6}{-19}\,\si{\coulomb}$; $m_e=\pot{9.11}{-31}\,\si{\kilogram}$}
\resposta{$F=\pot{8.2}{-8}\,\si{\newton}$;
$a=\pot{9.0}{22}\,\si{\meter\per\second\squared}$}


\end{questao}
\begin{questao}[2]{}

Duas esferas metálicas idênticas com cargas $+Q$ e $+3Q$, separadas por $d$,
repelem-se com força $F$. Elas são postas em contato e depois afastadas até
$2d$. Determine a nova força em função de $F$.
\resposta{$F'=F/3$}


\end{questao}
\begin{questao}[2]{}

Compare a força entre duas cargas de $\SI{1.0}{\micro\coulomb}$ separadas por
$\SI{1.0}{\centi\meter}$ com o peso de um corpo de $\SI{1.0}{\kilogram}$, e
interprete.
\dados{$g=\SI{9.8}{\meter\per\second\squared}$}
\resposta{$F=\SI{90}{\newton}$ contra $P=\SI{9.8}{\newton}$; razão $\approx 9{,}2$}


\end{questao}

\blocoaprofundamento
\begin{questao}[3]{}

Duas cargas fixas sobre o eixo $x$: $q_1=\SI{+4}{\micro\coulomb}$ em $x=0$ e
$q_2=\SI{-1}{\micro\coulomb}$ em $x=\SI{30}{\centi\meter}$. Determine todos os
pontos do eixo onde a força sobre uma carga de prova seria nula.
\resposta{Um único ponto: $x=\SI{60}{\centi\meter}$}


\end{questao}
\begin{questao}[3]{}

Seis cargas iguais $Q=\SI{2}{\micro\coulomb}$ ocupam os vértices de um hexágono
regular de lado $a=\SI{10}{\centi\meter}$, com uma carga
$q=\SI{1}{\micro\coulomb}$ no centro.
\begin{itensq}
\item Qual é a força resultante sobre $q$?
\item Uma das seis cargas é removida. Qual passa a ser a força sobre $q$
      (módulo e direção)?
\end{itensq}
\resposta{(a) nula; (b) $\SI{1.8}{\newton}$, apontando \emph{para} o vértice vazio}


\end{questao}
\begin{questao}[3]{}

Três cargas ocupam os vértices de um triângulo retângulo:
$q_A=\SI{+2}{\micro\coulomb}$ na origem, $q_B=\SI{+3}{\micro\coulomb}$ em
$(\SI{30}{\centi\meter},0)$ e $q_C=\SI{+4}{\micro\coulomb}$ em
$(0,\SI{40}{\centi\meter})$. Determine a força resultante sobre $q_A$: módulo e
ângulo com o eixo $x$.
\resposta{$F=\SI{0.75}{\newton}$, a $\SI{216.9}{\degree}$ (isto é, $\SI{36.9}{\degree}$ abaixo
do semieixo $-x$)}


\end{questao}
\begin{questao}[3]{}

Duas pequenas esferas metálicas idênticas, separadas por $\SI{10}{\centi\meter}$,
\textbf{atraem-se} com força de $\SI{0.108}{\newton}$. Postas em contato e
recolocadas na mesma distância, passam a \textbf{repelir-se} com
$\SI{0.036}{\newton}$. Determine as cargas iniciais.
\resposta{$q_1=\SI{+0.6}{\micro\coulomb}$ e $q_2=\SI{-0.2}{\micro\coulomb}$
(ou a permuta)}


\end{questao}
\begin{questao}[3]{}

Um pequeno bloco de massa $m=\SI{20}{\gram}$, com carga $q=\SI{1}{\micro\coulomb}$,
repousa sobre um plano inclinado \textbf{liso} de $\SI{30}{\degree}$. Uma carga fixa
$Q=\SI{1}{\micro\coulomb}$ está presa na base do plano, e a força elétrica atua
ao longo da superfície. Determine a distância de equilíbrio $d$ medida sobre o
plano.
\dados{$g=\SI{9.8}{\meter\per\second\squared}$}
\resposta{$d\approx\SI{30.3}{\centi\meter}$}


\end{questao}
\begin{questao}[3]{}

Um dipolo é formado por $+q$ e $-q$, com $q=\SI{1}{\micro\coulomb}$, fixos em
$(\pm\SI{5}{\centi\meter},0)$. Uma carga $Q=\SI{2}{\micro\coulomb}$ é colocada
em $(0,\SI{12}{\centi\meter})$, sobre a mediatriz.
\begin{itensq}
\item Mostre que a resultante sobre $Q$ é paralela ao eixo do dipolo.
\item Calcule seu módulo.
\end{itensq}
\resposta{(a) as componentes verticais se cancelam; (b) $F=\SI{0.82}{\newton}$,
no sentido $-x$ (de $+q$ para $-q$)}


\end{questao}

\blocodesafio
\begin{questao}[4]{}

\textbf{(Modelagem experimental.)} Em uma reprodução da balança de torção, duas
esferas com cargas iguais $q=\SI{0.10}{\micro\coulomb}$ foram mantidas a várias
distâncias e a força foi medida:

\begin{center}
\begin{tabular}{cccccc}
\hline
$r$ (m) & 0,020 & 0,030 & 0,040 & 0,050 & 0,060\\
$F$ (N) & 0,223 & 0,101 & 0,0558 & 0,0362 & 0,0249\\
\hline
\end{tabular}
\end{center}

\begin{itensq}
\item Proponha uma linearização que permita testar a hipótese $F=C\,r^{n}$.
\item Estime $n$ e discuta se os dados sustentam a lei do inverso do quadrado.
\item Obtenha o valor experimental de $\ke$ e compare com o tabelado.
\end{itensq}
\resposta{(a) $\log F=\log C+n\log r$; (b) $n\approx-2{,}00$; (c)
$\ke\approx\pot{9.1}{9}\,\si{\newton\meter\squared\per\coulomb\squared}$
(erro $\approx 1\%$)}


\end{questao}
\begin{questao}[4]{}

\textbf{(Otimização com cálculo.)} Um anel fino de raio $R$, com carga total $Q$
uniformemente distribuída, está no plano $yz$. Uma carga puntiforme $q$ é
colocada sobre o eixo do anel, a uma distância $x$ do centro.
\begin{itensq}
\item Mostre que $F(x)=\dfrac{\ke Q q\,x}{(x^{2}+R^{2})^{3/2}}$.
\item Determine o valor de $x$ que \textbf{maximiza} a força.
\item Calcule $F_{\max}$ para $Q=\SI{10}{\micro\coulomb}$,
      $q=\SI{2}{\micro\coulomb}$ e $R=\SI{10}{\centi\meter}$.
\end{itensq}
\resposta{(b) $x=R/\sqrt{2}\approx\SI{7.07}{\centi\meter}$;
(c) $F_{\max}=\dfrac{2\ke Qq}{3\sqrt{3}\,R^{2}}=\SI{6.93}{\newton}$}


\end{questao}
\begin{questao}[4]{}

\textbf{(Distribuição contínua.)} Um bastão isolante retilíneo de comprimento
$L$ possui carga $Q$ uniformemente distribuída. Uma carga puntiforme $q$ está
sobre o prolongamento do bastão, a uma distância $d$ da extremidade mais
próxima.
\begin{itensq}
\item Deduza $F=\dfrac{\ke Qq}{d\,(d+L)}$.
\item Verifique o comportamento para $L\ll d$.
\item Calcule $F$ para $Q=\SI{4}{\micro\coulomb}$, $L=\SI{20}{\centi\meter}$,
      $q=\SI{1}{\micro\coulomb}$ e $d=\SI{10}{\centi\meter}$.
\end{itensq}
\resposta{(b) $F\to\ke Qq/d^{2}$; (c) $F=\SI{1.2}{\newton}$}


\end{questao}
\begin{questao}[4]{}

\textbf{(Demonstração --- estabilidade.)} Duas cargas $+Q$ estão fixas em
$x=-d$ e $x=+d$. Uma carga $q$ é colocada na origem.
\begin{itensq}
\item Mostre que a origem é posição de equilíbrio para qualquer $q$.
\item Para $q>0$, analise a estabilidade a um pequeno deslocamento
      \emph{longitudinal} ($\pm x$).
\item Repita para um deslocamento \emph{transversal} ($\pm y$) e conclua.
\end{itensq}
\resposta{(a) por simetria; (b) estável, com $F\approx-4\ke Qq\,x/d^{3}$;
(c) instável, com $F\approx+2\ke Qq\,y/d^{3}$ --- não há equilíbrio estável em
todas as direções}


\end{questao}
\begin{questao}[4]{}

\textbf{(Série infinita.)} Cargas puntiformes de mesmo módulo $Q$ e
\textbf{sinais alternados} ocupam as posições $x=a,\,2a,\,3a,\dots$ sobre um
eixo: $-Q$ em $a$, $+Q$ em $2a$, $-Q$ em $3a$ e assim por diante,
indefinidamente. Uma carga $+q$ está na origem.
\begin{itensq}
\item Escreva a força resultante sobre $q$ como uma série.
\item Some a série e obtenha a expressão fechada.
\end{itensq}
\dados{$\displaystyle\sum_{n=1}^{\infty}\frac{(-1)^{n+1}}{n^{2}}=\frac{\pi^{2}}{12}$}
\resposta{$F=\dfrac{\pi^{2}}{12}\cdot\dfrac{\ke Qq}{a^{2}}
\approx 0{,}822\,\dfrac{\ke Qq}{a^{2}}$, atrativa (sentido $+x$)}


\end{questao}
\begin{questao}[4]{}

\textbf{(Oscilações.)} Uma partícula de massa $m=\SI{1.0}{\gram}$ e carga
$-q=\SI{-2}{\micro\coulomb}$ é liberada em repouso sobre o eixo de um anel de
raio $R=\SI{10}{\centi\meter}$ e carga $Q=\SI{5}{\micro\coulomb}$, a uma
distância $x_0\ll R$ do centro.
\begin{itensq}
\item Mostre que, para $x\ll R$, o movimento é harmônico simples.
\item Obtenha $\omega$ e o período $T$.
\end{itensq}
\resposta{(b) $\omega=\sqrt{\ke Qq/(mR^{3})}=\SI{300}{\radian\per\second}$;
$T\approx\SI{20.9}{\milli\second}$}


\end{questao}
\begin{questao}[4]{}

\textbf{(Modelagem com taxas.)} Duas esferas idênticas de massa
$m=\SI{2.0}{\gram}$, cada uma com carga $q$, pendem de fios isolantes de
comprimento $L=\SI{60}{\centi\meter}$ presos ao mesmo ponto. Elas se afastam até
uma separação $x\ll L$. Por fuga de carga para o ar úmido, as esferas se
aproximam lentamente.
\begin{itensq}
\item Mostre que $x^{3}=\dfrac{2\ke q^{2}L}{mg}$.
\item Demonstre que
      $\dfrac{\mathrm{d} x}{\mathrm{d} t}=\dfrac{2x}{3q}\dfrac{\mathrm{d} q}{\mathrm{d} t}$.
\item Se em certo instante $x=\SI{6.0}{\centi\meter}$ e as esferas se aproximam
      a $\SI{1.0}{\milli\meter\per\second}$, calcule $\mathrm{d} q/\mathrm{d} t$.
\end{itensq}
\dados{$g=\SI{9.8}{\meter\per\second\squared}$}
\resposta{(c) $q\approx\SI{19.8}{\nano\coulomb}$ e
$\mathrm{d} q/\mathrm{d} t\approx\SI{-0.49}{\nano\coulomb\per\second}$}


\end{questao}
\begin{questao}[4]{}

\textbf{(Simetria e integração.)} Um fio semicircular de raio
$R=\SI{5.0}{\centi\meter}$ possui carga $Q=\SI{3}{\micro\coulomb}$ distribuída
uniformemente. Uma carga $q=\SI{1}{\micro\coulomb}$ ocupa o centro do
semicírculo.
\begin{itensq}
\item Justifique por que a resultante é dirigida ao longo do eixo de simetria.
\item Mostre que $F=\dfrac{2\ke Qq}{\pi R^{2}}$ e calcule seu valor.
\end{itensq}
\resposta{$F=\SI{6.9}{\newton}$, ao longo do eixo de simetria}


\end{questao}
\begin{questao}[4]{}

\textbf{(Síntese --- equilíbrio de um sistema.)} Quatro cargas iguais
$Q=\SI{2}{\micro\coulomb}$ ocupam os vértices de um quadrado de lado $a$. Uma
quinta carga $q$ é colocada no centro.
\begin{itensq}
\item Determine $q$ (sinal e módulo, em função de $Q$) para que \textbf{cada}
      carga dos vértices fique em equilíbrio.
\item Calcule $q$ numericamente.
\item O sistema completo está em equilíbrio estável? Justifique.
\end{itensq}
\resposta{(a) $q=-\dfrac{Q(2\sqrt{2}+1)}{4}$; (b) $q\approx\SI{-1.91}{\micro\coulomb}$;
(c) não --- o equilíbrio é instável}


\end{questao}

\gabarito[1.2 Lei de Coulomb]

\subsection{Campo elétrico}

\blocofixacao
\begin{questao}[1]{}

Qual é a intensidade do campo elétrico gerado por uma carga de
$\SI{-15}{\micro\coulomb}$, a $\SI{10}{\milli\meter}$ dela, em um meio de constante
$\ke = \SI{8e9}{\newton\meter\squared\per\coulomb\squared}$?
\resposta{$E = \pot{1.2}{9}\,\si{\newton\per\coulomb}$}


\end{questao}
\begin{questao}[1]{}

O módulo do campo elétrico de uma carga puntiforme $q$, em um ponto P a uma dada
distância, vale $E$. Afastando-se a carga de modo que a distância a P
\textbf{dobre}, o campo em P passa a ser:
\begin{alternativas}[3]
\item $E$
\item $E/2$
\item $E/4$
\item $2E$
\item $4E$
\end{alternativas}
\resposta{(c)}


\end{questao}

\blocoaplicacao
\begin{questao}[2]{}

Uma carga de prova de $\SI{e-9}{\coulomb}$, colocada em um ponto P de um campo
elétrico, fica sujeita a uma força de $\SI{e-2}{\newton}$, vertical e para baixo.
Determine:
\begin{itensq}
\item a intensidade, direção e sentido do campo em P;
\item a força que atuaria sobre uma carga de $\SI{3}{\milli\coulomb}$ colocada em P.
\end{itensq}
\resposta{(a) $E=\pot{1.0}{7}\,\si{\newton\per\coulomb}$, vertical para baixo;
(b) $F = \pot{3}{4}\,\si{\newton}$}


\end{questao}
\begin{questao}[2]{}

Uma carga positiva $Q = \SI{4.0}{\micro\coulomb}$ está no ar. Determine a
intensidade, a direção e o sentido do campo elétrico em um ponto M, a
$\SI{20}{\centi\meter}$ da carga.
\dados{$\ke = \SI{9e9}{\newton\meter\squared\per\coulomb\squared}$.}
\resposta{$E = \pot{9.0}{5}\,\si{\newton\per\coulomb}$, radial, afastando-se de $Q$}


\end{questao}
\begin{questao}[2]{}

Duas cargas puntiformes, $Q$ e $4Q$ (de mesmo sinal), estão fixas no vácuo. O
módulo do campo elétrico resultante é \textbf{nulo} em um ponto sobre a reta que
as une. Esse ponto está tal que sua distância à carga menor, $d_1$, e à carga
maior, $d_2$, satisfazem:
\begin{alternativas}[2]
\item $d_2 = d_1$
\item $d_2 = 2d_1$
\item $d_2 = 4d_1$
\item $d_1 = 2d_2$
\item $d_1 = 4d_2$
\end{alternativas}
\resposta{(b)}


\end{questao}
\begin{questao}[2]{}

Duas cargas de mesmo módulo $Q$ estão fixas, separadas por $\SI{20}{\centi\meter}$.
Analise o campo elétrico resultante no ponto médio M do segmento que as une, em
dois casos:

\circuitoq{
\begin{tikzpicture}[scale=0.9,line width=0.8pt,>=Stealth]
 
 \begin{scope}
   \shade[ball color=Vcol!25] (-1.6,0) circle (0.25); \node at (-1.6,0){\footnotesize$+$};
   \shade[ball color=Icol!25] (1.6,0) circle (0.25);   \node at (1.6,0){\footnotesize$-$};
   \fill[laranja] (0,0) circle (0.06); \node[above,font=\footnotesize] at (0,0.12){M};
   \draw[Vcol,->,thick] (0,-0.5)--(1.1,-0.5); \node[Vcol,font=\footnotesize] at (0.55,-0.8){$\vec E_1$};
   \draw[Vcol,->,thick] (0,-1.1)--(1.1,-1.1); \node[Vcol,font=\footnotesize] at (0.55,-1.4){$\vec E_2$};
   \node[font=\footnotesize] at (0,1.0){Caso I: $+Q$ e $-Q$};
 \end{scope}
\end{tikzpicture}}{Campo no ponto médio entre duas cargas.}

No \textbf{Caso I} ($+Q$ e $-Q$) e no \textbf{Caso II} (ambas $+Q$), o campo
resultante em M é, respectivamente:
\begin{alternativas}
\item nulo em I e nulo em II.
\item não-nulo em I e nulo em II.
\item nulo em I e não-nulo em II.
\item não-nulo em ambos.
\item nulo em ambos.
\end{alternativas}
\resposta{(b)}


\end{questao}

\blocoaprofundamento
\begin{questao}[3]{}

As cargas $q_1 = \SI{20}{\coulomb}$ e $q_2 = \SI{64}{\coulomb}$ estão fixas no
vácuo nos pontos A e B, separados por $\SI{9}{\meter}$. Determine a posição, sobre
o segmento AB, onde o campo elétrico resultante é nulo.
\dados{$\ke = \SI{9e9}{\newton\meter\squared\per\coulomb\squared}$.}
\resposta{A $\SI{4}{\meter}$ de A (e $\SI{5}{\meter}$ de B)}


\end{questao}
\begin{questao}[3]{}

A intensidade do campo elétrico de uma carga puntiforme positiva, em função da
distância $d$, obedece ao gráfico $E \times d$ abaixo (curva do tipo $E = c/d^2$).
Sabendo que a $\SI{3}{\meter}$ o campo vale $\SI{80}{\newton\per\coulomb}$,
determine o campo a $\SI{6}{\meter}$.

\circuitoqq{
\begin{tikzpicture}
 \begin{axis}[width=9cm,height=5cm,xlabel={$d$ (m)},ylabel={$E$ (N/C)},
   xmin=0,xmax=8,ymin=0,ymax=90,xtick={0,2,3,4,6,8},ytick={0,20,40,60,80},
   grid=both,axis lines=left,domain=1.6:8,samples=100]
   \addplot[Vcol,very thick] {720/(x^2)};
   \addplot[only marks,mark=*,Vcol] coordinates {(3,80)(6,20)};
   \node[Vcol,anchor=south west,font=\footnotesize] at (axis cs:3,80) {$(3,\,80)$};
   \node[Vcol,anchor=south west,font=\footnotesize] at (axis cs:6,20) {$(6,\,20)$};
 \end{axis}
\end{tikzpicture}}{Campo elétrico em função da distância ($E \propto 1/d^2$).}

\resposta{$E = \SI{20}{\newton\per\coulomb}$}


\end{questao}
\begin{questao}[3]{}

Duas cargas $q_1 = \SI{+9}{\micro\coulomb}$ e $q_2 = \SI{+4}{\micro\coulomb}$
estão fixas e separadas por $\SI{10}{\centi\meter}$. Determine a que distância da
carga $q_1$, sobre a reta que as une, o campo elétrico resultante é nulo.
\resposta{A $\SI{6}{\centi\meter}$ de $q_1$}


\end{questao}
\begin{questao}[3]{}

Duas cargas puntiformes, $+4Q$ e $-Q$, estão fixas no vácuo, separadas por uma
distância $d$.
\begin{itensq}
\item Mostre que existe um ponto sobre a reta que as une onde o campo elétrico
      resultante é \textbf{nulo}, e determine sua posição.
\item Explique por que esse ponto está \emph{fora} do segmento que une as cargas,
      e do lado da carga de \emph{menor} módulo.
\item Nesse mesmo ponto, o potencial elétrico é nulo? Justifique.
\end{itensq}
\resposta{(a) a $2d$ da carga $+4Q$ (isto é, a $d$ além de $-Q$);
(b) e (c) ver comentário}


\end{questao}
\begin{questao}[3]{}

Um elétron entra com velocidade horizontal $v_0=\SI{2e7}{\meter\per\second}$ na
região entre duas placas paralelas de comprimento $L=\SI{5}{\centi\meter}$, onde
existe campo elétrico uniforme $E=\SI{e4}{\volt\per\meter}$, perpendicular à
velocidade.

\circuitoq{
\begin{tikzpicture}[scale=1,line width=0.8pt,>=Stealth]
 \draw[cinza!70,line width=2.2pt] (0,1.3)--(4.2,1.3);
 \draw[cinza!70,line width=2.2pt] (0,-1.3)--(4.2,-1.3);
 \node[above,font=\footnotesize] at (2.1,1.35) {$+$\ placa};
 \node[below,font=\footnotesize] at (2.1,-1.35) {$-$\ placa};
 \foreach \x in {0.6,1.5,2.4,3.3}
   \draw[->,Vcol!70] (\x,1.2)--(\x,-1.2);
 \node[Vcol,font=\footnotesize] at (3.85,0.55) {$\vec E$};
 \draw[->,laranja,line width=1.2pt] (-1.4,0)--(0,0)
       node[midway,above,font=\footnotesize]{$v_0$};
 \draw[laranja,line width=1.2pt,domain=0:4.2,samples=40,smooth]
       plot (\x,{0.055*\x*\x});
 \draw[->,laranja,line width=1.2pt] (4.2,{0.055*4.2*4.2})--(5.3,{0.055*4.2*4.2+0.55});
 \draw[<->,cinza] (0,-1.75)--(4.2,-1.75) node[midway,below,font=\footnotesize]{$L$};
 \draw[<->,cinza] (4.5,0)--(4.5,{0.055*4.2*4.2}) node[midway,right,font=\footnotesize]{$y$};
\end{tikzpicture}}{Deflexão de um elétron em campo elétrico uniforme.}

\begin{itensq}
\item Calcule a aceleração do elétron.
\item Determine o tempo de permanência entre as placas e o desvio vertical $y$.
\item Calcule o ângulo de saída em relação à direção inicial.
\end{itensq}
\dados{$e=\SI{1.6e-19}{\coulomb}$; $m_e=\SI{9.11e-31}{\kilogram}$.}
\resposta{(a) $a\approx\SI{1.76e15}{\meter\per\second\squared}$;
(b) $t=\SI{2.5}{\nano\second}$, $y\approx\SI{5.5}{\milli\meter}$;
(c) $\theta\approx\SI{12.4}{\degree}$}


\end{questao}
\begin{questao}[3]{}

Quatro cargas de módulo $q=\SI{1}{\micro\coulomb}$ ocupam os vértices de um
quadrado de lado $a=\SI{10}{\centi\meter}$. Compare o \textbf{campo elétrico no
centro} em três configurações:
\begin{itensq}
\item todas positivas;
\item duas positivas e duas negativas, alternadas ($+,-,+,-$);
\item duas positivas e duas negativas, adjacentes ($+,+,-,-$).
\end{itensq}
\dados{$\ke=\SI{9e9}{}$; $\sqrt2\approx1{,}41$.}
\resposta{(a) $\vec E=0$; (b) $\vec E=0$; (c) $E\approx\pot{5.1}{6}\,\si{\newton\per\coulomb}$}


\end{questao}
\begin{questao}[3]{}

Uma carga $-Q$ está na origem e uma carga $+4Q$ está em $x=d$.
\begin{itensq}
\item Determine o ponto sobre o eixo $x$ onde o campo resultante é nulo.
\item Explique por que ele fica \emph{fora} do segmento e do lado da carga de
      menor módulo.
\item Nesse ponto, quanto vale o potencial elétrico?
\end{itensq}
\resposta{(a) $x=-d$ (a uma distância $d$ à esquerda de $-Q$);
(b) e (c) ver comentário; $V=\dfrac{\ke Q}{d}$}


\end{questao}
\begin{questao}[3]{}

Um anel isolante de raio $R$ está uniformemente carregado com carga total $Q$.
Considere um ponto P sobre o eixo do anel, a uma distância $x$ do centro.
\begin{itensq}
\item Explique, por simetria, por que o campo em P é paralelo ao eixo.
\item Sabendo que $E(x)=\dfrac{\ke\,Q\,x}{\left(x^{2}+R^{2}\right)^{3/2}}$,
      determine $E$ no centro do anel e no limite $x\gg R$.
\item Mostre que o campo é \emph{máximo} em $x = \dfrac{R}{\sqrt2}$.
\end{itensq}
\resposta{(b) $E(0)=0$ e $E\to\ke Q/x^{2}$; (c) demonstração}


\end{questao}

\blocodesafio
\begin{questao}[4]{}

Uma haste isolante fina de comprimento $L$ está uniformemente carregada com carga
total $Q$ (densidade linear $\lambda=Q/L$). Considere um ponto P sobre o
prolongamento da haste, a uma distância $a$ da extremidade mais próxima.
\begin{itensq}
\item Monte a integral que fornece o campo elétrico em P.
\item Mostre que $E=\dfrac{\ke\,Q}{a\,(a+L)}$.
\item Verifique o comportamento no limite $a \gg L$ e interprete.
\end{itensq}
\resposta{(b) demonstração; (c) $E\to \ke Q/a^{2}$ (carga puntiforme)}


\end{questao}
\begin{questao}[4]{}

Um dipolo elétrico é formado por cargas $+q$ e $-q$ separadas por uma pequena
distância $2\ell$. Define-se o \textbf{momento de dipolo} $p=q\,(2\ell)$.
\begin{itensq}
\item Mostre que, sobre a mediatriz do dipolo, a uma distância $r\gg\ell$ do
      centro, o campo vale $E\approx\dfrac{\ke\,p}{r^{3}}$.
\item Por que o campo do dipolo decai mais rápido que o de uma carga isolada?
\item Um dipolo colocado em campo uniforme sofre força resultante? E torque?
\end{itensq}
\resposta{(a) demonstração; (b) e (c) ver comentário}


\end{questao}

\blocofixacao
\begin{questao}[1]{}

Determine o módulo do campo elétrico produzido por $Q=+\SI{3}{\micro\coulomb}$
a $\SI{30}{\centi\meter}$ de distância.
\dados{$\ke=\SI{9e9}{\newton\meter\squared\per\coulomb\squared}$}
\resposta{$E=\SI{300}{\kilo\newton\per\coulomb}$}


\end{questao}
\begin{questao}[1]{}

Determine o campo produzido por $Q=-\SI{5}{\micro\coulomb}$ a
$\SI{50}{\centi\meter}$, indicando o sentido.
\resposta{$E=\SI{180}{\kilo\newton\per\coulomb}$, apontando \emph{para} a carga}


\end{questao}
\begin{questao}[1]{}

Uma carga $q=+\SI{2}{\micro\coulomb}$ é colocada em um ponto onde
$E=\SI{5000}{\newton\per\coulomb}$. Determine a força sobre ela.
\resposta{$F=\SI{10}{\milli\newton}$, no sentido do campo}


\end{questao}
\begin{questao}[1]{}

Determine a carga que produz $E=\SI{2000}{\newton\per\coulomb}$ a
$\SI{30}{\centi\meter}$.
\resposta{$Q=\SI{20}{\nano\coulomb}$}


\end{questao}
\begin{questao}[1]{}

A que distância de $Q=+\SI{8}{\micro\coulomb}$ o campo vale
$\SI{200}{\kilo\newton\per\coulomb}$?
\resposta{$r=\SI{0.6}{\meter}$}


\end{questao}
\begin{questao}[1]{}

O campo elétrico em um ponto, medido com carga de prova $q_0$, vale $E$. Se a
carga de prova for \textbf{dobrada}, o valor medido:
\begin{alternativas}[2]
\item dobra
\item cai à metade
\item não se altera
\item depende do sinal de $q_0$
\end{alternativas}
\resposta{(c)}


\end{questao}
\begin{questao}[1]{}

Mostre que $\si{\newton\per\coulomb}$ e $\si{\volt\per\meter}$ são a mesma
unidade.
\resposta{$\si{\volt}=\si{\joule\per\coulomb}$ e
$\si{\joule}=\si{\newton\meter}$}


\end{questao}
\begin{questao}[1]{}

Uma carga negativa é colocada em uma região de campo elétrico uniforme. Sobre
seu movimento inicial, é correto afirmar que ela:
\begin{alternativas}[2]
\item acelera no sentido do campo
\item acelera no sentido oposto ao campo
\item permanece em repouso
\item move-se perpendicularmente ao campo
\end{alternativas}
\resposta{(b)}


\end{questao}

\blocoaplicacao
\begin{questao}[2]{}

Uma carga $q_0=-\SI{1}{\micro\coulomb}$ é colocada em um campo uniforme de
$E=\SI{2000}{\newton\per\coulomb}$.
\begin{itensq}
\item Determine a força.
\item Determine sua aceleração, se a massa for $\SI{4}{\gram}$.
\end{itensq}
\resposta{$F=\SI{2}{\milli\newton}$, oposta a $\vec E$;
$a=\SI{0.5}{\meter\per\second\squared}$}


\end{questao}
\begin{questao}[2]{}

Um próton é abandonado em um campo uniforme de
$E=\SI{1e4}{\newton\per\coulomb}$. Determine sua aceleração.
\dados{$e=\pot{1.6}{-19}\,\si{\coulomb}$; $m_p=\pot{1.67}{-27}\,\si{\kilogram}$}
\resposta{$a=\pot{9.6}{11}\,\si{\meter\per\second\squared}$}


\end{questao}
\begin{questao}[2]{}

Duas cargas estão sobre o eixo $x$: $+\SI{4}{\micro\coulomb}$ em $x=0$ e
$-\SI{2}{\micro\coulomb}$ em $x=\SI{0.30}{\meter}$. Determine o campo em
$x=\SI{0.60}{\meter}$.
\resposta{$E=\SI{100}{\kilo\newton\per\coulomb}$, no sentido $-x$}


\end{questao}
\begin{questao}[2]{}

Duas cargas iguais $Q=+\SI{2}{\micro\coulomb}$ estão em
$(\pm\SI{0.30}{\meter},\,0)$. Determine o campo no ponto
$(0,\ \SI{0.40}{\meter})$.
\resposta{$E=\SI{115.2}{\kilo\newton\per\coulomb}$, ao longo de $+y$}


\end{questao}
\begin{questao}[2]{}

Duas placas paralelas separadas por $\SI{2}{\milli\meter}$ estão submetidas a
$\SI{600}{\volt}$. Determine o campo entre elas.
\resposta{$E=\SI{300}{\kilo\volt\per\meter}$}


\end{questao}
\begin{questao}[2]{}

Uma placa condutora extensa tem densidade superficial de carga
$\sigma=\SI{2}{\micro\coulomb\per\meter\squared}$. Determine o campo próximo à
sua superfície.
\dados{$\varepsilon_0=\pot{8.85}{-12}\,\si{\farad\per\meter}$}
\resposta{$E=\SI{226}{\kilo\newton\per\coulomb}$}


\end{questao}
\begin{questao}[2]{}

Compare o campo produzido por $Q_1=+\SI{9}{\micro\coulomb}$ a
$\SI{30}{\centi\meter}$ com o de $Q_2=+\SI{4}{\micro\coulomb}$ a
$\SI{20}{\centi\meter}$.
\resposta{$E_1=E_2=\SI{900}{\kilo\newton\per\coulomb}$ --- iguais}


\end{questao}
\begin{questao}[2]{}

Uma carga em um ponto sofre campos de $\SI{3}{\kilo\newton\per\coulomb}$ na
direção $x$ e $\SI{4}{\kilo\newton\per\coulomb}$ na direção $y$, produzidos por
duas fontes distintas. Determine o campo resultante.
\resposta{$E=\SI{5}{\kilo\newton\per\coulomb}$, a $\ang{53.1}$ do eixo $x$}


\end{questao}
\begin{questao}[2]{}

Em uma representação por linhas de campo, o que indica uma região onde as
linhas estão mais próximas entre si?
\begin{alternativas}[2]
\item potencial maior
\item campo mais intenso
\item carga maior naquele ponto
\item ausência de campo
\end{alternativas}
\resposta{(b)}


\end{questao}
\begin{questao}[2]{}

Uma carga produz campo de $\SI{8}{\kilo\newton\per\coulomb}$ no vácuo. O
conjunto é imerso em um dielétrico de $\varepsilon_r=4$. Determine o novo campo.
\resposta{$E=\SI{2}{\kilo\newton\per\coulomb}$}


\end{questao}
\begin{questao}[2]{}

Uma gota de massa $\SI{2}{\milli\gram}$ e carga $\SI{1}{\nano\coulomb}$ deve
ficar suspensa em repouso. Determine o campo elétrico necessário.
\dados{$g=\SI{9.8}{\meter\per\second\squared}$}
\resposta{$E=\SI{19.6}{\kilo\newton\per\coulomb}$}


\end{questao}

\blocoaprofundamento
\begin{questao}[3]{}

Duas cargas iguais $Q=+\SI{1}{\micro\coulomb}$ ocupam
$(\pm\SI{0.20}{\meter},\,0)$.
\begin{itensq}
\item Determine o campo em $(0,\ \SI{0.20}{\meter})$.
\item Determine o campo no ponto médio e justifique.
\end{itensq}
\resposta{(a) $E=\SI{159}{\kilo\newton\per\coulomb}$ ao longo de $+y$;
(b) $\vec E=\vec0$}


\end{questao}
\begin{questao}[3]{}

Um dipolo é formado por $+q$ e $-q$ separadas por $2a$. Considere um ponto do
eixo perpendicular à linha que as une, a distância $y$ do centro.
\begin{itensq}
\item Deduza a expressão exata do campo.
\item Obtenha a forma assintótica para $y\gg a$.
\item Compare a queda com a de uma carga isolada.
\end{itensq}
\resposta{$E=\dfrac{2\ke qa}{(y^{2}+a^{2})^{3/2}}\to\dfrac{2\ke qa}{y^{3}}$}


\end{questao}
\begin{questao}[3]{}

Duas placas condutoras extensas e paralelas têm densidades superficiais
$+\sigma$ e $-\sigma$, com
$\sigma=\SI{1}{\micro\coulomb\per\meter\squared}$.
\begin{itensq}
\item Determine o campo entre as placas.
\item Determine o campo fora delas.
\item Justifique fisicamente.
\end{itensq}
\resposta{Entre: $E=\SI{113}{\kilo\newton\per\coulomb}$; fora: $E=0$}


\end{questao}
\begin{questao}[3]{}

Uma esfera condutora de raio $R=\SI{10}{\centi\meter}$ tem carga
$Q=\SI{5}{\micro\coulomb}$.
\begin{itensq}
\item Determine o campo no interior, na superfície e a $\SI{30}{\centi\meter}$
      do centro.
\item Avalie se a configuração é fisicamente sustentável no ar.
\end{itensq}
\resposta{$0$; $\SI{4.5}{\mega\volt\per\meter}$;
$\SI{500}{\kilo\volt\per\meter}$ --- \textbf{insustentável}}


\end{questao}
\begin{questao}[3]{}

Um elétron parte do repouso em um campo uniforme de
$E=\SI{1000}{\volt\per\meter}$ e percorre $\SI{2}{\centi\meter}$.
\begin{itensq}
\item Determine a aceleração e o tempo de percurso.
\item Determine a velocidade final e a energia adquirida.
\end{itensq}
\dados{$e=\pot{1.6}{-19}\,\si{\coulomb}$; $m_e=\pot{9.11}{-31}\,\si{\kilogram}$}
\resposta{$a=\pot{1.76}{14}\,\si{\meter\per\second\squared}$;
$t=\SI{15.1}{\nano\second}$; $v=\pot{2.65}{6}\,\si{\meter\per\second}$;
$\SI{20}{\electronvolt}$}


\end{questao}
\begin{questao}[3]{}

Compare a queda do campo com a distância para três configurações: carga
puntiforme, fio infinito carregado e plano infinito carregado.
\begin{itensq}
\item Escreva a dependência de cada uma.
\item Explique a origem geométrica das diferenças.
\end{itensq}
\resposta{$1/r^{2}$, $1/r$ e constante}


\end{questao}
\begin{questao}[3]{}

Uma haste isolante fina é carregada de modo que produz, em certo ponto, campo
de $\SI{600}{\newton\per\coulomb}$. A haste é então partida ao meio, e apenas
uma das metades permanece.
\begin{itensq}
\item O campo cai à metade? Justifique.
\item Determine em que condição isso valeria.
\end{itensq}
\resposta{Não necessariamente --- só vale se a metade removida contribuísse com
exatamente metade do campo}


\end{questao}

\blocodesafio
\begin{questao}[4]{}

\textbf{(Campo nulo com potencial não nulo.)} Construa uma configuração em que
$\vec E=\vec0$ em um ponto onde $V\ne0$.
\begin{itensq}
\item Proponha o arranjo mais simples e determine os valores.
\item Explique por que isso é possível.
\item Determine se o inverso ($V=0$ com $\vec E\ne\vec0$) também ocorre.
\end{itensq}
\resposta{Duas cargas iguais $+Q$; no ponto médio $\vec E=\vec0$ e
$V=2\ke Q/a$}


\end{questao}
\begin{questao}[4]{}

\textbf{(Linhas de campo.)} Analise as linhas de campo de um dipolo elétrico.
\begin{itensq}
\item Descreva seu traçado.
\item Demonstre que linhas de campo nunca se cruzam.
\item Explique o que ocorre em pontos de campo nulo.
\end{itensq}
\resposta{Saem de $+q$ e entram em $-q$; não se cruzam porque $\vec E$ tem
direção única em cada ponto}


\end{questao}
\begin{questao}[4]{}

\textbf{(Projeto de acelerador de placas.)} Projete um arranjo de placas
paralelas que imprima aceleração de
$a=\pot{2}{12}\,\si{\meter\per\second\squared}$ a um elétron.
\begin{itensq}
\item Determine o campo necessário.
\item Determine a tensão para separação de $\SI{1}{\centi\meter}$.
\item Determine o tempo de trânsito e a velocidade de saída.
\end{itensq}
\dados{$e=\pot{1.6}{-19}\,\si{\coulomb}$; $m_e=\pot{9.11}{-31}\,\si{\kilogram}$}
\resposta{$E=\SI{11.4}{\volt\per\meter}$; $U=\SI{0.114}{\volt}$;
$t=\SI{100}{\nano\second}$; $v=\pot{2}{5}\,\si{\meter\per\second}$}


\end{questao}
\begin{questao}[4]{}

\textbf{(Cavidade em condutor.)} Uma carga puntiforme $+q$ é colocada no
interior de uma cavidade vazia dentro de um condutor \emph{neutro} e isolado.
\begin{itensq}
\item Determine a carga induzida em cada superfície.
\item Determine o campo dentro do metal e fora do condutor.
\item Analise o que muda se o condutor for aterrado.
\end{itensq}
\resposta{$-q$ na parede da cavidade, $+q$ na superfície externa; campo nulo no
metal e não nulo fora}


\end{questao}
\begin{questao}[4]{}

\textbf{(Esfera isolante uniformemente carregada.)} Uma esfera isolante maciça
de raio $R$ tem carga $Q$ distribuída uniformemente em seu \emph{volume}.
\begin{itensq}
\item Determine o campo para $r<R$ e para $r>R$.
\item Identifique onde o campo é máximo.
\item Compare com a esfera \emph{condutora} de mesma carga.
\end{itensq}
\resposta{$E=\dfrac{\ke Qr}{R^{3}}$ dentro e $\dfrac{\ke Q}{r^{2}}$ fora;
máximo na superfície}


\end{questao}
\begin{questao}[4]{}

\textbf{(Traçado numérico de linhas de campo.)} Descreva um procedimento
computacional para traçar linhas de campo a partir de uma expressão conhecida de
$\vec E(x,y)$.
\begin{itensq}
\item Formule o algoritmo.
\item Indique como escolher os pontos de partida.
\item Aponte as dificuldades numéricas e como contorná-las.
\end{itensq}
\resposta{Integrar $\dfrac{\mathrm{d}\vec r}{\mathrm{d}s}
=\dfrac{\vec E}{|\vec E|}$ a partir de pontos distribuídos em torno das cargas}


\end{questao}
\begin{questao}[4]{}

\textbf{(Limite de campo em capacitor.)} Um capacitor de placas paralelas com
$d=\SI{0.1}{\milli\meter}$ usa dielétrico de rigidez
$\SI{20}{\kilo\volt\per\milli\meter}$ e $\varepsilon_r=3$.
\begin{itensq}
\item Determine a tensão máxima e o campo correspondente.
\item Determine a densidade superficial de carga no limite.
\item Discuta o compromisso entre capacitância e tensão.
\end{itensq}
\dados{$\varepsilon_0=\pot{8.85}{-12}\,\si{\farad\per\meter}$}
\resposta{$U_{\max}=\SI{2}{\kilo\volt}$;
$E=\SI{2e7}{\volt\per\meter}$;
$\sigma=\SI{531}{\micro\coulomb\per\meter\squared}$}


\end{questao}
\begin{questao}[4]{}

\textbf{(Superposição vetorial em anel.)} Um anel de raio $R$ tem carga
$Q$ uniformemente distribuída.
\begin{itensq}
\item Explique por que o campo é nulo no centro.
\item Determine, sem integrar, o campo em um ponto muito distante no eixo.
\item Discuta o que ocorre se metade do anel for removida.
\end{itensq}
\resposta{Centro: $\vec E=\vec0$ por simetria; longe: $E\to\ke Q/x^{2}$;
com meio anel, campo não nulo no centro}


\end{questao}

\gabarito[1.3 Campo elétrico]

\subsection{Potencial elétrico e superfícies equipotenciais}

\blocofixacao
\begin{questao}[1]{}

Sobre superfícies equipotenciais, é correto afirmar que:
\begin{alternativas}
\item o campo elétrico é tangente a elas.
\item o campo elétrico é perpendicular a elas.
\item o potencial varia ao longo de uma mesma superfície.
\item é necessário realizar trabalho para mover uma carga sobre a superfície.
\item elas sempre coincidem com as linhas de campo.
\end{alternativas}
\resposta{(b)}


\end{questao}
\begin{questao}[1]{}

Considere dois pontos A e B de um campo elétrico uniforme de intensidade
$E = \SI{e4}{\newton\per\coulomb}$, separados por $\SI{5}{\centi\meter}$ ao longo
da direção do campo. Calcule a ddp entre A e B.
\resposta{$U_{AB} = \SI{500}{\volt}$}


\end{questao}

\blocoaplicacao
\begin{questao}[2]{}

A figura representa três superfícies equipotenciais de um campo elétrico
uniforme, com $V_A = \SI{60}{\volt}$, $V_B = \SI{40}{\volt}$ e
$V_C = \SI{20}{\volt}$, igualmente espaçadas de $\SI{2}{\centi\meter}$.

\circuitoq{
\begin{tikzpicture}[scale=1,line width=0.8pt,>=Stealth]
 \foreach \x/\v in {0/A,2/B,4/C}{
   \draw[azulprof,thick] (\x,-1.3) -- (\x,1.3);
   \node[above,azulprof] at (\x,1.3) {$V_{\v}$};}
 \node[below] at (0,-1.3) {\footnotesize$\SI{60}{\volt}$};
 \node[below] at (2,-1.3) {\footnotesize$\SI{40}{\volt}$};
 \node[below] at (4,-1.3) {\footnotesize$\SI{20}{\volt}$};
 
 \foreach \y in {-0.7,0,0.7}
   \draw[->,Vcol,thick] (0.3,\y) -- (3.7,\y);
 \node[Vcol] at (2,1.0) {$\vec{E}$};
 \draw[<->,cinza] (0,-1.9) -- (2,-1.9) node[midway,below]{\footnotesize$\SI{2}{\centi\meter}$};
\end{tikzpicture}}{Superfícies equipotenciais em campo uniforme.}

\begin{itensq}
\item Determine a intensidade do campo elétrico.
\item Qual é o trabalho para levar uma carga de $\SI{5}{\micro\coulomb}$ de A até C?
\end{itensq}
\resposta{(a) $E = \SI{1000}{\volt\per\meter}$; (b) $W = \SI{2e-4}{\joule}$}


\end{questao}
\begin{questao}[2]{}

A figura mostra linhas equipotenciais no campo gerado por duas cargas de mesmo
módulo e sinais opostos. São dados $V_A = \SI{30}{\volt}$, $V_B = \SI{10}{\volt}$
e $V_C = \SI{-10}{\volt}$. Calcule a ddp entre A e B, e entre B e C.

\circuitoq{
\begin{tikzpicture}[scale=0.85,line width=0.8pt]
 
 \shade[ball color=Vcol!25] (-3,0) circle (0.3); \node at (-3,0){\footnotesize$+$};
 \shade[ball color=Icol!25] (3,0) circle (0.3);  \node at (3,0){\footnotesize$-$};
 
 \foreach \r/\lab/\x in {1.2/A/-1.6, 2.0/B/0, 1.2/C/1.6}{}
 \draw[azulprof,thick] (-1.6,-1.6) .. controls (-2.2,0) .. (-1.6,1.6);
 \node[azulprof] at (-1.6,1.85) {$A$};
 \draw[azulprof,thick] (0,-1.9) -- (0,1.9);
 \node[azulprof] at (0,2.15) {$B$};
 \draw[azulprof,thick] (1.6,-1.6) .. controls (2.2,0) .. (1.6,1.6);
 \node[azulprof] at (1.6,1.85) {$C$};
 \node[below,font=\footnotesize] at (-1.6,-1.7) {$\SI{30}{\volt}$};
 \node[below,font=\footnotesize] at (0,-2.0) {$\SI{10}{\volt}$};
 \node[below,font=\footnotesize] at (1.6,-1.7) {$\SI{-10}{\volt}$};
\end{tikzpicture}}{Linhas equipotenciais de um dipolo elétrico.}

\resposta{$U_{AB} = \SI{20}{\volt}$; $U_{BC} = \SI{20}{\volt}$}


\end{questao}
\begin{questao}[2]{}

As linhas cheias da figura representam linhas de força de um campo eletrostático
e as tracejadas, linhas equipotenciais, com $V_A > V_B > V_C$. Uma partícula de
carga $q = \SI{2e-6}{\coulomb}$ é levada de A até C.

\circuitoq{
\begin{tikzpicture}[scale=1,line width=0.8pt,>=Stealth]
 
 \foreach \x/\lab/\v in {0/A/{V_A},2.2/B/{V_B},4.4/C/{V_C}}{
   \draw[azulprof,dashed,thick] (\x,-1.4) -- (\x,1.4);
   \node[azulprof,above] at (\x,1.4) {$\v$};
   \fill[laranja] (\x,0) circle (0.07);
   \node[below,font=\footnotesize] at (\x,-0.15) {\lab};}
 
 \foreach \y in {-1,0,1}
   \draw[Vcol,->,thick] (-0.6,\y) -- (5.0,\y);
 \node[Vcol] at (5.2,0.7) {$\vec{E}$};
\end{tikzpicture}}{Linhas de força (cheias) e equipotenciais (tracejadas).}

Sobre o trabalho da força elétrica nesse deslocamento, é correto afirmar que:
\begin{alternativas}
\item é nulo, pois A e C estão na mesma linha de força.
\item é positivo, pois a carga se move no sentido do campo.
\item é negativo, pois a carga se move contra o campo.
\item depende do caminho seguido entre A e C.
\item é máximo se a carga for levada de A até B.
\end{alternativas}
\resposta{(b)}


\end{questao}
\begin{questao}[2]{}

A figura representa as linhas equipotenciais e as linhas de força do campo gerado
por uma carga puntiforme $Q$ fixa no vácuo. As equipotenciais são círculos
concêntricos.

\circuitoq{
\begin{tikzpicture}[scale=0.9,line width=0.8pt,>=Stealth]
 \shade[ball color=Vcol!30] (0,0) circle (0.28);
 \node at (0,0) {\footnotesize$Q$};
 \foreach \r in {1,1.7,2.4}
   \draw[azulprof,dashed] (0,0) circle (\r);
 \foreach \a in {0,45,...,315}
   \draw[Vcol,->,thick] (\a:0.35) -- (\a:2.7);
 \fill[laranja] (0:1) circle (0.06) node[above right,font=\footnotesize]{A};
 \fill[laranja] (0:2.4) circle (0.06) node[above right,font=\footnotesize]{C};
\end{tikzpicture}}{Campo e equipotenciais de uma carga puntiforme positiva.}

Como as linhas de força \emph{apontam para fora} de $Q$, pode-se concluir que:
\begin{alternativas}
\item $Q$ é negativa e o potencial aumenta com a distância.
\item $Q$ é positiva e o potencial diminui com a distância.
\item $Q$ é positiva e o potencial aumenta com a distância.
\item o potencial é o mesmo em A e C.
\item o campo é uniforme.
\end{alternativas}
\resposta{(b)}


\end{questao}
\begin{questao}[2]{}

A figura mostra a configuração das linhas de força do campo gerado por duas
partículas carregadas A e B. As linhas \emph{saem} de A e \emph{entram} em B.

\circuitoq{
\begin{tikzpicture}[scale=0.85,line width=0.7pt,>=Stealth]
 \shade[ball color=Vcol!25] (-2,0) circle (0.3); \node at (-2,0){\footnotesize A};
 \shade[ball color=Icol!25] ( 2,0) circle (0.3); \node at (2,0){\footnotesize B};
 \foreach \a in {30,60,90,120,150,-30,-60,-90,-120,-150}
   \draw[Vcol,->] (-2,0) .. controls (\a:1.3) and ({180-\a}:1.3) .. (2,0);
 \draw[Vcol,->] (-2,0)--(-3.3,0);
 \draw[Vcol,->] (3.3,0)--(2,0);
\end{tikzpicture}}{Linhas de força de duas cargas A e B.}

Os sinais das cargas A e B são, respectivamente:
\begin{alternativas}[2]
\item positivo e positivo.
\item positivo e negativo.
\item negativo e positivo.
\item negativo e negativo.
\item ambos nulos.
\end{alternativas}
\resposta{(b)}


\end{questao}
\begin{questao}[2]{}

Uma carga elétrica cria, em um ponto P situado a $\SI{20}{\centi\meter}$ dela, um
campo de intensidade $\SI{900}{\newton\per\coulomb}$. O módulo do potencial
elétrico em P é:
\begin{alternativas}[3]
\item \SI{45}{\volt}
\item \SI{90}{\volt}
\item \SI{180}{\volt}
\item \SI{450}{\volt}
\item \SI{900}{\volt}
\end{alternativas}
\resposta{(c)}


\end{questao}
\begin{questao}[2]{}

Próximo ao solo, em um dia limpo, a atmosfera apresenta um campo elétrico
aproximadamente uniforme, vertical, apontando para baixo, de intensidade
$\SI{120}{\volt\per\meter}$. Considere dois pontos M (mais alto) e N (mais baixo),
separados verticalmente por $\SI{10}{\meter}$.

\circuitoq{
\begin{tikzpicture}[scale=0.9,line width=0.8pt,>=Stealth]
 \fill[cinza!25] (-2.5,-1.6) rectangle (2.5,-1.2);
 \node[cinza,font=\footnotesize] at (0,-1.4) {solo};
 \foreach \x in {-2,-1,0,1,2}
   \draw[Vcol,->,thick] (\x,1.4) -- (\x,-1.1);
 \node[Vcol] at (2.5,0.4) {$\vec{E}$};
 \fill[laranja] (-1,1.1) circle (0.08) node[left,font=\footnotesize]{M};
 \fill[laranja] (-1,-0.9) circle (0.08) node[left,font=\footnotesize]{N};
 \draw[<->,cinza] (0.6,1.1)--(0.6,-0.9) node[midway,right,font=\footnotesize]{$\SI{10}{\meter}$};
\end{tikzpicture}}{Campo elétrico atmosférico próximo ao solo.}

A diferença de potencial entre M e N vale:
\begin{alternativas}[2]
\item \SI{12}{\volt}
\item \SI{120}{\volt}
\item \SI{1200}{\volt}
\item \SI{12000}{\volt}
\item zero
\end{alternativas}
\resposta{(c)}


\end{questao}

\blocoaprofundamento
\begin{questao}[3]{}

Duas cargas $Q_1 = \SI{+2}{\micro\coulomb}$ e $Q_2 = \SI{-2}{\micro\coulomb}$
estão fixas, separadas por $\SI{20}{\centi\meter}$. Determine o potencial elétrico
no ponto médio M do segmento que as une.
\dados{$\ke = \SI{9e9}{\newton\meter\squared\per\coulomb\squared}$.}
\resposta{$V_M = 0$}


\end{questao}
\begin{questao}[3]{}

Uma carga puntiforme $Q = \SI{+5}{\micro\coulomb}$ está fixa no vácuo. Uma segunda
carga $q = \SI{+2}{\micro\coulomb}$ é trazida do infinito até um ponto a
$\SI{30}{\centi\meter}$ de $Q$.
\begin{itensq}
\item Qual é o potencial criado por $Q$ nesse ponto?
\item Qual é o trabalho realizado pela força externa para trazer $q$ (supondo-a
      trazida lentamente)?
\end{itensq}
\dados{$\ke = \SI{9e9}{\newton\meter\squared\per\coulomb\squared}$;
$V_\infty = 0$.}
\resposta{(a) $V = \SI{1.5e5}{\volt}$; (b) $W = \SI{0.3}{\joule}$}


\end{questao}
\begin{questao}[3]{}

Um elétron ($q = \SI{-1.6e-19}{\coulomb}$, $m = \SI{9.1e-31}{\kilogram}$) parte do
repouso e é acelerado por uma ddp de $\SI{200}{\volt}$. Determine a energia
cinética adquirida e a velocidade final do elétron.
\resposta{$E_c = \SI{3.2e-17}{\joule}$; $v \approx \SI{8.4e6}{\meter\per\second}$}


\end{questao}
\begin{questao}[3]{}

O diagrama representa o potencial elétrico $V$ criado por uma carga puntiforme em
função da distância $d$ até ela.

\circuitoqq{
\begin{tikzpicture}
 \begin{axis}[width=9cm,height=5cm,xlabel={$d$ (m)},ylabel={$V$ (V)},
   xmin=0,xmax=6,ymin=0,ymax=90,xtick={0,1,2,3,4},ytick={0,20,40,80},
   grid=both,axis lines=left,domain=0.9:6,samples=100]
   \addplot[Vcol,very thick] {80/x};
   \addplot[only marks,mark=*,Vcol] coordinates {(1,80)(2,40)(4,20)};
   \node[Vcol,anchor=south west,font=\footnotesize] at (axis cs:2,40){$(2,\,40)$};
 \end{axis}
\end{tikzpicture}}{Potencial de uma carga puntiforme em função da distância.}

Sabendo que, a $\SI{2}{\meter}$, o potencial vale $\SI{40}{\volt}$, o potencial a
$\SI{4}{\meter}$ e o campo elétrico a $\SI{2}{\meter}$ valem, respectivamente:
\begin{alternativas}[2]
\item \SI{20}{\volt} e \SI{20}{\volt\per\meter}
\item \SI{20}{\volt} e \SI{40}{\volt\per\meter}
\item \SI{10}{\volt} e \SI{10}{\volt\per\meter}
\item \SI{40}{\volt} e \SI{40}{\volt\per\meter}
\item \SI{10}{\volt} e \SI{20}{\volt\per\meter}
\end{alternativas}
\resposta{(a)}


\end{questao}
\begin{questao}[3]{}

Na figura, uma partícula de carga $Q > 0$ está fixa em P. As superfícies A, B e C
são equipotenciais, e um ponto D situa-se sobre uma linha de força entre B e C.

\circuitoq{
\begin{tikzpicture}[scale=0.9,line width=0.8pt,>=Stealth]
 \shade[ball color=Vcol!30] (0,0) circle (0.25); \node at (0,0){\footnotesize P};
 \foreach \r/\lab in {1.1/A,1.8/B,2.5/C}{
   \draw[azulprof,dashed] (0,0) circle (\r);
   \node[azulprof,font=\footnotesize] at (\r,0.28) {\lab};}
 \draw[Vcol,->,thick] (20:0.3)--(20:2.7);
 \fill[laranja] (20:2.15) circle (0.07) node[right,font=\footnotesize]{D};
\end{tikzpicture}}{Equipotenciais de uma carga puntiforme positiva.}

Sobre os potenciais, é correto afirmar que:
\begin{alternativas}
\item $V_A < V_B < V_C$.
\item $V_A > V_B > V_C$.
\item $V_A = V_B = V_C$.
\item $V_C$ é o maior, por estar mais longe.
\item nada se pode afirmar sem o valor de $Q$.
\end{alternativas}
\resposta{(b)}


\end{questao}
\begin{questao}[3]{}

Uma esfera condutora de raio $R = \SI{10}{\centi\meter}$ está carregada com
$Q = \SI{6}{\micro\coulomb}$, isolada no vácuo. Os gráficos mostram o campo
elétrico e o potencial em função da distância $d$ ao centro.

\circuitoqq{
\begin{tikzpicture}
 \begin{axis}[width=6.2cm,height=4.2cm,xlabel={$d$ (m)},ylabel={$E$ (kN/C)},
   xmin=0,xmax=0.45,ymin=0,ymax=6200,xtick={0,0.1,0.2,0.3,0.4},
   ytick={0,2000,4000,5400},yticklabels={0,2000,4000,5400},
   grid=both,axis lines=left,domain=0.1:0.45,samples=80,
   scaled y ticks=false,y tick label style={/pgf/number format/fixed}]
   \addplot[Vcol,very thick] {9e9*6e-6/(x^2)/1000};
   \addplot[Vcol,very thick] coordinates {(0,0)(0.1,0)};
   \addplot[Vcol,dashed] coordinates {(0.1,0)(0.1,5400)};
 \end{axis}
\end{tikzpicture}\hspace{4mm}
\begin{tikzpicture}
 \begin{axis}[width=6.2cm,height=4.2cm,xlabel={$d$ (m)},ylabel={$V$ (kV)},
   xmin=0,xmax=0.45,ymin=0,ymax=620,xtick={0,0.1,0.2,0.3,0.4},
   ytick={0,180,540},grid=both,axis lines=left,domain=0.1:0.45,samples=80]
   \addplot[Icol,very thick] {9e9*6e-6/x/1000};
   \addplot[Icol,very thick] coordinates {(0,540)(0.1,540)};
 \end{axis}
\end{tikzpicture}}{Campo e potencial de uma esfera condutora carregada.}

\begin{itensq}
\item Calcule $E$ e $V$ na superfície da esfera.
\item Determine $E$ e $V$ a $\SI{30}{\centi\meter}$ do centro.
\item Explique por que, \emph{dentro} da esfera, $E = 0$ mas $V \neq 0$.
\end{itensq}
\dados{$\ke = \SI{9e9}{\newton\meter\squared\per\coulomb\squared}$.}
\resposta{(a) $E=\SI{5.4e6}{\newton\per\coulomb}$, $V=\SI{540}{\kilo\volt}$;
(b) $E=\SI{6e5}{\newton\per\coulomb}$, $V=\SI{180}{\kilo\volt}$; (c) ver comentário}


\end{questao}
\begin{questao}[3]{}

Três cargas puntiformes iguais, $q = \SI{1}{\micro\coulomb}$, são trazidas
\emph{uma a uma} desde o infinito até ocuparem os vértices de um triângulo
equilátero de lado $L = \SI{10}{\centi\meter}$.
\begin{itensq}
\item Calcule o trabalho externo necessário para trazer cada carga.
\item Determine a energia potencial total da configuração.
\item Se as cargas forem liberadas simultaneamente, qual é a energia cinética total
      quando estiverem infinitamente afastadas?
\end{itensq}
\dados{$\ke = \SI{9e9}{\newton\meter\squared\per\coulomb\squared}$.}
\resposta{(a) $0$, \SI{0.09}{\joule}, \SI{0.18}{\joule};
(b) $U = \SI{0.27}{\joule}$; (c) \SI{0.27}{\joule}}


\end{questao}
\begin{questao}[3]{}

Duas esferas condutoras, de raios $R$ e $3R$, estão distantes entre si e são
ligadas por um fio condutor longo e fino. A carga total do sistema é $Q$.
\begin{itensq}
\item Determine a carga final de cada esfera.
\item Compare os campos elétricos nas superfícies das duas esferas.
\item Explique o "poder das pontas" a partir desse resultado.
\end{itensq}
\resposta{(a) $Q_1=Q/4$, $Q_2=3Q/4$; (b) $E_1/E_2=3$; (c) ver comentário}


\end{questao}

\blocodesafio
\begin{questao}[4]{}

Quatro cargas iguais $q=\SI{1}{\micro\coulomb}$ ocupam os vértices de um quadrado
de lado $a=\SI{10}{\centi\meter}$.
\begin{itensq}
\item Calcule a energia potencial elétrica total da configuração.
\item Quanto trabalho externo seria necessário para trazer uma quinta carga $q$
      desde o infinito até o centro do quadrado?
\item Se todas fossem liberadas, qual seria a energia cinética total no infinito?
\end{itensq}
\dados{$\ke=\SI{9e9}{}$; $\sqrt2\approx1{,}41$.}
\resposta{(a) $U\approx\SI{0.49}{\joule}$; (b) $W\approx\SI{0.51}{\joule}$;
(c) $\approx\SI{0.49}{\joule}$}


\end{questao}
\begin{questao}[4]{}

Uma esfera condutora oca de raio interno $a$ e raio externo $b$ envolve, no
centro, uma carga puntiforme $+Q$. A casca está inicialmente \emph{neutra} e
isolada.
\begin{itensq}
\item Determine as cargas induzidas nas superfícies interna e externa da casca.
\item Esboce o comportamento de $E(r)$ nas quatro regiões ($r<a$, $a<r<b$,
      $r>b$).
\item Se a casca for \textbf{aterrada}, o que muda?
\end{itensq}
\resposta{(a) $-Q$ na interna, $+Q$ na externa; (b) e (c) ver comentário}


\end{questao}

\blocofixacao
\begin{questao}[1]{}

Determine o potencial a $\SI{30}{\centi\meter}$ de uma carga
$Q=+\SI{2}{\micro\coulomb}$, adotando $V(\infty)=0$.
\dados{$\ke=\SI{9e9}{\newton\meter\squared\per\coulomb\squared}$}
\resposta{$V=\SI{60}{\kilo\volt}$}


\end{questao}
\begin{questao}[1]{}

Determine o potencial a $\SI{60}{\centi\meter}$ de uma carga
$Q=-\SI{4}{\micro\coulomb}$.
\resposta{$V=\SI{-60}{\kilo\volt}$}


\end{questao}
\begin{questao}[1]{}

A que distância de $Q=+\SI{3}{\micro\coulomb}$ o potencial vale
$\SI{45}{\kilo\volt}$?
\resposta{$r=\SI{0.6}{\meter}$}


\end{questao}
\begin{questao}[1]{}

Determine a carga que produz $\SI{-18}{\kilo\volt}$ a
$\SI{0.5}{\meter}$ de distância.
\resposta{$Q=-\SI{1}{\micro\coulomb}$}


\end{questao}
\begin{questao}[1]{}

Sobre o potencial criado por várias cargas em um ponto, é correto afirmar:
\begin{alternativas}[2]
\item soma-se vetorialmente
\item soma-se algebricamente, com sinal
\item soma-se em módulo
\item depende do caminho até o ponto
\end{alternativas}
\resposta{(b)}


\end{questao}
\begin{questao}[1]{}

Uma carga $q=+\SI{3}{\micro\coulomb}$ é colocada em um ponto onde
$V=\SI{200}{\volt}$. Determine sua energia potencial elétrica.
\resposta{$U=\SI{0.6}{\milli\joule}$}


\end{questao}

\blocoaplicacao
\begin{questao}[2]{}

Uma carga $q=+\SI{2}{\micro\coulomb}$ desloca-se de um ponto a
$\SI{50}{\volt}$ para outro a $\SI{20}{\volt}$.
\begin{itensq}
\item Determine o trabalho realizado pela força elétrica.
\item Determine a variação da energia potencial.
\end{itensq}
\resposta{$W=\SI{60}{\micro\joule}$; $\Delta U=\SI{-60}{\micro\joule}$}


\end{questao}
\begin{questao}[2]{}

Duas cargas estão fixas: $q_1=+\SI{3}{\micro\coulomb}$ a
$\SI{30}{\centi\meter}$ do ponto $P$, e $q_2=-\SI{2}{\micro\coulomb}$ a
$\SI{20}{\centi\meter}$ de $P$. Determine o potencial em $P$.
\resposta{$V=0$}


\end{questao}
\begin{questao}[2]{}

Em um campo uniforme de $E=\SI{500}{\volt\per\meter}$, dois pontos estão
separados por $\SI{4}{\centi\meter}$ ao longo da direção do campo. Determine a
diferença de potencial.
\resposta{$U=\SI{20}{\volt}$}


\end{questao}
\begin{questao}[2]{}

Um elétron é acelerado através de uma diferença de potencial de
$\SI{500}{\volt}$.
\begin{itensq}
\item Determine a energia adquirida, em joules e em elétron-volts.
\item Determine sua velocidade final, partindo do repouso.
\end{itensq}
\dados{$e=\pot{1.6}{-19}\,\si{\coulomb}$;
$m_e=\pot{9.11}{-31}\,\si{\kilogram}$}
\resposta{$\SI{500}{\electronvolt}=\pot{8}{-17}\,\si{\joule}$;
$v=\pot{1.33}{7}\,\si{\meter\per\second}$}


\end{questao}
\begin{questao}[2]{}

Duas cargas, $q_1=+\SI{4}{\micro\coulomb}$ em $x=0$ e
$q_2=-\SI{1}{\micro\coulomb}$ em $x=\SI{0.5}{\meter}$, estão fixas. Determine o
ponto do segmento entre elas onde o potencial é nulo.
\resposta{$x=\SI{0.4}{\meter}$}


\end{questao}

\blocoaprofundamento
\begin{questao}[3]{}

Duas cargas $+Q$ e $-Q$, com $Q=\SI{2}{\micro\coulomb}$, estão separadas por
$\SI{40}{\centi\meter}$.
\begin{itensq}
\item Determine o potencial no ponto médio.
\item Determine o campo no ponto médio.
\item Explique a aparente contradição.
\end{itensq}
\resposta{$V=0$; $E=\SI{900}{\kilo\volt\per\meter}$}


\end{questao}
\begin{questao}[3]{}

O potencial ao longo de um eixo é dado por
$V(x)=\SI{200}{}-\SI{50}{}\,x$ (com $V$ em volts e $x$ em metros), no trecho
$0\le x\le\SI{4}{\meter}$.
\begin{itensq}
\item Determine o campo elétrico.
\item Determine o trabalho para levar $q=+\SI{2}{\micro\coulomb}$ de
      $x=1$ a $x=\SI{3}{\meter}$.
\item Identifique as superfícies equipotenciais.
\end{itensq}
\resposta{$E=\SI{50}{\volt\per\meter}$ no sentido $+x$;
$W=\SI{200}{\micro\joule}$}


\end{questao}
\begin{questao}[3]{}

Uma carga $q=+\SI{1}{\micro\coulomb}$ é trazida do infinito até
$\SI{20}{\centi\meter}$ de uma carga fixa $Q=+\SI{5}{\micro\coulomb}$.
\begin{itensq}
\item Determine o trabalho realizado por um agente externo.
\item Determine a energia potencial do par.
\item Interprete o sinal.
\end{itensq}
\resposta{$W_{ext}=U=\SI{0.225}{\joule}$}


\end{questao}
\begin{questao}[3]{}

Compare, para duas cargas positivas iguais, o lugar geométrico dos pontos de
potencial nulo e o dos pontos de campo nulo.
\begin{itensq}
\item Determine onde $V=0$.
\item Determine onde $\vec E=\vec 0$.
\item Explique a diferença.
\end{itensq}
\resposta{$V$ nunca se anula (exceto no infinito); $\vec E$ se anula no ponto
médio}


\end{questao}

\blocodesafio
\begin{questao}[4]{}

\textbf{(Demonstração --- trabalho e independência do caminho.)} Prove que o
trabalho da força elétrica entre dois pontos independe do caminho, e deduza a
relação com a diferença de potencial.
\begin{itensq}
\item Calcule o trabalho para o campo de uma carga puntiforme.
\item Generalize para uma distribuição qualquer.
\item Explique por que isso permite \emph{definir} o potencial.
\end{itensq}
\resposta{$W_{AB}=q(V_A-V_B)$, com
$V=\ke Q/r$}


\end{questao}
\begin{questao}[4]{}

\textbf{(Projeto --- $V=0$ com $E\neq0$.)} Projete um par de cargas de sinais
opostos e módulos \emph{diferentes} tal que exista um ponto do segmento onde
$V=0$ mas $\vec E\neq\vec0$.
\begin{itensq}
\item Determine a condição geral.
\item Escolha valores e localize o ponto.
\item Verifique que o campo não é nulo ali.
\end{itensq}
\resposta{Condição: $|q_1|/r_1=|q_2|/r_2$; com $+4\,\si{\micro\coulomb}$ e
$-1\,\si{\micro\coulomb}$ a $\SI{0.5}{\meter}$, o ponto é
$x=\SI{0.4}{\meter}$}


\end{questao}
\begin{questao}[4]{}

\textbf{(Estabilidade em um mínimo de potencial.)} Uma curva
$V(x)$ apresenta mínimo local em $x_0$.
\begin{itensq}
\item Analise a estabilidade de uma carga positiva colocada em $x_0$.
\item Repita para uma carga negativa.
\item Relacione com o teorema de Earnshaw.
\end{itensq}
\resposta{Positiva: \textbf{instável}; negativa: estável em $x$ --- mas nenhuma
é estável em três dimensões}


\end{questao}
\begin{questao}[4]{}

\textbf{(Condutor em equilíbrio.)} Demonstre que toda a superfície de um
condutor em equilíbrio eletrostático é equipotencial.
\begin{itensq}
\item Prove a partir da condição de equilíbrio.
\item Mostre que o interior também está no mesmo potencial.
\item Deduza a orientação do campo na superfície.
\end{itensq}
\resposta{Campo interno nulo $\Rightarrow$ $V$ constante em todo o condutor; o
campo externo é \textbf{perpendicular} à superfície}


\end{questao}
\begin{questao}[4]{}

\textbf{(Do potencial ao campo.)} O potencial em uma região é
$V(x,y)=\SI{100}{}-\SI{20}{}\,x+\SI{5}{}\,y^{2}$ (volts, com $x,y$ em metros).
\begin{itensq}
\item Determine $\vec E(x,y)$.
\item Avalie em $(x,y)=(2;\,1)$.
\item Determine a forma das equipotenciais.
\end{itensq}
\resposta{$\vec E=(20;\,-10y)\,\si{\volt\per\meter}$; em $(2;1)$,
$\vec E=(20;\,-10)$ com $|E|=\SI{22.4}{\volt\per\meter}$}


\end{questao}
\begin{questao}[4]{}

\textbf{(Equipotenciais e linhas de campo.)} Compare os dois conjuntos de
curvas para duas cargas positivas iguais.
\begin{itensq}
\item Descreva as equipotenciais perto de cada carga e muito longe do par.
\item Descreva as linhas de campo, incluindo o ponto médio.
\item Estabeleça as regras gerais que relacionam os dois conjuntos.
\end{itensq}
\resposta{Perto: esferas em torno de cada carga; longe: esferas em torno do par;
o ponto médio é ponto de sela do potencial}


\end{questao}

\gabarito[1.4 Potencial elétrico e superfícies equipotenciais]

\section{Conceitos Iniciais: Leis de Ohm, Potência, Fontes e Kirchhoff}

\subsection{Corrente elétrica e carga}

\blocofixacao
\begin{questao}[1]{}

Se uma corrente de $\SI{3}{\ampere}$ percorre um fio durante $\SI{2}{\minute}$, a
carga que atravessa sua seção reta, em coulombs, vale:
\begin{alternativas}[3]
\item 60
\item 110
\item 360
\item 220
\item 180
\end{alternativas}
\resposta{(c)}


\end{questao}
\begin{questao}[1]{}

Uma carga de $\SI{120}{\coulomb}$ passa uniformemente pela seção de um fio durante
$\SI{1}{\minute}$. A intensidade da corrente é:
\begin{alternativas}[3]
\item $\tfrac{1}{30}\,\si{\ampere}$
\item $\tfrac{1}{2}\,\si{\ampere}$
\item \SI{2}{\ampere}
\item \SI{30}{\ampere}
\item \SI{120}{\ampere}
\end{alternativas}
\resposta{(c)}


\end{questao}
\begin{questao}[1]{}

Uma lâmpada permanece acesa durante $\SI{5}{\minute}$ sob corrente de
$\SI{2}{\ampere}$. A carga total que circula, em coulombs, é:
\begin{alternativas}[3]
\item 0,4
\item 2,5
\item 10
\item 150
\item 600
\end{alternativas}
\resposta{(e)}


\end{questao}

\blocoaplicacao
\begin{questao}[2]{}

Um condutor metálico é percorrido por uma corrente contínua e constante de
$\SI{32}{\milli\ampere}$. Determine:
\begin{itensq}
\item a carga que atravessa uma seção reta por segundo;
\item o número de elétrons correspondente.
\end{itensq}
\dados{$e = \SI{1.6e-19}{\coulomb}$.}
\resposta{(a) \SI{32}{\milli\coulomb}; (b) $n = \pot{2.0}{17}$ elétrons}


\end{questao}
\begin{questao}[2]{}

O gráfico mostra a intensidade da corrente $i$ em um condutor em função do tempo.
Determine a carga que atravessa uma seção do condutor nos cinco primeiros
segundos.

\circuitoqq{
\begin{tikzpicture}
 \begin{axis}[width=9cm,height=5cm,xlabel={$t$ (s)},ylabel={$i$ (A)},
   xmin=0,xmax=5.5,ymin=0,ymax=5,xtick={0,1,2,3,4,5},ytick={0,1,2,3,4},
   grid=both,axis lines=left]
   \addplot[Icol,very thick] coordinates {(0,4)(5,4)};
   \addplot[Icol!30] coordinates {(0,0)(0,4)};
 \end{axis}
\end{tikzpicture}}{Corrente constante em função do tempo.}

\resposta{$Q = \SI{20}{\coulomb}$}


\end{questao}
\begin{questao}[2]{}

O gráfico representa a corrente elétrica $i$ que atravessa um condutor em função
do tempo. Determine a carga total que passa por uma seção reta entre $t=0$ e
$t=\SI{6}{\second}$.

\circuitoqq{
\begin{tikzpicture}
 \begin{axis}[width=9.5cm,height=5cm,xlabel={$t$ (s)},ylabel={$i$ (A)},
   xmin=0,xmax=6.5,ymin=0,ymax=5,xtick={0,1,2,3,4,5,6},ytick={0,2,4},
   grid=both,axis lines=left]
   \addplot[Icol,very thick] coordinates {(0,4)(2,4)(6,0)};
   \addplot[Icol!12] coordinates {(0,0)(0,4)};
 \end{axis}
\end{tikzpicture}}{Corrente variável em função do tempo.}

\resposta{$Q = \SI{16}{\coulomb}$}


\end{questao}
\begin{questao}[2]{}

Qual é a carga elétrica que atravessa, durante $\SI{10}{\hour}$, a seção de um
condutor percorrido por corrente constante de $\SI{5}{\ampere}$?
\begin{alternativas}[3]
\item \SI{50}{\coulomb}
\item \SI{1000}{\coulomb}
\item \SI{90000}{\coulomb}
\item \SI{180000}{\coulomb}
\item \SI{360000}{\coulomb}
\end{alternativas}
\resposta{(d)}


\end{questao}

\blocoaprofundamento
\begin{questao}[3]{}

O gráfico mostra a corrente em um condutor variando linearmente de
$\SI{6}{\ampere}$ a $\SI{2}{\ampere}$ entre $t=0$ e $t=\SI{8}{\second}$.

\circuitoqq{
\begin{tikzpicture}
 \begin{axis}[width=9cm,height=4.6cm,xlabel={$t$ (s)},ylabel={$i$ (A)},
   xmin=0,xmax=9,ymin=0,ymax=7,xtick={0,2,4,6,8},ytick={0,2,4,6},
   grid=both,axis lines=left]
   \addplot[Icol,very thick] coordinates {(0,6)(8,2)};
 \end{axis}
\end{tikzpicture}}{Corrente decrescente em função do tempo.}

\begin{itensq}
\item Determine a carga total que atravessa o condutor nesse intervalo.
\item Calcule a corrente \emph{média} e verifique a coerência com o item (a).
\item Em que instante a corrente instantânea iguala a média?
\end{itensq}
\resposta{(a) $Q = \SI{32}{\coulomb}$; (b) $i_m = \SI{4}{\ampere}$;
(c) $t = \SI{4}{\second}$}


\end{questao}
\begin{questao}[3]{}

Um fio de cobre de seção $A = \SI{1}{\milli\meter\squared}$ conduz corrente
constante de $\SI{10}{\ampere}$. O cobre possui aproximadamente
$n = \SI{8.5e28}{}$ elétrons livres por metro cúbico.
\begin{itensq}
\item Determine a \emph{velocidade de deriva} dos elétrons.
\item Quanto tempo um elétron leva para percorrer $\SI{1}{\meter}$ de fio?
\item Como conciliar esse resultado com o fato de uma lâmpada acender
      \emph{instantaneamente} ao acionarmos o interruptor?
\end{itensq}
\dados{$e = \SI{1.6e-19}{\coulomb}$.}
\resposta{(a) $v_d \approx \SI{0.74}{\milli\meter\per\second}$;
(b) $\approx\SI{1360}{\second}$ ($\approx\SI{23}{\minute}$); (c) ver comentário}


\end{questao}

\blocofixacao
\begin{questao}[1]{}

Uma corrente constante de $\SI{1.5}{\ampere}$ circula durante
$\SI{40}{\second}$. Determine a carga transportada.
\resposta{$Q=\SI{60}{\coulomb}$}


\end{questao}
\begin{questao}[1]{}

Uma carga de $\SI{90}{\coulomb}$ atravessa a seção de um condutor em
$\SI{45}{\second}$. Determine a corrente média.
\resposta{$I=\SI{2}{\ampere}$}


\end{questao}
\begin{questao}[1]{}

Determine o tempo necessário para transportar $\SI{60}{\coulomb}$ com corrente
de $\SI{250}{\milli\ampere}$.
\resposta{$t=\SI{240}{\second}=\SI{4}{\minute}$}


\end{questao}
\begin{questao}[1]{}

Determine o número de elétrons que atravessa uma seção quando circula
$\SI{16}{\coulomb}$.
\dados{$e=\pot{1.6}{-19}\,\si{\coulomb}$}
\resposta{$n=\pot{1}{20}$ elétrons}


\end{questao}
\begin{questao}[1]{}

Sobre o sentido convencional da corrente, é correto afirmar:
\begin{alternativas}[2]
\item coincide com o movimento dos elétrons em um metal
\item é oposto ao movimento dos elétrons em um metal
\item existe apenas em semicondutores
\item depende da intensidade da corrente
\end{alternativas}
\resposta{(b)}


\end{questao}

\blocoaplicacao
\begin{questao}[2]{}

Por um condutor passam $\SI{240}{\coulomb}$ em $\SI{120}{\second}$.
\begin{itensq}
\item Determine a corrente média.
\item Determine o número de elétrons.
\end{itensq}
\resposta{$I=\SI{2}{\ampere}$; $n=\pot{1.5}{21}$}


\end{questao}
\begin{questao}[2]{}

Uma bateria de celular tem capacidade de $\SI{2000}{\milli\ampere\hour}$.
\begin{itensq}
\item Converta para coulombs.
\item Determine por quanto tempo ela alimenta uma carga de
      $\SI{250}{\milli\ampere}$.
\end{itensq}
\resposta{$Q=\SI{7200}{\coulomb}$; $t=\SI{8}{\hour}$}


\end{questao}
\begin{questao}[2]{}

Uma corrente pulsada vale $\SI{5}{\ampere}$ durante
$\SI{20}{\milli\second}$ e zero durante os $\SI{80}{\milli\second}$
seguintes, repetindo-se.
\begin{itensq}
\item Determine a carga por ciclo.
\item Determine a corrente média.
\end{itensq}
\resposta{$Q=\SI{100}{\milli\coulomb}$; $I_{\text{méd}}=\SI{1}{\ampere}$}


\end{questao}
\begin{questao}[2]{}

Um condutor de seção $\SI{1.5}{\milli\meter\squared}$ conduz
$\SI{15}{\ampere}$. Determine a densidade de corrente.
\resposta{$J=\SI{10}{\ampere\per\milli\meter\squared}=\pot{1}{7}\,
\si{\ampere\per\meter\squared}$}


\end{questao}
\begin{questao}[2]{}

Um capacitor de $\SI{100}{\micro\farad}$, inicialmente descarregado, recebe
corrente constante de $\SI{2}{\milli\ampere}$ durante $\SI{5}{\second}$.
\begin{itensq}
\item Determine a carga acumulada.
\item Determine a tensão final.
\end{itensq}
\resposta{$Q=\SI{10}{\milli\coulomb}$; $U=\SI{100}{\volt}$}


\end{questao}
\begin{questao}[2]{}

Em uma solução de $\mathrm{NaCl}$, íons $\mathrm{Na^{+}}$ movem-se para a
direita transportando $\SI{3}{\coulomb}$ por segundo, enquanto
$\mathrm{Cl^{-}}$ movem-se para a esquerda transportando
$\SI{2}{\coulomb}$ por segundo em módulo. Determine a corrente resultante.
\resposta{$I=\SI{5}{\ampere}$ para a direita}


\end{questao}
\begin{questao}[2]{}

Uma bateria descarregada de $\SI{1200}{\milli\ampere\hour}$ é recarregada com
corrente constante de $\SI{400}{\milli\ampere}$.
\begin{itensq}
\item Determine o tempo teórico de recarga.
\item Explique por que o tempo real é maior.
\end{itensq}
\resposta{$t=\SI{3}{\hour}$ teórico; na prática, $20$ a $40\%$ mais}


\end{questao}
\begin{questao}[2]{}

Uma corrente de $\SI{8}{\ampere}$ atravessa um fusível cuja especificação é
$I^{2}t=\SI{50}{\ampere\squared\second}$. Determine o tempo até a fusão.
\resposta{$t\approx\SI{0.78}{\second}$}


\end{questao}

\blocoaprofundamento
\begin{questao}[3]{}

A corrente em um circuito varia segundo $i(t)=1+2t$ (ampères, com $t$ em
segundos), no intervalo $0\le t\le\SI{2}{\second}$.
\begin{itensq}
\item Determine a carga total transferida.
\item Determine a corrente média e compare com a média dos extremos.
\end{itensq}
\resposta{$Q=\SI{6}{\coulomb}$; $I_{\text{méd}}=\SI{3}{\ampere}$}


\end{questao}
\begin{questao}[3]{}

Um pulso de corrente cresce linearmente de zero a $\SI{6}{\ampere}$ em
$\SI{2}{\milli\second}$ e retorna a zero em mais $\SI{3}{\milli\second}$.
\begin{itensq}
\item Determine a carga transferida.
\item Determine a corrente média no pulso.
\end{itensq}
\resposta{$Q=\SI{15}{\milli\coulomb}$; $I_{\text{méd}}=\SI{3}{\ampere}$}


\end{questao}
\begin{questao}[3]{}

Uma corrente trapezoidal sobe de $0$ a $\SI{4}{\ampere}$ em
$\SI{1}{\second}$, permanece constante por $\SI{3}{\second}$ e cai a zero em
$\SI{2}{\second}$.
\begin{itensq}
\item Determine a carga total.
\item Determine a corrente média.
\end{itensq}
\resposta{$Q=\SI{18}{\coulomb}$; $I_{\text{méd}}=\SI{3}{\ampere}$}


\end{questao}
\begin{questao}[3]{}

Uma corrente alterna entre $+\SI{4}{\ampere}$ por $\SI{3}{\second}$ e
$-\SI{6}{\ampere}$ por $\SI{2}{\second}$, repetindo-se.
\begin{itensq}
\item Determine a carga líquida por ciclo.
\item Determine a carga total \emph{transportada} em módulo.
\item Interprete a diferença.
\end{itensq}
\resposta{Líquida: $\SI{0}{\coulomb}$; total: $\SI{24}{\coulomb}$}


\end{questao}
\begin{questao}[3]{}

Uma corrente de descarga vale $i(t)=I_0e^{-t/\tau}$, com
$I_0=\SI{2}{\ampere}$ e $\tau=\SI{50}{\milli\second}$.
\begin{itensq}
\item Determine a carga total transferida.
\item Determine a fração transferida no primeiro $\tau$.
\end{itensq}
\resposta{$Q=\SI{0.1}{\coulomb}$; $63{,}2\%$}


\end{questao}
\begin{questao}[3]{}

Uma corrente de $\SI{2}{\ampere}$ atravessa uma célula eletrolítica de cobre
durante $\SI{30}{\minute}$.
\begin{itensq}
\item Determine a carga transferida.
\item Determine a massa de cobre depositada.
\end{itensq}
\dados{$M_{Cu}=\SI{63.5}{\gram\per\mole}$; valência $z=2$;
$F=\SI{96500}{\coulomb\per\mole}$}
\resposta{$Q=\SI{3600}{\coulomb}$; $m=\SI{1.18}{\gram}$}


\end{questao}
\begin{questao}[3]{}

Um condutor de $\SI{2.5}{\milli\meter\squared}$ deve conduzir
$\SI{20}{\ampere}$.
\begin{itensq}
\item Determine a densidade de corrente.
\item Compare com o limite usual de $\SI{5}{\ampere\per\milli\meter\squared}$ e
      determine a seção adequada.
\end{itensq}
\resposta{$J=\SI{8}{\ampere\per\milli\meter\squared}$ --- excessiva; adotar
$\SI{4}{\milli\meter\squared}$ ou mais}


\end{questao}
\begin{questao}[3]{}

Um capacitor de $\SI{470}{\micro\farad}$ carregado a $\SI{100}{\volt}$
apresenta corrente de fuga aproximadamente constante de
$\SI{20}{\micro\ampere}$.
\begin{itensq}
\item Determine a carga armazenada.
\item Estime o tempo para descarga completa.
\item Comente a validade da hipótese de corrente constante.
\end{itensq}
\resposta{$Q=\SI{47}{\milli\coulomb}$; $t\approx\SI{39}{\minute}$}


\end{questao}
\begin{questao}[3]{}

Compare a carga transportada e o aquecimento produzido por duas correntes de
mesma média: (A) $\SI{2}{\ampere}$ constantes; (B) $\SI{10}{\ampere}$ durante
$20\%$ do tempo e zero no restante.
\begin{itensq}
\item Determine a carga em $\SI{10}{\second}$ para cada uma.
\item Determine o valor eficaz de cada uma.
\item Compare o aquecimento em um resistor de $\SI{1}{\ohm}$.
\end{itensq}
\resposta{Mesma carga ($\SI{20}{\coulomb}$); $I_{rms}=2$ e
$\SI{4.47}{\ampere}$; aquecimento $5$ vezes maior em (B)}


\end{questao}
\begin{questao}[3]{}

Um circuito recebe pulsos de $\SI{50}{\ampere}$ com duração de
$\SI{100}{\micro\second}$, a uma taxa de $\SI{200}{\per\second}$.
\begin{itensq}
\item Determine a carga por pulso e a corrente média.
\item Determine o ciclo de trabalho.
\item Determine o valor eficaz.
\end{itensq}
\resposta{$\SI{5}{\milli\coulomb}$ por pulso; $I_{\text{méd}}=\SI{1}{\ampere}$;
$D=2\%$; $I_{rms}=\SI{7.07}{\ampere}$}


\end{questao}

\blocodesafio
\begin{questao}[4]{}

\textbf{(Autonomia sob corrente variável.)} Uma bateria de
$\SI{2.4}{\ampere\hour}$ alimenta uma carga cuja corrente alterna entre
$\SI{0.4}{\ampere}$ por $\SI{20}{\minute}$ e $\SI{1.2}{\ampere}$ por
$\SI{5}{\minute}$, repetidamente.
\begin{itensq}
\item Determine a carga consumida por ciclo e a corrente média.
\item Determine a autonomia.
\item Discuta por que a autonomia real será menor.
\end{itensq}
\resposta{$\SI{0.233}{\ampere\hour}$ por ciclo; $I_{\text{méd}}=\SI{0.56}{\ampere}$;
autonomia $\approx\SI{4.3}{\hour}$}


\end{questao}
\begin{questao}[4]{}

\textbf{(Erro de estimativa pelo pico.)} Um técnico estima a carga de um pulso
triangular ($0\to\SI{6}{\ampere}\to0$ em $\SI{5}{\milli\second}$) como
$Q\approx I_{pico}T$.
\begin{itensq}
\item Determine o erro cometido.
\item Generalize para formas de onda quaisquer.
\item Proponha uma regra prática.
\end{itensq}
\resposta{Erro de $+100\%$ --- o dobro do valor real}


\end{questao}
\begin{questao}[4]{}

\textbf{(Projeto de forma de onda.)} Projete uma corrente em \textbf{dois
patamares} que transfira exatamente $\SI{10}{\coulomb}$ em
$\SI{4}{\second}$, com valores comerciais simples.
\begin{itensq}
\item Estabeleça a condição geral.
\item Proponha duas soluções e compare o valor eficaz.
\item Indique qual é preferível e por quê.
\end{itensq}
\resposta{Condição: $I_1t_1+I_2t_2=10$ com $t_1+t_2=4$; a solução mais
uniforme minimiza o aquecimento}


\end{questao}
\begin{questao}[4]{}

\textbf{(Corrente bipolar.)} Uma corrente periódica é positiva em parte do
ciclo e negativa no restante.
\begin{itensq}
\item Formule a condição para carga líquida nula.
\item Determine se isso implica ausência de efeitos.
\item Cite duas aplicações em que essa condição é imposta deliberadamente.
\end{itensq}
\resposta{$\int_0^{T}i\,\mathrm{d}t=0$; os efeitos térmicos e mecânicos
permanecem}


\end{questao}
\begin{questao}[4]{}

\textbf{(Medição --- média contra instantâneo.)} Dois instrumentos medem a
mesma corrente pulsada: um indica a média temporal, outro o valor instantâneo.
\begin{itensq}
\item Explique por que os resultados diferem.
\item Determine qual usar para estimar carga, aquecimento e pico.
\item Projete uma medição que forneça os três.
\end{itensq}
\resposta{Média para carga, eficaz para aquecimento, instantâneo para pico}


\end{questao}
\begin{questao}[4]{}

\textbf{(Integração de sinal ruidoso.)} Uma corrente medida contém ruído
aditivo de média zero.
\begin{itensq}
\item Explique por que integrar melhora a estimativa da carga.
\item Determine como o erro decresce com o tempo de integração.
\item Indique em que situação integrar \emph{piora} o resultado.
\end{itensq}
\resposta{O erro relativo decresce como $1/\sqrt{N}$; integrar piora se houver
\emph{deriva} (ruído de média não nula)}


\end{questao}
\begin{questao}[4]{}

\textbf{(Por que a LKC é tão precisa.)} A lei dos nós afirma que a carga não se
acumula em um nó. Quantifique o quanto ela poderia, em princípio, ser violada.
\begin{itensq}
\item Estime a capacitância parasita de um nó típico e relacione carga acumulada
      com potencial.
\item Determine o desequilíbrio de corrente necessário para elevar o nó em
      $\SI{1}{\volt}$.
\item Conclua sobre o caráter da lei.
\end{itensq}
\resposta{Com $C\approx\SI{1}{\pico\farad}$, basta $\SI{1}{\pico\coulomb}$ para
$\SI{1}{\volt}$ --- um desequilíbrio de $\SI{1}{\micro\ampere}$ por
$\SI{1}{\micro\second}$}


\end{questao}
\begin{questao}[4]{}

\textbf{(Deriva contra propagação.)} Um fio de cobre de
$\SI{2.5}{\milli\meter\squared}$ conduz $\SI{10}{\ampere}$, com densidade de
portadores $n=\pot{8.5}{28}\,\si{\per\meter\cubed}$.
\begin{itensq}
\item Determine a velocidade de deriva dos elétrons.
\item Determine o tempo para um elétron percorrer $\SI{1}{\meter}$.
\item Compare com o tempo de propagação do sinal e explique o paradoxo
      aparente.
\end{itensq}
\resposta{$v_d=\SI{0.29}{\milli\meter\per\second}$; $\SI{57}{\minute}$ para
$\SI{1}{\meter}$; o sinal leva $\SI{5}{\nano\second}$}


\end{questao}

\gabarito[2.1 Corrente elétrica e carga]

\subsection{Lei de Ohm e resistividade}

\blocofixacao
\begin{questao}[1]{}

De acordo com a segunda lei de Ohm, a resistência de um fio condutor:
\begin{alternativas}
\item aumenta com a área da seção e diminui com o comprimento.
\item aumenta com o comprimento e diminui com a área da seção.
\item depende apenas do material.
\item independe da temperatura.
\item é sempre igual para todos os materiais.
\end{alternativas}
\resposta{(b)}


\end{questao}
\begin{questao}[1]{}

O gráfico mostra a resistência de um fio em função de seu comprimento, mantida a
área constante.

\circuitoqq{
\begin{tikzpicture}
 \begin{axis}[width=9cm,height=5cm,xlabel={Comprimento (m)},
   ylabel={$R$ ($\Omega$)},xmin=0,xmax=4,ymin=0,ymax=7,
   xtick={0,1,2,3},ytick={0,2,4,6},grid=major,axis lines=left]
   \addplot[Rcol,very thick,domain=0:3] {2*x};
   \addplot[only marks,mark=*,Rcol] coordinates {(1,2)(2,4)(3,6)};
 \end{axis}
\end{tikzpicture}}{Resistência proporcional ao comprimento.}

\begin{itensq}
\item Qual é a constante de proporcionalidade $k$ da reta $R = kL$?
\item Qual seria a resistência de um fio de $\SI{5}{\meter}$ do mesmo tipo?
\item Que lei essa proporcionalidade confirma?
\end{itensq}
\resposta{(a) $k=\SI{2}{\ohm\per\meter}$; (b) \SI{10}{\ohm}; (c) 2ª lei de Ohm}


\end{questao}

\blocoaplicacao
\begin{questao}[2]{}

Um estudante quer construir um resistor de $\SI{10}{\ohm}$ com fio de
níquel-cromo, de resistividade $\rho = \SI{1.1e-6}{\ohm\meter}$ e seção
$A = \SI{0.5}{\milli\meter\squared}$. Qual é o comprimento necessário?
\resposta{$L \approx \SI{4.55}{\meter}$}


\end{questao}
\begin{questao}[2]{}

Um fio condutor de comprimento $L$, diâmetro $D$ e resistência $R$ tem seu
comprimento \textbf{e} seu diâmetro \emph{duplicados}. A nova resistência será:
\begin{alternativas}[3]
\item $R$
\item $2R$
\item $R/2$
\item $4R$
\item $R/4$
\end{alternativas}
\resposta{(c)}


\end{questao}
\begin{questao}[2]{}

O gráfico representa a diferença de potencial $U$ entre dois pontos de um fio em
função da corrente $i$. A resistência do trecho, em ohms, vale:

\circuitoqq{
\begin{tikzpicture}
 \begin{axis}[width=9cm,height=5cm,xlabel={$i$ (A)},ylabel={$U$ (V)},
   xmin=0,xmax=5,ymin=0,ymax=22,xtick={0,1,2,3,4},ytick={0,4,8,12,16,20},
   grid=major,axis lines=left]
   \addplot[Vcol,very thick,domain=0:4] {4*x};
   \addplot[only marks,mark=*,Vcol] coordinates {(1,4)(2,8)(3,12)(4,16)};
 \end{axis}
\end{tikzpicture}}{Curva característica $U \times i$ de um resistor ôhmico.}

\begin{alternativas}[3]
\item 0,05
\item 4
\item 20
\item 80
\item 160
\end{alternativas}
\resposta{(b)}


\end{questao}
\begin{questao}[2]{}

Dois fios do mesmo material têm resistências $R_1 = \SI{6}{\ohm}$ e
$R_2 = \SI{3}{\ohm}$; o fio 1 tem o dobro do comprimento do fio 2. Qual é a relação
entre as áreas das seções transversais, $A_1/A_2$?
\resposta{$A_1/A_2 = 1$ (áreas iguais)}


\end{questao}
\begin{questao}[2]{}

A tabela apresenta a resistividade de alguns materiais a $\SI{20}{\celsius}$.

\begin{table}[H]
\centering\small\sffamily
\begin{tabular}{@{}l c@{}}
\toprule
\textbf{Material} & \textbf{$\rho$ ($\times10^{-8}\,\si{\ohm\meter}$)}\\
\midrule
Prata      & 1,6\\
Cobre      & 1,7\\
Alumínio   & 2,8\\
Ferro      & 10\\
Níquel-cromo & 110\\
\bottomrule
\end{tabular}
\caption{Resistividade de materiais condutores.}
\end{table}

Dois fios de mesmo comprimento e mesma seção, um de cobre e outro de alumínio,
têm resistências $R_{\text{Cu}}$ e $R_{\text{Al}}$. A razão
$R_{\text{Al}}/R_{\text{Cu}}$ vale aproximadamente:
\begin{alternativas}[3]
\item 0,6
\item 1,0
\item 1,6
\item 2,8
\item 4,4
\end{alternativas}
\resposta{(c)}


\end{questao}
\begin{questao}[2]{}

Um fio de cobre tem comprimento $L = \SI{2}{\meter}$ e seção transversal
$A = \SI{1}{\milli\meter\squared}$. Calcule sua resistência.
\dados{$\rho_{\text{Cu}} = \SI{1.7e-8}{\ohm\meter}$.}
\resposta{$R = \SI{34}{\milli\ohm}$}


\end{questao}
\begin{questao}[2]{Apostila original revisada}

Um fio possui resistência elétrica $R$. Mantendo o mesmo material, seu comprimento é
duplicado e sua área de seção transversal é reduzida à metade. Determine a nova
resistência em função de $R$.
\resposta{$R'=4R$}


\end{questao}

\blocoaprofundamento
\begin{questao}[3]{}

Um fio de aço e um de cobre, de mesmo comprimento e mesma área, são ligados
individualmente à mesma tensão. Sabendo que
$\rho_{\text{aço}} = 10\,\rho_{\text{cobre}}$, compare:
\begin{itensq}
\item as correntes em cada fio;
\item as potências dissipadas em cada um.
\end{itensq}
\resposta{(a) $I_{\text{cobre}} = 10\,I_{\text{aço}}$;
(b) $P_{\text{cobre}} = 10\,P_{\text{aço}}$}


\end{questao}
\begin{questao}[3]{}

Dois fios de \emph{mesmo comprimento} e \emph{mesma seção}, um de cobre
($\rho_{\text{Cu}}=\SI{1.7e-8}{\ohm\meter}$) e outro de alumínio
($\rho_{\text{Al}}=\SI{2.8e-8}{\ohm\meter}$), são emendados em série e ligados a
$\SI{12}{\volt}$.
\begin{itensq}
\item Determine a razão entre as tensões em cada trecho.
\item Calcule a tensão em cada fio.
\item Em qual dos dois há maior dissipação de calor? Justifique.
\end{itensq}
\resposta{(a) $U_{\text{Al}}/U_{\text{Cu}}\approx1{,}65$;
(b) $\SI{4.53}{\volt}$ e $\SI{7.47}{\volt}$; (c) no alumínio}


\end{questao}
\begin{questao}[3]{}

A resistência de um condutor metálico varia com a temperatura segundo
$R = R_0\big[1+\alpha(T-T_0)\big]$, onde $\alpha$ é o coeficiente de temperatura.
Para o cobre, $\alpha \approx \SI{4e-3}{\per\celsius}$.

Um enrolamento de cobre de um motor mede $\SI{10}{\ohm}$ a $\SI{20}{\celsius}$
(motor frio). Após operar em regime, sua resistência sobe para
$\SI{12.4}{\ohm}$.
\begin{itensq}
\item Determine a temperatura do enrolamento em operação.
\item Se a tensão aplicada for mantida em $\SI{220}{\volt}$, qual é a variação
      percentual da potência dissipada entre o motor frio e o quente?
\item Explique como esse efeito é usado, na prática, para medir a temperatura
      interna de máquinas elétricas sem sensores.
\end{itensq}
\resposta{(a) $T = \SI{80}{\celsius}$; (b) a potência cai $\approx\SI{19}{\percent}$;
(c) ver comentário}


\end{questao}
\begin{questao}[3]{Apostila original revisada}

A resistividade de um material varia aproximadamente de forma linear com a
temperatura. Foram medidos os valores
$\rho(\SI{0}{\celsius})=\SI{1.5e-8}{\ohm\metre}$ e
$\rho(\SI{100}{\celsius})=\SI{2.5e-8}{\ohm\metre}$.
\begin{itensq}
\item Determine a taxa de variação da resistividade por grau Celsius.
\item Estime a resistividade a $\SI{150}{\celsius}$, mantendo o modelo linear.
\end{itensq}
\resposta{(a) $\SI{1.0e-10}{\ohm\metre\per\celsius}$; (b) $\SI{3.0e-8}{\ohm\metre}$}


\end{questao}

\blocofixacao
\begin{questao}[1]{}

Um resistor de $\SI{470}{\ohm}$ é percorrido por $\SI{15}{\milli\ampere}$.
Determine a tensão em seus terminais.
\resposta{$U=\SI{7.05}{\volt}$}


\end{questao}
\begin{questao}[1]{}

Uma tensão de $\SI{9}{\volt}$ produz corrente de
$\SI{30}{\milli\ampere}$ em um componente ôhmico. Determine sua resistência.
\resposta{$R=\SI{300}{\ohm}$}


\end{questao}
\begin{questao}[1]{}

Determine a corrente em um resistor de $\SI{2.2}{\kilo\ohm}$ submetido a
$\SI{12}{\volt}$.
\resposta{$I=\SI{5.45}{\milli\ampere}$}


\end{questao}
\begin{questao}[1]{}

Um condutor de cobre ($\rho=\pot{1.7}{-8}\,\si{\ohm\meter}$) tem
$\SI{50}{\meter}$ e seção de $\SI{2.5}{\milli\meter\squared}$. Determine sua
resistência.
\resposta{$R=\SI{0.34}{\ohm}$}


\end{questao}
\begin{questao}[1]{}

Determine a condutância de um resistor de $\SI{25}{\ohm}$.
\resposta{$G=\SI{0.04}{\ensuremath{\mathrm{S}}}$}


\end{questao}
\begin{questao}[1]{}

Segundo a segunda lei de Ohm, dobrar o \textbf{comprimento} de um fio, mantida
a seção, faz a resistência:
\begin{alternativas}[2]
\item dobrar
\item cair à metade
\item quadruplicar
\item não se alterar
\end{alternativas}
\resposta{(a)}


\end{questao}
\begin{questao}[1]{}

Dobrar a \textbf{área} da seção de um fio, mantido o comprimento, faz a
resistência:
\begin{alternativas}[2]
\item dobrar
\item cair à metade
\item quadruplicar
\item não se alterar
\end{alternativas}
\resposta{(b)}


\end{questao}
\begin{questao}[1]{}

Dois fios de mesmo material e mesmo comprimento têm diâmetros $d$ e $2d$.
Determine a razão entre suas resistências.
\resposta{$R_1/R_2=4$}


\end{questao}

\blocoaplicacao
\begin{questao}[2]{}

Determine o comprimento de um fio de níquel-cromo
($\rho=\pot{1.1}{-6}\,\si{\ohm\meter}$) de seção
$\SI{0.2}{\milli\meter\squared}$ necessário para obter
$\SI{22}{\ohm}$.
\resposta{$L=\SI{4}{\meter}$}


\end{questao}
\begin{questao}[2]{}

Um cabo de alumínio ($\rho=\pot{2.8}{-8}\,\si{\ohm\meter}$) de
$\SI{200}{\meter}$ deve ter resistência de $\SI{0.5}{\ohm}$. Determine a seção
e o diâmetro.
\resposta{$A=\SI{11.2}{\milli\meter\squared}$; $d=\SI{3.78}{\milli\meter}$}


\end{questao}
\begin{questao}[2]{}

Mediu-se $R=\SI{2.5}{\ohm}$ em um fio de $\SI{30}{\meter}$ e seção
$\SI{0.5}{\milli\meter\squared}$. Determine a resistividade e identifique o
material provável.
\resposta{$\rho=\pot{4.2}{-8}\,\si{\ohm\meter}$ --- compatível com latão ou
zinco}


\end{questao}
\begin{questao}[2]{}

Um fio cilíndrico é \textbf{esticado} até dobrar seu comprimento, conservando o
volume e o material. Determine a nova resistência.
\resposta{$R'=4R$}


\end{questao}
\begin{questao}[2]{}

Um resistor de fio tem $\SI{50}{\ohm}$ a $\SI{20}{\celsius}$, com
$\alpha=\SI{0.0039}{\per\celsius}$. Determine sua resistência a
$\SI{80}{\celsius}$.
\resposta{$R=\SI{61.7}{\ohm}$}


\end{questao}
\begin{questao}[2]{}

Um estudante mediu, em um componente, os pares
$(\SI{2}{\volt};\SI{20}{\milli\ampere})$,
$(\SI{4}{\volt};\SI{40}{\milli\ampere})$ e
$(\SI{6}{\volt};\SI{57}{\milli\ampere})$. O componente é ôhmico?
\resposta{Não --- a razão $U/I$ cresce de $100$ para $\SI{105}{\ohm}$}


\end{questao}
\begin{questao}[2]{}

Compare a seção necessária para obter a mesma resistência com cobre
($\rho=\pot{1.7}{-8}$) e com alumínio ($\pot{2.8}{-8}\,\si{\ohm\meter}$),
mantido o comprimento.
\resposta{$A_{Al}/A_{Cu}=1{,}65$}


\end{questao}
\begin{questao}[2]{}

Um fio de cobre de $\SI{100}{\meter}$ e $\SI{1.5}{\milli\meter\squared}$
conduz $\SI{10}{\ampere}$. Determine a resistência, a queda de tensão e a
potência dissipada.
\resposta{$R=\SI{1.13}{\ohm}$; $\Delta U=\SI{11.3}{\volt}$;
$P=\SI{113}{\watt}$}


\end{questao}

\blocoaprofundamento
\begin{questao}[3]{}

Projete um resistor de fio de $\SI{10}{\ohm}$ usando material de
$\rho=\pot{5}{-7}\,\si{\ohm\meter}$ e fio de $\SI{0.5}{\milli\meter}$ de
diâmetro.
\begin{itensq}
\item Determine o comprimento necessário.
\item Determine a corrente máxima para $\SI{5}{\watt}$ de dissipação.
\item Comente a construção.
\end{itensq}
\resposta{(a) $L=\SI{3.93}{\meter}$; (b) $I=\SI{0.707}{\ampere}$}


\end{questao}
\begin{questao}[3]{}

Um resistor nominal de $\SI{100}{\ohm}$ tem tolerância de $\SI{5}{\percent}$ e
coeficiente térmico de $400\,\mathrm{ppm/^{\circ}C}$, especificado a
$\SI{20}{\celsius}$. Determine a faixa total de resistência na operação entre
$\SI{0}{\celsius}$ e $\SI{70}{\celsius}$.
\resposta{de $\SI{94.2}{\ohm}$ a $\SI{107.1}{\ohm}$ ($-5{,}8\%$ a
$+7{,}1\%$)}


\end{questao}
\begin{questao}[3]{}

Uma lâmpada incandescente apresenta $\SI{2.00}{\ampere}$ a
$\SI{12.0}{\volt}$ e $\SI{2.02}{\ampere}$ a $\SI{12.5}{\volt}$.
\begin{itensq}
\item Determine a resistência estática no primeiro ponto.
\item Determine a resistência dinâmica.
\item Explique a diferença.
\end{itensq}
\resposta{$R_{est}=\SI{6}{\ohm}$; $r_{din}=\SI{25}{\ohm}$}


\end{questao}
\begin{questao}[3]{}

Um cabo de cobre deve alimentar uma carga de $\SI{20}{\ampere}$ a
$\SI{220}{\volt}$, a $\SI{40}{\meter}$ de distância, com queda de tensão
inferior a $\SI{3}{\percent}$.
\begin{itensq}
\item Determine a resistência máxima admissível do cabo.
\item Determine a seção mínima.
\item Escolha a bitola comercial.
\end{itensq}
\resposta{$R\le\SI{0.33}{\ohm}$; $A\ge\SI{4.12}{\milli\meter\squared}$;
adotar $\SI{6}{\milli\meter\squared}$}


\end{questao}
\begin{questao}[3]{}

Um fio é esticado até que seu comprimento fique $n$ vezes maior, conservando o
volume.
\begin{itensq}
\item Deduza a nova resistência em função de $n$.
\item Avalie para $n=3$ e $n=5$.
\item Determine a redução do diâmetro em cada caso.
\end{itensq}
\resposta{$R'=n^{2}R$; $9R$ e $25R$; $d'/d=1/\sqrt{n}$}


\end{questao}
\begin{questao}[3]{}

Uma barra condutora de comprimento $L=\SI{1}{\meter}$ tem seção que varia
linearmente segundo $A(x)=A_0\left(1+x/L\right)$, com
$A_0=\SI{1}{\milli\meter\squared}$ e $\rho=\pot{1.7}{-8}\,\si{\ohm\meter}$.
\begin{itensq}
\item Determine a resistência total.
\item Compare com uma barra uniforme de seção $A_0$ e com uma de seção média.
\end{itensq}
\resposta{$R=\dfrac{\rho L\ln2}{A_0}=\SI{11.8}{\milli\ohm}$}


\end{questao}
\begin{questao}[3]{}

Um sensor Pt100 tem $R_0=\SI{100}{\ohm}$ a $\SI{0}{\celsius}$ e
$\alpha=\SI{0.00385}{\per\celsius}$, alimentado por corrente constante de
$\SI{1}{\milli\ampere}$.
\begin{itensq}
\item Determine a tensão a $\SI{0}{\celsius}$ e a sensibilidade em
      $\si{\milli\volt\per\celsius}$.
\item Explique a vantagem da alimentação por corrente constante.
\item Verifique o autoaquecimento.
\end{itensq}
\resposta{$U_0=\SI{100}{\milli\volt}$;
$\SI{0.385}{\milli\volt\per\celsius}$}


\end{questao}
\begin{questao}[3]{}

Mede-se um resistor de $\SI{0.1}{\ohm}$ com um ohmímetro cujos cabos têm
$\SI{0.2}{\ohm}$ cada.
\begin{itensq}
\item Determine a leitura pelo método de \textbf{dois fios}.
\item Determine o erro.
\item Explique como o método de \textbf{quatro fios} o elimina.
\end{itensq}
\resposta{Leitura de $\SI{0.5}{\ohm}$ --- erro de $400\%$}


\end{questao}
\begin{questao}[3]{}

Dois resistores têm o mesmo valor nominal de $\SI{100}{\ohm}$, mas potências de
$\frac14\,\si{\watt}$ e $\SI{5}{\watt}$.
\begin{itensq}
\item Determine a tensão e a corrente máximas de cada um.
\item Ambos obedecem igualmente à lei de Ohm?
\item Explique o que a especificação de potência realmente limita.
\end{itensq}
\resposta{$\SI{5}{\volt}$/$\SI{50}{\milli\ampere}$ e
$\SI{22.4}{\volt}$/$\SI{224}{\milli\ampere}$; ambos são ôhmicos}


\end{questao}
\begin{questao}[3]{}

Um resistor de $\SI{100}{\ohm}$ com $\alpha=\SI{+0.004}{\per\celsius}$ dissipa
$\SI{1}{\watt}$ e tem resistência térmica de
$\SI{100}{\celsius\per\watt}$.
\begin{itensq}
\item Determine a elevação de temperatura e a nova resistência.
\item Analise o efeito sob tensão constante e sob corrente constante.
\end{itensq}
\resposta{$\Delta T=\SI{100}{\celsius}$; $R=\SI{140}{\ohm}$}


\end{questao}
\begin{questao}[3]{}

Um resistor de $\SI{100}{\ohm}$ e $\frac14\,\si{\watt}$ deve ser testado quanto
à ohmicidade. Projete o ensaio.
\begin{itensq}
\item Determine a faixa segura de tensão.
\item Proponha os pontos de medição e o critério de decisão.
\item Indique os cuidados experimentais.
\end{itensq}
\resposta{Até $\SI{5}{\volt}$; usar $20\%$ da potência nominal como teto
prático}


\end{questao}

\blocodesafio
\begin{questao}[4]{}

\textbf{(Demonstração --- deformação e resistência.)} Um fio de comprimento $L$
e seção $A$ sofre pequena deformação axial $\varepsilon=\Delta L/L$, a volume
constante.
\begin{itensq}
\item Deduza $\Delta R/R$ em primeira ordem.
\item Determine o fator de sensibilidade.
\item Discuta o efeito de a resistividade também variar com a deformação.
\end{itensq}
\resposta{$\Delta R/R\approx2\varepsilon$; fator $\approx2$}


\end{questao}
\begin{questao}[4]{}

\textbf{(Integração --- seção variável.)} Uma barra tem seção
$A(x)=A_0\left(1+x/L\right)$.
\begin{itensq}
\item Deduza $R$ por integração.
\item Generalize para $A(x)=A_0(1+kx/L)$.
\item Analise os limites $k\to0$ e $k\to\infty$.
\end{itensq}
\resposta{$R=\dfrac{\rho L}{A_0}\cdot\dfrac{\ln(1+k)}{k}$; para $k=1$,
$\ln2$}


\end{questao}
\begin{questao}[4]{}

\textbf{(Ponto de operação com fonte não ideal.)} Um resistor de
$\SI{100}{\ohm}$ é ligado a uma fonte cuja característica é
$U=24-8I$ (com $U$ em volts e $I$ em ampères).
\begin{itensq}
\item Determine o ponto de operação.
\item Determine a potência no resistor e o rendimento.
\item Repita para um resistor de $\SI{8}{\ohm}$ e compare.
\end{itensq}
\resposta{$I=\SI{0.222}{\ampere}$, $U=\SI{22.2}{\volt}$;
$P=\SI{4.94}{\watt}$, $\eta=92{,}6\%$}


\end{questao}
\begin{questao}[4]{}

\textbf{(Estabilidade térmica.)} Um resistor de resistência $R$, coeficiente
$\alpha$ e resistência térmica $R_{\text{térm}}$ é alimentado por corrente constante
$I$.
\begin{itensq}
\item Deduza a condição de estabilidade térmica.
\item Avalie para $R=\SI{100}{\ohm}$, $\alpha=\SI{0.004}{\per\celsius}$,
      $R_{\text{térm}}=\SI{100}{\celsius\per\watt}$ e $I=\SI{100}{\milli\ampere}$.
\item Determine a corrente crítica.
\end{itensq}
\resposta{Condição: $\alpha R I^{2}R_{\text{térm}}<1$; aqui $0{,}4$ --- estável;
$I_c=\SI{158}{\milli\ampere}$}


\end{questao}
\begin{questao}[4]{}

\textbf{(Dois critérios de dimensionamento.)} Um cabo deve alimentar
$\SI{30}{\ampere}$ a $\SI{220}{\volt}$, com queda máxima de
$\SI{4}{\percent}$.
\begin{itensq}
\item Determine a seção exigida pelo critério de queda de tensão, para
      $\SI{20}{\meter}$ e para $\SI{80}{\meter}$.
\item Sabendo que $\SI{6}{\milli\meter\squared}$ suporta cerca de
      $\SI{41}{\ampere}$, determine qual critério domina em cada caso.
\item Formule a regra geral.
\end{itensq}
\resposta{$\SI{3.87}{\milli\meter\squared}$ e
$\SI{15.5}{\milli\meter\squared}$; a corrente domina no cabo curto, a queda no
longo}


\end{questao}
\begin{questao}[4]{}

\textbf{(Sensor a corrente constante.)} Compare duas formas de converter a
variação de um sensor resistivo em tensão.
\begin{itensq}
\item Obtenha $U(R)$ para excitação por corrente constante e para divisor de
      tensão.
\item Compare a linearidade das duas.
\item Determine a sensibilidade de cada uma no ponto de operação
      $R=R_{fixo}$.
\end{itensq}
\resposta{Corrente constante: $U=IR$, exatamente linear; divisor:
$U=U_0R/(R+R_f)$, não linear}


\end{questao}
\begin{questao}[4]{}

\textbf{(Resistividade e temperatura.)} A resistividade de um metal varia como
$\rho(T)=\rho_0\left[1+\alpha(T-T_0)\right]$.
\begin{itensq}
\item Mostre que a mesma lei vale para $R$, e identifique a hipótese oculta.
\item Determine o erro cometido ao ignorar a dilatação térmica das dimensões.
\item Avalie para cobre entre $\SI{20}{\celsius}$ e
      $\SI{120}{\celsius}$.
\end{itensq}
\resposta{A hipótese oculta é dimensões constantes; o erro é da ordem de
$\SI{0.05}{\percent}$ --- desprezível}


\end{questao}
\begin{questao}[4]{}

\textbf{(Medição a quatro fios.)} Quantifique a vantagem do método de Kelvin.
\begin{itensq}
\item Deduza o erro relativo do método de dois fios em função de
      $R_{cabo}/R_x$.
\item Determine a partir de qual valor de $R_x$ o método de dois fios dá erro
      inferior a $\SI{1}{\percent}$, com cabos de $\SI{0.2}{\ohm}$ cada.
\item Explique por que o método de quatro fios não elimina \emph{todos} os
      erros.
\end{itensq}
\resposta{Erro $=2R_{cabo}/R_x$; é preciso $R_x\ge\SI{40}{\ohm}$}


\end{questao}
\begin{questao}[4]{}

\textbf{(Condutor composto.)} Um cabo de alta tensão tem alma de aço
($\rho_{\text{aço}}=\pot{1.5}{-7}\,\si{\ohm\meter}$, $A=\SI{20}{\milli\meter\squared}$)
envolvida por alumínio ($\rho_{Al}=\pot{2.8}{-8}\,\si{\ohm\meter}$,
$A=\SI{100}{\milli\meter\squared}$), com $\SI{1}{\kilo\meter}$ de comprimento.
\begin{itensq}
\item Determine a resistência de cada material e a do conjunto.
\item Determine a fração de corrente em cada um.
\item Explique a função da alma de aço.
\end{itensq}
\resposta{$R_{\text{aço}}=\SI{7.5}{\ohm}$, $R_{Al}=\SI{0.28}{\ohm}$,
$R_{eq}=\SI{0.270}{\ohm}$; o aço conduz $3{,}6\%$}


\end{questao}
\begin{questao}[4]{}

\textbf{(Limites da lei de Ohm.)} Discuta situações em que a relação
$U=RI$ deixa de descrever o comportamento de um condutor.
\begin{itensq}
\item Enumere ao menos quatro mecanismos, com exemplos.
\item Explique por que a lei de Ohm não é uma lei fundamental.
\item Indique como se procede quando ela falha.
\end{itensq}
\resposta{Autoaquecimento, não linearidade intrínseca, efeitos de alta
frequência e ruptura do meio}


\end{questao}

\gabarito[2.2 Lei de Ohm e resistividade]

\subsection{Potência elétrica e energia}

\blocofixacao
\begin{questao}[1]{}

A tensão nos terminais de um condutor é $\SI{12}{\volt}$ e ele dissipa
$\SI{24}{\watt}$.
\begin{itensq}
\item Qual é a corrente que o percorre?
\item Qual é o valor de sua resistência?
\end{itensq}
\resposta{(a) \SI{2}{\ampere}; (b) \SI{6}{\ohm}}


\end{questao}
\begin{questao}[1]{}

Um aparelho opera em $\SI{220}{\volt}$ e consome $\SI{440}{\watt}$.
\begin{itensq}
\item Qual é a corrente que ele consome?
\item Qual sua resistência equivalente?
\end{itensq}
\resposta{(a) \SI{2}{\ampere}; (b) \SI{110}{\ohm}}


\end{questao}
\begin{questao}[1]{Apostila original revisada}

Uma fonte fornece $\SI{100}{\watt}$ a um circuito, mas somente
$\SI{95}{\watt}$ são aproveitados pela carga. Determine o rendimento do circuito
e a potência dissipada em perdas.
\resposta{$\eta=95\%$; perdas de $\SI{5}{\watt}$}


\end{questao}

\blocoaplicacao
\begin{questao}[2]{}

Uma lâmpada de especificação $\SI{60}{\watt}$/$\SI{220}{\volt}$ é ligada, por
engano, a uma rede de $\SI{110}{\volt}$. Sua potência dissipada passa a ser:
\begin{alternativas}
\item inalterada.
\item reduzida à metade.
\item duplicada.
\item reduzida à quarta parte.
\item aumentada quatro vezes.
\end{alternativas}
\resposta{(d)}


\end{questao}
\begin{questao}[2]{}

Um chuveiro elétrico tem especificações $\SI{220}{\volt}$/$\SI{4400}{\watt}$. Se
for ligado a uma rede de $\SI{110}{\volt}$, a potência dissipada será:
\begin{alternativas}[2]
\item \SI{550}{\watt}
\item \SI{1100}{\watt}
\item \SI{2200}{\watt}
\item \SI{4400}{\watt}
\item \SI{8800}{\watt}
\end{alternativas}
\resposta{(b)}


\end{questao}
\begin{questao}[2]{}

Dois resistores $R_1 = \SI{10}{\ohm}$ e $R_2 = \SI{20}{\ohm}$ estão em série sob
uma tensão total de $\SI{60}{\volt}$.
\begin{itensq}
\item Qual é a corrente do circuito?
\item Qual é a potência dissipada em cada resistor?
\end{itensq}
\resposta{(a) \SI{2}{\ampere}; (b) $P_1=\SI{40}{\watt}$, $P_2=\SI{80}{\watt}$}


\end{questao}
\begin{questao}[2]{}

Um chuveiro elétrico de $\SI{5500}{\watt}$ opera em $\SI{220}{\volt}$.
\begin{itensq}
\item Qual é a corrente que ele exige da rede?
\item Qual é a resistência da sua resistência de aquecimento?
\item O disjuntor do circuito é de $\SI{20}{\ampere}$. Ele suporta esse chuveiro?
\end{itensq}
\resposta{(a) \SI{25}{\ampere}; (b) \SI{8.8}{\ohm}; (c) não}


\end{questao}
\begin{questao}[2]{Apostila original revisada}

Em uma residência, a tarifa média foi obtida de uma conta de
R\$~80,00 para um consumo de $\SI{200}{\kilo\watt\hour}$. Um chuveiro de
$\SI{3}{\kilo\watt}$ é utilizado durante $\SI{30}{\minute}$ por dia, ao longo de
30 dias. Calcule a energia mensal consumida pelo chuveiro e seu custo.
\resposta{$\SI{45}{\kilo\watt\hour}$; R\$~18,00}


\end{questao}

\blocoaprofundamento
\begin{questao}[3]{}

Em uma residência permanecem ligados, durante $\SI{2}{\hour}$, um ferro elétrico
de $\SI{1440}{\watt}$/$\SI{120}{\volt}$ e duas lâmpadas de
$\SI{60}{\watt}$/$\SI{120}{\volt}$. Mantida a tensão em $\SI{120}{\volt}$,
determine a corrente total no fusível e a energia elétrica consumida.
\resposta{$I = \SI{13}{\ampere}$; $E = \SI{3.12}{\kilo\watt\hour}$}


\end{questao}
\begin{questao}[3]{}

Uma residência utiliza, por dia: uma lâmpada de $\SI{100}{\watt}$ acesa
$\SI{5}{\hour}$, um chuveiro de $\SI{4400}{\watt}$ ligado $\SI{0.5}{\hour}$ e uma
geladeira de $\SI{200}{\watt}$ que funciona, em média, $\SI{8}{\hour}$ por dia.
Calcule o consumo mensal de energia (30 dias) e o custo, considerando a tarifa de
$\SI{0.80}{\text{R\$}\per\kilo\watt\hour}$.
\resposta{$E \approx \SI{127.5}{\kilo\watt\hour}$; custo $\approx$ R\$\,102,00}


\end{questao}
\begin{questao}[3]{}

Duas lâmpadas incandescentes de especificações $\SI{60}{\watt}$/$\SI{220}{\volt}$
e $\SI{100}{\watt}$/$\SI{220}{\volt}$ são ligadas \textbf{em série} à rede de
$\SI{220}{\volt}$.
\begin{itensq}
\item Calcule a resistência de cada lâmpada (suposta constante).
\item Determine a potência dissipada por cada uma nessa ligação.
\item \textbf{Qual delas brilha mais?} O resultado contraria a intuição?
      Explique.
\end{itensq}
\resposta{(a) $\SI{806.7}{\ohm}$ e $\SI{484}{\ohm}$;
(b) $\SI{23.4}{\watt}$ e $\SI{14.1}{\watt}$; (c) a de \SI{60}{\watt}}


\end{questao}
\begin{questao}[3]{}

Um chuveiro elétrico de $\SI{220}{\volt}$ possui uma resistência de aquecimento
com uma \emph{derivação central}, permitindo duas posições: na posição "verão",
usa-se a resistência inteira ($\SI{8.8}{\ohm}$); na posição "inverno", apenas
\emph{metade} dela.
\begin{itensq}
\item Calcule a potência em cada posição.
\item Determine a corrente em cada caso e verifique se um disjuntor de
      $\SI{40}{\ampere}$ é adequado.
\item Sabendo que a água entra a $\SI{20}{\celsius}$ com vazão de
      $\SI{3}{\liter\per\minute}$, estime a temperatura de saída no inverno.
\end{itensq}
\dados{$c_{\text{água}}=\SI{4180}{\joule\per\kilogram\per\celsius}$;
$\rho_{\text{água}}=\SI{1}{\kilogram\per\liter}$.}
\resposta{(a) \SI{5500}{\watt} e \SI{11000}{\watt}; (b) \SI{25}{\ampere} e
\SI{50}{\ampere} --- disjuntor insuficiente; (c) $\approx\SI{72}{\celsius}$}


\end{questao}

\blocofixacao
\begin{questao}[1]{}

Um resistor de $\SI{8}{\ohm}$ é percorrido por $\SI{2.5}{\ampere}$. Determine a
potência dissipada.
\resposta{$P=\SI{50}{\watt}$}


\end{questao}
\begin{questao}[1]{}

Um resistor de $\SI{44}{\ohm}$ é submetido a $\SI{220}{\volt}$. Determine a
potência.
\resposta{$P=\SI{1100}{\watt}$}


\end{questao}
\begin{questao}[1]{}

Converta $\SI{250}{\watt\hour}$ para joules.
\resposta{$\SI{900}{\kilo\joule}$}


\end{questao}
\begin{questao}[1]{}

Conhecendo apenas a tensão sobre um resistor e o valor de sua resistência, a
expressão adequada para a potência é:
\begin{alternativas}[2]
\item $P=UI$
\item $P=RI^{2}$
\item $P=U^{2}/R$
\item $P=U/R$
\end{alternativas}
\resposta{(c)}


\end{questao}
\begin{questao}[1]{}

Um dispositivo de $\SI{60}{\watt}$ opera por $\SI{5}{\minute}$. Determine a
energia consumida em joules.
\resposta{$E=\SI{18}{\kilo\joule}$}


\end{questao}

\blocoaplicacao
\begin{questao}[2]{}

Um chuveiro de $\SI{5500}{\watt}$ opera em $\SI{220}{\volt}$. Determine a
corrente e a resistência.
\resposta{$I=\SI{25}{\ampere}$; $R=\SI{8.8}{\ohm}$}


\end{questao}
\begin{questao}[2]{}

Projete a resistência de um aquecedor de $\SI{1200}{\watt}$ para
$\SI{220}{\volt}$ e determine a corrente nominal e o disjuntor adequado.
\resposta{$R=\SI{40.3}{\ohm}$; $I=\SI{5.45}{\ampere}$; disjuntor de
$\SI{10}{\ampere}$}


\end{questao}
\begin{questao}[2]{}

Duas lâmpadas incandescentes de $\SI{60}{\watt}$ são fabricadas para
$\SI{127}{\volt}$ e para $\SI{220}{\volt}$. Determine a resistência de cada
uma.
\resposta{$\SI{268.8}{\ohm}$ e $\SI{806.7}{\ohm}$}


\end{questao}
\begin{questao}[2]{}

Um aparelho opera em dois modos: $\SI{1500}{\watt}$ durante
$\SI{20}{\minute}$ e $\SI{600}{\watt}$ durante $\SI{40}{\minute}$, todos os
dias. Determine o consumo mensal ($\SI{30}{\day}$) e o custo a
$\mathrm{R\$}\,0{,}95$ por $\si{\kilo\watt\hour}$.
\resposta{$\SI{27}{\kilo\watt\hour}$; $\mathrm{R\$}\,25{,}65$}


\end{questao}
\begin{questao}[2]{}

Um resistor é submetido a pulsos de $\SI{50}{\watt}$ com ciclo de trabalho de
$\SI{10}{\percent}$. Determine a potência média.
\resposta{$P_{\text{média}}=\SI{5}{\watt}$}


\end{questao}
\begin{questao}[2]{}

Uma fonte com rendimento de $\SI{90}{\percent}$ alimenta um conversor de
$\SI{85}{\percent}$. Determine o rendimento global e a potência de entrada
necessária para entregar $\SI{100}{\watt}$ úteis.
\resposta{$\eta=76{,}5\%$; $P_{entrada}=\SI{130.7}{\watt}$}


\end{questao}
\begin{questao}[2]{}

Uma carga alterna entre $\SI{500}{\watt}$ por $\SI{5}{\minute}$ e
$\SI{100}{\watt}$ por $\SI{15}{\minute}$, em ciclos de
$\SI{20}{\minute}$. Determine a potência média e a energia por ciclo.
\resposta{$P_{\text{média}}=\SI{200}{\watt}$; $E=\SI{66.7}{\watt\hour}$}


\end{questao}

\blocoaprofundamento
\begin{questao}[3]{}

Um resistor de $\SI{10}{\ohm}$ conduz $\SI{4}{\ampere}$ durante
$\SI{25}{\percent}$ do tempo e nada nos $\SI{75}{\percent}$ restantes.
\begin{itensq}
\item Determine a corrente média e a potência calculada a partir dela.
\item Determine a potência média verdadeira.
\item Explique a discrepância.
\end{itensq}
\resposta{(a) $\SI{1}{\ampere}$ e $\SI{10}{\watt}$ (errado); (b)
$\SI{40}{\watt}$; (c) a potência depende de $\langle i^{2}\rangle$}


\end{questao}
\begin{questao}[3]{}

Um resistor com potência nominal de $\SI{10}{\watt}$ recebe pulsos de
$\SI{50}{\watt}$ com ciclo de trabalho de $\SI{10}{\percent}$.
\begin{itensq}
\item Determine a potência média e compare com o nominal.
\item Estabeleça o critério para decidir se o componente sobrevive.
\item Analise dois casos: pulsos de $\SI{1}{\milli\second}$ e de
      $\SI{10}{\second}$.
\end{itensq}
\resposta{$P_{\text{média}}=\SI{5}{\watt}$; a decisão depende da constante térmica}


\end{questao}
\begin{questao}[3]{}

Um sistema tem uma fonte de $\SI{92}{\percent}$ em cascata com um conversor de
$\SI{85}{\percent}$.
\begin{itensq}
\item Determine o rendimento global.
\item Compare o ganho de melhorar o estágio pior (de $85$ para
      $\SI{92}{\percent}$) com o de melhorar o melhor (de $92$ para
      $\SI{97}{\percent}$).
\item Formule a regra geral.
\end{itensq}
\resposta{(a) $78{,}2\%$; (b) $84{,}6\%$ contra $82{,}5\%$ --- melhorar o pior
rende mais}


\end{questao}
\begin{questao}[3]{}

Um aparelho consome $\SI{4}{\watt}$ em espera, permanentemente.
\begin{itensq}
\item Determine o consumo anual e o custo a
      $\mathrm{R\$}\,0{,}95/\si{\kilo\watt\hour}$.
\item Estime o impacto de dez aparelhos semelhantes.
\item Proponha uma decisão de engenharia.
\end{itensq}
\resposta{$\SI{35}{\kilo\watt\hour}$ e $\mathrm{R\$}\,33{,}29$ por aparelho;
$\mathrm{R\$}\,333$ para dez}


\end{questao}
\begin{questao}[3]{}

Um aparelho projetado para $\SI{220}{\volt}$ e $\SI{1000}{\watt}$ é ligado por
engano em $\SI{127}{\volt}$.
\begin{itensq}
\item Determine a potência resultante, supondo resistência constante.
\item Analise o caso inverso: aparelho de $\SI{127}{\volt}$ em
      $\SI{220}{\volt}$.
\end{itensq}
\resposta{(a) $\SI{333}{\watt}$ ($33\%$); (b) potência \emph{triplicada} ---
destruição}


\end{questao}
\begin{questao}[3]{}

Um resistor de $\SI{4}{\ohm}$ é percorrido por
$i(t)=\SI{5}{\ampere}\cdot e^{-t/\tau}$, com $\tau=\SI{20}{\milli\second}$.
\begin{itensq}
\item Determine a potência inicial.
\item Determine a energia total dissipada.
\item Determine o instante em que metade da energia já foi dissipada.
\end{itensq}
\resposta{(a) $\SI{100}{\watt}$; (b) $\SI{1}{\joule}$; (c)
$t=\SI{6.93}{\milli\second}$}


\end{questao}
\begin{questao}[3]{}

Compare uma lâmpada incandescente de $\SI{60}{\watt}$ com uma de LED de
$\SI{9}{\watt}$ e mesmo fluxo luminoso, usadas $\SI{5}{\hour}$ por dia.
\begin{itensq}
\item Determine o consumo e o custo anual de cada uma, a
      $\mathrm{R\$}\,0{,}95/\si{\kilo\watt\hour}$.
\item Determine o tempo de retorno se a lâmpada de LED custar
      $\mathrm{R\$}\,15{,}00$ a mais.
\end{itensq}
\resposta{(a) $\mathrm{R\$}\,104{,}02$ contra $\mathrm{R\$}\,15{,}60$;
(b) cerca de $2 meses$}


\end{questao}
\begin{questao}[3]{}

Uma instalação alimenta uma carga de $\SI{10}{\kilo\watt}$ por um cabo de
resistência total $\SI{0.2}{\ohm}$, em $\SI{220}{\volt}$.
\begin{itensq}
\item Determine a corrente e a perda no cabo.
\item Determine a fração da energia perdida.
\item Repita para $\SI{440}{\volt}$ e comente.
\end{itensq}
\resposta{(a) $\SI{45.5}{\ampere}$, $\SI{414}{\watt}$; (b) $4{,}1\%$;
(c) $\SI{103}{\watt}$, $1{,}0\%$}


\end{questao}

\blocodesafio
\begin{questao}[4]{}

\textbf{(Demonstração --- potência média e valor eficaz.)} Prove que a potência
média em um resistor depende do valor \emph{eficaz} da corrente.
\begin{itensq}
\item Deduza $P_{\text{média}}=R\,I_{rms}^{2}$ a partir da definição.
\item Mostre que $I_{rms}\ge I_{\text{média}}$, com igualdade apenas para corrente
      constante.
\item Aplique a uma onda quadrada de $\SI{0}{\ampere}$ e
      $\SI{4}{\ampere}$ com ciclo de $\SI{25}{\percent}$.
\end{itensq}
\resposta{$I_{rms}=\SI{2}{\ampere}$ contra $I_{\text{média}}=\SI{1}{\ampere}$;
$P=R I_{rms}^{2}$}


\end{questao}
\begin{questao}[4]{}

\textbf{(Integração --- energia em transitório.)} Um resistor $R$ conduz
$i(t)=I_0e^{-t/\tau}$.
\begin{itensq}
\item Deduza a energia total dissipada.
\item Determine a "potência média equivalente" ao longo de $5\tau$.
\item Compare com a energia de um pulso retangular de mesma amplitude e
      duração $\tau$.
\end{itensq}
\resposta{$E=\dfrac{RI_0^{2}\tau}{2}$; equivale à potência inicial mantida por
$\tau/2$}


\end{questao}
\begin{questao}[4]{}

\textbf{(Projeto de proteção.)} Um resistor de $\SI{100}{\ohm}$ e
$\frac14\,\si{\watt}$ deve ser protegido contra sobrecarga.
\begin{itensq}
\item Determine a tensão e a corrente máximas admissíveis.
\item Projete um resistor limitador em série, para uma fonte de
      $\SI{24}{\volt}$.
\item Verifique a potência do limitador.
\end{itensq}
\resposta{$U_{\max}=\SI{5}{\volt}$, $I_{\max}=\SI{50}{\milli\ampere}$;
limitador de $\SI{380}{\ohm}$}


\end{questao}
\begin{questao}[4]{}

\textbf{(Análise de retorno.)} Dois equipamentos entregam a mesma energia útil
de $\SI{1000}{\kilo\watt\hour}$ por ano: o modelo A tem rendimento de
$\SI{80}{\percent}$ e o B, de $\SI{92}{\percent}$.
\begin{itensq}
\item Determine a energia de entrada de cada um e a diferença.
\item Determine o tempo de retorno se B custar $\mathrm{R\$}\,800$ a mais, com
      tarifa de $\mathrm{R\$}\,0{,}95/\si{\kilo\watt\hour}$.
\item Discuta os fatores que a análise simples ignora.
\end{itensq}
\resposta{(a) $1250$ e $\SI{1087}{\kilo\watt\hour}$, diferença de
$\SI{163}{\kilo\watt\hour}$; (b) cerca de $\SI{5.2}{year}$}


\end{questao}
\begin{questao}[4]{}

\textbf{(Restrição orçamentária.)} Deseja-se limitar a
$\mathrm{R\$}\,90{,}00$ o custo mensal de um aquecedor de
$\SI{1.5}{\kilo\watt}$, com tarifa de
$\mathrm{R\$}\,1{,}00/\si{\kilo\watt\hour}$.
\begin{itensq}
\item Determine o tempo de uso admissível.
\item Determine a potência que permitiria uso de $\SI{4}{\hour}$ diárias dentro
      do mesmo orçamento.
\item Discuta qual das duas soluções é preferível.
\end{itensq}
\resposta{(a) $\SI{60}{\hour}$ por mês, ou $\SI{2}{\hour}$ por dia;
(b) $\SI{750}{\watt}$}


\end{questao}
\begin{questao}[4]{}

\textbf{(Otimização econômica de condutor.)} Um cabo alimenta uma carga de
$\SI{40}{\ampere}$ por $\SI{4000}{\hour}$ ao ano. O custo do cabo é
proporcional à seção, e a resistência é inversamente proporcional a ela.
\begin{itensq}
\item Escreva o custo total como função da seção.
\item Determine a condição de mínimo.
\item Interprete o resultado.
\end{itensq}
\resposta{No ótimo, o custo anualizado do cabo iguala o custo anual da energia
perdida}


\end{questao}
\begin{questao}[4]{}

\textbf{(Balanço energético de um chuveiro.)} Um chuveiro de
$\SI{5500}{\watt}$ aquece água de $\SI{20}{\celsius}$ a
$\SI{40}{\celsius}$.
\begin{itensq}
\item Determine a vazão máxima, supondo rendimento de $\SI{100}{\percent}$.
\item Determine o custo de um banho de $\SI{8}{\minute}$ a
      $\mathrm{R\$}\,0{,}95/\si{\kilo\watt\hour}$.
\item Discuta por que o rendimento é tão alto neste caso.
\end{itensq}
\dados{$c_{\text{água}}=\SI{4180}{\joule\per\kilogram\per\kelvin}$}
\resposta{(a) $\SI{0.0658}{\kilogram\per\second}\approx
\SI{3.9}{\liter\per\minute}$; (b) $\mathrm{R\$}\,0{,}70$}


\end{questao}
\begin{questao}[4]{}

\textbf{(Diagnóstico por consumo.)} A conta de energia de uma residência subiu
de $\SI{180}{\kilo\watt\hour}$ para $\SI{320}{\kilo\watt\hour}$ mensais, sem
mudança de hábitos.
\begin{itensq}
\item Determine a potência contínua equivalente ao acréscimo.
\item Levante hipóteses compatíveis com esse valor.
\item Proponha um procedimento de investigação.
\end{itensq}
\resposta{Acréscimo de $\SI{140}{\kilo\watt\hour}$, equivalente a
$\SI{194}{\watt}$ contínuos}


\end{questao}

\gabarito[2.3 Potência elétrica e energia]

\subsection{Fontes reais e Leis de Kirchhoff}

\blocofixacao
\begin{questao}[1]{}

Um circuito tem duas fontes ideais de $\SI{6}{\volt}$ e $\SI{4}{\volt}$ ligadas em
série e em \textbf{oposição}.
\begin{itensq}
\item Qual é a tensão total fornecida ao circuito?
\item Se a resistência total for $\SI{5}{\ohm}$, qual é a corrente?
\end{itensq}
\resposta{(a) \SI{2}{\volt}; (b) \SI{0.4}{\ampere}}


\end{questao}
\begin{questao}[1]{}

No nó da figura, as correntes $I_1 = \SI{5}{\ampere}$ e $I_2 = \SI{3}{\ampere}$
entram, e $I_3$ sai. Pela lei dos nós de Kirchhoff, $I_3$ vale:

\circuitoq{
\begin{circuitikz}[scale=1,transform shape,line width=0.8pt]
 \coordinate (o) at (0,0);
 \draw[->,Icol,line width=1.1pt] (-2,0.6) -- (o) node[midway,above]{$I_1$};
 \draw[->,Icol,line width=1.1pt] (-2,-0.6) -- (o) node[midway,below]{$I_2$};
 \draw[->,Icol,line width=1.1pt] (o) -- (2,0) node[midway,above]{$I_3$};
 \fill (o) circle (0.06);
\end{circuitikz}}

\resposta{$I_3 = \SI{8}{\ampere}$}


\end{questao}

\blocoaplicacao
\begin{questao}[2]{}

Uma bateria de fem $\EMF = \SI{12}{\volt}$ e resistência interna
$r = \SI{0.5}{\ohm}$ alimenta um resistor externo $R = \SI{5.5}{\ohm}$. Determine
a corrente e a tensão nos terminais da bateria.
\resposta{$I = \SI{2}{\ampere}$; $U = \SI{11}{\volt}$}


\end{questao}
\begin{questao}[2]{}

Uma bateria de fem $\EMF$ e resistência interna $r = \SI{0.4}{\ohm}$ apresenta
tensão de $\SI{12}{\volt}$ em seus terminais quando fornece
$\SI{5}{\ampere}$. Determine a fem e a tensão em vazio (corrente nula).
\resposta{$\EMF = \SI{14}{\volt}$; em vazio, $U = \SI{14}{\volt}$}


\end{questao}
\begin{questao}[2]{}

No nó da figura, são conhecidas quatro das cinco correntes. Determine $I_5$,
indicando se entra ou sai do nó.

\circuitoq{
\begin{tikzpicture}[scale=1,transform shape,line width=0.9pt,>=Stealth]
 \coordinate (o) at (0,0);
 \draw[->,Icol] (-2,0.8) -- (o) node[midway,above]{$\SI{6}{\ampere}$};
 \draw[->,Icol] (-2,-0.8) -- (o) node[midway,below]{$\SI{4}{\ampere}$};
 \draw[->,Icol] (o) -- (2,0.8) node[midway,above]{$\SI{3}{\ampere}$};
 \draw[->,Icol] (o) -- (2,-0.8) node[midway,below]{$\SI{5}{\ampere}$};
 \draw[->,Vcol,thick] (o) -- (0,-1.8) node[midway,right]{$I_5$};
 \fill (o) circle (0.07);
\end{tikzpicture}}{Nó com cinco ramos.}

\resposta{$I_5 = \SI{2}{\ampere}$, saindo do nó}


\end{questao}

\blocoaprofundamento
\begin{questao}[3]{}

No circuito de duas malhas da figura, aplique as leis de Kirchhoff para determinar
as correntes $I_1$, $I_2$ e a corrente $I_3$ no resistor central.

\circuitoq{
\begin{circuitikz}[scale=1,transform shape,line width=0.8pt]
 
 \draw (0,0) to[battery1,l_={$E_1=\SI{18}{\volt}$}] (0,3)
       to[R,l={$R_1=\SI{2}{\ohm}$}] (3,3) coordinate(top);
 \draw (top) to[R,l={$R_2=\SI{6}{\ohm}$},i>_={$I_3$}] (3,0) coordinate(bot) -- (0,0);
 
 \draw (top) to[R,l={$R_3=\SI{6}{\ohm}$}] (6,3)
       to[battery1,l={$E_2=\SI{6}{\volt}$},invert] (6,0) -- (bot);
 \draw[->,Icol,line width=1pt] (0.6,3.4)--(1.5,3.4) node[midway,above,font=\footnotesize]{$I_1$};
 \draw[->,Icol,line width=1pt] (5.4,3.4)--(4.5,3.4) node[midway,above,font=\footnotesize]{$I_2$};
\end{circuitikz}}

\resposta{$I_1 = \SI{3}{\ampere}$, $I_2 = \SI{1}{\ampere}$, $I_3 = \SI{2}{\ampere}$}


\end{questao}
\begin{questao}[3]{}

No circuito de duas malhas da figura, determine as três correntes de ramo
($I_1$, $I_2$, $I_3$) aplicando as leis de Kirchhoff.

\circuitoq{
\begin{circuitikz}[scale=1,transform shape,line width=0.8pt]
 \draw (0,0) to[battery1,l_={$E_1=\SI{12}{\volt}$}] (0,3)
       to[R,l={$R_1=\SI{2}{\ohm}$}] (3,3) coordinate(top);
 \draw (top) to[R,l={$R_2=\SI{4}{\ohm}$},i>_={$I_3$}] (3,0) coordinate(bot) -- (0,0);
 \draw (top) to[R,l={$R_3=\SI{2}{\ohm}$}] (6,3)
       to[battery1,l={$E_2=\SI{6}{\volt}$}] (6,0) -- (bot);
 \draw[->,Icol,line width=1pt] (0.6,3.4)--(1.5,3.4) node[midway,above,font=\footnotesize]{$I_1$};
 \draw[->,Icol,line width=1pt] (5.4,3.4)--(4.5,3.4) node[midway,above,font=\footnotesize]{$I_2$};
\end{circuitikz}}

\resposta{$I_1 = \SI{2.4}{\ampere}$, $I_2 = \SI{-0.6}{\ampere}$, $I_3 = \SI{1.8}{\ampere}$}


\end{questao}
\begin{questao}[3]{}

No circuito da figura, quatro fontes ideais de tensão estão interligadas, e os
terminais $a$, $b$ e $c$ estão \textbf{abertos} (nenhuma carga conectada).
Observe atentamente a \emph{polaridade} indicada dentro de cada fonte.

\circuitoq{
\begin{circuitikz}[scale=1.05,transform shape,line width=0.9pt]
 
 \coordinate (N1) at (0,4.2);      
 \coordinate (N2) at (0,2.2);      
 \draw (N1) -- (0,4.2);
 
 \draw (N2) to[\fonteV,l_={$\SI{5}{\volt}$},invert] (N1);
 
 \draw (N1) -- (0.9,4.2) to[\fonteV,l={$\SI{4}{\volt}$}] (2.9,4.2)
       to[short,-o] (4.0,4.2) node[right,font=\Large]{$a$};
 
 \draw (N2) -- (0,0.2);
 
 \draw (0,1.4) -- (0.9,1.4) to[\fonteV,l_={$\SI{5}{\volt}$},invert] (2.9,1.4)
       to[short,-o] (4.0,1.4) node[right,font=\Large]{$b$};
 \fill (0,1.4) circle (0.07);
 
 \draw (0,0.2) -- (0.9,0.2) to[\fonteV,l_={$\SI{3}{\volt}$}] (2.9,0.2)
       to[short,-o] (4.0,0.2) node[right,font=\Large]{$c$};
\end{circuitikz}}

\begin{itensq}
\item Determine $V_{ac}$.
\item Determine também $V_{ab}$ e $V_{bc}$, e verifique a consistência
      ($V_{ab}+V_{bc}=V_{ac}$).
\item Por que os valores \emph{não} dependem das resistências dos fios?
\end{itensq}
\resposta{(a) $V_{ac}=\SI{6}{\volt}$; (b) $V_{ab}=\SI{14}{\volt}$,
$V_{bc}=\SI{-8}{\volt}$; (c) ver comentário}


\end{questao}

\blocodesafio
\begin{questao}[4]{}

Duas baterias de fem $\EMF_1 = \SI{12}{\volt}$ ($r_1=\SI{0.2}{\ohm}$) e
$\EMF_2 = \SI{12.6}{\volt}$ ($r_2=\SI{0.1}{\ohm}$) são ligadas em
\textbf{paralelo} para alimentar uma carga $R_L = \SI{2}{\ohm}$ --- situação
comum em bancos de baterias e em sistemas com gerador e bateria em paralelo.

\circuitoq{
\begin{circuitikz}[scale=0.95,transform shape,line width=0.8pt]
 \draw (0,0) to[battery1,l_={$\EMF_1$}] (0,1.6) to[R,l_={$r_1$},color=Vcol] (0,3);
 \draw (2.4,0) to[battery1,l_={$\EMF_2$}] (2.4,1.6) to[R,l_={$r_2$},color=Vcol] (2.4,3);
 \draw (0,3) -- (2.4,3) -- (5,3) to[R,l={$R_L=\SI{2}{\ohm}$}] (5,0) -- (2.4,0) -- (0,0);
 \node[above,Vcol] at (3.6,3.1) {$V$};
 \fill (2.4,3) circle (0.07); \fill (2.4,0) circle (0.07);
\end{circuitikz}}

\begin{itensq}
\item Determine a tensão $V$ sobre a carga (use análise nodal).
\item Calcule a corrente fornecida por \emph{cada} bateria.
\item Interprete o sinal obtido e explique o risco prático dessa associação.
\end{itensq}
\resposta{(a) $V = \SI{12}{\volt}$; (b) $I_1 = 0$, $I_2 = \SI{6}{\ampere}$;
(c) ver comentário}


\end{questao}

\blocofixacao
\begin{questao}[1]{}

Uma fonte de $\EMF=\SI{6}{\volt}$ e $r=\SI{1}{\ohm}$ fornece
$\SI{1}{\ampere}$. Determine a tensão em seus terminais.
\resposta{$V=\SI{5}{\volt}$}


\end{questao}
\begin{questao}[1]{}

Uma fonte de $\EMF=\SI{9}{\volt}$ apresenta $\SI{8.4}{\volt}$ nos terminais
quando fornece $\SI{2}{\ampere}$. Determine sua resistência interna.
\resposta{$r=\SI{0.3}{\ohm}$}


\end{questao}
\begin{questao}[1]{}

Em um nó entram $\SI{2.5}{\ampere}$ e $\SI{1.5}{\ampere}$, e sai uma única
corrente. Determine seu valor.
\resposta{$I=\SI{4}{\ampere}$}


\end{questao}
\begin{questao}[1]{}

Uma malha simples contém uma fonte ideal de $\SI{12}{\volt}$ e resistores de
$\SI{4}{\ohm}$ e $\SI{8}{\ohm}$. Determine a corrente pela lei das malhas.
\resposta{$I=\SI{1}{\ampere}$}


\end{questao}
\begin{questao}[1]{}

Sobre a diferença entre fonte ideal e fonte real, é correto afirmar:
\begin{alternativas}[2]
\item a ideal mantém a tensão constante para qualquer corrente
\item a real mantém a tensão constante para qualquer corrente
\item ambas mantêm a corrente constante
\item a ideal tem resistência interna maior
\end{alternativas}
\resposta{(a)}


\end{questao}
\begin{questao}[1]{}

Uma pilha de $\EMF=\SI{1.5}{\volt}$ tem $r=\SI{0.25}{\ohm}$. Determine sua
corrente de curto-circuito.
\resposta{$I_{cc}=\SI{6}{\ampere}$}


\end{questao}

\blocoaplicacao
\begin{questao}[2]{}

Uma bateria de $\EMF=\SI{12}{\volt}$ e $r=\SI{0.6}{\ohm}$ alimenta
$R=\SI{5.4}{\ohm}$.
\begin{itensq}
\item Determine a corrente e a tensão terminal.
\item Determine as potências e o rendimento.
\end{itensq}
\resposta{$I=\SI{2}{\ampere}$, $V=\SI{10.8}{\volt}$;
$P_R=\SI{21.6}{\watt}$, $P_r=\SI{2.4}{\watt}$, $\eta=90\%$}


\end{questao}
\begin{questao}[2]{}

Um carregador de $\SI{14}{\volt}$ é ligado a uma bateria de
$\SI{12}{\volt}$, com resistência total de $\SI{0.4}{\ohm}$ no circuito.
\begin{itensq}
\item Determine a corrente de carga.
\item Determine a tensão nos terminais da bateria.
\end{itensq}
\resposta{$I=\SI{5}{\ampere}$; $V=\SI{14}{\volt}$}


\end{questao}
\begin{questao}[2]{}

Em um nó, entram $\SI{3}{\ampere}$, $\SI{4}{\ampere}$ e $\SI{2}{\ampere}$;
sai $\SI{6}{\ampere}$ por um ramo e $I_5$ por outro. Determine $I_5$.
\resposta{$I_5=\SI{3}{\ampere}$, saindo}


\end{questao}
\begin{questao}[2]{}

Uma malha contém fontes ideais de $\SI{24}{\volt}$ e $\SI{6}{\volt}$ em
oposição, com resistores somando $\SI{6}{\ohm}$. Determine a corrente e
identifique qual fonte absorve energia.
\resposta{$I=\SI{3}{\ampere}$; a de $\SI{6}{\volt}$ absorve}


\end{questao}
\begin{questao}[2]{}

Duas pilhas de $\SI{1.5}{\volt}$ e $r=\SI{0.2}{\ohm}$ cada são associadas.
Determine a f.e.m. e a resistência interna equivalentes em série e em paralelo.
\resposta{Série: $\SI{3}{\volt}$ e $\SI{0.4}{\ohm}$; paralelo:
$\SI{1.5}{\volt}$ e $\SI{0.1}{\ohm}$}


\end{questao}
\begin{questao}[2]{}

Determine o número de equações independentes de LKC e de LKT para um circuito
com $n=5$ nós e $b=8$ ramos.
\resposta{$4$ de LKC e $4$ de LKT, totalizando $8$}


\end{questao}
\begin{questao}[2]{}

Uma fonte de $\EMF=\SI{12}{\volt}$ e $r=\SI{0.5}{\ohm}$ alimenta duas cargas
em paralelo, de $\SI{12}{\ohm}$ e $\SI{6}{\ohm}$. Determine a tensão terminal e
a corrente de cada carga.
\resposta{$V=\SI{10.67}{\volt}$; $\SI{0.89}{\ampere}$ e
$\SI{1.78}{\ampere}$}


\end{questao}
\begin{questao}[2]{}

Uma bateria em uso apresenta $\EMF=\SI{1.5}{\volt}$ e $r=\SI{0.3}{\ohm}$
quando nova, e $\EMF=\SI{1.3}{\volt}$ com $r=\SI{1.2}{\ohm}$ ao final da vida.
Compare a tensão entregue a uma carga de $\SI{3}{\ohm}$.
\resposta{$\SI{1.36}{\volt}$ contra $\SI{0.93}{\volt}$ --- queda de $32\%$}


\end{questao}
\begin{questao}[2]{}

Verifique o balanço de potência de um circuito com fonte de
$\EMF=\SI{20}{\volt}$, $r=\SI{1}{\ohm}$ e carga de $\SI{9}{\ohm}$.
\resposta{Fornecido $\SI{40}{\watt}$; dissipado $4+\SI{36}{\watt}$}


\end{questao}

\blocoaprofundamento
\begin{questao}[3]{}

Duas fontes reais alimentam a mesma carga: $\EMF_1=\SI{12}{\volt}$ com
$r_1=\SI{1}{\ohm}$ e $\EMF_2=\SI{6}{\volt}$ com $r_2=\SI{1}{\ohm}$, ambas
ligadas a $R=\SI{4}{\ohm}$.
\begin{itensq}
\item Determine a tensão comum e as três correntes.
\item Identifique o comportamento de cada fonte.
\end{itensq}
\resposta{$V=\SI{8}{\volt}$; $I_1=\SI{4}{\ampere}$,
$I_2=\SI{-2}{\ampere}$, $I_R=\SI{2}{\ampere}$ --- a fonte 2 é
\textbf{carregada}}


\end{questao}
\begin{questao}[3]{}

Um banco é formado por quatro pilhas de $\SI{1.5}{\volt}$ e
$r=\SI{0.2}{\ohm}$.
\begin{itensq}
\item Compare série e paralelo quanto a $\EMF$, $r$ e corrente de
      curto.
\item Determine a potência máxima disponível em cada arranjo.
\item Estabeleça o critério de escolha.
\end{itensq}
\resposta{Série: $\SI{6}{\volt}$/$\SI{0.8}{\ohm}$; paralelo:
$\SI{1.5}{\volt}$/$\SI{0.05}{\ohm}$; $P_{\max}=\SI{11.25}{\watt}$ nos
\textbf{dois}}


\end{questao}
\begin{questao}[3]{}

Aplique a LKC a uma \textbf{superfície fechada} que envolva dois nós de um
circuito, e explique por que o resultado é válido.
\begin{itensq}
\item Enuncie a forma generalizada.
\item Justifique a partir da conservação de carga.
\item Dê um exemplo de uso.
\end{itensq}
\resposta{A soma das correntes que atravessam qualquer superfície fechada é
nula}


\end{questao}
\begin{questao}[3]{}

Uma fonte foi ensaiada sob diferentes cargas, obtendo-se:
\begin{center}
\begin{tabular}{cccc}
\hline
$I$ ($\si{\ampere}$) & $0{,}50$ & $1{,}00$ & $2{,}00$\\
$V$ ($\si{\volt}$) & $11{,}75$ & $11{,}50$ & $11{,}00$\\
\hline
\end{tabular}
\end{center}
\begin{itensq}
\item Determine $\EMF$ e $r$.
\item Verifique a linearidade do modelo.
\item Estime a corrente de curto e discuta a extrapolação.
\end{itensq}
\resposta{$\EMF=\SI{12}{\volt}$, $r=\SI{0.5}{\ohm}$;
$I_{cc}=\SI{24}{\ampere}$}


\end{questao}
\begin{questao}[3]{}

Um circuito tem $n=6$ nós, $b=9$ ramos e $2$ fontes de corrente ideais.
\begin{itensq}
\item Determine o número de equações de Kirchhoff necessárias.
\item Explique como as fontes de corrente reduzem o trabalho.
\item Compare com a contagem para análise nodal.
\end{itensq}
\resposta{$5$ LKC $+$ $4$ LKT $=9$; as fontes de corrente fixam duas correntes
de ramo, restando $7$ incógnitas}


\end{questao}
\begin{questao}[3]{}

Uma fonte tem resistência interna que \emph{cresce} com a corrente, segundo
$r=0{,}2+0{,}1\,I$ (com $r$ em $\si{\ohm}$ e $I$ em $\si{\ampere}$), com
$\EMF=\SI{6}{\volt}$.
\begin{itensq}
\item Determine a corrente para uma carga de $\SI{1}{\ohm}$.
\item Compare com a previsão do modelo linear com $r=\SI{0.2}{\ohm}$.
\end{itensq}
\resposta{$I=\SI{3.80}{\ampere}$ contra $\SI{5}{\ampere}$ previstos}


\end{questao}
\begin{questao}[3]{}

Uma bateria de $\EMF=\SI{12}{\volt}$ e $r=\SI{0.5}{\ohm}$ deve entregar tensão
terminal de pelo menos $\SI{11.4}{\volt}$.
\begin{itensq}
\item Determine a corrente máxima admissível.
\item Determine a menor carga que pode ser conectada.
\item Determine a potência entregue nessa condição.
\end{itensq}
\resposta{$I\le\SI{1.2}{\ampere}$; $R\ge\SI{9.5}{\ohm}$;
$P=\SI{13.68}{\watt}$}


\end{questao}
\begin{questao}[3]{}

Um circuito contém uma fonte \textbf{dependente} de tensão de valor
$3\,I_x$ (em volts, com $I_x$ em ampères) em série com $\SI{2}{\ohm}$, além de
uma fonte independente de $\SI{20}{\volt}$ e um resistor de $\SI{8}{\ohm}$, em
malha única. A corrente de controle $I_x$ é a própria corrente da malha.
\begin{itensq}
\item Escreva a LKT e resolva.
\item Explique por que Kirchhoff continua válido.
\end{itensq}
\resposta{$I=\SI{1.54}{\ampere}$ (fonte dependente em oposição)}


\end{questao}
\begin{questao}[3]{}

Uma fonte alimenta uma carga através de um cabo de $\SI{0.3}{\ohm}$ (ida e
volta). A fonte tem $\EMF=\SI{24}{\volt}$ e $r=\SI{0.2}{\ohm}$; a carga vale
$\SI{9.5}{\ohm}$.
\begin{itensq}
\item Determine a corrente, a tensão na carga e a tensão nos terminais da
      fonte.
\item Determine onde a energia é perdida.
\end{itensq}
\resposta{$I=\SI{2.4}{\ampere}$; $V_{carga}=\SI{22.8}{\volt}$;
$V_{fonte}=\SI{23.52}{\volt}$}


\end{questao}

\blocodesafio
\begin{questao}[4]{}

\textbf{(Ensaio de caracterização.)} Descreva como obter $\EMF$ e $r$ de uma
fonte por ensaios em vazio e em curto-circuito.
\begin{itensq}
\item Deduza as relações.
\item Analise os riscos do ensaio em curto.
\item Proponha uma alternativa segura e compare.
\end{itensq}
\resposta{$\EMF=V_{oc}$ e $r=V_{oc}/I_{cc}$; o curto é destrutivo em fontes de
baixa impedância}


\end{questao}
\begin{questao}[4]{}

\textbf{(Demonstração --- balanço de potência.)} Prove que a soma algébrica das
potências de todos os elementos de um circuito é nula.
\begin{itensq}
\item Estabeleça a convenção de sinais.
\item Demonstre a partir das leis de Kirchhoff.
\item Explique por que o resultado independe da natureza dos elementos.
\end{itensq}
\resposta{$\sum_k v_ki_k=0$ --- consequência direta de LKC e LKT}


\end{questao}
\begin{questao}[4]{}

\textbf{(Formulação com fonte ideal entre dois nós.)} Um circuito tem dois nós
não referenciais, $A$ e $B$, com uma fonte ideal de tensão $\EMF$ ligada
diretamente entre eles.
\begin{itensq}
\item Explique por que a LKC não pode ser escrita isoladamente em $A$ nem em
      $B$.
\item Formule o sistema corretamente.
\item Determine a corrente na fonte após resolver.
\end{itensq}
\resposta{Usar supernó: uma LKC envolvendo $A$ e $B$, mais a restrição
$V_A-V_B=\EMF$}


\end{questao}
\begin{questao}[4]{}

\textbf{(Projeto de verificação experimental.)} Elabore um ensaio de laboratório
para verificar a LKC em um nó com três ramos, considerando incertezas dos
instrumentos.
\begin{itensq}
\item Descreva o procedimento.
\item Propague as incertezas e estabeleça o critério de aceitação.
\item Interprete um resíduo de $\SI{0.35}{\ampere}$ em correntes de alguns
      ampères.
\end{itensq}
\resposta{Critério: resíduo compatível com $u_c\approx\sqrt{\sum u_k^{2}}$;
$\SI{0.35}{\ampere}$ indica erro real}


\end{questao}
\begin{questao}[4]{}

\textbf{(Baterias em série desequilibradas.)} Quatro células de
$\SI{1.5}{\volt}$ estão em série alimentando uma carga, mas uma delas está
gasta, com $\EMF=\SI{0.9}{\volt}$ e $r=\SI{1.5}{\ohm}$ (as demais mantêm
$r=\SI{0.2}{\ohm}$). A carga vale $\SI{5}{\ohm}$.
\begin{itensq}
\item Determine a corrente e a tensão em cada célula.
\item Identifique o que ocorre com a célula gasta.
\item Discuta a implicação prática.
\end{itensq}
\resposta{$I=\SI{0.76}{\ampere}$; a célula gasta apresenta
$\SI{-0.24}{\volt}$ --- está sendo \textbf{invertida}}


\end{questao}
\begin{questao}[4]{}

\textbf{(Modelagem --- fonte com limitação de corrente.)} Uma fonte de bancada
tem $\EMF=\SI{20}{\volt}$ e $r=\SI{0.1}{\ohm}$, com limitação eletrônica em
$\SI{2}{\ampere}$.
\begin{itensq}
\item Descreva a característica $V\times I$ completa.
\item Determine o ponto de operação para cargas de $\SI{50}{\ohm}$,
      $\SI{10}{\ohm}$ e $\SI{2}{\ohm}$.
\item Explique por que o modelo de fonte real deixa de valer.
\end{itensq}
\resposta{Duas regiões: fonte de tensão até $\SI{2}{\ampere}$, fonte de corrente
depois}


\end{questao}
\begin{questao}[4]{}

\textbf{(Sistema completo por Kirchhoff.)} Um circuito de duas malhas tem
$\EMF_1=\SI{18}{\volt}$ com $r_1=\SI{1}{\ohm}$ no primeiro ramo,
$\EMF_2=\SI{12}{\volt}$ com $r_2=\SI{2}{\ohm}$ no segundo, e
$R=\SI{6}{\ohm}$ no ramo comum.
\begin{itensq}
\item Escreva as equações de Kirchhoff e resolva.
\item Verifique o balanço de potência.
\item Identifique o papel de cada fonte.
\end{itensq}
\resposta{$I_1=\SI{3.6}{\ampere}$, $I_2=\SI{-1.2}{\ampere}$,
$I_R=\SI{2.4}{\ampere}$}


\end{questao}

\gabarito[2.4 Fontes reais e Leis de Kirchhoff]

\section{Associação de Resistores e Transformações}

\subsection{Circuitos em série}

\blocofixacao
\begin{questao}[1]{}

Em uma associação de resistores \textbf{em série}, é correto afirmar que:
\begin{alternativas}
\item a tensão é a mesma em todos os resistores.
\item a corrente é a mesma em todos os resistores.
\item a resistência equivalente é menor que a menor das resistências.
\item a corrente se divide proporcionalmente às resistências.
\item a potência dissipada é a mesma em todos os resistores.
\end{alternativas}
\resposta{(b)}


\end{questao}
\begin{questao}[1]{}

Três resistores de $\SI{4}{\ohm}$, $\SI{6}{\ohm}$ e $\SI{10}{\ohm}$ são ligados em
série a uma fonte de $\SI{40}{\volt}$ de resistência interna desprezível.
Determine:
\begin{itensq}
\item a resistência equivalente;
\item a corrente do circuito;
\item a tensão em cada resistor.
\end{itensq}
\resposta{(a) \SI{20}{\ohm}; (b) \SI{2}{\ampere}; (c) \SI{8}{\volt}, \SI{12}{\volt} e \SI{20}{\volt}}


\end{questao}
\begin{questao}[1]{}

Para cada circuito, calcule a resistência equivalente e a corrente $I$.

\circuitoq{
\subcirc{a}{
\begin{circuitikz}[scale=0.78,transform shape,line width=0.8pt]
 \draw (0,0) to[battery1,l_={$E=\SI{10}{\volt}$}] (0,3)
       to[R,l={$R_1=\SI{3}{\ohm}$}] (3,3)
       to[R,l={$R_2=\SI{2}{\ohm}$}] (3,0) -- (0,0);
 \draw[->,Icol,line width=1pt] (0.7,3.4) -- (1.7,3.4)
       node[midway,above,font=\footnotesize]{$I$};
\end{circuitikz}}
\hfill
\subcirc{b}{
\begin{circuitikz}[scale=0.78,transform shape,line width=0.8pt]
 \draw (0,0) to[battery1,l_={$E=\SI{15}{\volt}$}] (0,3)
       to[R,l={$R_1=\SI{4}{\ohm}$}] (3,3)
       to[R,l={$R_2=\SI{6}{\ohm}$}] (3,0) -- (0,0);
 \draw[->,Icol,line width=1pt] (0.7,3.4) -- (1.7,3.4)
       node[midway,above,font=\footnotesize]{$I$};
\end{circuitikz}}

\vspace{1.1em}

\subcirc{c}{
\begin{circuitikz}[scale=0.78,transform shape,line width=0.8pt]
 \draw (0,0) to[battery1,l_={$E=\SI{25}{\volt}$}] (0,3)
       to[R,l={$R_1=\SI{4.5}{\ohm}$}] (3.4,3)
       to[R,l={$R_2=\SI{3}{\ohm}$}] (3.4,0)
       to[R,l={$R_3=\SI{2.5}{\ohm}$}] (0,0);
 \draw[->,Icol,line width=1pt] (0.7,3.4) -- (1.7,3.4)
       node[midway,above,font=\footnotesize]{$I$};
\end{circuitikz}}
\hfill
\subcirc{d}{
\begin{circuitikz}[scale=0.78,transform shape,line width=0.8pt]
 \draw (0,0) to[battery1,l_={$E=\SI{30}{\volt}$}] (0,4.4)
       to[R,l={$R_1=\SI{2.5}{\ohm}$}] (3.6,4.4)
       to[R,l={$R_2=\SI{4}{\ohm}$}] (3.6,2.2)
       to[R,l={$R_3=\SI{1}{\ohm}$}] (3.6,0)
       to[R,l={$R_4=\SI{2.5}{\ohm}$}] (0,0);
 \draw[->,Icol,line width=1pt] (0.7,4.8) -- (1.7,4.8)
       node[midway,above,font=\footnotesize]{$I$};
\end{circuitikz}}}

\resposta{(a) \SI{5}{\ohm}, \SI{2}{\ampere}; (b) \SI{10}{\ohm}, \SI{1.5}{\ampere};
(c) \SI{10}{\ohm}, \SI{2.5}{\ampere}; (d) \SI{10}{\ohm}, \SI{3}{\ampere}}


\end{questao}
\begin{questao}[1]{}

Uma associação em série de $n$ resistores iguais de $\SI{15}{\ohm}$ apresenta
resistência equivalente de $\SI{105}{\ohm}$. O valor de $n$ é:
\begin{alternativas}[3]
\item 3
\item 5
\item 6
\item 7
\item 9
\end{alternativas}
\resposta{(d)}


\end{questao}

\blocoaplicacao
\begin{questao}[2]{}

No circuito da figura, duas fontes ideais estão ligadas em \textbf{oposição}.

\circuitoq{
\begin{circuitikz}[scale=0.9,transform shape,line width=0.8pt]
 \draw (0,0) to[battery1,l_={$E_1=\SI{30}{\volt}$}] (0,3.6)
       to[R,l={$R_1=\SI{2.5}{\ohm}$}] (4.5,3.6)
       to[battery1,l={$E_2=\SI{10}{\volt}$},invert] (4.5,1.8)
       to[R,l={$R_2=\SI{1.5}{\ohm}$}] (4.5,0)
       to[R,l={$R_3=\SI{1}{\ohm}$}] (0,0);
 \draw[->,Icol,line width=1pt] (1.0,4.0) -- (2.2,4.0)
       node[midway,above,font=\footnotesize]{$I$};
\end{circuitikz}}

Determine a corrente do circuito e a tensão sobre $R_1$.
\resposta{$I = \SI{4}{\ampere}$; $U_{R_1} = \SI{10}{\volt}$}


\end{questao}
\begin{questao}[2]{}

Uma bateria de força eletromotriz $\EMF = \SI{12}{\volt}$ e resistência interna
$r = \SI{0.5}{\ohm}$ alimenta um resistor $R = \SI{5.5}{\ohm}$.

\circuitoq{
\begin{circuitikz}[scale=0.9,transform shape,line width=0.8pt]
 \draw (0,0) to[battery1,l_={$\EMF=\SI{12}{\volt}$}] (0,2)
       to[R,l_={$r=\SI{0.5}{\ohm}$},color=Vcol] (0,3.4) -- (3.6,3.4)
       to[R,l={$R=\SI{5.5}{\ohm}$}] (3.6,0) -- (0,0);
 \draw[dashed,cinza] (-1.15,-0.5) rectangle (0.95,3.9);
 \node[cinza,font=\footnotesize\sffamily,above] at (-0.1,3.95) {bateria real};
 \draw[->,Icol,line width=1pt] (1.2,3.75) -- (2.3,3.75)
       node[midway,above,font=\footnotesize]{$I$};
\end{circuitikz}}

Determine:
\begin{itensq}
\item a corrente do circuito;
\item a tensão nos terminais da bateria;
\item a potência dissipada internamente.
\end{itensq}
\resposta{(a) \SI{2}{\ampere}; (b) \SI{11}{\volt}; (c) \SI{2}{\watt}}


\end{questao}
\begin{questao}[2]{}

Seis lâmpadas idênticas, de resistência $\SI{40}{\ohm}$ cada, são ligadas
\textbf{em série} a uma fonte de $\SI{120}{\volt}$.
\begin{itensq}
\item Calcule a corrente e a tensão em cada lâmpada.
\item Uma das lâmpadas queima (filamento aberto). O que acontece com as demais?
      Justifique.
\end{itensq}
\resposta{(a) $I=\SI{0.5}{\ampere}$ e $U=\SI{20}{\volt}$ por lâmpada;
(b) todas apagam}


\end{questao}
\begin{questao}[2]{}

Três resistores estão em série sob $\SI{120}{\volt}$: $R_1=\SI{20}{\ohm}$,
$R_2$ desconhecido e $R_3=\SI{30}{\ohm}$. Deseja-se que a tensão sobre $R_2$
seja de $\SI{20}{\volt}$. O valor de $R_2$ deve ser:
\begin{alternativas}[3]
\item \SI{5}{\ohm}
\item \SI{10}{\ohm}
\item \SI{15}{\ohm}
\item \SI{20}{\ohm}
\item \SI{25}{\ohm}
\end{alternativas}
\resposta{(b)}


\end{questao}

\blocoaprofundamento
\begin{questao}[3]{}

Uma carga de $\SI{9.6}{\ohm}$ é alimentada por uma fonte de $\SI{24}{\volt}$ por
meio de um cabo de dois condutores, cada um com resistência de
$\SI{0.2}{\ohm}$.

\circuitoq{
\begin{circuitikz}[scale=0.9,transform shape,line width=0.8pt]
 \draw (0,0) to[battery1,l_={$E=\SI{24}{\volt}$}] (0,3)
       to[R,l=$\SI{0.2}{\ohm}$,color=Vcol] (4.4,3)
       to[R,l={$R_c=\SI{9.6}{\ohm}$}] (4.4,0)
       to[R,l=$\SI{0.2}{\ohm}$,color=Vcol] (0,0);
 \node[Vcol,font=\footnotesize\sffamily] at (2.2,3.85) {condutores do cabo};
 \draw[->,Icol,line width=1pt] (0.5,2.55) -- (1.5,2.55)
       node[midway,below,font=\footnotesize]{$I$};
\end{circuitikz}}

Determine:
\begin{itensq}
\item a corrente na linha;
\item a queda de tensão total no cabo;
\item a tensão efetivamente entregue à carga;
\item o rendimento da transmissão, $\eta = P_{\text{carga}}/P_{\text{total}}$.
\end{itensq}
\resposta{(a) \SI{2.4}{\ampere}; (b) \SI{0.96}{\volt}; (c) \SI{23.04}{\volt};
(d) \SI{96}{\percent}}


\end{questao}
\begin{questao}[3]{}

Dispõe-se de resistores de $\SI{100}{\ohm}$, todos idênticos, e de uma fonte
ideal de $\SI{50}{\volt}$. Deseja-se obter, a partir dessa fonte, uma saída de
exatamente $\SI{20}{\volt}$ em vazio (sem carga conectada à saída), usando apenas
uma associação série desses resistores como divisor de tensão.
\begin{itensq}
\item Qual é o menor número de resistores capaz de fornecer essa tensão?
\item Nessa configuração, quanto vale a corrente que circula pelo divisor?
\end{itensq}
\resposta{(a) 5 resistores (saída sobre 2 deles); (b) \SI{0.1}{\ampere}}


\end{questao}
\begin{questao}[3]{}

Uma linha de transmissão de $\SI{2}{\kilo\meter}$ alimenta uma carga de
$\SI{10}{\ohm}$ a partir de uma fonte de $\SI{240}{\volt}$. O cabo é de cobre
($\rho=\SI{1.7e-8}{\ohm\meter}$) e deve-se escolher sua seção $A$ de modo que a
queda de tensão na linha não ultrapasse $\SI{4}{\percent}$ da tensão da fonte.
\begin{itensq}
\item Determine a resistência máxima admissível para o par de condutores.
\item Calcule a seção mínima do cabo.
\item Determine o rendimento da transmissão nessa condição e a potência perdida.
\end{itensq}
\resposta{(a) $R_{\text{linha}} \le \SI{0.417}{\ohm}$;
(b) $A \ge \SI{163}{\milli\meter\squared}$; (c) $\eta = \SI{96}{\percent}$,
$P_{\text{perdida}}\approx\SI{221}{\watt}$}


\end{questao}

\blocofixacao
\begin{questao}[1]{}

Determine a resistência equivalente de $R_1=\SI{1.2}{\kilo\ohm}$,
$R_2=\SI{470}{\ohm}$ e $R_3=\SI{2.2}{\kilo\ohm}$ em série.
\resposta{$R_{eq}=\SI{3.87}{\kilo\ohm}$}


\end{questao}
\begin{questao}[1]{}

Uma associação série de três resistores tem $R_{eq}=\SI{75}{\ohm}$. Sabendo que
$R_1=\SI{22}{\ohm}$ e $R_2=\SI{33}{\ohm}$, determine $R_3$.
\resposta{$R_3=\SI{20}{\ohm}$}


\end{questao}
\begin{questao}[1]{}

Uma fonte ideal de $\SI{24}{\volt}$ alimenta $R_1=\SI{4}{\ohm}$ em série com
$R_2=\SI{8}{\ohm}$. Determine a corrente do circuito.
\resposta{$I=\SI{2}{\ampere}$}


\end{questao}
\begin{questao}[1]{}

Em uma cadeia série percorrida por $\SI{0.5}{\ampere}$, determine a queda de
tensão em um resistor de $\SI{68}{\ohm}$.
\resposta{$U=\SI{34}{\volt}$}


\end{questao}
\begin{questao}[1]{}

Acrescentando-se mais um resistor \textbf{em série} a um circuito alimentado
por fonte ideal, a corrente total:
\begin{alternativas}[2]
\item aumenta
\item diminui
\item não se altera
\item pode aumentar ou diminuir, conforme o valor
\end{alternativas}
\resposta{(b)}


\end{questao}
\begin{questao}[1]{}

Dois resistores de $\SI{15}{\ohm}$ e $\SI{45}{\ohm}$ são ligados em série a uma
fonte. Trocando-se a \textbf{ordem} deles no circuito, o que muda?
\begin{alternativas}[2]
\item a resistência equivalente
\item a corrente do circuito
\item a queda em cada resistor
\item nada
\end{alternativas}
\resposta{(d)}


\end{questao}

\blocoaplicacao
\begin{questao}[2]{}

$R_1=\SI{5}{\ohm}$, $R_2=\SI{10}{\ohm}$ e $R_3=\SI{15}{\ohm}$ estão em série
sob $\SI{60}{\volt}$. Determine a corrente e as três quedas de tensão.
\resposta{$I=\SI{2}{\ampere}$; $U_1=\SI{10}{\volt}$, $U_2=\SI{20}{\volt}$,
$U_3=\SI{30}{\volt}$}


\end{questao}
\begin{questao}[2]{}

No circuito da questão anterior, calcule a potência dissipada em cada resistor
e confronte com a potência fornecida pela fonte.
\resposta{$\SI{20}{\watt}$, $\SI{40}{\watt}$, $\SI{60}{\watt}$; total
$\SI{120}{\watt}$}


\end{questao}
\begin{questao}[2]{}

Em uma cadeia série de três resistores, mediu-se $\SI{18}{\volt}$ sobre o
resistor de $\SI{90}{\ohm}$. Os outros dois valem $\SI{30}{\ohm}$ e
$\SI{60}{\ohm}$. Determine a tensão da fonte.
\resposta{$V=\SI{36}{\volt}$}


\end{questao}
\begin{questao}[2]{}

Duas lâmpadas incandescentes, de resistências $\SI{240}{\ohm}$ e
$\SI{960}{\ohm}$, são ligadas \textbf{em série} a $\SI{120}{\volt}$. Qual
brilha mais e por quê?
\resposta{A de $\SI{960}{\ohm}$; $P=\SI{9.6}{\watt}$ contra $\SI{2.4}{\watt}$}


\end{questao}
\begin{questao}[2]{}

Deseja-se limitar a $\SI{150}{\milli\ampere}$ a corrente de uma carga de
$\SI{60}{\ohm}$ alimentada por $\SI{24}{\volt}$. Determine o resistor em série
necessário e sua potência.
\resposta{$R_s=\SI{100}{\ohm}$; $P_s=\SI{2.25}{\watt}$}


\end{questao}
\begin{questao}[2]{}

Duas montagens têm a mesma $R_{eq}=\SI{100}{\ohm}$ sob $\SI{50}{\volt}$:
(A) dois resistores de $\SI{50}{\ohm}$; (B) um de $\SI{10}{\ohm}$ e outro de
$\SI{90}{\ohm}$. Compare corrente total, potência total e a potência do
resistor mais solicitado.
\resposta{Corrente e potência totais iguais ($\SI{0.5}{\ampere}$,
$\SI{25}{\watt}$); resistor mais solicitado: $\SI{12.5}{\watt}$ em (A) e
$\SI{22.5}{\watt}$ em (B)}


\end{questao}
\begin{questao}[2]{}

Uma cadeia série de $\SI{48}{\ohm}$ é alimentada por $\SI{12}{\volt}$ e
protegida por fusível. Determine a corrente normal de operação e justifique a
escolha de um fusível de $\SI{500}{\milli\ampere}$.
\resposta{$I=\SI{250}{\milli\ampere}$; o fusível deve ficar acima da corrente
normal e abaixo da corrente de falha}


\end{questao}
\begin{questao}[2]{}

Um reostato de $\SI{0}{\ohm}$ a $\SI{100}{\ohm}$ está em série com uma carga de
$\SI{50}{\ohm}$, sob $\SI{30}{\volt}$. Determine a faixa de variação da
corrente.
\resposta{de $\SI{0.2}{\ampere}$ a $\SI{0.6}{\ampere}$}


\end{questao}
\begin{questao}[2]{}

Em uma cadeia série sob tensão fixa, se um dos resistores tiver seu valor
\textbf{dobrado}, a queda sobre ele:
\begin{alternativas}[2]
\item dobra
\item aumenta, mas menos que o dobro
\item não se altera
\item diminui
\end{alternativas}
\resposta{(b)}


\end{questao}
\begin{questao}[2]{}

Uma cadeia série dissipa $\SI{300}{\watt}$ e opera $\SI{6}{\hour}$ por dia
durante $\SI{30}{\day}$. Calcule a energia consumida e o custo, a
$\mathrm{R\$}\,0{,}80$ por $\si{\kilo\watt\hour}$.
\resposta{$E=\SI{54}{\kilo\watt\hour}$; custo $=\mathrm{R\$}\,43{,}20$}


\end{questao}
\begin{questao}[2]{}

Em uma cadeia série sob $\SI{100}{\volt}$ mediram-se as quedas
$\SI{22}{\volt}$, $\SI{35}{\volt}$ e $\SI{41}{\volt}$. A medição é consistente?
Se a corrente for $\SI{50}{\milli\ampere}$, determine os três resistores.
\resposta{Sim ($22+35+41=98\approx100$, dentro de $2\%$);
$\SI{440}{\ohm}$, $\SI{700}{\ohm}$ e $\SI{820}{\ohm}$}


\end{questao}

\blocoaprofundamento
\begin{questao}[3]{}

$R_1=\SI{10}{\ohm}$ e $R_2=\SI{20}{\ohm}$ estão em série sob
$\SI{30}{\volt}$.
\begin{itensq}
\item Calcule a potência de cada um.
\item Especifique a potência nominal comercial de cada resistor, com fator de
      segurança $2$.
\item Explique por que especificar ambos igualmente seria desperdício ou risco.
\end{itensq}
\resposta{(a) $\SI{10}{\watt}$ e $\SI{20}{\watt}$; (b) $\SI{25}{\watt}$ e
$\SI{50}{\watt}$}


\end{questao}
\begin{questao}[3]{}

Um resistor de fio tem $\SI{50}{\ohm}$ a $\SI{20}{\celsius}$ e coeficiente
térmico $\alpha=\SI{0.004}{\per\celsius}$. Ele está em série com um resistor de
filme de $\SI{50}{\ohm}$ e $\alpha\approx0$, sob $\SI{124}{\volt}$.
\begin{itensq}
\item Determine a resistência do resistor de fio a $\SI{80}{\celsius}$.
\item Calcule a corrente a $\SI{20}{\celsius}$ e a $\SI{80}{\celsius}$.
\item Comente a variação percentual.
\end{itensq}
\resposta{(a) $\SI{62}{\ohm}$; (b) $\SI{1.24}{\ampere}$ e
$\SI{1.107}{\ampere}$; (c) queda de $10{,}7\%$}


\end{questao}
\begin{questao}[3]{}

$R_1=\SI{100}{\ohm}\pm5\%$ e $R_2=\SI{200}{\ohm}\pm5\%$ estão em série.
\begin{itensq}
\item Determine $R_{eq}$ e seus limites de pior caso.
\item Expresse a tolerância resultante em porcentagem.
\item Repita supondo $R_2$ com tolerância de $1\%$.
\end{itensq}
\resposta{(a) $\SI{300}{\ohm}$, entre $\SI{285}{\ohm}$ e $\SI{315}{\ohm}$;
(b) $\pm5\%$; (c) $\pm2{,}33\%$}


\end{questao}
\begin{questao}[3]{}

Na cadeia $R_1=\SI{10}{\ohm}$, $R_2=\SI{20}{\ohm}$, $R_3=\SI{30}{\ohm}$ sob
$\SI{60}{\volt}$, o resistor $R_2$ sofre \textbf{curto-circuito}. Compare
corrente e potências antes e depois.
\resposta{$I$: $\SI{1}{\ampere}\to\SI{1.5}{\ampere}$;
$P_1$: $10\to\SI{22.5}{\watt}$; $P_3$: $30\to\SI{67.5}{\watt}$}


\end{questao}
\begin{questao}[3]{}

Na mesma cadeia ($\SI{10}{\ohm}$, $\SI{20}{\ohm}$, $\SI{30}{\ohm}$ sob
$\SI{60}{\volt}$), agora $R_2$ \textbf{abre}. Determine a corrente e a leitura
de um voltímetro ideal colocado sobre cada resistor.
\resposta{$I=0$; $\SI{0}{\volt}$ em $R_1$ e $R_3$; $\SI{60}{\volt}$ em $R_2$}


\end{questao}
\begin{questao}[3]{}

Um LED com queda direta de $\SI{2.0}{\volt}$ deve operar com
$\SI{15}{\milli\ampere}$ a partir de uma fonte de $\SI{5.0}{\volt}$.
\begin{itensq}
\item Determine o resistor em série ideal.
\item Escolha o valor comercial (série E12) e recalcule a corrente real.
\item Determine a potência do resistor.
\end{itensq}
\resposta{(a) $\SI{200}{\ohm}$; (b) $\SI{220}{\ohm}\Rightarrow
\SI{13.6}{\milli\ampere}$; (c) $\SI{41}{\milli\watt}$}


\end{questao}
\begin{questao}[3]{}

Uma bateria de $\EMF=\SI{12}{\volt}$ e resistência interna $r=\SI{0.5}{\ohm}$
alimenta uma carga de $\SI{9.5}{\ohm}$.
\begin{itensq}
\item Determine a corrente e a tensão nos terminais.
\item Calcule a regulação de tensão percentual.
\item Determine a fração da potência perdida internamente.
\end{itensq}
\resposta{(a) $\SI{1.2}{\ampere}$, $\SI{11.4}{\volt}$; (b) $5{,}26\%$;
(c) $5\%$}


\end{questao}
\begin{questao}[3]{}

Duas cadeias têm $R_{eq}=\SI{120}{\ohm}$ sob $\SI{120}{\volt}$: (A) quatro
resistores de $\SI{30}{\ohm}$; (B) um de $\SI{100}{\ohm}$ e dois de
$\SI{10}{\ohm}$. Compare a potência do componente mais solicitado e a potência
total, e escolha a montagem preferível.
\resposta{Total $\SI{120}{\watt}$ em ambas; pior componente:
$\SI{30}{\watt}$ em (A) e $\SI{100}{\watt}$ em (B) --- (A) é preferível}


\end{questao}
\begin{questao}[3]{}

Um termistor NTC de $\SI{120}{\ohm}$ (frio) e $\SI{5}{\ohm}$ (quente) é
colocado em série com uma carga de $\SI{20}{\ohm}$, sob $\SI{100}{\volt}$.
Determine a corrente no instante da energização e em regime, e explique a
função do componente.
\resposta{$\SI{0.71}{\ampere}$ na partida e $\SI{4.0}{\ampere}$ em regime}


\end{questao}
\begin{questao}[3]{}

Em uma cadeia de três resistores sob $\SI{90}{\volt}$, sabe-se que
$R_2=2R_1$, $R_3=3R_1$ e que a corrente é $\SI{0.5}{\ampere}$. Determine os
três valores e as quedas.
\resposta{$R_1=\SI{30}{\ohm}$, $R_2=\SI{60}{\ohm}$, $R_3=\SI{90}{\ohm}$;
quedas $\SI{15}{\volt}$, $\SI{30}{\volt}$, $\SI{45}{\volt}$}


\end{questao}
\begin{questao}[3]{}

Uma cadeia de $n$ resistores iguais de $\SI{47}{\ohm}$ deve limitar a corrente
a no máximo $\SI{40}{\milli\ampere}$ sob $\SI{24}{\volt}$. Determine o menor
$n$ e a corrente resultante.
\resposta{$n=13$; $I=\SI{39.3}{\milli\ampere}$}


\end{questao}
\begin{questao}[3]{}

Em uma cadeia série de dois resistores sob tensão fixa $V$, demonstre que a
soma das potências é $P=\dfrac{V^{2}}{R_1+R_2}$ e explique por que aumentar
\emph{qualquer} resistor reduz a potência total, embora possa \emph{aumentar} a
potência do outro.
\resposta{Aumentar $R_k$ reduz $P_{tot}$; a potência do outro sempre diminui
--- é a do próprio $R_k$ que pode subir}


\end{questao}

\blocodesafio
\begin{questao}[4]{}

\textbf{(Projeto.)} Uma cadeia série de três resistores deve dividir
$\SI{24}{\volt}$ em quedas de $\SI{6}{\volt}$, $\SI{8}{\volt}$ e
$\SI{10}{\volt}$, com corrente de $\SI{20}{\milli\ampere}$.
\begin{itensq}
\item Determine os três resistores e suas potências.
\item Escolha valores comerciais E24 e recalcule as quedas reais.
\item Avalie se o erro é aceitável para alimentar cargas com tolerância de
      $\pm5\%$.
\end{itensq}
\resposta{(a) $\SI{300}{\ohm}$, $\SI{400}{\ohm}$, $\SI{500}{\ohm}$;
$\SI{0.12}{\watt}$, $\SI{0.16}{\watt}$, $\SI{0.20}{\watt}$;
(b) $300/390/510\ \si{\ohm}$ $\Rightarrow$ $5{,}98/7{,}77/10{,}16\ \si{\volt}$;
(c) sim, erro máximo de $2{,}9\%$}


\end{questao}
\begin{questao}[4]{}

\textbf{(Propagação de incerteza.)} Uma cadeia é formada por $N$ resistores
nominalmente iguais, cada um com tolerância relativa $t$.
\begin{itensq}
\item Obtenha a tolerância de $R_{eq}$ no critério de \textbf{pior caso}.
\item Obtenha-a no critério \textbf{estatístico} (erros independentes).
\item Compare para $N=5$ e $t=5\%$, e discuta qual usar.
\end{itensq}
\resposta{(a) $t$; (b) $t/\sqrt{N}$; (c) $5\%$ contra $2{,}24\%$}


\end{questao}
\begin{questao}[4]{}

\textbf{(Projeto com compensação térmica.)} Deseja-se uma associação série cuja
resistência total seja \emph{independente} da temperatura, combinando um
resistor de $R_1=\SI{100}{\ohm}$ com $\alpha_1=+4000\,\mathrm{ppm/^{\circ}C}$ e
outro com $\alpha_2=-1000\,\mathrm{ppm/^{\circ}C}$.
\begin{itensq}
\item Deduza a condição de compensação.
\item Determine $R_2$ e a resistência total.
\item Discuta a limitação prática dessa compensação.
\end{itensq}
\resposta{(a) $R_1\alpha_1+R_2\alpha_2=0$; (b) $R_2=\SI{400}{\ohm}$,
$R_{eq}=\SI{500}{\ohm}$; (c) só é exata em primeira ordem e numa temperatura}


\end{questao}
\begin{questao}[4]{}

\textbf{(Projeto sob incerteza.)} Um LED deve ser alimentado por uma fonte
entre $\SI{9}{\volt}$ e $\SI{12}{\volt}$. Sua queda direta varia entre
$\SI{1.8}{\volt}$ e $\SI{2.2}{\volt}$, e a corrente não pode exceder
$\SI{20}{\milli\ampere}$.
\begin{itensq}
\item Identifique a combinação de pior caso e determine o menor $R$ admissível.
\item Determine a corrente mínima resultante e avalie o impacto no brilho.
\item Especifique a potência do resistor.
\end{itensq}
\resposta{(a) $\SI{12}{\volt}$ com $\SI{1.8}{\volt}$ $\Rightarrow$
$R\ge\SI{510}{\ohm}$; (b) $\SI{13.3}{\milli\ampere}$;
(c) $\SI{0.204}{\watt}\Rightarrow\frac{1}{2}\,\si{\watt}$}


\end{questao}
\begin{questao}[4]{}

\textbf{(Sensibilidade.)} Em uma cadeia $I=\dfrac{V}{R_1+R_2}$, com
$V=\SI{12}{\volt}$, $R_1=\SI{100}{\ohm}$ e $R_2=\SI{200}{\ohm}$.
\begin{itensq}
\item Obtenha $\partial I/\partial R_1$ e avalie-a numericamente.
\item Obtenha a sensibilidade \emph{relativa}
      $S=\dfrac{\partial I/I}{\partial R_1/R_1}$.
\item Interprete o resultado para $R_1\ll R_2$ e $R_1\gg R_2$.
\end{itensq}
\resposta{(a) $-\pot{1.33}{-4}\,\si{\ampere\per\ohm}$;
(b) $S=-R_1/(R_1+R_2)=-1/3$; (c) $S\to0$ e $S\to-1$}


\end{questao}
\begin{questao}[4]{}

\textbf{(Análise de falhas.)} Dez resistores de $\SI{10}{\ohm}$ formam uma
cadeia sob $\SI{100}{\volt}$.
\begin{itensq}
\item Determine corrente, queda e potência por resistor em operação normal.
\item Repita supondo que \emph{um} resistor entre em curto.
\item Compare com o caso de um resistor \emph{aberto} e discuta qual falha é
      mais perigosa.
\end{itensq}
\resposta{(a) $\SI{1}{\ampere}$, $\SI{10}{\volt}$, $\SI{10}{\watt}$;
(b) $\SI{1.11}{\ampere}$, $\SI{11.1}{\volt}$, $\SI{12.3}{\watt}$;
(c) o curto é mais perigoso}


\end{questao}
\begin{questao}[4]{}

\textbf{(Método experimental.)} Para determinar $\EMF$ e $r$ de uma bateria,
mediu-se a tensão de terminal em duas condições de carga:
$(\SI{1}{\ampere};\ \SI{11.5}{\volt})$ e
$(\SI{3}{\ampere};\ \SI{10.5}{\volt})$.
\begin{itensq}
\item Obtenha $r$ e $\EMF$.
\item Explique a vantagem de usar dois pontos em vez de medir a tensão em
      vazio.
\item Estime a corrente de curto-circuito e comente sua validade.
\end{itensq}
\resposta{(a) $r=\SI{0.5}{\ohm}$, $\EMF=\SI{12}{\volt}$;
(c) $I_{cc}=\SI{24}{\ampere}$, valor apenas indicativo}


\end{questao}
\begin{questao}[4]{}

\textbf{(Projeto com componentes disponíveis.)} Deseja-se uma resistência de
$\SI{1}{\kilo\ohm}$ capaz de dissipar $\SI{2}{\watt}$, dispondo apenas de
resistores de $\SI{200}{\ohm}$ e $\frac{1}{2}\,\si{\watt}$.
\begin{itensq}
\item Proponha a associação e verifique o valor.
\item Determine a potência de cada resistor e a margem de segurança.
\item Calcule a tensão e a corrente nominais da associação.
\end{itensq}
\resposta{(a) cinco em série; (b) $\SI{0.4}{\watt}$ cada, margem de $25\%$;
(c) $\SI{44.7}{\volt}$ e $\SI{44.7}{\milli\ampere}$}


\end{questao}
\begin{questao}[4]{}

\textbf{(Instrumentação.)} Cinco sensores resistivos de $\SI{1}{\kilo\ohm}$
estão em série, alimentados por uma fonte de \textbf{corrente constante} de
$\SI{1}{\milli\ampere}$. Mede-se a tensão de cada nó em relação ao terra.
\begin{itensq}
\item Determine as cinco tensões nodais em condição nominal.
\item Um sensor passa a $\SI{1.2}{\kilo\ohm}$. Mostre como identificar
      \emph{qual} deles mudou.
\item Explique por que a fonte de corrente é preferível a uma fonte de tensão.
\end{itensq}
\resposta{(a) $1$, $2$, $3$, $4$ e $\SI{5}{\volt}$; (b) pela quebra do degrau
uniforme de $\SI{1}{\volt}$; (c) desacopla os sensores entre si}


\end{questao}
\begin{questao}[4]{}

\textbf{(Otimização.)} Uma cadeia série sob tensão fixa $V$ contém um resistor
ajustável $R$ e um bloco fixo de resistência total $R_f$.
\begin{itensq}
\item Obtenha $P_R(R)$ e mostre que ela é máxima em $R=R_f$.
\item Calcule $P_{R,\max}$ e o rendimento nesse ponto.
\item Compare com a condição de máxima potência \emph{total} e explique a
      diferença.
\end{itensq}
\resposta{(a) $R=R_f$; (b) $P_{\max}=V^{2}/(4R_f)$, rendimento $50\%$;
(c) a potência total é máxima em $R\to0$}


\end{questao}

\gabarito[3.1 Circuitos em série]

\subsection{Circuitos em paralelo}

\blocofixacao
\begin{questao}[1]{}

Sobre uma associação de resistores \textbf{em paralelo}, é correto afirmar que:
\begin{alternativas}
\item a corrente é a mesma em todos os ramos.
\item a resistência equivalente é sempre maior que a maior das resistências.
\item a tensão é a mesma em todos os ramos.
\item a tensão se divide proporcionalmente às resistências.
\item a resistência equivalente é a soma das resistências.
\end{alternativas}
\resposta{(c)}


\end{questao}
\begin{questao}[1]{}

Para cada circuito, calcule a resistência equivalente, as correntes de ramo e a
corrente total $I_t$.

\circuitoq{
\subcirc{a}{
\begin{circuitikz}[scale=0.72,transform shape,line width=0.8pt]
 \draw (0,0) to[battery1,l_={$E=\SI{30}{\volt}$}] (0,3.4);
 \draw (0,3.4) -- (9,3.4);  \draw (0,0) -- (9,0);
 \draw (2.2,3.4) to[R,l=$\SI{4}{\ohm}$,*-*] (2.2,0);
 \draw (4.5,3.4) to[R,l=$\SI{3}{\ohm}$,*-*] (4.5,0);
 \draw (6.8,3.4) to[R,l=$\SI{6}{\ohm}$,*-*] (6.8,0);
 \draw (9,3.4)   to[R,l=$\SI{12}{\ohm}$] (9,0);
 \draw[->,Icol,line width=1pt] (0.5,3.75) -- (1.6,3.75)
       node[midway,above,font=\footnotesize]{$I_t$};
\end{circuitikz}}
\hfill
\subcirc{b}{
\begin{circuitikz}[scale=0.72,transform shape,line width=0.8pt]
 \draw (0,0) to[battery1,l_={$E=\SI{24}{\volt}$}] (0,3.4);
 \draw (0,3.4) -- (6.6,3.4);  \draw (0,0) -- (6.6,0);
 \draw (2.2,3.4) to[R,l=$\SI{12}{\ohm}$,*-*] (2.2,0);
 \draw (4.4,3.4) to[R,l=$\SI{12}{\ohm}$,*-*] (4.4,0);
 \draw (6.6,3.4) to[R,l=$\SI{12}{\ohm}$] (6.6,0);
 \draw[->,Icol,line width=1pt] (0.5,3.75) -- (1.6,3.75)
       node[midway,above,font=\footnotesize]{$I_t$};
\end{circuitikz}}

\vspace{1.1em}

\subcirc{c}{
\begin{circuitikz}[scale=0.72,transform shape,line width=0.8pt]
 \draw (0,0) to[battery1,l_={$E=\SI{15}{\volt}$}] (0,3.4);
 \draw (0,3.4) -- (6.6,3.4);  \draw (0,0) -- (6.6,0);
 \draw (2.2,3.4) to[R,l=$\SI{12}{\ohm}$,*-*] (2.2,0);
 \draw (4.4,3.4) to[R,l=$\SI{12}{\ohm}$,*-*] (4.4,0);
 \draw (6.6,3.4) to[R,l=$\SI{3}{\ohm}$] (6.6,0);
 \draw[->,Icol,line width=1pt] (0.5,3.75) -- (1.6,3.75)
       node[midway,above,font=\footnotesize]{$I_t$};
\end{circuitikz}}
\hfill
\subcirc{d}{
\begin{circuitikz}[scale=0.72,transform shape,line width=0.8pt]
 \draw (0,0) to[battery1,l_={$E=\SI{9}{\volt}$}] (0,3.4);
 \draw (0,3.4) -- (4.4,3.4);  \draw (0,0) -- (4.4,0);
 \draw (2.2,3.4) to[R,l=$\SI{30}{\ohm}$,*-*] (2.2,0);
 \draw (4.4,3.4) to[R,l=$\SI{60}{\ohm}$] (4.4,0);
 \draw[->,Icol,line width=1pt] (0.5,3.75) -- (1.6,3.75)
       node[midway,above,font=\footnotesize]{$I_t$};
\end{circuitikz}}}

\resposta{(a) \SI{1.2}{\ohm}, $I_t=\SI{25}{\ampere}$;
(b) \SI{4}{\ohm}, $I_t=\SI{6}{\ampere}$;
(c) \SI{2}{\ohm}, $I_t=\SI{7.5}{\ampere}$;
(d) \SI{20}{\ohm}, $I_t=\SI{0.45}{\ampere}$}


\end{questao}
\begin{questao}[1]{}

$n$ resistores iguais de $\SI{60}{\ohm}$ são associados em paralelo, resultando
em $\SI{10}{\ohm}$. O valor de $n$ é:
\begin{alternativas}[3]
\item 2
\item 3
\item 4
\item 6
\item 10
\end{alternativas}
\resposta{(d)}


\end{questao}

\blocoaplicacao
\begin{questao}[2]{}

Um resistor de $\SI{20}{\ohm}$ é associado em paralelo com um resistor $R$
desconhecido, resultando em uma resistência equivalente de $\SI{12}{\ohm}$.
Determine $R$.
\resposta{$R = \SI{30}{\ohm}$}


\end{questao}
\begin{questao}[2]{}

Uma corrente total de $\SI{12}{\ampere}$ chega a um nó e se divide entre dois
resistores em paralelo, de $\SI{4}{\ohm}$ e $\SI{12}{\ohm}$. As correntes em cada
resistor valem, respectivamente:
\begin{alternativas}[2]
\item \SI{3}{\ampere} e \SI{9}{\ampere}
\item \SI{9}{\ampere} e \SI{3}{\ampere}
\item \SI{6}{\ampere} e \SI{6}{\ampere}
\item \SI{4}{\ampere} e \SI{8}{\ampere}
\item \SI{8}{\ampere} e \SI{4}{\ampere}
\end{alternativas}
\resposta{(b)}


\end{questao}
\begin{questao}[2]{}

Três lâmpadas incandescentes de $\SI{60}{\watt}$/$\SI{120}{\volt}$ são ligadas
\textbf{em paralelo} à rede de $\SI{120}{\volt}$.
\begin{itensq}
\item Calcule a resistência de cada lâmpada e a resistência equivalente.
\item Determine a corrente total fornecida pela rede.
\item Se uma das lâmpadas queimar, o que acontece com o brilho das outras duas?
\end{itensq}
\resposta{(a) \SI{240}{\ohm} cada, $\Req=\SI{80}{\ohm}$;
(b) \SI{1.5}{\ampere}; (c) nada muda}


\end{questao}
\begin{questao}[2]{}

No trecho da figura, um fio ideal (sem resistência) é ligado em paralelo com o
resistor $R_2$.

\circuitoq{
\begin{circuitikz}[scale=0.95,transform shape,line width=0.8pt]
 \draw (0,0) to[battery1,l_={$E=\SI{20}{\volt}$}] (0,3)
       to[R,l={$R_1=\SI{10}{\ohm}$}] (3,3)
       to[R,l={$R_2=\SI{20}{\ohm}$}] (6,3) -- (6,0) -- (0,0);
 \draw[Vcol,line width=1.1pt] (3,3) -- (3,4.2) -- (6,4.2) -- (6,3);
 \node[Vcol,font=\footnotesize\sffamily,above] at (4.5,4.25) {fio ideal};
 \fill (3,3) circle (0.06); \fill (6,3) circle (0.06);
\end{circuitikz}}

A corrente fornecida pela fonte vale:
\begin{alternativas}[3]
\item \SI{0.67}{\ampere}
\item \SI{1}{\ampere}
\item \SI{1.5}{\ampere}
\item \SI{2}{\ampere}
\item \SI{3}{\ampere}
\end{alternativas}
\resposta{(d)}


\end{questao}

\blocoaprofundamento
\begin{questao}[3]{}

Um resistor de $\SI{100}{\ohm}$ suporta no máximo $\SI{2}{\watt}$ de dissipação.
Deseja-se construir, com resistores desse mesmo tipo, uma associação
\textbf{em paralelo} de $\SI{25}{\ohm}$.
\begin{itensq}
\item Quantos resistores são necessários?
\item Qual é a máxima tensão que pode ser aplicada à associação?
\item Qual é a potência máxima que a associação suporta?
\end{itensq}
\resposta{(a) 4 resistores; (b) $\approx\SI{14.1}{\volt}$;
(c) \SI{8}{\watt}}


\end{questao}
\begin{questao}[3]{}

Determine a resistência equivalente entre os terminais A e B do circuito da
figura, no qual dois trechos são interligados por fios ideais.

\circuitoq{
\begin{circuitikz}[scale=0.85,transform shape,line width=0.8pt]
 \draw (0,0) node[left,font=\bfseries] {A}
       to[R,l=$\SI{15}{\ohm}$,o-*] (3,0)
       to[R,l=$\SI{30}{\ohm}$,-*] (6,0)
       to[R,l=$\SI{30}{\ohm}$,-o] (9,0)
       node[right,font=\bfseries] {B};
 \draw[Vcol,line width=1.1pt] (3,0) -- (3,1.8) -- (6,1.8) -- (6,0);
 \node[Vcol,font=\footnotesize\sffamily,above] at (4.5,1.85) {fio ideal};
\end{circuitikz}}

\resposta{$R_{AB} = \SI{45}{\ohm}$}


\end{questao}

\blocodesafio
\begin{questao}[4]{}

$n$ resistores \emph{diferentes} $R_1, R_2, \dots, R_n$ são associados em
paralelo.
\begin{itensq}
\item Demonstre que a resistência equivalente é sempre \textbf{menor} que a menor
      das resistências.
\item Mostre que, se um dos resistores tende a zero, $\Req \to 0$, e interprete
      fisicamente.
\item Dois resistores em paralelo dão $\SI{6}{\ohm}$; em série, dariam
      $\SI{25}{\ohm}$. Determine seus valores.
\end{itensq}
\resposta{(a) e (b) demonstrações; (c) $\SI{10}{\ohm}$ e $\SI{15}{\ohm}$}


\end{questao}

\blocofixacao
\begin{questao}[1]{}

Determine a resistência equivalente de $R_1=\SI{30}{\ohm}$ e
$R_2=\SI{60}{\ohm}$ em paralelo.
\resposta{$R_{eq}=\SI{20}{\ohm}$}


\end{questao}
\begin{questao}[1]{}

Três resistores de $\SI{4}{\ohm}$, $\SI{8}{\ohm}$ e $\SI{8}{\ohm}$ estão em
paralelo. Calcule a condutância total e $R_{eq}$.
\resposta{$G=\SI{0.5}{\ensuremath{\mathrm{S}}}$; $R_{eq}=\SI{2}{\ohm}$}


\end{questao}
\begin{questao}[1]{}

Três resistores iguais de $\SI{90}{\ohm}$ são associados em paralelo. Determine
$R_{eq}$.
\resposta{$R_{eq}=\SI{30}{\ohm}$}


\end{questao}
\begin{questao}[1]{}

Acrescentando-se mais um ramo \textbf{em paralelo} a um circuito alimentado por
fonte ideal de tensão, a corrente total:
\begin{alternativas}[2]
\item aumenta
\item diminui
\item não se altera
\item pode aumentar ou diminuir, conforme o valor
\end{alternativas}
\resposta{(a)}


\end{questao}
\begin{questao}[1]{}

Um ramo de $\SI{48}{\ohm}$ pertence a uma associação paralela ligada a
$\SI{24}{\volt}$. Determine a corrente desse ramo.
\resposta{$I=\SI{0.5}{\ampere}$}


\end{questao}
\begin{questao}[1]{}

Uma associação paralela sob $\SI{40}{\volt}$ tem ramos que conduzem
$\SI{2}{\ampere}$, $\SI{3}{\ampere}$ e $\SI{5}{\ampere}$. Determine a corrente
total e $R_{eq}$.
\resposta{$I_t=\SI{10}{\ampere}$; $R_{eq}=\SI{4}{\ohm}$}


\end{questao}
\begin{questao}[1]{}

Em uma associação paralela, o ramo que conduz a \textbf{maior} corrente é:
\begin{alternativas}[2]
\item o de maior resistência
\item o de menor resistência
\item o mais próximo da fonte
\item todos conduzem igualmente
\end{alternativas}
\resposta{(b)}


\end{questao}

\blocoaplicacao
\begin{questao}[2]{}

$R_1=\SI{10}{\ohm}$, $R_2=\SI{20}{\ohm}$ e $R_3=\SI{20}{\ohm}$ estão em
paralelo sob $\SI{60}{\volt}$. Determine as correntes de ramo, a corrente total
e $R_{eq}$.
\resposta{$\SI{6}{\ampere}$, $\SI{3}{\ampere}$, $\SI{3}{\ampere}$;
$I_t=\SI{12}{\ampere}$; $R_{eq}=\SI{5}{\ohm}$}


\end{questao}
\begin{questao}[2]{}

No circuito da questão anterior, calcule a potência de cada ramo e confronte
com a potência total.
\resposta{$\SI{360}{\watt}$, $\SI{180}{\watt}$, $\SI{180}{\watt}$; total
$\SI{720}{\watt}$}


\end{questao}
\begin{questao}[2]{}

Uma corrente de $\SI{15}{\ampere}$ chega a um nó e se divide entre
$R_1=\SI{8}{\ohm}$ e $R_2=\SI{12}{\ohm}$. Determine as correntes de ramo pela
fórmula do divisor.
\resposta{$I_1=\SI{9}{\ampere}$, $I_2=\SI{6}{\ampere}$}


\end{questao}
\begin{questao}[2]{}

Três resistores em paralelo sob $\SI{100}{\volt}$ drenam $\SI{5}{\ampere}$ no
total. Dois deles valem $\SI{50}{\ohm}$ e $\SI{100}{\ohm}$. Determine o
terceiro.
\resposta{$R_3=\SI{50}{\ohm}$}


\end{questao}
\begin{questao}[2]{}

Duas lâmpadas de resistências $\SI{240}{\ohm}$ e $\SI{960}{\ohm}$ são ligadas
\textbf{em paralelo} a $\SI{120}{\volt}$. Qual brilha mais? Compare com o
resultado da mesma dupla em série.
\resposta{Em paralelo, a de $\SI{240}{\ohm}$ ($\SI{60}{\watt}$ contra
$\SI{15}{\watt}$) --- inverte-se em relação à série}


\end{questao}
\begin{questao}[2]{}

Um galvanômetro tem resistência interna $\SI{100}{\ohm}$ e deflexão total com
$\SI{1}{\milli\ampere}$. Determine o resistor \emph{shunt} que amplia a escala
para $\SI{11}{\milli\ampere}$.
\resposta{$R_{sh}=\SI{10}{\ohm}$}


\end{questao}
\begin{questao}[2]{}

Um circuito residencial de $\SI{220}{\volt}$ alimenta, em paralelo, cargas de
$\SI{4400}{\watt}$, $\SI{1100}{\watt}$ e $\SI{550}{\watt}$. Determine a
corrente total e escolha o disjuntor.
\resposta{$I_t=\SI{27.5}{\ampere}$; disjuntor de $\SI{32}{\ampere}$}


\end{questao}
\begin{questao}[2]{}

Dois condutores idênticos de $\SI{0.8}{\ohm}$ cada são ligados em paralelo para
alimentar uma carga. Determine a resistência resultante e explique a prática.
\resposta{$\SI{0.4}{\ohm}$ --- metade}


\end{questao}
\begin{questao}[2]{}

Um resistor de $\SI{1}{\mega\ohm}$ é ligado em paralelo com um de
$\SI{100}{\ohm}$. A resistência equivalente:
\begin{alternativas}[2]
\item fica praticamente igual a $\SI{100}{\ohm}$
\item fica praticamente igual a $\SI{1}{\mega\ohm}$
\item vale aproximadamente $\SI{500}{\kilo\ohm}$
\item vale exatamente $\SI{50}{\ohm}$
\end{alternativas}
\resposta{(a)}


\end{questao}
\begin{questao}[2]{}

Um chuveiro possui duas resistências de $\SI{22}{\ohm}$ que podem ser ligadas
isoladamente ou em paralelo, sob $\SI{220}{\volt}$. Determine a potência em
cada caso.
\resposta{$\SI{2200}{\watt}$ com uma; $\SI{4400}{\watt}$ com as duas}


\end{questao}
\begin{questao}[2]{}

Em uma associação paralela mediram-se as correntes de ramo $\SI{4}{\ampere}$,
$\SI{1}{\ampere}$ e $\SI{1}{\ampere}$, com corrente total de
$\SI{6}{\ampere}$ e tensão de $\SI{20}{\volt}$. Verifique a consistência e
determine as três resistências e a potência total.
\resposta{Consistente; $\SI{5}{\ohm}$, $\SI{20}{\ohm}$, $\SI{20}{\ohm}$;
$P=\SI{120}{\watt}$}


\end{questao}

\blocoaprofundamento
\begin{questao}[3]{}

Uma corrente total de $\SI{6}{\ampere}$ divide-se entre $\SI{5}{\ohm}$,
$\SI{20}{\ohm}$ e $\SI{20}{\ohm}$ em paralelo.
\begin{itensq}
\item Determine $R_{eq}$ e a tensão da associação.
\item Obtenha as três correntes pelo divisor em condutâncias.
\item Verifique por dois caminhos independentes.
\end{itensq}
\resposta{$R_{eq}=\SI{3.33}{\ohm}$, $V=\SI{20}{\volt}$;
$\SI{4}{\ampere}$, $\SI{1}{\ampere}$, $\SI{1}{\ampere}$}


\end{questao}
\begin{questao}[3]{}

Três ramos de $\SI{10}{\ohm}$, $\SI{20}{\ohm}$ e $\SI{20}{\ohm}$ operam sob
fonte \textbf{ideal} de $\SI{60}{\volt}$. O ramo de $\SI{10}{\ohm}$ se
\textbf{abre}. Determine o que muda nas correntes de ramo, na corrente total e
em $R_{eq}$.
\resposta{Ramos restantes inalterados ($\SI{3}{\ampere}$ cada);
$I_t$: $12\to\SI{6}{\ampere}$; $R_{eq}$: $5\to\SI{10}{\ohm}$}


\end{questao}
\begin{questao}[3]{}

Um dos ramos de uma associação paralela sofre \textbf{curto-circuito}. Analise
$R_{eq}$, a corrente total e o efeito sobre os demais ramos, e justifique a
necessidade de proteção.
\resposta{$R_{eq}\to0$; $I_t\to\infty$ (limitada só pela fonte e pela fiação);
os demais ramos ficam sem tensão}


\end{questao}
\begin{questao}[3]{}

Dispõe-se de resistores de $\SI{220}{\ohm}$ com dissipação máxima de
$\SI{1}{\watt}$. Projete uma associação paralela de $\SI{55}{\ohm}$ capaz de
dissipar pelo menos $\SI{4}{\watt}$, e determine a tensão máxima aplicável.
\resposta{Quatro em paralelo; $P_{\max}=\SI{4}{\watt}$;
$V_{\max}=\SI{14.8}{\volt}$}


\end{questao}
\begin{questao}[3]{}

Uma rede tem cinco ramos em paralelo: $\SI{10}{\ohm}$, $\SI{20}{\ohm}$,
$\SI{20}{\ohm}$, $\SI{40}{\ohm}$ e $\SI{40}{\ohm}$. Calcule $R_{eq}$ e explique
por que o cálculo em condutâncias é preferível.
\resposta{$G=\SI{0.25}{\ensuremath{\mathrm{S}}}$; $R_{eq}=\SI{4}{\ohm}$}


\end{questao}
\begin{questao}[3]{}

Um galvanômetro de $\SI{99}{\ohm}$ e $\SI{1}{\milli\ampere}$ de fundo de escala
deve medir até $\SI{100}{\milli\ampere}$.
\begin{itensq}
\item Determine o shunt.
\item Calcule a resistência do amperímetro resultante.
\item Explique por que essa resistência precisa ser pequena.
\end{itensq}
\resposta{(a) $R_{sh}=\SI{1}{\ohm}$; (b) $\SI{0.99}{\ohm}$}


\end{questao}
\begin{questao}[3]{}

Uma carga de $\SI{100}{\ohm}$ opera a $\SI{200}{\volt}$. A isolação
degradada cria um caminho de fuga de $\SI{10}{\kilo\ohm}$ em paralelo.
\begin{itensq}
\item Determine a corrente de fuga e a corrente total.
\item Calcule $R_{eq}$ e o erro relativo cometido ao ignorar a fuga.
\item Comente o risco.
\end{itensq}
\resposta{(a) $\SI{20}{\milli\ampere}$ e $\SI{2.02}{\ampere}$;
(b) $\SI{99.01}{\ohm}$, erro de $1\%$; (c) risco de choque}


\end{questao}
\begin{questao}[3]{}

$R_1=\SI{100}{\ohm}\pm5\%$ e $R_2=\SI{400}{\ohm}\pm5\%$ estão em paralelo.
Determine $R_{eq}$ nominal e seus extremos de pior caso.
\resposta{$R_{eq}=\SI{80}{\ohm}$, entre $\SI{76}{\ohm}$ e $\SI{84}{\ohm}$
($\pm5\%$)}


\end{questao}
\begin{questao}[3]{}

$R_1=\SI{100}{\ohm}$ com $\alpha=\SI{0.004}{\per\celsius}$ está em paralelo com
$R_2=\SI{100}{\ohm}$ de coeficiente desprezível. Determine $R_{eq}$ a
$\SI{20}{\celsius}$ e a $\SI{80}{\celsius}$, e calcule a variação percentual.
\resposta{$\SI{50}{\ohm}$ e $\SI{55.36}{\ohm}$; variação de $+10{,}7\%$}


\end{questao}
\begin{questao}[3]{}

Um resistor de $\SI{1000}{\ohm}$ está em paralelo com um de $\SI{10}{\ohm}$.
\begin{itensq}
\item Calcule $R_{eq}$.
\item Recalcule supondo que $R_1$ varie $+10\%$.
\item Interprete o resultado em termos de sensibilidade.
\end{itensq}
\resposta{(a) $\SI{9.901}{\ohm}$; (b) $\SI{9.910}{\ohm}$; (c) variação de
apenas $0{,}09\%$}


\end{questao}
\begin{questao}[3]{}

As cargas da questão residencial ($\SI{27.5}{\ampere}$ totais) são alimentadas
por um circuito cuja fiação tem $\SI{0.2}{\ohm}$ de ida e volta, a partir de um
quadro de $\SI{220}{\volt}$.
\begin{itensq}
\item Determine a queda de tensão na fiação e a tensão nas cargas.
\item Calcule a perda de potência na fiação.
\item Avalie se a queda atende ao limite usual de $4\%$.
\end{itensq}
\resposta{(a) $\SI{5.5}{\volt}$; $\SI{214.5}{\volt}$;
(b) $\SI{151.25}{\watt}$; (c) $2{,}5\%$ --- atende}


\end{questao}
\begin{questao}[3]{}

Deseja-se obter $\SI{75}{\ohm}$ dispondo de um resistor de $\SI{100}{\ohm}$.
Determine o resistor a ser associado em paralelo e discuta a viabilidade.
\resposta{$R=\SI{300}{\ohm}$}


\end{questao}
\begin{questao}[3]{}

Compare a mesma dupla $R_1=\SI{30}{\ohm}$ e $R_2=\SI{60}{\ohm}$ ligada em série
e em paralelo, sob $\SI{90}{\volt}$: resistência, corrente total, distribuição
de corrente e de potência.
\resposta{Série: $\SI{90}{\ohm}$, $\SI{1}{\ampere}$, $P=30$ e
$\SI{60}{\watt}$. Paralelo: $\SI{20}{\ohm}$, $\SI{4.5}{\ampere}$,
$P=270$ e $\SI{135}{\watt}$}


\end{questao}

\blocodesafio
\begin{questao}[4]{}

\textbf{(Projeto de instrumentação.)} Um galvanômetro tem $R_g=\SI{99}{\ohm}$ e
fundo de escala $I_g=\SI{1}{\milli\ampere}$. Deseja-se um amperímetro de três
escalas: $\SI{10}{\milli\ampere}$, $\SI{100}{\milli\ampere}$ e
$\SI{1}{\ampere}$.
\begin{itensq}
\item Determine os três shunts na configuração de shunts comutados.
\item Aponte o risco dessa configuração e descreva a alternativa (shunt de
      Ayrton).
\item Calcule a resistência do instrumento em cada escala.
\end{itensq}
\resposta{(a) $\SI{11}{\ohm}$, $\SI{1}{\ohm}$ e $\SI{0.0991}{\ohm}$;
(c) $\SI{9.9}{\ohm}$, $\SI{0.99}{\ohm}$ e $\SI{0.099}{\ohm}$}


\end{questao}
\begin{questao}[4]{}

\textbf{(Demonstração --- escalonamento de bancos.)} Considere $N$ resistores
iguais de valor $R$ e potência nominal $P_1$.
\begin{itensq}
\item Mostre que a associação paralela vale $R/N$ e suporta $NP_1$.
\item Mostre que a série vale $NR$ e também suporta $NP_1$.
\item Deduza a matriz $m\times n$ necessária para manter o valor $R$ e obter
      potência $mnP_1$.
\end{itensq}
\resposta{(c) $m$ colunas em paralelo de $n$ unidades em série, com $m=n$}


\end{questao}
\begin{questao}[4]{}

\textbf{(Sensibilidade.)} Para $R_{eq}=\dfrac{R_1R_2}{R_1+R_2}$:
\begin{itensq}
\item Obtenha $\dfrac{\partial R_{eq}}{\partial R_1}$ e mostre que vale
      $\left(\dfrac{R_2}{R_1+R_2}\right)^{2}$.
\item Obtenha a sensibilidade relativa e mostre que
      $S_1+S_2=1$.
\item Avalie para $R_1=\SI{100}{\ohm}$ e $R_2=\SI{400}{\ohm}$, e interprete.
\end{itensq}
\resposta{(a) $0{,}64$; (b) $S_1=R_2/(R_1+R_2)$, $S_2=R_1/(R_1+R_2)$;
(c) $S_1=0{,}8$ e $S_2=0{,}2$}


\end{questao}
\begin{questao}[4]{}

\textbf{(Divisão de corrente sob tolerância.)} Três shunts iguais de
$\SI{0.01}{\ohm}$ devem dividir igualmente $\SI{30}{\ampere}$.
\begin{itensq}
\item Determine a corrente ideal por shunt e a tensão sobre eles.
\item Se um deles for $2\%$ menor e os outros dois nominais, calcule o
      desequilíbrio de corrente.
\item Discuta a implicação térmica.
\end{itensq}
\resposta{(a) $\SI{10}{\ampere}$ e $\SI{0.1}{\volt}$;
(b) $\SI{10.13}{\ampere}$ contra $\SI{9.93}{\ampere}$ --- desvio de $1{,}3\%$;
(c) o mais "fraco" aquece mais}


\end{questao}
\begin{questao}[4]{}

\textbf{(Compensação térmica em paralelo.)} Deseja-se uma associação paralela
com resistência independente da temperatura, usando
$R_1=\SI{100}{\ohm}$ com $\alpha_1=+4000\,\mathrm{ppm/^{\circ}C}$ e um segundo
resistor com $\alpha_2=-1000\,\mathrm{ppm/^{\circ}C}$.
\begin{itensq}
\item Deduza a condição de compensação em termos de condutâncias.
\item Determine $R_2$ e $R_{eq}$.
\item Verifique a compensação e comente sua exatidão.
\end{itensq}
\resposta{(a) $G_1\alpha_1+G_2\alpha_2=0$; (b) $R_2=\SI{25}{\ohm}$,
$R_{eq}=\SI{20}{\ohm}$; (c) exata só em primeira ordem}


\end{questao}
\begin{questao}[4]{}

\textbf{(Falha de ramo com fonte real.)} Uma fonte de $\SI{24}{\volt}$ com
$r=\SI{2}{\ohm}$ alimenta três ramos de $\SI{6}{\ohm}$ em paralelo.
\begin{itensq}
\item Determine a tensão do barramento, a corrente total e a de cada ramo.
\item Um ramo se abre. Recalcule tudo.
\item Explique por que os ramos remanescentes \emph{mudam}, ao contrário do
      caso com fonte ideal.
\end{itensq}
\resposta{(a) $\SI{12}{\volt}$, $\SI{6}{\ampere}$, $\SI{2}{\ampere}$ por ramo;
(b) $\SI{14.4}{\volt}$, $\SI{4.8}{\ampere}$, $\SI{2.4}{\ampere}$ por ramo}


\end{questao}
\begin{questao}[4]{}

\textbf{(Limite do método de ajuste.)} Dispõe-se de um resistor fixo de
$\SI{100}{\ohm}$ e deseja-se obter valores-alvo por associação paralela.
\begin{itensq}
\item Deduza $R(alvo)$ e calcule para alvos de $\SI{90}{\ohm}$,
      $\SI{95}{\ohm}$ e $\SI{99}{\ohm}$.
\item Obtenha a sensibilidade do alvo em relação a $R$ e avalie no caso de
      $\SI{99}{\ohm}$.
\item Conclua sobre a viabilidade prática.
\end{itensq}
\resposta{(a) $\SI{900}{\ohm}$, $\SI{1900}{\ohm}$ e $\SI{9900}{\ohm}$;
(c) inviável para alvos próximos do resistor fixo}


\end{questao}
\begin{questao}[4]{}

\textbf{(Demonstração --- propagação de incerteza.)} Mostre que, se
\emph{todos} os resistores de uma rede puramente resistiva tiverem a mesma
tolerância relativa $t$, então $R_{eq}$ terá exatamente a tolerância $t$,
qualquer que seja a topologia.
\begin{itensq}
\item Demonstre a propriedade de homogeneidade.
\item Verifique em um exemplo série e em um paralelo.
\item Discuta o que muda quando as tolerâncias são diferentes.
\end{itensq}
\resposta{$R_{eq}$ é homogênea de grau 1: $R_{eq}(\lambda R_k)=\lambda
R_{eq}(R_k)$; logo a tolerância relativa se conserva}


\end{questao}
\begin{questao}[4]{}

\textbf{(Seletividade da proteção.)} Um barramento de $\SI{220}{\volt}$
alimenta três circuitos em paralelo, com correntes nominais de
$\SI{20}{\ampere}$, $\SI{5}{\ampere}$ e $\SI{2.5}{\ampere}$, protegidos
individualmente e por um disjuntor geral.
\begin{itensq}
\item Especifique os disjuntores de ramo e o geral.
\item Explique o requisito de seletividade e o que ocorre se ele falhar.
\item Discuta por que a soma dos disjuntores de ramo pode exceder o geral.
\end{itensq}
\resposta{(a) ramos de $25$, $10$ e $\SI{10}{\ampere}$; geral de
$\SI{32}{\ampere}$; (b) o ramo deve atuar antes do geral}


\end{questao}

\gabarito[3.2 Circuitos em paralelo]

\subsection{Circuitos mistos}

\blocofixacao
\begin{questao}[1]{}

A estratégia geral para reduzir um circuito misto é:
\begin{alternativas}
\item somar todas as resistências, independentemente da ligação.
\item resolver primeiro as associações série mais internas, depois os paralelos.
\item identificar os trechos em série e em paralelo e reduzi-los, dos mais internos
      para os mais externos, um de cada vez.
\item aplicar sempre a transformação estrela-triângulo.
\item calcular a corrente antes de reduzir o circuito.
\end{alternativas}
\resposta{(c)}


\end{questao}
\begin{questao}[1]{}

Para cada circuito, calcule a resistência equivalente vista pela fonte e a
corrente total.

\circuitoq{
\subcirc{a}{
\begin{circuitikz}[scale=0.8,transform shape,line width=0.8pt]
 \draw (0,0) to[battery1,l_=$\SI{24}{\volt}$] (0,3)
       to[R,l=$\SI{6}{\ohm}$] (2.6,3) coordinate(n);
 \draw (n) to[R,l=$\SI{6}{\ohm}$] (2.6,0) -- (0,0);
 \draw (n) -- (4.6,3) to[R,l=$\SI{3}{\ohm}$] (4.6,0) -- (2.6,0);
 \draw[->,Icol,line width=1pt] (0.6,3.4)--(1.5,3.4) node[midway,above,font=\footnotesize]{$I$};
\end{circuitikz}}
\hfill
\subcirc{b}{
\begin{circuitikz}[scale=0.8,transform shape,line width=0.8pt]
 \draw (0,0) to[battery1,l_=$\SI{20}{\volt}$] (0,3) coordinate(t);
 \draw (t) -- (2,3) to[R,l=$\SI{4}{\ohm}$] (2,0) coordinate(b);
 \draw (2,3) -- (4,3) to[R,l=$\SI{4}{\ohm}$] (4,0) -- (2,0);
 \draw (4,3) to[R,l=$\SI{8}{\ohm}$] (6,3) -- (6,0) -- (b) -- (0,0);
 \draw[->,Icol,line width=1pt] (0.5,3.4)--(1.4,3.4) node[midway,above,font=\footnotesize]{$I$};
\end{circuitikz}}}

\resposta{(a) \SI{8}{\ohm}, \SI{3}{\ampere}; (b) \SI{10}{\ohm}, \SI{2}{\ampere}}


\end{questao}

\blocoaplicacao
\begin{questao}[2]{}

No circuito da figura, determine a resistência equivalente vista pela fonte, a
corrente total e a corrente em cada resistor de \SI{20}{\ohm}.

\circuitoq{
\begin{circuitikz}[scale=0.9,transform shape,line width=0.8pt]
 \draw (0,0) to[battery1,l_={$E=\SI{40}{\volt}$}] (0,3)
       to[R,l=$\SI{10}{\ohm}$] (3,3) coordinate(n);
 \draw (n) to[R,l=$\SI{20}{\ohm}$] (3,0);
 \draw (n) -- (5.5,3) to[R,l=$\SI{20}{\ohm}$] (5.5,0) -- (3,0) -- (0,0);
 \draw[->,Icol,line width=1pt] (0.6,3.4)--(1.6,3.4) node[midway,above,font=\footnotesize]{$I_t$};
\end{circuitikz}}

\resposta{$\Req=\SI{20}{\ohm}$; $I_t=\SI{2}{\ampere}$; \SI{1}{\ampere} em cada \SI{20}{\ohm}}


\end{questao}
\begin{questao}[2]{}

Calcule a resistência equivalente entre A e B.

\circuitoq{
\begin{circuitikz}[scale=0.9,transform shape,line width=0.8pt]
 \draw (0,0) node[left,font=\bfseries]{A} to[short,o-*] (1.2,0) coordinate(a);
 
 \draw (a) to[R,l=$\SI{2}{\ohm}$] (3.2,0) coordinate(b);
 
 \draw (b) to[R,l=$\SI{6}{\ohm}$] (5.4,0) coordinate(c);
 \draw (b) ++(0,1.4) coordinate(b2) to[R,l=$\SI{3}{\ohm}$] ($(c)+(0,1.4)$) coordinate(c2);
 \draw (b) -- (b2);  \draw (c) -- (c2);
 
 \draw (c) to[R,l=$\SI{10}{\ohm}$] (7.6,0) coordinate(d);
 \draw (c) ++(0,1.4) coordinate(c3) to[R,l=$\SI{15}{\ohm}$] ($(d)+(0,1.4)$) coordinate(d2);
 \draw (c) -- (c3);  \draw (d) -- (d2);
 \draw (d) to[short,-o] (8.8,0) node[right,font=\bfseries]{B};
\end{circuitikz}}

\resposta{$R_{AB} = \SI{10}{\ohm}$}


\end{questao}

\blocoaprofundamento
\begin{questao}[3]{}

Determine a resistência equivalente entre A e B da rede em escada da figura.

\circuitoq{
\begin{circuitikz}[scale=0.9,transform shape,line width=0.8pt]
 \draw (0,2) node[left,font=\bfseries]{A} to[short,o-*] (1,2) coordinate(a);
 \draw (a) to[R,l=$\SI{12}{\ohm}$] (1,0) coordinate(b);
 \draw (a) to[R,l=$\SI{4}{\ohm}$] (4,2) coordinate(c);
 \draw (c) to[R,l=$\SI{6}{\ohm}$] (4,0) coordinate(d);
 \draw (c) to[R,l=$\SI{3}{\ohm}$] (6.5,2) coordinate(e);
 \draw (e) -- (6.5,0) -- (d) -- (b);
 \draw (b) to[short,*-o] (0,0) node[left,font=\bfseries]{B};
 \draw (e) -- (7,2) node[right]{};
\end{circuitikz}}

\resposta{$R_{AB} = \SI{4}{\ohm}$}


\end{questao}
\begin{questao}[3]{}

No circuito da figura, um fio ideal curto-circuita parte da rede. Determine a
resistência equivalente entre A e B.

\circuitoq{
\begin{circuitikz}[scale=0.9,transform shape,line width=0.8pt]
 \draw (0,1.5) node[left,font=\bfseries]{A} to[short,o-*] (1.5,1.5) coordinate(a);
 \draw (a) to[R,l=$\SI{10}{\ohm}$] (1.5,3.2) -- (5,3.2) coordinate(top);
 \draw (a) to[R,l=$\SI{20}{\ohm}$] (5,1.5) coordinate(m);
 \draw (top) -- (m);
 \draw[Vcol,line width=1.1pt] (a) -- (1.5,0) -- (5,0) -- (m);
 \node[Vcol,font=\footnotesize\sffamily,below] at (3.25,0) {fio ideal};
 \draw (m) to[R,l=$\SI{10}{\ohm}$] (7,1.5) to[short,-o] (7.8,1.5)
       node[right,font=\bfseries]{B};
\end{circuitikz}}

\resposta{$R_{AB} = \SI{10}{\ohm}$}


\end{questao}
\begin{questao}[3]{}

Doze resistores iguais de $\SI{6}{\ohm}$ são soldados formando as arestas de um
\textbf{cubo}. Determine a resistência equivalente entre dois vértices ligados por
uma \emph{aresta} (A e B na figura).

\circuitoq{
\begin{tikzpicture}[scale=1,line width=0.8pt]
 \def\a{2}
 
 \coordinate (A) at (0,0);
 \coordinate (B) at (\a,0);
 \coordinate (C) at (\a,\a);
 \coordinate (D) at (0,\a);
 \coordinate (E) at (0.7,0.6);
 \coordinate (F) at (\a+0.7,0.6);
 \coordinate (G) at (\a+0.7,\a+0.6);
 \coordinate (H) at (0.7,\a+0.6);
 \draw[Rcol] (A)--(B)--(C)--(D)--cycle;
 \draw[Rcol] (E)--(F)--(G)--(H)--cycle;
 \draw[Rcol] (A)--(E); \draw[Rcol] (B)--(F);
 \draw[Rcol] (C)--(G); \draw[Rcol] (D)--(H);
 \fill[Vcol] (A) circle (0.09) node[below left,font=\bfseries]{A};
 \fill[Vcol] (B) circle (0.09) node[below right,font=\bfseries]{B};
 \node[Rcol,font=\footnotesize] at (\a/2,-0.3){$\SI{6}{\ohm}$ (cada aresta)};
\end{tikzpicture}}{Cubo de resistores.}

\resposta{$R_{AB} = \SI{5}{\ohm}$}


\end{questao}
\begin{questao}[3]{}

Doze resistores iguais de $\SI{12}{\ohm}$ formam as arestas de um cubo.
Determine a resistência equivalente entre:
\begin{itensq}
\item dois vértices de uma mesma \emph{aresta};
\item dois vértices opostos de uma mesma \emph{face} (diagonal de face);
\item dois vértices \emph{diametralmente opostos} (diagonal principal).
\end{itensq}
\resposta{(a) \SI{7}{\ohm}; (b) \SI{9}{\ohm}; (c) \SI{10}{\ohm}}


\end{questao}
\begin{questao}[3]{}

No circuito da figura, todos os resistores valem $\SI{6}{\ohm}$, exceto o
resistor central da ponte, de $\SI{5}{\ohm}$. Determine $R_{AB}$ e justifique por
que o valor do resistor central é irrelevante.

\circuitoq{
\begin{circuitikz}[scale=1,transform shape,line width=0.8pt]
 \draw (0,0) node[left,font=\bfseries]{A} coordinate (A);
 \draw (5,0) node[right,font=\bfseries]{B} coordinate (B);
 \coordinate (C) at (2.5,1.6); \coordinate (D) at (2.5,-1.6);
 \draw (A) to[R,l={$\SI{4}{\ohm}$}] (C);
 \draw (C) to[R,l={$\SI{8}{\ohm}$}] (B);
 \draw (A) to[R,l_={$\SI{4}{\ohm}$}] (D);
 \draw (D) to[R,l_={$\SI{8}{\ohm}$}] (B);
 \draw (C) to[R,l={$\SI{5}{\ohm}$},color=Vcol] (D);
\end{circuitikz}}

\resposta{$R_{AB}=\SI{6}{\ohm}$; a ponte está equilibrada}


\end{questao}
\begin{questao}[3]{}

Dispõe-se de um número ilimitado de resistores de $\SI{10}{\ohm}$. Projete
associações que resultem em (a) $\SI{15}{\ohm}$; (b) $\SI{25}{\ohm}$;
(c) $\SI{4}{\ohm}$, usando o \textbf{menor número possível} de resistores em cada
caso.
\resposta{(a) 3 resistores; (b) 4 resistores; (c) 4 resistores}


\end{questao}

\blocodesafio
\begin{questao}[4]{}

A figura mostra uma rede em escada \textbf{infinita}, em que todos os resistores
valem $R$. Determine a resistência equivalente $\Req$ entre os terminais A e B.

\circuitoq{
\begin{circuitikz}[scale=0.9,transform shape,line width=0.8pt]
 \draw (0,2) node[left,font=\bfseries]{A} to[short,o-] (0.8,2);
 \draw (0.8,2) to[R,l=$R$] (2.8,2) coordinate(a1);
 \draw (a1) to[R,l=$R$] (2.8,0) coordinate(b1);
 \draw (a1) to[R,l=$R$] (4.8,2) coordinate(a2);
 \draw (a2) to[R,l=$R$] (4.8,0) coordinate(b2);
 \draw (a2) to[R,l=$R$] (6.8,2) coordinate(a3);
 \draw (a3) to[R,l=$R$] (6.8,0) coordinate(b3);
 \draw[dotted,line width=1.2pt] (6.9,2) -- (8.0,2);
 \draw[dotted,line width=1.2pt] (6.9,0) -- (8.0,0);
 \node[font=\footnotesize] at (8.4,1) {$\cdots\infty$};
 \draw (0.8,0) to[short,o-] (0,0) node[left,font=\bfseries]{B};
 \draw (0.8,2)--(0.8,2); \draw (0.8,0)--(2.8,0);
 \draw (2.8,0)--(4.8,0); \draw (4.8,0)--(6.8,0);
\end{circuitikz}}

\resposta{$\Req = \dfrac{R(1+\sqrt5)}{2} \approx 1{,}618\,R$}


\end{questao}
\begin{questao}[4]{}

Quantas seções da escada da questão anterior são necessárias para que a
resistência de entrada difira do valor infinito em menos de
$\SI{0.1}{\percent}$? Considere $R = \SI{1}{\ohm}$ e use a recorrência
$\Req^{(n)} = R + \big(R \parallel \Req^{(n-1)}\big)$, com
$\Req^{(1)} = R$.

\begin{table}[H]
\centering\small\sffamily
\begin{tabular}{@{}c c c@{}}
\toprule
\textbf{$n$ (seções)} & \textbf{$\Req$ ($\Omega$)} & \textbf{erro relativo}\\
\midrule
1  & 1,000000 & \SI{38}{\percent}\\
2  & 1,500000 & \SI{7.3}{\percent}\\
3  & 1,600000 & \SI{1.1}{\percent}\\
5  & 1,617647 & \SI{0.024}{\percent}\\
10 & 1,618034 & $<\SI{0.001}{\percent}$\\
\bottomrule
\end{tabular}
\caption{Convergência da escada finita para o valor infinito.}
\end{table}

\resposta{$n = 5$ seções já bastam}


\end{questao}
\begin{questao}[4]{}

Considere agora uma escada infinita em que os resistores \emph{em série} valem $R$
e os \emph{em paralelo} valem $2R$ --- a chamada rede \textbf{R-2R}, usada em
conversores digital-analógico (DAC).

\circuitoq{
\begin{circuitikz}[scale=0.9,transform shape,line width=0.8pt]
 \draw (0,2) node[left,font=\bfseries]{A} to[short,o-] (0.8,2);
 \draw (0.8,2) to[R,l=$R$] (2.8,2) coordinate(a1);
 \draw (a1) to[R,l=$2R$,color=Rcol] (2.8,0);
 \draw (a1) to[R,l=$R$] (4.8,2) coordinate(a2);
 \draw (a2) to[R,l=$2R$,color=Rcol] (4.8,0);
 \draw (a2) to[R,l=$R$] (6.8,2) coordinate(a3);
 \draw (a3) to[R,l=$2R$,color=Rcol] (6.8,0);
 \draw[dotted,line width=1.2pt] (6.9,2) -- (8.0,2);
 \draw[dotted,line width=1.2pt] (6.9,0) -- (8.0,0);
 \node[font=\footnotesize] at (8.4,1) {$\cdots\infty$};
 \draw (0.8,0) to[short,o-] (0,0) node[left,font=\bfseries]{B};
 \draw (0.8,0)--(6.8,0);
\end{circuitikz}}

\begin{itensq}
\item Determine $\Req$ entre A e B.
\item Explique por que esse resultado torna a rede R-2R especialmente útil.
\end{itensq}
\resposta{(a) $\Req = 2R$ (exato); (b) ver comentário}


\end{questao}
\begin{questao}[4]{}

Generalize: para uma escada infinita com resistores $R_s$ em série e $R_p$ em
paralelo, obtenha a expressão de $\Req$ e determine a condição sobre
$R_p$ para que $\Req$ seja um múltiplo \emph{racional} de $R_s$.
\resposta{$\Req = \dfrac{R_s + \sqrt{R_s^2+4R_sR_p}}{2}$; racional se
$R_s^2+4R_sR_p$ for quadrado perfeito}


\end{questao}
\begin{questao}[4]{}

A figura mostra uma rede \textbf{infinita} em que cada estágio $n$ é formado por
$(n+1)$ resistores \emph{em paralelo}, ligados entre dois barramentos
consecutivos. Os valores seguem o padrão do \textbf{triângulo de Pascal}: no
estágio $n$, os resistores valem $\dfrac{1}{\binom{n}{k}}\,\si{\ohm}$, com
$k = 0,1,\dots,n$.

\circuitoq{
\begin{circuitikz}[scale=0.82,transform shape,line width=0.75pt,font=\footnotesize]
 
 \foreach \x/\h in {0/0.6, 2.4/1.2, 4.8/1.8, 7.2/2.4}
   \draw[cinza] (\x,-\h) -- (\x,\h);
 \draw[cinza] (9.6,-3.0) -- (9.6,3.0);
 \node[left,font=\bfseries\normalsize] at (-0.35,0) {A};
 \draw (-0.35,0)--(0,0);
 
 \foreach \y in {-0.55,0.55}
   \draw (0,\y) to[R,l={$1$}] (2.4,\y);
 
 \foreach \y/\lab in {-1.1/{$1$}, 0/{$\frac12$}, 1.1/{$1$}}
   \draw (2.4,\y) to[R,l={\lab}] (4.8,\y);
 
 \foreach \y/\lab in {-1.65/{$1$}, -0.55/{$\frac13$}, 0.55/{$\frac13$}, 1.65/{$1$}}
   \draw (4.8,\y) to[R,l={\lab}] (7.2,\y);
 
 \foreach \y/\lab in {-2.2/{$1$}, -1.1/{$\frac14$}, 0/{$\frac16$}, 1.1/{$\frac14$}, 2.2/{$1$}}
   \draw (7.2,\y) to[R,l={\lab}] (9.6,\y);
 
 \foreach \y in {-2.6,-1.3,0,1.3,2.6}
   \draw[dotted,line width=1pt] (9.7,\y) -- (10.9,\y);
 \node[font=\normalsize] at (11.4,0) {$\infty$};
 \node[cinza,font=\scriptsize] at (1.2,-1.7) {estágio 1};
 \node[cinza,font=\scriptsize] at (3.6,-2.2) {2};
 \node[cinza,font=\scriptsize] at (6.0,-2.6) {3};
 \node[cinza,font=\scriptsize] at (8.4,-3.0) {4};
\end{circuitikz}}

Determine a resistência equivalente $R_{A\infty}$ vista do terminal A.
\resposta{$R_{A\infty} = \SI{1}{\ohm}$ (exatamente)}


\end{questao}

\blocofixacao
\begin{questao}[1]{}

$R_1=\SI{10}{\ohm}$ está em série com o paralelo de $R_2=\SI{20}{\ohm}$ e
$R_3=\SI{20}{\ohm}$. Determine $R_{eq}$.
\resposta{$R_{eq}=\SI{20}{\ohm}$}


\end{questao}
\begin{questao}[1]{}

Calcule $R_{eq}$ de $(\SI{10}{\ohm}\parallel\SI{10}{\ohm})$ em série com
$(\SI{30}{\ohm}\parallel\SI{60}{\ohm})$.
\resposta{$R_{eq}=\SI{25}{\ohm}$}


\end{questao}
\begin{questao}[1]{}

Dois resistores estão ligados \emph{aos mesmos dois nós} do circuito. Eles
estão:
\begin{alternativas}[2]
\item em série
\item em paralelo
\item em série ou paralelo, conforme o desenho
\item nem em série nem em paralelo
\end{alternativas}
\resposta{(b)}


\end{questao}
\begin{questao}[1]{}

A estratégia correta para reduzir uma rede mista é:
\begin{alternativas}[2]
\item começar sempre pelos resistores mais próximos da fonte
\item reduzir dos blocos mais internos para os mais externos
\item somar todas as resistências e depois corrigir
\item reduzir os paralelos antes de qualquer série, sempre
\end{alternativas}
\resposta{(b)}


\end{questao}
\begin{questao}[1]{}

Determine $R_{eq}$ de $(\SI{10}{\ohm}+\SI{20}{\ohm})$ em paralelo com
$(\SI{15}{\ohm}+\SI{15}{\ohm})$.
\resposta{$R_{eq}=\SI{15}{\ohm}$}


\end{questao}
\begin{questao}[1]{}

$R_1=\SI{12}{\ohm}$ está em série com o paralelo entre $R_2=\SI{6}{\ohm}$ e a
série $R_3=\SI{4}{\ohm}$ com $R_4=\SI{8}{\ohm}$. Determine $R_{eq}$.
\resposta{$R_{eq}=\SI{16}{\ohm}$}


\end{questao}
\begin{questao}[1]{}

Uma fonte de $\SI{24}{\volt}$ alimenta a rede
$\SI{10}{\ohm}+(\SI{20}{\ohm}\parallel\SI{20}{\ohm})$. Determine a corrente
total e a tensão sobre o bloco paralelo.
\resposta{$I=\SI{1.2}{\ampere}$; $U_{par}=\SI{12}{\volt}$}


\end{questao}
\begin{questao}[1]{}

Em uma rede mista, a corrente que atravessa o resistor imediatamente ligado ao
terminal da fonte é:
\begin{alternativas}[2]
\item igual à corrente total
\item metade da corrente total
\item dividida entre os ramos internos
\item indeterminada sem mais dados
\end{alternativas}
\resposta{(a)}


\end{questao}

\blocoaplicacao
\begin{questao}[2]{}

Uma fonte de $\SI{24}{\volt}$ alimenta $R_1=\SI{10}{\ohm}$ em série com
$R_2=\SI{20}{\ohm}$ e $R_3=\SI{20}{\ohm}$ em paralelo. Determine todas as
correntes e tensões.
\resposta{$I_t=\SI{1.2}{\ampere}$; $U_1=\SI{12}{\volt}$;
$U_{23}=\SI{12}{\volt}$; $I_2=I_3=\SI{0.6}{\ampere}$}


\end{questao}
\begin{questao}[2]{}

No circuito anterior, $R_3$ é trocado por $\SI{60}{\ohm}$. Recalcule
$R_{eq}$, a corrente total e as correntes de ramo.
\resposta{$R_{eq}=\SI{25}{\ohm}$; $I_t=\SI{0.96}{\ampere}$;
$I_2=\SI{0.72}{\ampere}$, $I_3=\SI{0.24}{\ampere}$}


\end{questao}
\begin{questao}[2]{}

Na rede $\SI{10}{\ohm}+(\SI{20}{\ohm}\parallel\SI{20}{\ohm})$ sob
$\SI{24}{\volt}$, calcule a potência de cada resistor e confronte com a
potência da fonte.
\resposta{$\SI{14.4}{\watt}$; $\SI{7.2}{\watt}$; $\SI{7.2}{\watt}$; total
$\SI{28.8}{\watt}$}


\end{questao}
\begin{questao}[2]{}

Determine $R_{eq}$ de $(\SI{5}{\ohm}+\SI{15}{\ohm})$ em paralelo com
$(\SI{10}{\ohm}+\SI{10}{\ohm})$, e a corrente de cada ramo sob
$\SI{30}{\volt}$.
\resposta{$R_{eq}=\SI{10}{\ohm}$; $I_t=\SI{3}{\ampere}$; $\SI{1.5}{\ampere}$
em cada ramo}


\end{questao}
\begin{questao}[2]{}

Uma rede é formada por $R_1=\SI{8}{\ohm}$ em série com o paralelo de
$\SI{24}{\ohm}$ e $R$. Sabendo que $R_{eq}=\SI{24}{\ohm}$, determine $R$.
\resposta{$R=\SI{48}{\ohm}$}


\end{questao}
\begin{questao}[2]{}

Determine $R_{eq}$ de $R_1=\SI{10}{\ohm}$ em série com o paralelo entre
$(\SI{20}{\ohm}+\SI{20}{\ohm})$ e $\SI{10}{\ohm}$.
\resposta{$R_{eq}=\SI{18}{\ohm}$}


\end{questao}
\begin{questao}[2]{}

Na rede $\SI{10}{\ohm}+(\SI{20}{\ohm}\parallel\SI{20}{\ohm})$ sob
$\SI{24}{\volt}$, o nó entre $R_1$ e o bloco paralelo é chamado $A$, e o
terminal negativo da fonte é o terra. Determine $V_A$.
\resposta{$V_A=\SI{12}{\volt}$}


\end{questao}
\begin{questao}[2]{}

Na mesma rede, um dos resistores de $\SI{20}{\ohm}$ sofre
\textbf{curto-circuito}. Recalcule $R_{eq}$, a corrente total e a potência de
$R_1$.
\resposta{$R_{eq}=\SI{10}{\ohm}$; $I_t=\SI{2.4}{\ampere}$;
$P_1=\SI{57.6}{\watt}$}


\end{questao}
\begin{questao}[2]{}

Na mesma rede, agora um dos resistores de $\SI{20}{\ohm}$ \textbf{abre}.
Recalcule $R_{eq}$, a corrente total e a potência do resistor remanescente.
\resposta{$R_{eq}=\SI{30}{\ohm}$; $I_t=\SI{0.8}{\ampere}$;
$P=\SI{12.8}{\watt}$}


\end{questao}
\begin{questao}[2]{}

Duas redes têm $R_{eq}=\SI{20}{\ohm}$: (A) $\SI{10}{\ohm}$ em série com
$\SI{20}{\ohm}\parallel\SI{20}{\ohm}$; (B) $\SI{30}{\ohm}\parallel\SI{60}{\ohm}$.
Sob $\SI{60}{\volt}$, compare a potência total e a do componente mais
solicitado.
\resposta{Total $\SI{180}{\watt}$ em ambas; pior componente:
$\SI{90}{\watt}$ em (A) e $\SI{120}{\watt}$ em (B)}


\end{questao}
\begin{questao}[2]{}

Uma fonte de $\EMF=\SI{24}{\volt}$ e $r=\SI{2}{\ohm}$ alimenta a rede mista de
$\SI{10}{\ohm}$. Determine a corrente, a tensão nos terminais e a potência
entregue à rede.
\resposta{$I=\SI{2}{\ampere}$; $V=\SI{20}{\volt}$; $P=\SI{40}{\watt}$}


\end{questao}
\begin{questao}[2]{}

Em uma rede mista de $R_{eq}=\SI{16}{\ohm}$ sob $\SI{32}{\volt}$, verificou-se
$P_{total}=\SI{64}{\watt}$. A medição é consistente? Justifique de duas
maneiras.
\resposta{Sim: $P=V^{2}/R_{eq}=VI=\SI{64}{\watt}$}


\end{questao}
\begin{questao}[2]{}

Na rede $R_1=\SI{12}{\ohm}$ em série com $R_2=\SI{6}{\ohm}\parallel
(R_3=\SI{4}{\ohm}+R_4=\SI{8}{\ohm})$, sob $\SI{32}{\volt}$, determine a
corrente em $R_3$.
\resposta{$I_3=\SI{0.667}{\ampere}$}


\end{questao}

\blocoaprofundamento
\begin{questao}[3]{}

Uma rede consiste em $(R_1+R_2)$ em paralelo com $(R_3+R_4)$.
\begin{itensq}
\item Obtenha a expressão de $R_{eq}$.
\item Avalie para $R_1=\SI{8}{\ohm}$, $R_2=\SI{16}{\ohm}$,
      $R_3=R_4=\SI{12}{\ohm}$.
\item Analise o caso-limite $R_3+R_4\gg R_1+R_2$.
\end{itensq}
\resposta{(a) $R_{eq}=\dfrac{(R_1+R_2)(R_3+R_4)}{R_1+R_2+R_3+R_4}$;
(b) $\SI{12}{\ohm}$; (c) $R_{eq}\to R_1+R_2$}


\end{questao}
\begin{questao}[3]{}

Um divisor formado por $R_1=\SI{140}{\ohm}$ e $R_2=\SI{100}{\ohm}$ é alimentado
por $\SI{12}{\volt}$.
\begin{itensq}
\item Determine a tensão de saída em vazio.
\item Recalcule com uma carga de $\SI{1}{\kilo\ohm}$ conectada à saída.
\item Determine o erro percentual introduzido pela carga.
\end{itensq}
\resposta{(a) $\SI{5.00}{\volt}$; (b) $\SI{4.72}{\volt}$; (c) $-5{,}5\%$}


\end{questao}
\begin{questao}[3]{}

Uma rede tem $R_1=\SI{6}{\ohm}$ em série com $(\SI{12}{\ohm}\parallel
\SI{12}{\ohm})$, e este em série com $(\SI{8}{\ohm}\parallel\SI{8}{\ohm})$, sob
$\SI{32}{\volt}$. Determine a potência de cada um dos cinco resistores.
\resposta{$\SI{24}{\watt}$ em $R_1$; $\SI{12}{\watt}$ em cada
$\SI{12}{\ohm}$; $\SI{8}{\watt}$ em cada $\SI{8}{\ohm}$}


\end{questao}
\begin{questao}[3]{}

Um reostato de $\SI{0}{\ohm}$ a $\SI{60}{\ohm}$ está em paralelo com
$\SI{30}{\ohm}$, e o conjunto em série com $\SI{20}{\ohm}$, sob
$\SI{60}{\volt}$.
\begin{itensq}
\item Determine a faixa de $R_{eq}$.
\item Determine a faixa da corrente total.
\item Explique por que a faixa é limitada.
\end{itensq}
\resposta{(a) de $\SI{20}{\ohm}$ a $\SI{40}{\ohm}$; (b) de $\SI{1.5}{\ampere}$
a $\SI{3}{\ampere}$}


\end{questao}
\begin{questao}[3]{}

Na rede $\SI{10}{\ohm}+(\SI{20}{\ohm}\parallel\SI{20}{\ohm})$, todos os
resistores têm potência nominal de $\SI{10}{\watt}$. Determine a tensão máxima
aplicável e identifique o componente crítico.
\resposta{$V_{\max}=\SI{20}{\volt}$; o crítico é $R_1$}


\end{questao}
\begin{questao}[3]{}

Obtenha $\SI{24}{\ohm}$ dispondo apenas de resistores de $\SI{36}{\ohm}$.
Determine a associação de \textbf{menor} número de unidades e prove que não há
solução com duas.
\resposta{Três unidades: $(\SI{36}{\ohm}+\SI{36}{\ohm})\parallel\SI{36}{\ohm}$}


\end{questao}
\begin{questao}[3]{}

Em uma rede mista alimentada por fonte ideal, um resistor de um ramo paralelo
abre. Descreva um procedimento sistemático para localizar a falha usando apenas
um voltímetro, e aplique-o à rede
$\SI{10}{\ohm}+(\SI{20}{\ohm}\parallel\SI{20}{\ohm})$ sob $\SI{24}{\volt}$.
\resposta{Comparar as tensões medidas com as previstas; o bloco defeituoso
apresenta tensão \emph{maior} que a esperada}


\end{questao}
\begin{questao}[3]{}

Uma rede mista de $R_{eq}=\SI{20}{\ohm}$ é alimentada por fonte de
$\EMF=\SI{60}{\volt}$ com $r=\SI{5}{\ohm}$.
\begin{itensq}
\item Determine a corrente, a tensão na rede e o rendimento.
\item Determine o valor de $R_{eq}$ que maximizaria a potência na rede.
\item Compare o rendimento nas duas situações.
\end{itensq}
\resposta{(a) $\SI{2.4}{\ampere}$, $\SI{48}{\volt}$, $80\%$;
(b) $R_{eq}=\SI{5}{\ohm}$; (c) $80\%$ contra $50\%$}


\end{questao}
\begin{questao}[3]{}

Uma fonte de \textbf{corrente} de $\SI{3}{\ampere}$ alimenta $R_1=\SI{20}{\ohm}$
em paralelo com o ramo formado por $R_2=\SI{10}{\ohm}$ em série com
$(\SI{20}{\ohm}\parallel\SI{20}{\ohm})$.
\begin{itensq}
\item Determine $R_{eq}$ e a tensão da associação.
\item Obtenha a corrente em cada um dos cinco resistores.
\item Verifique o balanço de potência.
\end{itensq}
\resposta{(a) $\SI{10}{\ohm}$ e $\SI{30}{\volt}$;
(b) $\SI{1.5}{\ampere}$ em $R_1$ e em $R_2$, $\SI{0.75}{\ampere}$ em cada
$\SI{20}{\ohm}$ do bloco; (c) $\SI{90}{\watt}$}


\end{questao}

\blocodesafio
\begin{questao}[4]{}

\textbf{(Projeto de divisor carregado.)} Deseja-se obter $\SI{5}{\volt}$ a
partir de $\SI{12}{\volt}$ para alimentar uma carga fixa de
$\SI{1}{\kilo\ohm}$.
\begin{itensq}
\item Deduza o critério de "corrente de sangria" e determine $R_1,R_2$ para
      erro inferior a $2\%$.
\item Projete o divisor \emph{compensado}, que entregue exatamente
      $\SI{5}{\volt}$ com a carga presente.
\item Compare as duas soluções em potência dissipada e em robustez.
\end{itensq}
\resposta{(a) $R_2\approx\SI{25}{\ohm}$, $R_1=\SI{35}{\ohm}$ (erro $1{,}4\%$);
(b) $R_1=\SI{140}{\ohm}$, $R_2=\SI{111}{\ohm}$;
(c) a compensada dissipa muito menos, mas só é exata para aquela carga}


\end{questao}
\begin{questao}[4]{}

\textbf{(Enumeração de topologias.)} Dispõe-se de quatro resistores iguais de
$\SI{100}{\ohm}$.
\begin{itensq}
\item Determine todos os valores distintos de $R_{eq}$ obteníveis usando os
      quatro.
\item Identifique os dois arranjos que resultam em $\SI{100}{\ohm}$ e
      compare a potência do componente mais solicitado sob
      $\SI{100}{\volt}$.
\item Justifique por que os extremos são $\SI{25}{\ohm}$ e
      $\SI{400}{\ohm}$.
\end{itensq}
\resposta{(a) $25$, $40$, $60$, $75$, $100$, $133{,}3$, $166{,}7$, $250$ e
$\SI{400}{\ohm}$; (b) $\SI{25}{\watt}$ em ambos}


\end{questao}
\begin{questao}[4]{}

\textbf{(Projeto sob duas restrições.)} Uma rede deve apresentar
$R_{eq}=\SI{20}{\ohm}$ e dissipar $\SI{45}{\watt}$.
\begin{itensq}
\item Determine a tensão e a corrente nominais.
\item Compare três topologias quanto à potência do componente mais solicitado.
\item Escolha a melhor e justifique.
\end{itensq}
\resposta{(a) $\SI{30}{\volt}$ e $\SI{1.5}{\ampere}$; (b) pior componente:
$\SI{30}{\watt}$, $\SI{22.5}{\watt}$ e $\SI{11.25}{\watt}$;
(c) a matriz $2\times2$ de $\SI{20}{\ohm}$}


\end{questao}
\begin{questao}[4]{}

\textbf{(Diagnóstico quantitativo.)} Uma rede
$R_1=\SI{10}{\ohm}+(R_2=\SI{20}{\ohm}\parallel R_3=\SI{20}{\ohm})$ opera sob
$\SI{24}{\volt}$. Mediu-se $I_t=\SI{2.4}{\ampere}$.
\begin{itensq}
\item Mostre que a rede está defeituosa e identifique o tipo de falha.
\item Localize o componente e prove que a localização é única.
\item Proponha uma medição de confirmação.
\end{itensq}
\resposta{$R_{eq}$ medida $=\SI{10}{\ohm}$ contra $\SI{20}{\ohm}$ prevista;
o bloco paralelo está em curto}


\end{questao}
\begin{questao}[4]{}

\textbf{(Redução por simetria.)} Seis resistores de $\SI{6}{\ohm}$ formam as
arestas de um \textbf{tetraedro} regular (quatro vértices, todos ligados dois a
dois).
\begin{itensq}
\item Determine $R_{AB}$ entre dois vértices quaisquer, usando simetria.
\item Determine a corrente em cada aresta para $I_t=\SI{4}{\ampere}$.
\item Recalcule $R_{AB}$ se a aresta $AB$ for removida.
\end{itensq}
\resposta{(a) $R_{AB}=\SI{3}{\ohm}$; (c) $\SI{6}{\ohm}$}


\end{questao}

\gabarito[3.3 Circuitos mistos]

\subsection{Transformação estrela-triângulo}

\blocofixacao
\begin{questao}[1]{}

Considere uma configuração \textbf{estrela} com $R_1=\SI{3}{\ohm}$,
$R_2=\SI{6}{\ohm}$ e $R_3=\SI{9}{\ohm}$. Determine os resistores da configuração
\textbf{triângulo} equivalente.
\resposta{$R_{AB}=\SI{11}{\ohm}$, $R_{BC}=\SI{33}{\ohm}$, $R_{CA}=\SI{16.5}{\ohm}$}


\end{questao}
\begin{questao}[1]{}

Um circuito \textbf{triângulo} possui $R_{AB}=\SI{10}{\ohm}$,
$R_{BC}=\SI{5}{\ohm}$ e $R_{CA}=\SI{15}{\ohm}$. Calcule os resistores da
\textbf{estrela} equivalente.
\resposta{$R_A=\SI{5}{\ohm}$, $R_B=\SI{1.67}{\ohm}$, $R_C=\SI{2.5}{\ohm}$}


\end{questao}
\begin{questao}[1]{}

Explique, com suas palavras, por que a transformação estrela-triângulo é útil e
cite um exemplo prático de circuito em que ela é necessária.
\resposta{Resposta dissertativa --- ver comentário}


\end{questao}

\blocoaplicacao
\begin{questao}[2]{}

No triângulo da figura, com $R_{AB}=\SI{9}{\ohm}$ (lado direto entre A e B),
$R_{AC}=\SI{6}{\ohm}$ e $R_{CB}=\SI{3}{\ohm}$, determine a resistência
equivalente entre A e B.

\circuitoq{
\begin{circuitikz}[scale=1,transform shape,line width=0.8pt]
 \draw (0,0) node[below left,font=\bfseries]{A} coordinate(A)
       to[R,l=$\SI{6}{\ohm}$] (2,2.4) node[above,font=\bfseries]{C} coordinate(C)
       to[R,l=$\SI{3}{\ohm}$] (4,0) node[below right,font=\bfseries]{B} coordinate(B);
 \draw (A) to[R,l_=$\SI{9}{\ohm}$] (B);
\end{circuitikz}}

\resposta{$R_{AB} = \SI{4.5}{\ohm}$}


\end{questao}
\begin{questao}[2]{}

Três resistores iguais de $\SI{2}{\ohm}$ estão ligados em \textbf{estrela}.
Convertendo-os em \textbf{triângulo} e aplicando uma fonte de $\SI{12}{\volt}$
entre dois dos terminais (A e B), determine a corrente total fornecida pela fonte.
\resposta{$I = \SI{3}{\ampere}$}


\end{questao}
\begin{questao}[2]{}

Um triângulo entre A, B e C tem $R_{AB}=\SI{12}{\ohm}$, $R_{BC}=\SI{6}{\ohm}$ e
$R_{CA}=\SI{6}{\ohm}$.
\begin{itensq}
\item Converta para estrela.
\item Calcule a resistência equivalente entre A e B após a conversão.
\end{itensq}
\resposta{(a) $R_A=\SI{3}{\ohm}$, $R_B=\SI{3}{\ohm}$, $R_C=\SI{1.5}{\ohm}$;
(b) $R_{AB}=\SI{6}{\ohm}$}


\end{questao}

\blocoaprofundamento
\begin{questao}[3]{}

A figura mostra uma \textbf{ponte de Wheatstone}. Verifique se ela está
equilibrada e, caso esteja, determine a resistência equivalente entre A e B.

\circuitoq{
\begin{circuitikz}[scale=1,transform shape,line width=0.8pt]
 \draw (0,0) node[left,font=\bfseries]{A} coordinate(A);
 \draw (5,0) node[right,font=\bfseries]{B} coordinate(B);
 \coordinate (C) at (2.5,1.7);
 \coordinate (D) at (2.5,-1.7);
 \draw (A) to[R,l=$\SI{10}{\ohm}$] (C);
 \draw (C) to[R,l=$\SI{20}{\ohm}$] (B);
 \draw (A) to[R,l_=$\SI{15}{\ohm}$] (D);
 \draw (D) to[R,l_=$\SI{30}{\ohm}$] (B);
 \draw (C) to[R,l=$\SI{8}{\ohm}$] (D);
 \node[font=\footnotesize] at (3.1,-0.27) {galv.};
\end{circuitikz}}

\resposta{Equilibrada; $R_{AB} = \SI{18}{\ohm}$}


\end{questao}
\begin{questao}[3]{}

Na ponte da figura, $R_{AC}=\SI{6}{\ohm}$, $R_{CB}=\SI{18}{\ohm}$,
$R_{AD}=\SI{18}{\ohm}$, $R_{DB}=\SI{6}{\ohm}$ e o resistor central (ponte) vale
$R_{CD}=\SI{6}{\ohm}$. Determine a resistência equivalente entre A e B.

\circuitoq{
\begin{circuitikz}[scale=1,transform shape,line width=0.8pt]
 \draw (0,0) node[left,font=\bfseries]{A} coordinate(A);
 \draw (5,0) node[right,font=\bfseries]{B} coordinate(B);
 \coordinate (C) at (2.5,1.7);
 \coordinate (D) at (2.5,-1.7);
 \draw (A) to[R,l=$\SI{6}{\ohm}$] (C);
 \draw (C) to[R,l=$\SI{18}{\ohm}$] (B);
 \draw (A) to[R,l_=$\SI{18}{\ohm}$] (D);
 \draw (D) to[R,l_=$\SI{6}{\ohm}$] (B);
 \draw (C) to[R,l=$\SI{6}{\ohm}$] (D);
\end{circuitikz}}

\resposta{$R_{AB} = \SI{10}{\ohm}$}


\end{questao}
\begin{questao}[3]{}

Três resistores formam um triângulo entre os nós X, Y e Z, com
$R_{XY}=\SI{9}{\ohm}$, $R_{YZ}=\SI{12}{\ohm}$ e $R_{ZX}=\SI{15}{\ohm}$. Além
disso, cada nó é ligado ao terra por um resistor de $\SI{6}{\ohm}$. Descreva o
procedimento para obter a resistência entre o nó X e o terra e monte a expressão
final (sem exigir o valor numérico fechado).
\resposta{Ver roteiro comentado}


\end{questao}
\begin{questao}[3]{}

Uma ponte de Wheatstone é usada para medir uma resistência desconhecida $R_x$.
Os braços conhecidos são $R_1 = \SI{100}{\ohm}$, $R_2 = \SI{200}{\ohm}$ e
$R_3$ é uma década ajustável. O equilíbrio é obtido com
$R_3 = \SI{347}{\ohm}$.
\begin{itensq}
\item Determine $R_x$.
\item Se $R_1$ e $R_2$ têm tolerância de $\SI{0.1}{\percent}$ cada e $R_3$ de
      $\SI{0.05}{\percent}$, estime a incerteza relativa de $R_x$.
\item Por que a ponte é mais precisa que medir $R_x$ com um ohmímetro comum?
\end{itensq}
\resposta{(a) $R_x = \SI{694}{\ohm}$; (b) $\approx\SI{0.25}{\percent}$;
(c) ver comentário}


\end{questao}

\blocofixacao
\begin{questao}[1]{}

Um triângulo equilátero tem três lados de $\SI{30}{\ohm}$. Determine a estrela
equivalente.
\resposta{Três braços de $\SI{10}{\ohm}$}


\end{questao}
\begin{questao}[1]{}

Uma estrela equilibrada tem três braços de $\SI{5}{\ohm}$. Determine o
triângulo equivalente.
\resposta{Três lados de $\SI{15}{\ohm}$}


\end{questao}
\begin{questao}[1]{}

Escreva, \textbf{sem calcular}, a expressão do braço $R_B$ de uma estrela
equivalente a um triângulo de lados $R_{AB}$, $R_{BC}$ e $R_{CA}$.
\resposta{$R_B=\dfrac{R_{AB}\,R_{BC}}{R_{AB}+R_{BC}+R_{CA}}$}


\end{questao}
\begin{questao}[1]{}

Compare estrela e triângulo quanto ao número de elementos e de nós.
\begin{alternativas}[2]
\item ambos têm 3 resistores; a estrela tem um nó interno adicional
\item a estrela tem 3 e o triângulo tem 4 resistores
\item ambos têm o mesmo número de nós
\item o triângulo tem um nó interno adicional
\end{alternativas}
\resposta{(a)}


\end{questao}
\begin{questao}[1]{}

Identifique qual das expressões é a conversão triângulo-estrela:
\begin{alternativas}[2]
\item $R=\dfrac{\text{produto de dois}}{\text{soma dos três}}$
\item $R=\dfrac{\text{soma dos produtos}}{\text{oposto}}$
\item $R=\dfrac{\text{soma dos três}}{3}$
\item $R=\dfrac{\text{produto dos três}}{\text{soma dos três}}$
\end{alternativas}
\resposta{(a)}


\end{questao}

\blocoaplicacao
\begin{questao}[2]{}

Um triângulo tem $R_{AB}=\SI{20}{\ohm}$, $R_{BC}=\SI{30}{\ohm}$ e
$R_{CA}=\SI{50}{\ohm}$. Determine a estrela equivalente.
\resposta{$R_A=\SI{10}{\ohm}$, $R_B=\SI{6}{\ohm}$, $R_C=\SI{15}{\ohm}$}


\end{questao}
\begin{questao}[2]{}

Converta de volta a estrela obtida ($10$, $6$ e $\SI{15}{\ohm}$) para
triângulo e verifique a reversibilidade.
\resposta{$20$, $30$ e $\SI{50}{\ohm}$ --- os valores originais}


\end{questao}
\begin{questao}[2]{}

Para o triângulo $20/30/\SI{50}{\ohm}$ e sua estrela equivalente, calcule a
resistência entre $A$ e $B$ (com $C$ aberto) nas duas configurações.
\resposta{$\SI{16}{\ohm}$ em ambas}


\end{questao}
\begin{questao}[2]{}

Uma estrela tem $R_A=\SI{4}{\ohm}$, $R_B=\SI{8}{\ohm}$ e
$R_C=\SI{12}{\ohm}$. Determine o triângulo equivalente e identifique o maior
lado.
\resposta{$R_{AB}=\SI{14.67}{\ohm}$, $R_{BC}=\SI{44}{\ohm}$,
$R_{CA}=\SI{22}{\ohm}$; o maior é $R_{BC}$}


\end{questao}
\begin{questao}[2]{}

Em um triângulo, o lado $R_{BC}$ é substituído por um \textbf{circuito aberto}.
Determine a estrela equivalente e interprete.
\resposta{$R_A=\dfrac{R_{AB}R_{CA}}{R_{AB}+R_{CA}}$, $R_B=R_C=0$ --- ou seja,
apenas dois resistores em série}


\end{questao}
\begin{questao}[2]{}

Em um triângulo, o lado $R_{AB}$ é substituído por um \textbf{curto-circuito}.
Determine a estrela equivalente.
\resposta{$R_A=R_B=0$ e $R_C=\dfrac{R_{BC}R_{CA}}{R_{BC}+R_{CA}}$}


\end{questao}
\begin{questao}[2]{}

Verifique numericamente, para o triângulo $20/30/\SI{50}{\ohm}$, a identidade
$R_{AB}R_{BC}+R_{BC}R_{CA}+R_{CA}R_{AB}
=\left(\sum R_\Delta\right)\left(\sum R_Y\right)$.
\resposta{$3100=100\times31$ \checkmark}


\end{questao}
\begin{questao}[2]{}

Em uma rede, três resistores formam um triângulo entre os nós $P$, $Q$ e $R$,
e há mais resistores ligados a esses três nós. Explique por que convém
converter o triângulo em estrela, e não o contrário.
\resposta{A estrela cria um nó interno isolado, liberando associações
série/paralelo}


\end{questao}
\begin{questao}[2]{}

Uma estrela tem braços $R_A=R_B=\SI{6}{\ohm}$ e $R_C=\SI{3}{\ohm}$. Determine o
triângulo equivalente e verifique pela resistência entre $A$ e $B$.
\resposta{$R_{AB}=\SI{24}{\ohm}$, $R_{BC}=R_{CA}=\SI{12}{\ohm}$}


\end{questao}

\blocoaprofundamento
\begin{questao}[3]{}

Uma ponte é alimentada por $\SI{84}{\volt}$ entre o nó $S$ e o terra. Do nó
$S$ partem $\SI{6}{\ohm}$ até $A$ e $\SI{12}{\ohm}$ até $B$; de $A$ sai
$\SI{12}{\ohm}$ ao terra e de $B$ sai $\SI{6}{\ohm}$ ao terra; entre $A$ e
$B$ há $\SI{6}{\ohm}$.
\begin{itensq}
\item Verifique que a ponte está desequilibrada.
\item Converta o triângulo $S$--$A$--$B$ em estrela e determine $R_{eq}$.
\item Determine a corrente total.
\end{itensq}
\resposta{(b) $R_{eq}=\SI{8.4}{\ohm}$; (c) $I=\SI{10}{\ampere}$}


\end{questao}
\begin{questao}[3]{}

Resolva a mesma ponte por \textbf{análise nodal} e confronte com o resultado da
conversão.
\resposta{$V_A=\SI{48}{\volt}$, $V_B=\SI{36}{\volt}$; $R_{eq}=\SI{8.4}{\ohm}$
--- confere}


\end{questao}
\begin{questao}[3]{}

Um triângulo tem um lado muito menor que os outros dois:
$R_{AB}=\SI{0.1}{\ohm}$, $R_{BC}=R_{CA}=\SI{100}{\ohm}$.
\begin{itensq}
\item Determine a estrela equivalente.
\item Compare $R_A+R_B$ com $R_{AB}$ e interprete.
\end{itensq}
\resposta{$R_A=R_B=\SI{49.98}{\milli\ohm}$, $R_C=\SI{49.98}{\ohm}$;
$R_A+R_B\approx R_{AB}$}


\end{questao}
\begin{questao}[3]{}

Uma estrela tem um braço muito pequeno: $R_A=R_B=\SI{10}{\ohm}$ e
$R_C=\SI{0.1}{\ohm}$.
\begin{itensq}
\item Determine o triângulo equivalente.
\item Analise o comportamento quando $R_C\to0$.
\end{itensq}
\resposta{$R_{AB}=\SI{1020}{\ohm}$, $R_{BC}=R_{CA}=\SI{10.2}{\ohm}$;
$R_{AB}\to\infty$}


\end{questao}
\begin{questao}[3]{}

Discuta o condicionamento das fórmulas de conversão para resistores positivos.
\begin{itensq}
\item Identifique se há risco de cancelamento catastrófico.
\item Determine onde, na prática, a precisão realmente se perde.
\end{itensq}
\resposta{As fórmulas são bem condicionadas; a perda de precisão ocorre na
\emph{redução posterior} da rede}


\end{questao}
\begin{questao}[3]{}

Uma conversão produz um triângulo de $36{,}67$, $110$ e $\SI{55}{\ohm}$.
Escolha valores comerciais E24 e avalie o erro.
\begin{itensq}
\item Indique os valores adotados.
\item Converta de volta e compare com a estrela original de $10$, $20$ e
      $\SI{30}{\ohm}$.
\end{itensq}
\resposta{$36$, $110$ e $\SI{56}{\ohm}$; erros de $-0{,}2\%$, $-2{,}0\%$ e
$+1{,}7\%$ nos braços}


\end{questao}
\begin{questao}[3]{}

Compare o esforço de resolver a ponte desequilibrada por conversão
estrela-triângulo e por análise nodal.
\begin{itensq}
\item Enumere as etapas de cada método.
\item Indique o critério de escolha.
\end{itensq}
\resposta{Conversão: 3 fórmulas e reduções; nodal: sistema $2\times2$. A
conversão dá só $R_{eq}$; a nodal dá tudo}


\end{questao}
\begin{questao}[3]{}

Uma rede de três terminais é formada por resistores positivos em uma topologia
qualquer. Discuta se ela sempre admite equivalente em estrela.
\begin{itensq}
\item Estabeleça a condição.
\item Dê um exemplo em que a estrela equivalente teria braço negativo.
\end{itensq}
\resposta{A condição é a desigualdade triangular sobre as resistências entre
pares}


\end{questao}

\blocodesafio
\begin{questao}[4]{}

\textbf{(Dedução --- triângulo para estrela.)} Deduza as fórmulas de conversão a
partir das resistências entre terminais.
\begin{itensq}
\item Escreva as três equações de equivalência.
\item Resolva o sistema para $R_A$, $R_B$ e $R_C$.
\item Obtenha a forma final em função dos lados do triângulo.
\end{itensq}
\resposta{$R_A=\dfrac{R_{AB}R_{CA}}{R_{AB}+R_{BC}+R_{CA}}$ e análogas}


\end{questao}
\begin{questao}[4]{}

\textbf{(Dedução --- estrela para triângulo.)} Obtenha a transformação inversa.
\begin{itensq}
\item Deduza $R_{AB}$ em função dos braços.
\item Verifique que a composição das duas transformações é a identidade.
\item Explique por que o denominador é o braço \emph{oposto}.
\end{itensq}
\resposta{$R_{AB}=\dfrac{R_AR_B+R_BR_C+R_CR_A}{R_C}$}


\end{questao}
\begin{questao}[4]{}

\textbf{(Positividade.)} Prove que as transformações preservam a positividade
dos elementos.
\begin{itensq}
\item Mostre que $\Delta\to Y$ com lados positivos gera braços positivos.
\item Mostre que $Y\to\Delta$ com braços positivos gera lados positivos.
\item Determine o que ocorre quando um elemento tende a zero ou a infinito.
\end{itensq}
\resposta{Ambas as fórmulas são quocientes de quantidades positivas ---
positividade garantida}


\end{questao}
\begin{questao}[4]{}

\textbf{(Alcance da equivalência.)} Analise \emph{o que} a transformação
preserva e o que ela destrói.
\begin{itensq}
\item Prove que as resistências entre os três terminais são preservadas.
\item Determine se correntes e potências internas são preservadas.
\item Estabeleça o procedimento correto quando grandezas internas são
      necessárias.
\end{itensq}
\resposta{Preserva o comportamento nos três terminais; \emph{não} preserva nada
do interior}


\end{questao}
\begin{questao}[4]{}

\textbf{(Projeto com valores comerciais.)} Uma estrela de $10$, $20$ e
$\SI{30}{\ohm}$ deve ser substituída por um triângulo montado com resistores
E24.
\begin{itensq}
\item Determine o triângulo exato.
\item Escolha valores comerciais e avalie o erro na resistência entre $A$ e
      $B$.
\item Proponha uma estratégia para reduzir o erro.
\end{itensq}
\resposta{Exato: $36{,}67$, $110$ e $\SI{55}{\ohm}$; com $36/110/56$, o erro em
$R_{AB}$ é de cerca de $1\%$}


\end{questao}
\begin{questao}[4]{}

\textbf{(Ponte quase equilibrada.)} Uma ponte tem $R_1=\SI{100}{\ohm}$,
$R_2=\SI{100}{\ohm}$, $R_3=\SI{100}{\ohm}$ e
$R_4=\SI{101}{\ohm}$, com $R_g=\SI{100}{\ohm}$ e alimentação de
$\SI{10}{\volt}$.
\begin{itensq}
\item Estime a corrente do galvanômetro por análise nodal.
\item Discuta a viabilidade de obter esse valor por conversão
      estrela-triângulo.
\end{itensq}
\resposta{$I_g\approx\SI{-124}{\micro\ampere}$ (de $B$ para $A$); a conversão é
inadequada para esta grandeza}


\end{questao}
\begin{questao}[4]{}

\textbf{(Síntese.)} Uma rede de três terminais deve apresentar
$R_{AB}=\SI{20}{\ohm}$, $R_{BC}=\SI{30}{\ohm}$ e $R_{CA}=\SI{40}{\ohm}$
(medidas com o terceiro terminal aberto).
\begin{itensq}
\item Verifique a realizabilidade e determine a estrela que atende.
\item Determine o triângulo equivalente.
\item Escolha valores comerciais e avalie qual das duas topologias erra menos.
\end{itensq}
\resposta{$R_A=\SI{15}{\ohm}$, $R_B=\SI{5}{\ohm}$, $R_C=\SI{25}{\ohm}$;
triângulo $23$, $38{,}33$ e $\SI{115}{\ohm}$; a estrela é preferível}


\end{questao}
\begin{questao}[4]{}

\textbf{(Generalização.)} A transformação estrela-triângulo é o caso $n=3$ de
uma família mais ampla.
\begin{itensq}
\item Descreva a transformação estrela-polígono para $n$ terminais.
\item Determine quantos elementos há em cada configuração.
\item Explique por que a transformação inversa nem sempre existe para $n>3$.
\end{itensq}
\resposta{Estrela: $n$ braços; polígono: $n(n-1)/2$ lados. A inversa só existe
para $n=3$}


\end{questao}

\gabarito[3.4 Transformação estrela-triângulo]

\section{Divisor de Tensão e de Corrente}

\subsection{Divisor de tensão}

\blocofixacao
\begin{questao}[1]{}

Em um divisor de tensão formado por dois resistores em série, a maior parcela de
tensão aparece:
\begin{alternativas}
\item sobre o menor resistor.
\item sobre o maior resistor.
\item igualmente sobre os dois, sempre.
\item sobre o resistor mais próximo do terminal positivo.
\item sobre nenhum: a tensão fica toda na fonte.
\end{alternativas}
\resposta{(b)}


\end{questao}
\begin{questao}[1]{}

No circuito, três resistores em série de $\SI{3}{\ohm}$, $\SI{4}{\ohm}$ e
$\SI{5}{\ohm}$ estão ligados a uma fonte de $\SI{120}{\volt}$. Calcule a tensão
sobre cada resistor.

\circuitoq{
\begin{circuitikz}[scale=0.95,transform shape,line width=0.8pt]
 \draw (0,0) to[battery1,l_={$E=\SI{120}{\volt}$}] (0,3)
       to[R,l={$R_1=\SI{3}{\ohm}$}] (3.2,3)
       to[R,l={$R_2=\SI{4}{\ohm}$}] (3.2,0)
       to[R,l={$R_3=\SI{5}{\ohm}$}] (0,0);
 \draw[->,Icol,line width=1pt] (0.7,3.4)--(1.7,3.4)
       node[midway,above,font=\footnotesize]{$I$};
\end{circuitikz}}

\resposta{$U_1=\SI{30}{\volt}$, $U_2=\SI{40}{\volt}$, $U_3=\SI{50}{\volt}$}


\end{questao}
\begin{questao}[1]{}

Uma fonte de $\SI{100}{\volt}$ alimenta dois resistores em série,
$R_1=\SI{60}{\ohm}$ e $R_2=\SI{40}{\ohm}$. A tensão de saída, tomada sobre $R_2$,
vale:
\begin{alternativas}[3]
\item \SI{20}{\volt}
\item \SI{40}{\volt}
\item \SI{50}{\volt}
\item \SI{60}{\volt}
\item \SI{100}{\volt}
\end{alternativas}
\resposta{(b)}


\end{questao}
\begin{questao}[1]{}

Três resistores iguais estão em série sob $\SI{60}{\volt}$. A tensão sobre cada um
vale:
\begin{alternativas}[3]
\item \SI{10}{\volt}
\item \SI{15}{\volt}
\item \SI{20}{\volt}
\item \SI{30}{\volt}
\item \SI{60}{\volt}
\end{alternativas}
\resposta{(c)}


\end{questao}

\blocoaplicacao
\begin{questao}[2]{}

Deseja-se obter $\SI{5}{\volt}$ a partir de uma fonte de $\SI{15}{\volt}$ usando
um divisor de dois resistores. Sabendo que $R_1 = \SI{2}{\kilo\ohm}$ (ligado ao
terminal positivo), determine $R_2$ (a saída é tomada sobre $R_2$).
\resposta{$R_2 = \SI{1}{\kilo\ohm}$}


\end{questao}
\begin{questao}[2]{}

No circuito da figura, calcule a tensão de saída $U_{\text{out}}$ (a) em vazio e
(b) quando uma carga $R_L = \SI{1}{\kilo\ohm}$ é ligada aos terminais de saída.
Comente o resultado.

\circuitoq{
\begin{circuitikz}[scale=0.95,transform shape,line width=0.8pt]
 \draw (0,0) to[battery1,l_={$E=\SI{12}{\volt}$}] (0,3.4)
       to[R,l={$R_1=\SI{1}{\kilo\ohm}$}] (3.4,3.4) coordinate(a);
 \draw (a) to[R,l={$R_2=\SI{1}{\kilo\ohm}$}] (3.4,0) coordinate(b) -- (0,0);
 \draw (a) to[short,-o] (5,3.4) node[right]{$+$};
 \draw (b) to[short,-o] (5,0) node[right]{$-$};
 \draw[->] (5.3,3.0) -- (5.3,0.4) node[midway,right,font=\footnotesize]{$U_{\text{out}}$};
\end{circuitikz}}

\resposta{(a) \SI{6}{\volt}; (b) \SI{4}{\volt}}


\end{questao}
\begin{questao}[2]{}

Um potenciômetro de $\SI{10}{\kilo\ohm}$ é ligado como divisor de tensão a uma
fonte de $\SI{9}{\volt}$. Quando o cursor está a $\SI{30}{\percent}$ do curso (a
partir do terminal de referência), qual é a tensão de saída em vazio?
\resposta{\SI{2.7}{\volt}}


\end{questao}
\begin{questao}[2]{}

Para o circuito da figura (malha única com cinco resistores), calcule a corrente
e a tensão sobre o resistor de $\SI{8}{\ohm}$.

\circuitoq{
\begin{circuitikz}[scale=0.9,transform shape,line width=0.8pt]
 \draw (0,0) to[battery1,l_={$E=\SI{100}{\volt}$},invert] (0,3)
       to[R,l={$\SI{10}{\ohm}$}] (3,3)
       to[R,l={$\SI{6}{\ohm}$}] (6,3)
       to[R,l={$\SI{8}{\ohm}$}] (6,0)
       to[R,l={$\SI{12}{\ohm}$}] (3,0)
       to[R,l={$\SI{6}{\ohm}$}] (0,0);
 \draw[->,Icol,line width=1pt] (0.7,3.4)--(1.7,3.4)
       node[midway,above,font=\footnotesize]{$I$};
\end{circuitikz}}

\resposta{$I = \SI{2.38}{\ampere}$ (aprox.); $U_{8} \approx \SI{19}{\volt}$}


\end{questao}

\blocoaprofundamento
\begin{questao}[3]{}

Um divisor de tensão deve fornecer $\SI{5}{\volt}$ a uma carga que consome
$\SI{10}{\milli\ampere}$, a partir de uma fonte de $\SI{15}{\volt}$. Como regra
de projeto, adota-se a \emph{corrente de sangria} (a que circula pelo divisor sem
passar pela carga) igual a $\SI{10}{\milli\ampere}$. Determine $R_1$ e $R_2$.

\circuitoq{
\begin{circuitikz}[scale=0.95,transform shape,line width=0.8pt]
 \draw (0,0) to[battery1,l_={$E=\SI{15}{\volt}$}] (0,3.4)
       to[R,l={$R_1$},i={$I_1$}] (3.4,3.4) coordinate(a);
 \draw (a) to[R,l={$R_2$},i={$I_2$}] (3.4,0) coordinate(b) -- (0,0);
 \draw (a) to[short,-o] (5.2,3.4) coordinate(o1);
 \draw (b) to[short,-o] (5.2,0) coordinate(o2);
 \draw (o1) to[R,l={carga},i={$\SI{10}{\milli\ampere}$}] (o2);
\end{circuitikz}}

\resposta{$R_2 = \SI{500}{\ohm}$; $R_1 = \SI{500}{\ohm}$}


\end{questao}
\begin{questao}[3]{}

Um divisor de tensão com $R_1 = \SI{4}{\kilo\ohm}$ e $R_2 = \SI{2}{\kilo\ohm}$ é
alimentado por $\SI{12}{\volt}$; a saída é tomada sobre $R_2$. Determine a faixa
de valores da carga $R_L$ para que a tensão de saída não caia mais que
$\SI{10}{\percent}$ em relação ao valor em vazio.
\resposta{$R_L \gtrsim \SI{12}{\kilo\ohm}$}


\end{questao}
\begin{questao}[3]{}

Um sistema de alimentação deve entregar $\SI{5}{\volt}$ a uma carga de
$\SI{100}{\ohm}$, a partir de uma fonte de $\SI{15}{\volt}$, usando um divisor
resistivo. Especifica-se que a tensão na carga não deve cair mais de
$\SI{2}{\percent}$ se a carga variar de $\SI{100}{\ohm}$ a $\SI{200}{\ohm}$.
\begin{itensq}
\item Explique por que um divisor puramente resistivo tem dificuldade em atender
      a essa especificação.
\item Determine a ordem de grandeza que $R_2$ (resistor em paralelo com a carga)
      deveria ter e discuta o custo dessa escolha.
\end{itensq}
\resposta{Ver análise; $R_2 \ll R_L$, com alto consumo}


\end{questao}

\blocofixacao
\begin{questao}[1]{}

Três resistores de $\SI{1}{\kilo\ohm}$, $\SI{2}{\kilo\ohm}$ e
$\SI{3}{\kilo\ohm}$ estão em série sob $\SI{12}{\volt}$. Determine a tensão
sobre o resistor \emph{intermediário}.
\resposta{$U=\SI{4}{\volt}$}


\end{questao}
\begin{questao}[1]{}

Um divisor deve entregar $25\%$ da tensão de entrada. Determine a razão
$R_1/R_2$.
\resposta{$R_1/R_2=3$}


\end{questao}
\begin{questao}[1]{}

Em um divisor em vazio, ambos os resistores são \textbf{dobrados}. A tensão de
saída:
\begin{alternativas}[2]
\item dobra
\item cai à metade
\item não se altera
\item depende dos valores originais
\end{alternativas}
\resposta{(c)}


\end{questao}
\begin{questao}[1]{}

Analise os casos-limite de um divisor: (a) $R_2\to0$; (b) $R_2\to\infty$.
\resposta{(a) $U_o\to0$; (b) $U_o\to U_{in}$}


\end{questao}

\blocoaplicacao
\begin{questao}[2]{}

Um divisor de $\SI{12}{\volt}$ usa três resistores em série de
$\SI{1}{\kilo\ohm}$, $\SI{1}{\kilo\ohm}$ e $\SI{2}{\kilo\ohm}$, com duas
tomadas intermediárias.
\begin{itensq}
\item Determine as duas tensões de saída em relação ao terra.
\item Determine a corrente de repouso.
\end{itensq}
\resposta{(a) $\SI{9}{\volt}$ e $\SI{6}{\volt}$; (b) $\SI{3}{\milli\ampere}$}


\end{questao}
\begin{questao}[2]{}

Uma fonte de $\SI{30}{\volt}$ alimenta dois resistores iguais em série, e o
ponto médio é adotado como referência (terra). Determine as tensões dos dois
extremos.
\resposta{$\SI{+15}{\volt}$ e $\SI{-15}{\volt}$}


\end{questao}
\begin{questao}[2]{}

Um divisor de $\SI{24}{\volt}$ tem $R_1=\SI{8}{\kilo\ohm}$ e
$R_2=\SI{4}{\kilo\ohm}$. Determine a saída, a corrente de repouso e a potência
total dissipada.
\resposta{$\SI{8}{\volt}$; $\SI{2}{\milli\ampere}$; $\SI{48}{\milli\watt}$}


\end{questao}
\begin{questao}[2]{}

Um divisor tem $R_1=\SI{2}{\kilo\ohm}$ fixo e $R_2$ ajustável de
$\SI{0}{\ohm}$ a $\SI{8}{\kilo\ohm}$, sob $\SI{10}{\volt}$. Determine a faixa
da tensão de saída.
\resposta{de $\SI{0}{\volt}$ a $\SI{8}{\volt}$}


\end{questao}
\begin{questao}[2]{}

Um divisor deve entregar $\SI{4.5}{\volt}$ com $R_2=\SI{3}{\kilo\ohm}$.
Determine $U_{in}$ para $R_1=\SI{5}{\kilo\ohm}$.
\resposta{$U_{in}=\SI{12}{\volt}$}


\end{questao}
\begin{questao}[2]{}

Um termistor NTC vale $\SI{10}{\kilo\ohm}$ a $\SI{25}{\celsius}$ e
$\SI{4}{\kilo\ohm}$ a $\SI{50}{\celsius}$. Ele ocupa a posição de $R_2$ em um
divisor com $R_1=\SI{10}{\kilo\ohm}$ sob $\SI{5}{\volt}$. Determine a saída nas
duas temperaturas.
\resposta{$\SI{2.5}{\volt}$ e $\SI{1.43}{\volt}$}


\end{questao}
\begin{questao}[2]{}

Um LDR varia de $\SI{1}{\kilo\ohm}$ (claro) a $\SI{100}{\kilo\ohm}$ (escuro).
Ele é colocado como $R_1$ de um divisor com $R_2=\SI{10}{\kilo\ohm}$ sob
$\SI{5}{\volt}$. Determine a faixa de saída e o sentido da variação.
\resposta{de $\SI{4.55}{\volt}$ (claro) a $\SI{0.45}{\volt}$ (escuro)}


\end{questao}
\begin{questao}[2]{}

Um divisor projetado com $R_1=\SI{8.7}{\kilo\ohm}$ e $R_2=\SI{3.3}{\kilo\ohm}$
para $12\to\SI{3.3}{\volt}$ deve usar valores E24. Adotando
$R_1=\SI{9.1}{\kilo\ohm}$, determine a saída real e o erro.
\resposta{$U_o=\SI{3.19}{\volt}$; erro de $-3{,}2\%$}


\end{questao}

\blocoaprofundamento
\begin{questao}[3]{}

Um divisor $R_1=\SI{6}{\kilo\ohm}$, $R_2=\SI{3}{\kilo\ohm}$ é alimentado por
uma fonte \textbf{real} de $\SI{12}{\volt}$ com $r=\SI{1}{\kilo\ohm}$.
\begin{itensq}
\item Determine a saída, considerando $r$.
\item Compare com a previsão que ignora $r$.
\item Determine a resistência de Thévenin vista na saída.
\end{itensq}
\resposta{(a) $\SI{3.6}{\volt}$; (b) previsão de $\SI{4}{\volt}$, erro de
$10\%$; (c) $R_{Th}=\SI{2.1}{\kilo\ohm}$}


\end{questao}
\begin{questao}[3]{}

Dois divisores idênticos ($R_1=\SI{3}{\kilo\ohm}$, $R_2=\SI{1}{\kilo\ohm}$)
são alimentados pela mesma fonte de $\SI{8}{\volt}$, mas em um deles $R_2$ é
trocado por $\SI{1.1}{\kilo\ohm}$.
\begin{itensq}
\item Determine as duas saídas e a diferença entre elas.
\item Determine a sensibilidade da diferença a variações de $R_2$.
\end{itensq}
\resposta{(a) $\SI{2}{\volt}$ e $\SI{2.146}{\volt}$; diferença de
$\SI{146}{\milli\volt}$}


\end{questao}
\begin{questao}[3]{}

Um divisor $R_1=R_2=\SI{100}{\kilo\ohm}$ sob $\SI{10}{\volt}$ é medido por um
voltímetro cuja resistência de entrada é $\SI{1}{\mega\ohm}$.
\begin{itensq}
\item Determine a leitura do instrumento.
\item Determine o erro percentual.
\item Determine a resistência mínima do voltímetro para erro inferior a
      $1\%$.
\end{itensq}
\resposta{(a) $\SI{4.76}{\volt}$; (b) $-4{,}8\%$; (c)
$\ge\SI{5}{\mega\ohm}$}


\end{questao}
\begin{questao}[3]{}

Um divisor $R_1=\SI{4}{\kilo\ohm}$, $R_2=\SI{1}{\kilo\ohm}$ sob
$\SI{10}{\volt}$ fornece referência a um conversor A/D de $10$ bits com fundo
de escala de $\SI{2.5}{\volt}$.
\begin{itensq}
\item Determine a saída e a margem em relação ao fundo de escala.
\item Determine o degrau de quantização.
\item Compare o erro de quantização com o erro de resistores de $1\%$.
\end{itensq}
\resposta{(a) $\SI{2}{\volt}$, $80\%$ do fundo; (b) $\SI{2.44}{\milli\volt}$;
(c) resistores dominam ($\approx\SI{16}{\milli\volt}$)}


\end{questao}
\begin{questao}[3]{}

Um divisor de $\SI{48}{\volt}$ para $\SI{12}{\volt}$ deve alimentar uma carga
de $\SI{100}{\ohm}$.
\begin{itensq}
\item Determine $R_1$ e $R_2$ ignorando a carga, com corrente de repouso de
      $\SI{120}{\milli\ampere}$.
\item Determine a saída real com a carga conectada.
\item Avalie a viabilidade.
\end{itensq}
\resposta{(a) $R_1=\SI{300}{\ohm}$, $R_2=\SI{100}{\ohm}$;
(b) $\SI{6.86}{\volt}$; (c) inviável}


\end{questao}
\begin{questao}[3]{}

Um divisor com $R_1=\SI{10}{\kilo\ohm}$ e $R_2=\SI{10}{\kilo\ohm}$ opera sob
$\SI{200}{\volt}$.
\begin{itensq}
\item Determine a potência de cada resistor e a tensão sobre cada um.
\item Especifique os resistores considerando potência \emph{e} tensão máxima de
      trabalho.
\end{itensq}
\resposta{$\SI{1}{\watt}$ e $\SI{100}{\volt}$ em cada; exige resistor de
potência com tensão de trabalho adequada}


\end{questao}
\begin{questao}[3]{}

Um divisor deve fornecer $\SI{5}{\volt}$ a partir de $\SI{15}{\volt}$ para uma
entrada de microcontrolador com corrente de entrada de
$\SI{1}{\micro\ampere}$.
\begin{itensq}
\item Determine $R_1$ e $R_2$ para corrente de repouso de
      $\SI{10}{\micro\ampere}$.
\item Verifique o erro introduzido pela corrente de entrada.
\item Comente o compromisso com o consumo.
\end{itensq}
\resposta{(a) $R_1=\SI{1}{\mega\ohm}$, $R_2=\SI{500}{\kilo\ohm}$;
(b) erro de $\approx3\%$}


\end{questao}
\begin{questao}[3]{}

Determine, em um divisor, a expressão da resistência de Thévenin vista na saída
e avalie-a para $R_1=\SI{6}{\kilo\ohm}$ e $R_2=\SI{3}{\kilo\ohm}$.
\begin{itensq}
\item Deduza $R_{Th}$.
\item Mostre que $R_{Th}<\min(R_1,R_2)$.
\item Explique seu papel no carregamento.
\end{itensq}
\resposta{$R_{Th}=R_1\parallel R_2=\SI{2}{\kilo\ohm}$}


\end{questao}
\begin{questao}[3]{}

Um divisor alimenta uma carga que \textbf{varia} entre $\SI{10}{\kilo\ohm}$ e
$\SI{100}{\kilo\ohm}$. O divisor tem $R_{Th}=\SI{2}{\kilo\ohm}$ e
$U_{Th}=\SI{5}{\volt}$.
\begin{itensq}
\item Determine a faixa de tensão de saída.
\item Determine a regulação percentual.
\item Indique como melhorá-la.
\end{itensq}
\resposta{(a) de $\SI{4.17}{\volt}$ a $\SI{4.90}{\volt}$;
(b) $\approx15\%$}


\end{questao}

\blocodesafio
\begin{questao}[4]{}

\textbf{(Condição de carregamento.)} Deduza a condição para que uma carga
$R_L$ altere a saída de um divisor em menos de um erro relativo $\varepsilon$.
\begin{itensq}
\item Obtenha a expressão exata do erro.
\item Deduza a condição aproximada para $\varepsilon$ pequeno.
\item Construa a tabela de $R_L/R_{Th}$ para
      $\varepsilon=10\%$, $1\%$ e $0{,}1\%$.
\end{itensq}
\resposta{$\varepsilon=\dfrac{R_{Th}}{R_L+R_{Th}}$; aproximadamente
$R_L\ge R_{Th}/\varepsilon$}


\end{questao}
\begin{questao}[4]{}

\textbf{(Projeto sob restrição de potência.)} Projete um divisor que reduza
$\SI{100}{\volt}$ a $\SI{5}{\volt}$, dissipando no máximo
$\SI{0.5}{\watt}$.
\begin{itensq}
\item Determine $R_1$ e $R_2$.
\item Determine a potência e a tensão de cada resistor.
\item Determine a maior carga que pode ser atendida com erro de até $1\%$.
\end{itensq}
\resposta{(a) $R_1=\SI{19}{\kilo\ohm}$, $R_2=\SI{1}{\kilo\ohm}$;
(b) $0{,}475$ e $\SI{0.025}{\watt}$; (c) $R_L\ge\SI{95}{\kilo\ohm}$}


\end{questao}
\begin{questao}[4]{}

\textbf{(Potenciômetro carregado.)} Um potenciômetro de $\SI{10}{\kilo\ohm}$
é usado como divisor sob $\SI{10}{\volt}$, com carga
$R_L=\SI{10}{\kilo\ohm}$ na saída. Seja $x$ a fração do curso.
\begin{itensq}
\item Deduza $U_o(x)$ com a carga.
\item Avalie em $x=0{,}25$, $0{,}5$ e $0{,}75$ e compare com o caso ideal.
\item Determine onde o desvio é máximo e como reduzi-lo.
\end{itensq}
\resposta{Em $x=0{,}5$: $\SI{4.00}{\volt}$ contra $\SI{5}{\volt}$ ideal ---
desvio de $20\%$}


\end{questao}
\begin{questao}[4]{}

\textbf{(Tolerância de razão.)} Um divisor usa resistores de $1\%$ com entrada
exata.
\begin{itensq}
\item Deduza a expressão do pior erro relativo da saída.
\item Avalie para as razões $12\to\SI{6}{\volt}$, $12\to\SI{5}{\volt}$ e
      $100\to\SI{5}{\volt}$.
\item Explique por que resistores \emph{casados} melhoram drasticamente o
      resultado.
\end{itensq}
\resposta{$\varepsilon=\dfrac{R_1}{R_1+R_2}(t_1+t_2)$; $1{,}00\%$,
$1{,}17\%$ e $1{,}90\%$}


\end{questao}
\begin{questao}[4]{}

\textbf{(Sonda atenuadora.)} Uma sonda $10{:}1$ é formada por
$\SI{9}{\mega\ohm}$ na ponta, em série com a entrada de
$\SI{1}{\mega\ohm}$ do osciloscópio.
\begin{itensq}
\item Verifique a razão de atenuação e a resistência de entrada total.
\item Determine o erro ao medir uma fonte com $R_{Th}=\SI{100}{\kilo\ohm}$,
      com e sem a sonda.
\item Explique o que se perde ao usar a sonda.
\end{itensq}
\resposta{(a) $1/10$ e $\SI{10}{\mega\ohm}$; (b) $1{,}0\%$ com sonda,
$9{,}1\%$ sem}


\end{questao}
\begin{questao}[4]{}

\textbf{(Otimização de sensor.)} Um termistor de resistência $R_t$ ocupa a
posição inferior de um divisor com resistor fixo $R$ e alimentação $U$.
\begin{itensq}
\item Obtenha $\dfrac{\mathrm{d}U_o}{\mathrm{d}R_t}$.
\item Determine o valor de $R$ que maximiza a sensibilidade em um dado ponto de
      operação.
\item Avalie para $U=\SI{5}{\volt}$, $R_t=\SI{10}{\kilo\ohm}$ e coeficiente de
      $\SI{-400}{\ohm\per\celsius}$.
\end{itensq}
\resposta{(b) $R=R_t$; (c) $\SI{-50}{\milli\volt\per\celsius}$}


\end{questao}
\begin{questao}[4]{}

\textbf{(Divisor contra regulador.)} Compare as duas soluções para obter
$\SI{5}{\volt}$ a partir de $\SI{12}{\volt}$, para uma carga que varia de
$\SI{1}{\milli\ampere}$ a $\SI{10}{\milli\ampere}$.
\begin{itensq}
\item Dimensione o divisor com repouso de dez vezes a carga máxima e determine
      a regulação.
\item Determine o rendimento de cada solução.
\item Estabeleça o critério de escolha.
\end{itensq}
\resposta{(a) $R_1=\SI{70}{\ohm}$, $R_2=\SI{50}{\ohm}$; regulação
$\approx5{,}3\%$; (b) $3{,}8\%$ contra $42\%$}


\end{questao}
\begin{questao}[4]{}

\textbf{(Divisor com proteção.)} Projete a interface entre um sinal de entrada
que varia de $\SI{0}{\volt}$ a $\SI{30}{\volt}$ e um conversor A/D de
$\SI{3.3}{\volt}$ e impedância de entrada muito alta.
\begin{itensq}
\item Determine a razão e escolha os resistores.
\item Verifique o comportamento em sobretensão de $\SI{50}{\volt}$ e proponha
      proteção.
\item Determine o consumo e discuta o compromisso com a impedância exigida pelo
      conversor.
\end{itensq}
\resposta{(a) razão $\approx1/10$: $R_1=\SI{91}{\kilo\ohm}$,
$R_2=\SI{10}{\kilo\ohm}$; (c) $\SI{297}{\micro\ampere}$ no fundo de escala}


\end{questao}

\gabarito[4.1 Divisor de tensão]

\subsection{Divisor de corrente}

\blocofixacao
\begin{questao}[1]{}

Em um divisor de corrente de dois resistores em paralelo, a maior corrente
circula:
\begin{alternativas}
\item pelo maior resistor.
\item pelo menor resistor.
\item igualmente pelos dois, sempre.
\item pelo resistor mais próximo da fonte.
\item por nenhum: toda a corrente retorna à fonte.
\end{alternativas}
\resposta{(b)}


\end{questao}
\begin{questao}[1]{}

Uma corrente de $\SI{10}{\ampere}$ chega a um nó e se divide entre
$R_1=\SI{2}{\ohm}$ e $R_2=\SI{3}{\ohm}$ em paralelo. Calcule $i_1$ e $i_2$.

\circuitoq{
\begin{circuitikz}[scale=1,transform shape,line width=0.8pt,american]
 \draw (0,0) to[isource,l={$\SI{10}{\ampere}$}] (0,2.6)
       to[short,-*,i={$I_t$}] (2,2.6)
       to[R,l={$R_1=\SI{2}{\ohm}$},i>_={$i_1$}] (2,0) -- (0,0);
 \draw (2,2.6) -- (4,2.6)
       to[R,l={$R_2=\SI{3}{\ohm}$},i>_={$i_2$}] (4,0) to[short,-*] (2,0);
\end{circuitikz}}

\resposta{$i_1 = \SI{6}{\ampere}$; $i_2 = \SI{4}{\ampere}$}


\end{questao}
\begin{questao}[1]{}

Uma corrente de $\SI{15}{\ampere}$ divide-se entre dois resistores em paralelo,
$\SI{3}{\ohm}$ e $\SI{6}{\ohm}$. As correntes valem, respectivamente:
\begin{alternativas}[2]
\item \SI{5}{\ampere} e \SI{10}{\ampere}
\item \SI{10}{\ampere} e \SI{5}{\ampere}
\item \SI{7.5}{\ampere} e \SI{7.5}{\ampere}
\item \SI{9}{\ampere} e \SI{6}{\ampere}
\item \SI{6}{\ampere} e \SI{9}{\ampere}
\end{alternativas}
\resposta{(b)}


\end{questao}

\blocoaplicacao
\begin{questao}[2]{}

No circuito, uma fonte de corrente de $\SI{30}{\ampere}$ alimenta três resistores
em paralelo. Calcule a corrente em cada um.

\circuitoq{
\begin{circuitikz}[scale=1,transform shape,line width=0.8pt,american]
 \draw (0,0) to[isource,l={$\SI{30}{\ampere}$}] (0,2.6)
       to[short,-*,i={$I_t$}] (2,2.6)
       to[R,l={$\SI{2}{\ohm}$},i>_={$i_1$}] (2,0) -- (0,0);
 \draw (2,2.6) -- (4,2.6)
       to[R,l={$\SI{3}{\ohm}$},i>_={$i_2$}] (4,0) to[short,-*] (2,0);
 \draw (4,2.6) -- (6,2.6)
       to[R,l={$\SI{6}{\ohm}$},i>_={$i_3$}] (6,0) to[short,-*] (4,0);
\end{circuitikz}}

\resposta{$i_1=\SI{15}{\ampere}$, $i_2=\SI{10}{\ampere}$, $i_3=\SI{5}{\ampere}$}


\end{questao}
\begin{questao}[2]{}

No divisor da figura, com três ramos, determine a corrente em cada resistor.

\circuitoq{
\begin{circuitikz}[scale=1,transform shape,line width=0.8pt,american]
 \draw (0,0) to[isource,l={$\SI{25}{\ampere}$}] (0,2.6)
       to[short,-*,i={$I_t$}] (2,2.6)
       to[R,l={$\SI{12}{\ohm}$},i>_={$i_1$}] (2,0) -- (0,0);
 \draw (2,2.6) -- (4,2.6)
       to[R,l={$\SI{3}{\ohm}$},i>_={$i_2$}] (4,0) to[short,-*] (2,0);
 \draw (4,2.6) -- (6,2.6)
       to[R,l={$\SI{12}{\ohm}$},i>_={$i_3$}] (6,0) to[short,-*] (4,0);
\end{circuitikz}}

\resposta{$i_1=i_3\approx\SI{4.17}{\ampere}$; $i_2\approx\SI{16.7}{\ampere}$}


\end{questao}

\blocoaprofundamento
\begin{questao}[3]{}

Uma corrente de $\SI{35}{\ampere}$ alimenta três resistores em paralelo,
$\SI{5}{\ohm}$, $\SI{10}{\ohm}$ e $\SI{20}{\ohm}$.
\begin{itensq}
\item Calcule a corrente em cada ramo.
\item Verifique que a soma das potências dissipadas nos ramos é igual à potência
      entregue pela fonte.
\end{itensq}
\resposta{(a) \SI{20}{\ampere}, \SI{10}{\ampere}, \SI{5}{\ampere};
(b) \SI{3500}{\watt} nos dois lados}


\end{questao}
\begin{questao}[3]{}

Um miliamperímetro tem fundo de escala de $\SI{1}{\milli\ampere}$ e resistência
interna $R_g = \SI{50}{\ohm}$. Para medir correntes de até $\SI{1}{\ampere}$,
liga-se em paralelo com ele um resistor \emph{shunt} $R_s$, que desvia o excesso
de corrente. Determine $R_s$.

\circuitoq{
\begin{circuitikz}[scale=1,transform shape,line width=0.8pt,american]
 \draw (0,1.4) to[short,i={$I=\SI{1}{\ampere}$},-*] (2,1.4) coordinate(n1);
 \draw (n1) to[R,l={$R_g=\SI{50}{\ohm}$},i={$\SI{1}{\milli\ampere}$}] (5,1.4) coordinate(n2);
 \node[above,font=\footnotesize] at (3.5,1.75) {galvanômetro};
 \draw (n1) to[short] (2,0) to[R,l_={$R_s$},i_={$i_s$}] (5,0) to[short,-*] (n2);
 \draw (n2) to[short] (7,1.4);
\end{circuitikz}}

\resposta{$R_s \approx \SI{0.05}{\ohm}$}


\end{questao}
\begin{questao}[3]{}

Um galvanômetro tem fundo de escala $i_g = \SI{1}{\milli\ampere}$ e resistência
interna $R_g = \SI{50}{\ohm}$. Deseja-se convertê-lo em um amperímetro de
\textbf{três escalas}: \SI{100}{\milli\ampere}, \SI{1}{\ampere} e
\SI{10}{\ampere}, usando shunts comutáveis.
\begin{itensq}
\item Deduza a expressão geral do shunt para uma escala $I_{\max}$.
\item Calcule os três shunts.
\item Explique por que a resistência do amperímetro \emph{diminui} conforme se
      aumenta a escala, e por que isso é desejável.
\end{itensq}
\resposta{(a) $R_s = \dfrac{R_g\,i_g}{I_{\max}-i_g}$;
(b) \SI{0.505}{\ohm}, \SI{50.05}{\milli\ohm}, \SI{5.0}{\milli\ohm}; (c) ver comentário}


\end{questao}

\blocofixacao
\begin{questao}[1]{}

Três ramos com condutâncias $G_1=\SI{0.5}{\ensuremath{\mathrm{S}}}$,
$G_2=\SI{0.3}{\ensuremath{\mathrm{S}}}$ e $G_3=\SI{0.2}{\ensuremath{\mathrm{S}}}$ recebem
$I_T=\SI{10}{\ampere}$. Determine a corrente de cada ramo.
\resposta{$\SI{5}{\ampere}$, $\SI{3}{\ampere}$ e $\SI{2}{\ampere}$}


\end{questao}
\begin{questao}[1]{}

Dois ramos de $\SI{4}{\ohm}$ e $\SI{12}{\ohm}$ dividem uma corrente. Determine
a razão $I_1/I_2$.
\resposta{$I_1/I_2=3$}


\end{questao}
\begin{questao}[1]{}

Um ramo de $\SI{20}{\ohm}$ de uma associação paralela conduz
$\SI{1.5}{\ampere}$; o outro ramo vale $\SI{30}{\ohm}$. Determine a corrente
total.
\resposta{$I_T=\SI{2.5}{\ampere}$}


\end{questao}
\begin{questao}[1]{}

Um dos ramos de um divisor de corrente sofre \textbf{curto-circuito}. Determine
o que ocorre com as correntes.
\resposta{Toda a corrente passa pelo curto; os demais ramos ficam sem corrente}


\end{questao}
\begin{questao}[1]{}

Em um divisor de corrente alimentado por uma \textbf{fonte de corrente ideal},
um dos ramos é aberto. A corrente nos ramos restantes:
\begin{alternativas}[2]
\item permanece a mesma
\item aumenta
\item diminui
\item torna-se nula
\end{alternativas}
\resposta{(b)}


\end{questao}

\blocoaplicacao
\begin{questao}[2]{}

Uma corrente de $\SI{12}{\ampere}$ divide-se entre quatro ramos de
$\SI{5}{\ohm}$, $\SI{10}{\ohm}$, $\SI{20}{\ohm}$ e $\SI{20}{\ohm}$.
\begin{itensq}
\item Determine a tensão da associação.
\item Determine a corrente de cada ramo.
\end{itensq}
\resposta{$U=\SI{30}{\volt}$; $6$, $3$, $1{,}5$ e $\SI{1.5}{\ampere}$}


\end{questao}
\begin{questao}[2]{}

No circuito anterior, determine a fração percentual da corrente em cada ramo e
explique por que ela não depende de $I_T$.
\resposta{$50\%$, $25\%$, $12{,}5\%$ e $12{,}5\%$}


\end{questao}
\begin{questao}[2]{}

Uma corrente de $\SI{8}{\ampere}$ deve dividir-se de modo que
$\SI{25}{\percent}$ passe por um ramo $R$ e o restante por um ramo de
$\SI{15}{\ohm}$. Determine $R$.
\resposta{$R=\SI{45}{\ohm}$}


\end{questao}
\begin{questao}[2]{}

Uma fonte de $\SI{60}{\volt}$ com resistência em série $R_s=\SI{4}{\ohm}$ alimenta
dois ramos em paralelo de $\SI{10}{\ohm}$ e $\SI{15}{\ohm}$.
\begin{itensq}
\item Determine a corrente total.
\item Determine a corrente de cada ramo.
\end{itensq}
\resposta{$I_T=\SI{6}{\ampere}$; $\SI{3.6}{\ampere}$ e $\SI{2.4}{\ampere}$}


\end{questao}
\begin{questao}[2]{}

No circuito anterior, determine a potência dissipada em cada ramo e a tensão
sobre a associação.
\resposta{$U=\SI{36}{\volt}$; $\SI{129.6}{\watt}$ e $\SI{86.4}{\watt}$}


\end{questao}
\begin{questao}[2]{}

Uma corrente de $\SI{10}{\ampere}$ divide-se entre $\SI{1}{\ohm}$ e
$\SI{1}{\kilo\ohm}$.
\begin{itensq}
\item Determine as duas correntes.
\item Avalie o erro de simplesmente desprezar o ramo de alta resistência.
\end{itensq}
\resposta{$\SI{9.99}{\ampere}$ e $\SI{9.99}{\milli\ampere}$; erro de
$0{,}1\%$}


\end{questao}
\begin{questao}[2]{}

Uma fonte de corrente de $\SI{12}{\ampere}$ alimenta três ramos de
$\SI{6}{\ohm}$, $\SI{12}{\ohm}$ e $\SI{12}{\ohm}$. O ramo de
$\SI{6}{\ohm}$ é removido. Determine as novas correntes.
\resposta{Antes: $6$, $3$, $\SI{3}{\ampere}$. Depois: $\SI{6}{\ampere}$ em cada}


\end{questao}
\begin{questao}[2]{}

Um resistor de $\SI{10}{\ohm}$ está em paralelo com um \textbf{indutor} ideal,
e o conjunto recebe $\SI{4}{\ampere}$ em regime permanente CC. Determine a
corrente de cada ramo.
\resposta{$\SI{0}{\ampere}$ no resistor; $\SI{4}{\ampere}$ no indutor}


\end{questao}
\begin{questao}[2]{}

Um divisor de corrente tem três ramos e mediu-se $\SI{4}{\ampere}$,
$\SI{2}{\ampere}$ e $\SI{1.33}{\ampere}$, com tensão comum de
$\SI{20}{\volt}$. Verifique a consistência e determine as resistências.
\resposta{Consistente; $\SI{5}{\ohm}$, $\SI{10}{\ohm}$ e $\SI{15}{\ohm}$}


\end{questao}
\begin{questao}[2]{}

Explique por que a fórmula $I_1=I_T\dfrac{R_2}{R_1+R_2}$ \textbf{não} se
generaliza para três ramos, e mostre com um contraexemplo.
\resposta{A forma correta usa condutâncias; a versão com resistências vale só
para dois ramos}


\end{questao}

\blocoaprofundamento
\begin{questao}[3]{}

Três condutores em paralelo, de resistências $\SI{0.10}{\ohm}$,
$\SI{0.11}{\ohm}$ e $\SI{0.12}{\ohm}$, conduzem $\SI{300}{\ampere}$.
\begin{itensq}
\item Determine a corrente de cada um.
\item Determine o desvio percentual em relação à divisão ideal.
\item Comente a implicação térmica.
\end{itensq}
\resposta{$109{,}4$; $99{,}5$ e $\SI{91.2}{\ampere}$; desvios de
$+9{,}4\%$, $-0{,}5\%$ e $-8{,}8\%$}


\end{questao}
\begin{questao}[3]{}

Uma corrente de $\SI{12}{\ampere}$ deve dividir-se na razão $1{:}2{:}3$.
\begin{itensq}
\item Determine as três correntes.
\item Determine as resistências, adotando a tensão comum de
      $\SI{12}{\volt}$.
\item Determine a potência de cada ramo.
\end{itensq}
\resposta{$2$, $4$ e $\SI{6}{\ampere}$; $\SI{6}{\ohm}$, $\SI{3}{\ohm}$ e
$\SI{2}{\ohm}$; $24$, $48$ e $\SI{72}{\watt}$}


\end{questao}
\begin{questao}[3]{}

Dois ramos nominalmente iguais de $\SI{10}{\ohm}$, com tolerância de
$\pm5\%$, dividem $\SI{10}{\ampere}$.
\begin{itensq}
\item Determine o pior desequilíbrio de corrente.
\item Compare o desvio percentual da corrente com o dos resistores.
\end{itensq}
\resposta{$5{,}25$ e $\SI{4.75}{\ampere}$; desvio de $\pm5\%$}


\end{questao}
\begin{questao}[3]{}

Um ramo de $\SI{0.5}{\ohm}$ de um divisor adquire resistência de contato de
$\SI{0.05}{\ohm}$ em suas conexões. O outro ramo, também de
$\SI{0.5}{\ohm}$, permanece íntegro, e a corrente total é
$\SI{100}{\ampere}$.
\begin{itensq}
\item Determine as novas correntes.
\item Determine o desvio e a potência adicional no contato.
\end{itensq}
\resposta{$\SI{47.6}{\ampere}$ e $\SI{52.4}{\ampere}$; desvio de
$\mp4{,}8\%$; $\SI{113}{\watt}$ no contato}


\end{questao}
\begin{questao}[3]{}

Uma corrente de $\SI{10}{\ampere}$ divide-se entre $\SI{5}{\ohm}$ e
$\SI{20}{\ohm}$.
\begin{itensq}
\item Determine as correntes e as potências.
\item Determine qual ramo aquece mais e por qual fator.
\end{itensq}
\resposta{$8$ e $\SI{2}{\ampere}$; $320$ e $\SI{80}{\watt}$ --- o menor
resistor dissipa $4\times$ mais}


\end{questao}
\begin{questao}[3]{}

Três ramos de $\SI{10}{\ohm}$, $\SI{20}{\ohm}$ e $\SI{20}{\ohm}$ têm o
primeiro aberto. Compare o efeito sob dois tipos de alimentação:
\begin{itensq}
\item fonte de corrente ideal de $\SI{12}{\ampere}$;
\item fonte de tensão ideal de $\SI{60}{\volt}$.
\end{itensq}
\resposta{(a) tensão dobra ($60\to\SI{120}{\volt}$), ramos vão de $3$ a
$\SI{6}{\ampere}$; (b) tensão e ramos inalterados, total cai de $12$ a
$\SI{6}{\ampere}$}


\end{questao}
\begin{questao}[3]{}

Um ramo fixo de $\SI{10}{\ohm}$ está em paralelo com um reostato de
$\SI{0}{\ohm}$ a $\SI{40}{\ohm}$, sob corrente total de $\SI{5}{\ampere}$.
\begin{itensq}
\item Determine a faixa de corrente no ramo fixo.
\item Determine a posição do reostato para divisão igual.
\item Comente a linearidade do controle.
\end{itensq}
\resposta{(a) de $0$ a $\SI{4}{\ampere}$; (b) $R=\SI{10}{\ohm}$}


\end{questao}
\begin{questao}[3]{}

Deduza a expressão da corrente em um ramo $R_k$ quando todos os demais são
substituídos por sua resistência equivalente $R_p$, e verifique com o divisor de
$6$, $12$ e $\SI{12}{\ohm}$ sob $\SI{12}{\ampere}$.
\resposta{$I_k=I_T\dfrac{R_p}{R_k+R_p}$; para $R_k=\SI{6}{\ohm}$ e
$R_p=\SI{6}{\ohm}$, $I_k=\SI{6}{\ampere}$}


\end{questao}
\begin{questao}[3]{}

Um divisor de corrente tem um ramo contendo uma \textbf{fonte de tensão} de
$\SI{10}{\volt}$ em série com $\SI{5}{\ohm}$; o outro ramo é um resistor de
$\SI{5}{\ohm}$. A corrente total injetada é $\SI{4}{\ampere}$.
\begin{itensq}
\item Explique por que o divisor simples não se aplica.
\item Determine as correntes.
\end{itensq}
\resposta{$I_1=\SI{1}{\ampere}$ e $I_2=\SI{3}{\ampere}$}


\end{questao}

\blocodesafio
\begin{questao}[4]{}

\textbf{(Demonstração geral.)} Deduza a fórmula do divisor de corrente para
$n$ ramos em paralelo.
\begin{itensq}
\item Obtenha $I_k$ em função das condutâncias.
\item Mostre que a forma com resistências vale apenas para $n=2$.
\item Verifique que $\sum I_k=I_T$ identicamente.
\end{itensq}
\resposta{$I_k=I_T\dfrac{G_k}{\sum_j G_j}$}


\end{questao}
\begin{questao}[4]{}

\textbf{(Projeto de balanceamento.)} Deve-se dividir $\SI{60}{\ampere}$
igualmente entre três ramos, com perda total inferior a $\SI{10}{\watt}$.
\begin{itensq}
\item Determine a resistência de cada ramo e a queda de tensão.
\item Verifique a perda e determine a margem.
\item Discuta o compromisso entre precisão da divisão e perda.
\end{itensq}
\resposta{$R\le\SI{8.33}{\milli\ohm}$; queda de $\SI{0.167}{\volt}$;
$\SI{3.33}{\watt}$ por ramo no limite}


\end{questao}
\begin{questao}[4]{}

\textbf{(Amplificação de desequilíbrio.)} Em $n$ condutores nominalmente
iguais, um deles tem resistência $R(1-\delta)$.
\begin{itensq}
\item Deduza a fração de corrente que ele conduz.
\item Obtenha a expressão aproximada para $\delta$ pequeno e identifique o fator
      de amplificação.
\item Avalie para $n=3$ e $\delta=5\%$, $10\%$ e $20\%$.
\end{itensq}
\resposta{Desvio relativo $\approx\dfrac{n-1}{n}\,\delta$; para $n=3$:
$3{,}4\%$, $7{,}1\%$ e $15{,}4\%$}


\end{questao}
\begin{questao}[4]{}

\textbf{(Sensibilidade.)} Para um divisor de dois ramos:
\begin{itensq}
\item Obtenha $\dfrac{\partial I_1}{\partial R_1}$.
\item Obtenha a sensibilidade relativa e mostre que $|S_1|+|S_2|=1$.
\item Interprete os casos $R_1\ll R_2$ e $R_1\gg R_2$.
\end{itensq}
\resposta{$S_1=-\dfrac{R_1}{R_1+R_2}$, $S_2=+\dfrac{R_2}{R_1+R_2}$}


\end{questao}
\begin{questao}[4]{}

\textbf{(Divisor de baixa perda para medição.)} Deseja-se medir
$\SI{100}{\ampere}$ desviando uma fração conhecida para um instrumento, com
perda mínima.
\begin{itensq}
\item Projete um divisor de razão $1000{:}1$ com ramo principal de
      $\SI{0.1}{\milli\ohm}$.
\item Determine a corrente de medição, a queda e a perda total.
\item Compare com um resistor \emph{shunt} direto de $\SI{1}{\milli\ohm}$.
\end{itensq}
\resposta{Ramo de medição de $\SI{100}{\milli\ohm}$; $\SI{99.9}{\milli\ampere}$;
perda de $\SI{1}{\watt}$ contra $\SI{10}{\watt}$}


\end{questao}
\begin{questao}[4]{}

\textbf{(Estabilidade térmica da divisão.)} Dois ramos idênticos dividem uma
corrente elevada. Um deles aquece mais que o outro.
\begin{itensq}
\item Analise a evolução para material de coeficiente térmico \emph{positivo}.
\item Repita para coeficiente \emph{negativo}.
\item Determine a condição de estabilidade e cite exemplos.
\end{itensq}
\resposta{$\alpha>0$: estável (realimentação negativa); $\alpha<0$: instável
(fuga térmica)}


\end{questao}
\begin{questao}[4]{}

\textbf{(Divisor com elemento ativo.)} Um dos ramos de um divisor contém uma
fonte de corrente \textbf{dependente} que injeta $g\,U$, sendo $U$ a tensão
comum e $g=\SI{0.05}{\ensuremath{\mathrm{S}}}$. O outro ramo é um resistor de
$\SI{10}{\ohm}$, e a corrente total injetada é $\SI{6}{\ampere}$.
\begin{itensq}
\item Explique por que o divisor simples falha.
\item Determine a tensão comum e a corrente de cada ramo.
\item Determine o valor crítico de $g$.
\end{itensq}
\resposta{(b) $U=\SI{120}{\volt}$; $\SI{12}{\ampere}$ no resistor e
$\SI{6}{\ampere}$ injetados pelo ramo ativo; (c) $g_c=\SI{0.1}{\ensuremath{\mathrm{S}}}$}


\end{questao}
\begin{questao}[4]{}

\textbf{(Escolha de escala.)} Um divisor deve repartir $\SI{12}{\ampere}$ na
razão $1{:}2{:}3$. A razão fixa apenas as \emph{proporções}; resta escolher a
escala absoluta das resistências.
\begin{itensq}
\item Mostre que a perda total é proporcional à escala escolhida.
\item Determine a escala para perda de $\SI{14.4}{\watt}$ e para
      $\SI{144}{\watt}$.
\item Estabeleça o critério de escolha na prática.
\end{itensq}
\resposta{$P=I_T^{2}R_{eq}$, e $R_{eq}$ escala com o fator adotado; escalas de
$0{,}6/0{,}3/\SI{0.2}{\ohm}$ e $6/3/\SI{2}{\ohm}$}


\end{questao}

\gabarito[4.2 Divisor de corrente]

\section{Potência Elétrica e Máxima Transferência}

\blocofixacao
\begin{questao}[1]{}

Um resistor de $\SI{7}{\ohm}$ é percorrido por uma corrente de $\SI{3}{\ampere}$.
A potência dissipada é:
\begin{alternativas}[3]
\item \SI{21}{\watt}
\item \SI{49}{\watt}
\item \SI{63}{\watt}
\item \SI{147}{\watt}
\item \SI{441}{\watt}
\end{alternativas}
\resposta{(c)}


\end{questao}
\begin{questao}[1]{}

Na condição de máxima transferência de potência de uma fonte real para uma carga,
a relação entre a resistência de carga $R_L$ e a resistência interna
$R_{\text{Th}}$ da fonte é:
\begin{alternativas}
\item $R_L$ deve ser a menor possível.
\item $R_L$ deve ser a maior possível.
\item $R_L = R_{\text{Th}}$.
\item $R_L = 2R_{\text{Th}}$.
\item $R_L = R_{\text{Th}}/2$.
\end{alternativas}
\resposta{(c)}


\end{questao}
\begin{questao}[1]{}

Um dispositivo submetido a $\SI{12}{\volt}$ é percorrido por $\SI{3}{\ampere}$. A
potência que ele consome é:
\begin{alternativas}[3]
\item \SI{4}{\watt}
\item \SI{15}{\watt}
\item \SI{36}{\watt}
\item \SI{4}{\watt}
\item \SI{0.25}{\watt}
\end{alternativas}
\resposta{(c)}


\end{questao}

\blocoaplicacao
\begin{questao}[2]{}

Uma fonte de Thévenin tem $V_{\text{Th}} = \SI{12}{\volt}$ e
$R_{\text{Th}} = \SI{4}{\ohm}$.
\begin{itensq}
\item Qual é o valor de $R_L$ para máxima transferência de potência?
\item Qual é o valor dessa potência máxima?
\end{itensq}
\resposta{(a) \SI{4}{\ohm}; (b) \SI{9}{\watt}}


\end{questao}
\begin{questao}[2]{}

Uma fonte com resistência interna de $\SI{6}{\ohm}$ e fem de $\SI{18}{\volt}$
alimenta uma carga $R_L$. Determine $R_L$ para máxima transferência de potência e
o valor dessa potência.
\resposta{$R_L = \SI{6}{\ohm}$; $P_{\max} = \SI{13.5}{\watt}$}


\end{questao}
\begin{questao}[2]{}

No circuito da figura, a fonte de $\SI{60}{\volt}$ está \textbf{absorvendo}
potência (sendo carregada). Sobre a corrente $I$ indicada (que entra pelo terminal
positivo da fonte de $\SI{60}{\volt}$), pode-se afirmar que:

\circuitoq{
\begin{circuitikz}[scale=0.95,transform shape,line width=0.8pt]
 \draw (0,0) to[battery1,l_={$\SI{90}{\volt}$}] (0,3)
       to[R,l={$\SI{10}{\ohm}$}] (4,3)
       to[battery1,l={$\SI{60}{\volt}$},invert] (4,0) -- (0,0);
 \draw[->,Icol,line width=1pt] (0.8,3.4)--(2.0,3.4)
       node[midway,above,font=\footnotesize]{$I$};
\end{circuitikz}}

\begin{alternativas}
\item $I$ é positiva no sentido indicado.
\item $I$ é nula.
\item $I$ tem sentido tal que entra pelo $+$ da fonte de \SI{60}{\volt},
      confirmando a absorção.
\item a fonte de \SI{60}{\volt} nunca pode absorver potência.
\item nada se pode afirmar sem o valor do resistor.
\end{alternativas}
\resposta{(c)}


\end{questao}
\begin{questao}[2]{}

Uma bateria de fem $\SI{12}{\volt}$ e resistência interna $\SI{2}{\ohm}$ alimenta
uma carga $R_L = \SI{4}{\ohm}$.
\begin{itensq}
\item Qual é a potência entregue à carga?
\item Compare com a potência máxima teórica (obtida quando $R_L = R_{\text{int}}$).
\end{itensq}
\resposta{(a) \SI{16}{\watt}; (b) máximo é \SI{18}{\watt} em $R_L=\SI{2}{\ohm}$}


\end{questao}

\blocoaprofundamento
\begin{questao}[3]{}

Uma fonte real de $V_{\text{Th}} = \SI{20}{\volt}$ e $R_{\text{Th}} = \SI{10}{\ohm}$
alimenta uma carga variável $R_L$.
\begin{itensq}
\item Calcule a potência na carga para $R_L = \SI{5}{\ohm}$, $\SI{10}{\ohm}$ e
      $\SI{20}{\ohm}$.
\item Confirme que o máximo ocorre em $R_L = R_{\text{Th}}$.
\end{itensq}
\resposta{$P(5)\approx\SI{8.9}{\watt}$, $P(10)=\SI{10}{\watt}$,
$P(20)\approx\SI{8.9}{\watt}$; máximo em \SI{10}{\ohm}}


\end{questao}
\begin{questao}[3]{}

Explique por que, na condição de máxima transferência de potência
($R_L = R_{\text{Th}}$), o rendimento do circuito é de apenas
$\SI{50}{\percent}$, e por que, em sistemas de \emph{potência} (rede elétrica),
essa condição \textbf{não} é desejável.
\resposta{Ver comentário}


\end{questao}
\begin{questao}[3]{}

Uma fonte de $V_{\text{Th}}=\SI{24}{\volt}$ e $R_{\text{Th}}=\SI{4}{\ohm}$
alimenta uma carga variável $R_L$. Complete mentalmente a análise abaixo e
responda.

\begin{table}[H]
\centering\small\sffamily
\setlength{\arraystretch}{1.25}
\begin{tabular}{@{}c c c c@{}}
\toprule
$R_L$ ($\Omega$) & $I$ (A) & $P_L$ (W) & $\eta$ (\%)\\
\midrule
1  & 4,8 & 23,0 & 20\\
2  & 4,0 & 32,0 & 33\\
\textbf{4}  & \textbf{3,0} & \textbf{36,0} & \textbf{50}\\
8  & 2,0 & 32,0 & 67\\
16 & 1,2 & 23,0 & 80\\
\bottomrule
\end{tabular}
\caption{Potência na carga e rendimento em função de $R_L$.}
\end{table}

\begin{itensq}
\item Verifique os valores da linha $R_L=\SI{8}{\ohm}$.
\item Explique a \emph{simetria} observada na coluna $P_L$.
\item Se o objetivo for \emph{eficiência energética}, qual $R_L$ escolher? E se
      for \emph{máxima potência}? Justifique a diferença.
\end{itensq}
\resposta{(a) $I=\SI{2}{\ampere}$, $P_L=\SI{32}{\watt}$, $\eta=\SI{67}{\percent}$;
(b) e (c) ver comentário}


\end{questao}

\blocodesafio
\begin{questao}[4]{}

\textbf{Problema de síntese integrado.} Projete um circuito que satisfaça
simultaneamente:
\begin{itemize}[leftmargin=1.6em,itemsep=1pt,topsep=2pt]
\item alimentado por uma fonte de $\SI{24}{\volt}$ com resistência interna
      $\SI{2}{\ohm}$;
\item entregue \emph{máxima potência} a uma carga;
\item a corrente total drenada da fonte não ultrapasse $\SI{6}{\ampere}$.
\end{itemize}
Determine o valor da carga, a potência entregue e verifique a restrição de
corrente.
\resposta{$R_L = \SI{2}{\ohm}$; $P = \SI{72}{\watt}$; $I = \SI{6}{\ampere}$
(no limite)}


\end{questao}

\blocofixacao
\begin{questao}[1]{}

Uma fonte tem $V_{Th}=\SI{24}{\volt}$ e $R_{Th}=\SI{6}{\ohm}$. Determine
$R_L$ para máxima transferência e a potência correspondente.
\resposta{$R_L=\SI{6}{\ohm}$; $P_{\max}=\SI{24}{\watt}$}


\end{questao}
\begin{questao}[1]{}

Na condição anterior, determine a corrente e a tensão sobre a carga.
\resposta{$I=\SI{2}{\ampere}$; $U_L=\SI{12}{\volt}$}


\end{questao}
\begin{questao}[1]{}

Uma fonte entrega no máximo $\SI{50}{\watt}$, com $R_{Th}=\SI{8}{\ohm}$.
Determine $V_{Th}$.
\resposta{$V_{Th}=\SI{40}{\volt}$}


\end{questao}
\begin{questao}[1]{}

Uma fonte de $\SI{24}{\volt}$ entrega no máximo $\SI{24}{\watt}$. Determine
$R_{Th}$.
\resposta{$R_{Th}=\SI{6}{\ohm}$}


\end{questao}
\begin{questao}[1]{}

Se $R_L$ for ajustado para o \textbf{dobro} do valor ótimo, a potência
entregue:
\begin{alternativas}[2]
\item dobra
\item cai à metade
\item cai para $\SI{89}{\percent}$ do máximo
\item permanece máxima
\end{alternativas}
\resposta{(c)}


\end{questao}

\blocoaplicacao
\begin{questao}[2]{}

Uma fonte com $V_{Th}=\SI{24}{\volt}$ e $R_{Th}=\SI{6}{\ohm}$ alimenta
sucessivamente $R_L=\SI{3}{\ohm}$ e $R_L=\SI{12}{\ohm}$. Determine a potência
em cada caso e comente.
\resposta{$\SI{21.33}{\watt}$ nos dois casos}


\end{questao}
\begin{questao}[2]{}

No ponto de máxima transferência da fonte anterior, determine a potência total
fornecida, a entregue à carga e a dissipada internamente.
\resposta{$\SI{48}{\watt}$, $\SI{24}{\watt}$ e $\SI{24}{\watt}$}


\end{questao}
\begin{questao}[2]{}

Duas fontes têm a mesma potência máxima de $\SI{18}{\watt}$: a primeira com
$V_{Th}=\SI{12}{\volt}$ e a segunda com $V_{Th}=\SI{24}{\volt}$. Determine
$R_{Th}$ de cada uma.
\resposta{$\SI{2}{\ohm}$ e $\SI{8}{\ohm}$}


\end{questao}
\begin{questao}[2]{}

Uma carga fixa de $\SI{10}{\ohm}$ é alimentada, alternativamente, por duas
fontes de $\SI{20}{\volt}$: uma com $R_{Th}=\SI{2}{\ohm}$ e outra com
$R_{Th}=\SI{10}{\ohm}$. Compare a potência entregue.
\resposta{$\SI{27.8}{\watt}$ e $\SI{10}{\watt}$}


\end{questao}
\begin{questao}[2]{}

Uma fonte casada entrega $\SI{24}{\watt}$ durante $\SI{5}{\hour}$. Determine a
energia entregue à carga e a energia total consumida da fonte.
\resposta{$\SI{120}{\watt\hour}$ e $\SI{240}{\watt\hour}$}


\end{questao}
\begin{questao}[2]{}

Uma fonte de $\SI{24}{\volt}$ e $R_{Th}=\SI{6}{\ohm}$ alimenta
$R_L=\SI{18}{\ohm}$. Determine a potência na carga, o rendimento e compare com
o máximo disponível.
\resposta{$P_L=\SI{18}{\watt}$; $\eta=75\%$; $75\%$ do máximo}


\end{questao}
\begin{questao}[2]{}

Uma fonte tem $V_{oc}=\SI{30}{\volt}$ e $I_{sc}=\SI{5}{\ampere}$. Determine
$R_{Th}$ e a potência máxima disponível.
\resposta{$R_{Th}=\SI{6}{\ohm}$; $P_{\max}=\SI{37.5}{\watt}$}


\end{questao}
\begin{questao}[2]{}

Verifique o balanço de potência para a fonte de $\SI{24}{\volt}$ e
$R_{Th}=\SI{6}{\ohm}$ com $R_L=\SI{2}{\ohm}$.
\resposta{Fonte: $\SI{72}{\watt}$; carga: $\SI{18}{\watt}$; interna:
$\SI{54}{\watt}$}


\end{questao}

\blocoaprofundamento
\begin{questao}[3]{}

Para $V_{Th}=\SI{24}{\volt}$ e $R_{Th}=\SI{6}{\ohm}$, verifique a simetria dos
pares recíprocos de carga.
\begin{itensq}
\item Calcule $P_L$ para $R_L=\SI{2}{\ohm}$ e $R_L=\SI{18}{\ohm}$.
\item Repita para $\SI{1.5}{\ohm}$ e $\SI{24}{\ohm}$.
\item Identifique a relação entre os pares.
\end{itensq}
\resposta{(a) $\SI{18}{\watt}$ nos dois; (b) $\SI{15.36}{\watt}$ nos dois;
(c) $R_{L1}R_{L2}=R_{Th}^{2}$}


\end{questao}
\begin{questao}[3]{}

Determine a faixa de $R_L/R_{Th}$ que garante pelo menos $90\%$ da potência
máxima.
\begin{itensq}
\item Monte a equação e resolva.
\item Interprete a largura da faixa.
\end{itensq}
\resposta{$0{,}52\le x\le1{,}93$ --- razão de quase $4{:}1$}


\end{questao}
\begin{questao}[3]{}

Construa a tabela de rendimento e de fração de potência máxima para
$x=R_L/R_{Th}$ igual a $0{,}5$; $1$; $2$; $4$ e $9$, e interprete.
\resposta{$\eta$: $33$, $50$, $67$, $80$ e $90\%$; $P/P_{\max}$: $89$, $100$,
$89$, $64$ e $36\%$}


\end{questao}
\begin{questao}[3]{}

Uma carga fixa de $\SI{8}{\ohm}$ é alimentada por uma fonte de
$\SI{24}{\volt}$ cuja resistência interna pode ser reduzida por projeto.
\begin{itensq}
\item Calcule $P_L$ para $R_{Th}=16$, $8$, $4$, $2$ e $\SI{0}{\ohm}$.
\item Determine o valor ótimo de $R_{Th}$.
\item Explique por que a resposta não é $R_{Th}=R_L$.
\end{itensq}
\resposta{(b) $R_{Th}\to0$; a potência cresce monotonicamente conforme
$R_{Th}$ diminui}


\end{questao}
\begin{questao}[3]{}

Uma bateria de $\EMF=\SI{12}{\volt}$ e $r=\SI{0.5}{\ohm}$ alimenta um motor de
partida.
\begin{itensq}
\item Determine a potência máxima disponível e a carga correspondente.
\item Determine a corrente nessa condição e avalie sua viabilidade.
\item Determine a potência com $R_L=\SI{2}{\ohm}$ e compare.
\end{itensq}
\resposta{(a) $\SI{72}{\watt}$ com $R_L=\SI{0.5}{\ohm}$; (b)
$\SI{12}{\ampere}$; (c) $\SI{46.1}{\watt}$ com $\eta=80\%$}


\end{questao}
\begin{questao}[3]{}

Uma fonte alimenta a carga através de uma rede resistiva intermediária. A fonte
tem $\SI{36}{\volt}$ e $\SI{4}{\ohm}$; a rede é um resistor de
$\SI{12}{\ohm}$ em série e um de $\SI{6}{\ohm}$ em paralelo com a saída.
\begin{itensq}
\item Determine o equivalente de Thévenin visto pela carga.
\item Determine $R_L$ ótimo e a potência máxima.
\item Compare com a potência máxima disponível na fonte original.
\end{itensq}
\resposta{(a) $V_{Th}=\SI{9.82}{\volt}$, $R_{Th}=\SI{4.36}{\ohm}$;
(b) $R_L=\SI{4.36}{\ohm}$ e $P_{\max}=\SI{5.53}{\watt}$; (c) contra
$\SI{81}{\watt}$ disponíveis na fonte}


\end{questao}
\begin{questao}[3]{}

Uma carga só pode ser escolhida entre $\SI{12}{\ohm}$ e $\SI{40}{\ohm}$
(valores disponíveis em estoque). A fonte tem $V_{Th}=\SI{24}{\volt}$ e
$R_{Th}=\SI{6}{\ohm}$.
\begin{itensq}
\item Determine qual escolher para máxima potência.
\item Determine qual escolher para máximo rendimento.
\item Discuta o critério de decisão.
\end{itensq}
\resposta{(a) $\SI{12}{\ohm}$ ($\SI{21.3}{\watt}$); (b) $\SI{40}{\ohm}$
($\eta=87\%$)}


\end{questao}
\begin{questao}[3]{}

Uma fonte de $\SI{24}{\volt}$ e $R_{Th}=\SI{6}{\ohm}$ tem sua corrente limitada
eletronicamente a $\SI{1.5}{\ampere}$.
\begin{itensq}
\item Determine se o ponto ótimo teórico é atingível.
\item Determine a carga que maximiza a potência sob essa restrição.
\item Compare com o máximo irrestrito.
\end{itensq}
\resposta{(a) não --- exigiria $\SI{2}{\ampere}$; (b)
$R_L=\SI{10}{\ohm}$; (c) $\SI{22.5}{\watt}$ contra $\SI{24}{\watt}$}


\end{questao}
\begin{questao}[3]{}

Determine o rendimento necessário para que uma fonte entregue $\SI{64}{\percent}$
da potência máxima, e verifique a coerência com a curva $\eta(x)$.
\resposta{$\eta=80\%$ (com $x=4$) ou $\eta=20\%$ (com $x=0{,}25$)}


\end{questao}

\blocodesafio
\begin{questao}[4]{}

\textbf{(Demonstração completa.)} Prove a condição de máxima transferência de
potência.
\begin{itensq}
\item Derive $P_L(R_L)$ e obtenha o ponto crítico.
\item Confirme que é máximo pela segunda derivada ou pela análise dos limites.
\item Obtenha $P_{\max}$ e o rendimento correspondente.
\end{itensq}
\resposta{$R_L=R_{Th}$; $P_{\max}=\dfrac{V_{Th}^{2}}{4R_{Th}}$;
$\eta=50\%$}


\end{questao}
\begin{questao}[4]{}

\textbf{(Simetria dos pares recíprocos.)} Prove que $R_L$ e
$R_{Th}^{2}/R_L$ entregam a mesma potência.
\begin{itensq}
\item Demonstre algebricamente.
\item Mostre que os rendimentos correspondentes somam $100\%$.
\item Deduza um método para obter $R_{Th}$ a partir de duas medições.
\end{itensq}
\resposta{$P(x)=P(1/x)$; $\eta(x)+\eta(1/x)=1$;
$R_{Th}=\sqrt{R_{L1}R_{L2}}$}


\end{questao}
\begin{questao}[4]{}

\textbf{(Tolerância em torno do ótimo.)} Deduza a faixa de $x=R_L/R_{Th}$ que
garante uma fração $f$ da potência máxima.
\begin{itensq}
\item Obtenha a equação e sua solução geral.
\item Construa a tabela para $f=0{,}90$, $0{,}95$ e $0{,}99$.
\item Interprete o comportamento da largura da faixa.
\end{itensq}
\resposta{$x=\dfrac{2-f\pm2\sqrt{1-f}}{f}$; faixas de $3{,}7{:}1$,
$2{,}5{:}1$ e $1{,}5{:}1$}


\end{questao}
\begin{questao}[4]{}

\textbf{(Projeto de fonte.)} Projete uma fonte que entregue no máximo
$\SI{25}{\watt}$ a uma carga de $\SI{4}{\ohm}$.
\begin{itensq}
\item Determine $V_{Th}$ e $R_{Th}$.
\item Determine a corrente, a tensão de saída e a potência total.
\item Especifique a dissipação interna e discuta o dimensionamento.
\end{itensq}
\resposta{$V_{Th}=\SI{20}{\volt}$, $R_{Th}=\SI{4}{\ohm}$;
$I=\SI{2.5}{\ampere}$; $P_{tot}=\SI{50}{\watt}$}


\end{questao}
\begin{questao}[4]{}

\textbf{(Otimização com restrição de rendimento.)} Maximize a potência
entregue, com a restrição $\eta\ge80\%$, para
$V_{Th}=\SI{24}{\volt}$ e $R_{Th}=\SI{6}{\ohm}$.
\begin{itensq}
\item Traduza a restrição em uma condição sobre $x$.
\item Determine o ótimo restrito.
\item Compare com o ótimo irrestrito e interprete.
\end{itensq}
\resposta{$x\ge4$; $R_L=\SI{24}{\ohm}$ com $P_L=\SI{15.36}{\watt}$
($64\%$ do máximo)}


\end{questao}
\begin{questao}[4]{}

\textbf{(Potência sob dois graus de liberdade.)} Analise a maximização de
$P_L$ quando \emph{ambos} $R_L$ e $R_{Th}$ podem variar, com $V_{Th}$ fixo.
\begin{itensq}
\item Mostre que não existe máximo interior.
\item Determine o comportamento no contorno.
\item Explique a relevância prática dessa análise.
\end{itensq}
\resposta{Não há máximo interior; $P_L\to\infty$ com
$R_{Th}\to0$ e $R_L\to0$}


\end{questao}
\begin{questao}[4]{}

\textbf{(Análise de um sistema real.)} Um painel fotovoltaico apresenta, sob
certa irradiância, $V_{oc}=\SI{22}{\volt}$ e $I_{sc}=\SI{5}{\ampere}$, com
ponto de máxima potência em $\SI{18}{\volt}$ e $\SI{4.5}{\ampere}$.
\begin{itensq}
\item Compare a potência máxima real com a prevista pelo modelo de Thévenin.
\item Explique a discrepância.
\item Discuta a implicação para o rastreamento do ponto de máxima potência.
\end{itensq}
\resposta{$\SI{81}{\watt}$ reais contra $\SI{27.5}{\watt}$ previstos --- o
painel não é uma fonte linear}


\end{questao}

\gabarito[Capítulo 5 --- Potência e máxima transferência]

\section{Análise Nodal}

\blocofixacao
\begin{questao}[1]{}

Na análise nodal, o nó de referência (terra) é escolhido para:
\begin{alternativas}
\item ter o maior potencial do circuito.
\item servir como ponto de potencial zero, ao qual os demais são comparados.
\item ser sempre o nó com a fonte de maior tensão.
\item eliminar as fontes de corrente.
\item garantir que todas as correntes sejam positivas.
\end{alternativas}
\resposta{(b)}


\end{questao}
\begin{questao}[1]{}

No circuito da figura, uma fonte de corrente de $\SI{9}{\ampere}$ alimenta dois
resistores em paralelo ligados ao terra. Determine o potencial do nó $V$.

\circuitoq{
\begin{circuitikz}[scale=1,transform shape,line width=0.8pt,american]
 \draw (0,0) to[isource,l={$\SI{9}{\ampere}$}] (0,2.6)
       -- (1.5,2.6) coordinate(n);
 \node[above,Vcol] at (1.5,2.75) {$V$};
 \draw (n) to[R,l={$\SI{6}{\ohm}$}] (1.5,0) -- (0,0);
 \draw (n) -- (3.5,2.6) to[R,l={$\SI{3}{\ohm}$}] (3.5,0) -- (1.5,0);
 \node[ground] at (0,0){};
\end{circuitikz}}

\resposta{$V = \SI{18}{\volt}$}


\end{questao}

\blocoaplicacao
\begin{questao}[2]{}

Determine os potenciais nodais $V_1$ e $V_2$ do circuito da figura.

\circuitoq{
\begin{circuitikz}[scale=1,transform shape,line width=0.8pt,american]
 \draw (0,0) to[isource,l={$\SI{6}{\ampere}$}] (0,2.8) -- (1.5,2.8) coordinate(a);
 \node[above,Vcol] at (1.5,2.95) {$V_1$};
 \draw (a) to[R,l={$\SI{2}{\ohm}$}] (1.5,0) coordinate(g);
 \draw (a) to[R,l={$\SI{3}{\ohm}$}] (4.5,2.8) coordinate(b);
 \node[above,Vcol] at (4.5,2.95) {$V_2$};
 \draw (b) to[R,l={$\SI{6}{\ohm}$}] (4.5,0) -- (g) -- (0,0);
 \node[ground] at (1.5,0){};
\end{circuitikz}}

\resposta{$V_1 \approx \SI{9.82}{\volt}$; $V_2 \approx \SI{6.55}{\volt}$}


\end{questao}
\begin{questao}[2]{}

No nó da figura, uma fonte injeta $\SI{3}{\ampere}$ e outra retira
$\SI{1}{\ampere}$; dois resistores de $\SI{6}{\ohm}$ e $\SI{3}{\ohm}$ ligam o nó ao
terra. Determine o potencial $V$.

\circuitoq{
\begin{circuitikz}[scale=1,transform shape,line width=0.8pt,american]
 \draw (0,0) to[isource,l={$\SI{3}{\ampere}$}] (0,2.6) -- (1.5,2.6) coordinate(n);
 \node[above,Vcol] at (1.5,2.75){$V$};
 \draw (n) to[R,l={$\SI{6}{\ohm}$}] (1.5,0) -- (0,0);
 \draw (n) -- (3.5,2.6) to[R,l={$\SI{3}{\ohm}$}] (3.5,0) -- (1.5,0);
 \draw (n) -- (5,2.6) to[isource,l={$\SI{1}{\ampere}$},invert] (5,0) -- (3.5,0);
 \node[ground] at (1.5,0){};
\end{circuitikz}}

\resposta{$V = \SI{4}{\volt}$}


\end{questao}

\blocoaprofundamento
\begin{questao}[3]{}

No circuito da figura, há uma fonte de tensão e uma de corrente. Usando análise
nodal, determine o potencial $V_A$.

\circuitoq{
\begin{circuitikz}[scale=1,transform shape,line width=0.8pt,american]
 \draw (0,0) to[battery1,l_={$\SI{12}{\volt}$}] (0,3) -- (1.5,3)
       to[R,l={$\SI{4}{\ohm}$}] (3.5,3) coordinate(a);
 \node[above,Vcol] at (3.5,3.15) {$V_A$};
 \draw (a) to[R,l={$\SI{6}{\ohm}$}] (3.5,0) coordinate(g) -- (0,0);
 \draw (a) -- (5.5,3);
 \draw (5.5,3) to[isource,l={$\SI{2}{\ampere}$}] (5.5,0) -- (g);
 \node[ground] at (3.5,0){};
\end{circuitikz}}

\resposta{$V_A = \SI{12}{\volt}$}


\end{questao}
\begin{questao}[3]{}

Determine os potenciais nodais $V_1$, $V_2$ e $V_3$ do circuito da figura,
montando e resolvendo o sistema de equações.

\circuitoq{
\begin{circuitikz}[scale=1,transform shape,line width=0.8pt]
 \draw (0,0) to[\fonteV,l={$\SI{12}{\volt}$}] (0,3) -- (1.4,3)
       to[R,l={$\SI{2}{\ohm}$}] (3.2,3) coordinate(n1);
 \node[above,Vcol] at (3.2,3.2){$V_1$}; \fill (n1) circle (0.07);
 \draw (n1) to[R,l={$\SI{4}{\ohm}$}] (6,3) coordinate(n2);
 \node[above,Vcol] at (6,3.2){$V_2$}; \fill (n2) circle (0.07);
 \draw (n1) to[R,l_={$\SI{4}{\ohm}$}] (3.2,0) coordinate(n3);
 \node[below,Vcol] at (3.2,-0.25){$V_3$}; \fill (n3) circle (0.07);
 \draw (n2) to[R,l={$\SI{8}{\ohm}$}] (6,1.5) coordinate(m) -- (6,0) coordinate(b2);
 \draw (n3) to[R,l_={$\SI{4}{\ohm}$}] (0,0);
 \draw (n3) -- (b2);
 \draw (m) -- (7.6,1.5) to[R,l={$\SI{8}{\ohm}$}] (7.6,-1.2) -- (3.2,-1.2) -- (n3);
 \node[ground] at (3.2,-1.2){};
\end{circuitikz}}

\resposta{$V_1 = \dfrac{76}{9}\approx\SI{8.44}{\volt}$,
$V_2=\dfrac{16}{3}\approx\SI{5.33}{\volt}$,
$V_3=\dfrac{40}{9}\approx\SI{4.44}{\volt}$}


\end{questao}
\begin{questao}[3]{}

No circuito da figura, uma fonte de $\SI{6}{\volt}$ está \emph{flutuante} entre os
nós 1 e 2 (nenhum deles é o terra). Uma fonte de corrente de $\SI{2}{\ampere}$
injeta corrente no nó 1. Determine $V_1$ e $V_2$ pelo método do \textbf{supernó}.

\circuitoq{
\begin{circuitikz}[scale=1,transform shape,line width=0.8pt,american]
 
 \draw (0,0) to[isource,l={$\SI{2}{\ampere}$}] (0,3) -- (1.6,3) coordinate(n1);
 \node[above,Vcol] at (1.6,3.2){$V_1$};
 \fill (n1) circle (0.07);
 
 \draw (n1) to[R,l={$R_1=\SI{6}{\ohm}$}] (1.6,0) coordinate(g);
 
 \draw (n1) to[battery1,l={$\SI{6}{\volt}$}] (4.6,3) coordinate(n2);
 \node[above,Vcol] at (4.6,3.2){$V_2$};
 \fill (n2) circle (0.07);
 
 \draw (n2) to[R,l={$R_2=\SI{3}{\ohm}$}] (4.6,0) -- (g) -- (0,0);
 
 \draw[dashed,laranja,line width=1pt,rounded corners=8pt]
       (1.1,2.4) rectangle (5.1,3.75);
 \node[laranja,font=\footnotesize\sffamily,right] at (5.2,3.05) {supernó};
 \node[ground] at (1.6,0){};
\end{circuitikz}}

\resposta{$V_1 = \SI{8}{\volt}$; $V_2 = \SI{2}{\volt}$}


\end{questao}

\blocofixacao
\begin{questao}[1]{}

Em um circuito com cinco nós, quantas tensões nodais independentes existem após a escolha do nó de referência?
\resposta{4}


\end{questao}
\begin{questao}[1]{}

Um nó $A$ recebe $\SI{3}{\ampere}$ e está ligado ao terra por $\SI{10}{\ohm}$ e $\SI{15}{\ohm}$. Determine $V_A$.
\resposta{$V_A=\SI{18}{\volt}$}


\end{questao}
\begin{questao}[1]{}

Escreva a corrente, no sentido $A\to B$, em um resistor $R$ entre os nós $A$ e $B$.
\resposta{$i_{AB}=(V_A-V_B)/R$}


\end{questao}
\begin{questao}[1]{}

Uma fonte de corrente ideal de $I$ sai de um nó. Na forma $\mathbf G\mathbf V=\mathbf I$, qual é o sinal de sua contribuição no vetor de correntes injetadas?
\resposta{Negativo}


\end{questao}
\begin{questao}[1]{}

Converta $\SI{4}{\ohm}$, $\SI{5}{\ohm}$ e $\SI{20}{\ohm}$, todos ligados do mesmo nó ao terra, em uma única condutância equivalente.
\resposta{$G_{eq}=\SI{0.5}{\ensuremath{\mathrm{S}}}$}


\end{questao}
\begin{questao}[1]{}

Uma fonte ideal de tensão de $\SI{12}{\volt}$ tem o terminal negativo aterrado. Qual é a tensão do terminal positivo?
\resposta{$\SI{12}{\volt}$}


\end{questao}
\begin{questao}[1]{}

Quando uma fonte ideal de tensão liga dois nós desconhecidos e nenhum deles é o terra, qual técnica nodal é indicada?
\resposta{Supernó}


\end{questao}
\begin{questao}[1]{}

Após obter $V_A=\SI{18}{\volt}$, determine a corrente em um resistor de $\SI{15}{\ohm}$ entre $A$ e o terra.
\resposta{$\SI{1.2}{\ampere}$, de $A$ para o terra}


\end{questao}

\blocoaplicacao
\begin{questao}[2]{}

Determine $V$ e as correntes nos dois resistores.
\circuitoq{\begin{circuitikz}[american,scale=.95,transform shape]
\draw (0,0) to[isource,l_={$\SI{5}{\ampere}$}] (0,2.7) -- (1.8,2.7) coordinate(n);
\draw (n) to[R,l={$\SI{10}{\ohm}$}] (1.8,0) -- (0,0);
\draw (n) -- (3.8,2.7) to[R,l={$\SI{20}{\ohm}$}] (3.8,0) -- (1.8,0);
\node[above,Vcol] at (n){$V$}; \node[ground] at (1.8,0){};
\end{circuitikz}}
\resposta{$V=\SI{33.33}{\volt}$; $i_{10}=\SI{3.33}{\ampere}$; $i_{20}=\SI{1.67}{\ampere}$}


\end{questao}
\begin{questao}[2]{}

Determine a tensão nodal $V$.
\circuitoq{\begin{circuitikz}[american,scale=.9,transform shape]
\draw (0,0) to[isource,l_={$\SI{6}{\ampere}$}] (0,2.6) -- (1.4,2.6) coordinate(n);
\draw (n) to[R,l={$\SI{8}{\ohm}$}] (1.4,0) -- (0,0);
\draw (n) -- (3.3,2.6) to[R,l={$\SI{8}{\ohm}$}] (3.3,0) -- (1.4,0);
\draw (n) -- (5,2.6) to[isource,l={$\SI{2}{\ampere}$}] (5,0) -- (3.3,0);
\node[above,Vcol] at (n){$V$}; \node[ground] at (1.4,0){};
\end{circuitikz}}
\resposta{$V=\SI{16}{\volt}$}


\end{questao}
\begin{questao}[2]{}

Determine $V_A$ usando diretamente a LKC no nó.
\circuitoq{\begin{circuitikz}[american,scale=.9,transform shape]
\coordinate (g) at (2.2,0);
\draw (0,0) to[battery1,l_={$\SI{12}{\volt}$}] (0,2.7) to[R,l={$\SI{4}{\ohm}$}] (2.2,2.7) coordinate(a);
\draw (a) to[R,l={$\SI{12}{\ohm}$}] (g);
\draw (4.6,0) to[battery1,l_={$\SI{6}{\volt}$}] (4.6,2.7) to[R,l_={$\SI{6}{\ohm}$}] (a);
\draw (4.6,0) -- (g) -- (0,0);
\node[ground] at (g){}; \node[above,Vcol] at (a){$V_A$};
\end{circuitikz}}
\resposta{$V_A=\SI{8}{\volt}$}


\end{questao}
\begin{questao}[2]{}

Determine $V_1$ e $V_2$.
\circuitoq{\begin{circuitikz}[american,scale=.95,transform shape]
\draw (0,0) to[isource,l_={$\SI{3}{\ampere}$}] (0,2.8) -- (1.5,2.8) coordinate(n1);
\draw (n1) to[R,l={$\SI{10}{\ohm}$}] (1.5,0);
\draw (n1) to[R,l={$\SI{10}{\ohm}$}] (4.2,2.8) coordinate(n2);
\draw (n2) to[R,l={$\SI{10}{\ohm}$}] (4.2,0) -- (1.5,0) -- (0,0);
\node[above,Vcol] at (n1){$V_1$}; \node[above,Vcol] at (n2){$V_2$}; \node[ground] at (1.5,0){};
\end{circuitikz}}
\resposta{$V_1=\SI{20}{\volt}$; $V_2=\SI{10}{\volt}$}


\end{questao}
\begin{questao}[2]{}

Determine $V_A$. A fonte e o resistor de $\SI{10}{\ohm}$ formam uma fonte real de tensão.
\circuitoq{\begin{circuitikz}[american,scale=.95,transform shape]
\draw (0,0) to[battery1,l_={$\SI{30}{\volt}$}] (0,2.7) to[R,l={$\SI{10}{\ohm}$}] (2.7,2.7) coordinate(a);
\draw (a) to[R,l={$\SI{15}{\ohm}$}] (2.7,0) -- (0,0); \node[ground] at (2.7,0){}; \node[above,Vcol] at (a){$V_A$};
\end{circuitikz}}
\resposta{$V_A=\SI{18}{\volt}$}


\end{questao}
\begin{questao}[2]{}

Determine $V_A$ e a corrente no resistor de $\SI{5}{\ohm}$, positiva no sentido $A\to$ nó de $\SI{20}{\volt}$.
\circuitoq{\begin{circuitikz}[american,scale=.9,transform shape]
\draw (0,0) to[battery1,l_={$\SI{20}{\volt}$}] (0,2.7) to[R,l={$\SI{5}{\ohm}$}] (2.4,2.7) coordinate(a);
\draw (a) to[R,l={$\SI{20}{\ohm}$}] (2.4,0) -- (0,0);
\draw (4.4,0) to[isource,l_={$\SI{2}{\ampere}$}] (4.4,2.7) -- (a); \draw (4.4,0)--(2.4,0);
\node[ground] at (2.4,0){}; \node[above,Vcol] at (a){$V_A$};
\end{circuitikz}}
\resposta{$V_A=\SI{24}{\volt}$; $i=\SI{0.8}{\ampere}$}


\end{questao}
\begin{questao}[2]{}

Determine $V_1$ e $V_2$ na rede em escada.
\circuitoq{\begin{circuitikz}[american,scale=.95,transform shape]
\draw (0,0) to[isource,l_={$\SI{4}{\ampere}$}] (0,2.8) -- (1.4,2.8) coordinate(n1);
\draw (n1) to[R,l={$\SI{5}{\ohm}$}] (1.4,0);
\draw (n1) to[R,l={$\SI{10}{\ohm}$}] (4.1,2.8) coordinate(n2);
\draw (n2) to[R,l={$\SI{10}{\ohm}$}] (4.1,0) -- (1.4,0) -- (0,0);
\node[ground] at (1.4,0){}; \node[above,Vcol] at (n1){$V_1$}; \node[above,Vcol] at (n2){$V_2$};
\end{circuitikz}}
\resposta{$V_1=\SI{16}{\volt}$; $V_2=\SI{8}{\volt}$}


\end{questao}
\begin{questao}[2]{}

A fonte de $\SI{1}{\ampere}$ transfere corrente do nó 1 para o nó 2. Determine $V_1$ e $V_2$.
\circuitoq{\begin{circuitikz}[american,scale=.9,transform shape]
\draw (0,0) to[isource,l_={$\SI{5}{\ampere}$}] (0,3) -- (1.5,3) coordinate(n1);
\draw (n1) to[R,l={$\SI{6}{\ohm}$}] (1.5,0);
\draw (n1) to[R,l={$\SI{6}{\ohm}$}] (4.1,3) coordinate(n2);
\draw (n2) to[R,l={$\SI{3}{\ohm}$}] (4.1,0) -- (1.5,0)--(0,0);
\draw (n1) to[isource,l={$\SI{1}{\ampere}$}] (n2);
\node[ground] at (1.5,0){}; \node[above,Vcol] at (n1){$V_1$}; \node[above,Vcol] at (n2){$V_2$};
\end{circuitikz}}
\resposta{$V_1=\SI{15.6}{\volt}$; $V_2=\SI{7.2}{\volt}$}


\end{questao}
\begin{questao}[2]{}

Determine os três potenciais da escada resistiva.
\circuitoq{\begin{circuitikz}[american,scale=.85,transform shape]
\draw (0,0) to[isource,l_={$\SI{4}{\ampere}$}] (0,2.8) -- (1.2,2.8) coordinate(n1);
\draw (n1) to[R,l={$\SI{10}{\ohm}$}] (1.2,0);
\draw (n1) to[R,l={$\SI{10}{\ohm}$}] (3.7,2.8) coordinate(n2);
\draw (n2) to[R,l={$\SI{10}{\ohm}$}] (3.7,0);
\draw (n2) to[R,l={$\SI{10}{\ohm}$}] (6.2,2.8) coordinate(n3);
\draw (n3) to[R,l={$\SI{10}{\ohm}$}] (6.2,0) -- (1.2,0) -- (0,0);
\node[ground] at (1.2,0){}; \node[above,Vcol] at(n1){$V_1$};\node[above,Vcol] at(n2){$V_2$};\node[above,Vcol] at(n3){$V_3$};
\end{circuitikz}}
\resposta{$V_1=\SI{25}{\volt}$; $V_2=\SI{10}{\volt}$; $V_3=\SI{5}{\volt}$}


\end{questao}
\begin{questao}[2]{}

Duas fontes de corrente injetam corrente em nós diferentes. Determine $V_1$ e $V_2$.
\circuitoq{\begin{circuitikz}[american,scale=.9,transform shape]
\draw (0,0) to[isource,l_={$\SI{4}{\ampere}$}] (0,2.8) -- (1.5,2.8) coordinate(n1);
\draw (n1) to[R,l={$\SI{4}{\ohm}$}] (1.5,0);
\draw (n1) to[R,l={$\SI{8}{\ohm}$}] (4.5,2.8) coordinate(n2);
\draw (n2) to[R,l={$\SI{4}{\ohm}$}] (4.5,0);
\draw (6.2,0) to[isource,l_={$\SI{2}{\ampere}$}] (6.2,2.8) -- (n2); \draw (6.2,0)--(1.5,0)--(0,0);
\node[ground] at (1.5,0){};\node[above,Vcol] at(n1){$V_1$};\node[above,Vcol] at(n2){$V_2$};
\end{circuitikz}}
\resposta{$V_1=\SI{14}{\volt}$; $V_2=\SI{10}{\volt}$}


\end{questao}
\begin{questao}[2]{}

Determine $V_1$ e $V_2$.
\circuitoq{\begin{circuitikz}[american,scale=.9,transform shape]
\draw (0,0) to[battery1,l_={$\SI{18}{\volt}$}] (0,2.8) to[R,l={$\SI{6}{\ohm}$}] (2.4,2.8) coordinate(n1);
\draw (n1) to[R,l={$\SI{12}{\ohm}$}] (2.4,0);
\draw (n1) to[R,l={$\SI{6}{\ohm}$}] (5.1,2.8) coordinate(n2);
\draw (n2) to[R,l={$\SI{6}{\ohm}$}] (5.1,0);
\draw (6.7,0) to[isource,l_={$\SI{1}{\ampere}$}] (6.7,2.8)--(n2);\draw (6.7,0)--(0,0);
\node[ground] at (2.4,0){};\node[above,Vcol] at(n1){$V_1$};\node[above,Vcol] at(n2){$V_2$};
\end{circuitikz}}
\resposta{$V_1=\SI{10.5}{\volt}$; $V_2=\SI{8.25}{\volt}$}


\end{questao}
\begin{questao}[2]{}

A fonte horizontal força $\SI{2}{\ampere}$ do nó 1 para o nó 2. Determine as tensões nodais.
\circuitoq{\begin{circuitikz}[american,scale=.9,transform shape]
\draw (0,0) to[battery1,l_={$\SI{24}{\volt}$}] (0,2.8) to[R,l={$\SI{6}{\ohm}$}] (2.6,2.8) coordinate(n1);
\draw (n1) to[R,l={$\SI{12}{\ohm}$}] (2.6,0);
\draw (n1) to[isource,l={$\SI{2}{\ampere}$}] (5.2,2.8) coordinate(n2);
\draw (n2) to[R,l={$\SI{6}{\ohm}$}] (5.2,0)--(0,0);
\node[ground] at (2.6,0){};\node[above,Vcol] at(n1){$V_1$};\node[above,Vcol] at(n2){$V_2$};
\end{circuitikz}}
\resposta{$V_1=\SI{8}{\volt}$; $V_2=\SI{12}{\volt}$}


\end{questao}
\begin{questao}[2]{}

Determine $V_1$, $V_2$ e a corrente no resistor de $\SI{4}{\ohm}$ entre os nós.
\circuitoq{\begin{circuitikz}[american,scale=.9,transform shape]
\draw (0,0) to[battery1,l_={$\SI{10}{\volt}$}] (0,2.8) to[R,l={$\SI{2}{\ohm}$}] (2.4,2.8) coordinate(n1);
\draw (n1) to[R,l={$\SI{4}{\ohm}$}] (5,2.8) coordinate(n2);
\draw (n1) to[R,l={$\SI{4}{\ohm}$}] (2.4,0);
\draw (n2) to[R,l={$\SI{8}{\ohm}$}] (5,0)--(0,0);
\node[ground] at (2.4,0){};\node[above,Vcol] at(n1){$V_1$};\node[above,Vcol] at(n2){$V_2$};
\end{circuitikz}}
\resposta{$V_1=\SI{6}{\volt}$; $V_2=\SI{4}{\volt}$; $i_{12}=\SI{0.5}{\ampere}$}


\end{questao}

\blocoaprofundamento
\begin{questao}[3]{}

No circuito, $i_1$ e $i_2$ têm os sentidos indicados. Determine $i_1+i_2$.
\circuitoq{\begin{circuitikz}[american,scale=.92,transform shape]
\draw (0,0) to[isource,l_={$\SI{5}{\ampere}$}] (0,3) -- (1.8,3) coordinate(n);
\draw (n) to[R,l={$\SI{5}{\ohm}$},i>^={$i_2$}] (1.8,0) -- (0,0);
\draw (n) -- (3.9,3) to[cisource,l={$4i_2$},invert] (3.9,0) -- (1.8,0);
\draw (n) -- (6,3) to[R,l={$\SI{1}{\ohm}$},i>^={$i_1$}] (6,0) -- (3.9,0);
\node[ground] at (1.8,0){};
\end{circuitikz}}
\resposta{$i_1+i_2=\SI{15}{\ampere}$}


\end{questao}
\begin{questao}[3]{}

Determine $V_1$, $V_2$ e a corrente no resistor entre os nós.
\circuitoq{\begin{circuitikz}[american,scale=.9,transform shape]
\draw (0,0) to[isource,l_={$\SI{4}{\ampere}$}] (0,3)--(1.5,3) coordinate(n1);
\draw (n1) to[R,l={$\SI{6}{\ohm}$}] (1.5,0);
\draw (n1) to[R,l={$\SI{3}{\ohm}$}] (4.5,3) coordinate(n2);
\draw (n2) to[R,l={$\SI{6}{\ohm}$}] (4.5,0);
\draw (6.2,0) to[isource,l_={$\SI{1}{\ampere}$}] (6.2,3)--(n2); \draw (6.2,0)--(0,0);
\node[ground] at(1.5,0){}; \node[above,Vcol] at(n1){$V_1$}; \node[above,Vcol] at(n2){$V_2$};
\end{circuitikz}}
\resposta{$V_1=\SI{16.8}{\volt}$; $V_2=\SI{13.2}{\volt}$; $i_{12}=\SI{1.2}{\ampere}$ de 1 para 2}


\end{questao}
\begin{questao}[3]{}

Resolva o circuito por supernó e determine também a corrente na fonte de $\SI{12}{\volt}$, positiva de 1 para 2.
\circuitoq{\begin{circuitikz}[american,scale=.9,transform shape]
\draw (0,0) to[isource,l_={$\SI{4}{\ampere}$}] (0,3)--(1.5,3) coordinate(n1);
\draw (n1) to[R,l={$\SI{8}{\ohm}$}] (1.5,0);
\draw (n1) to[battery1,l={$\SI{12}{\volt}$}] (4.7,3) coordinate(n2);
\draw (n2) to[R,l={$\SI{4}{\ohm}$}] (4.7,0)--(0,0);
\node[ground] at(1.5,0){}; \node[above,Vcol] at(n1){$V_1$}; \node[above,Vcol] at(n2){$V_2$};
\end{circuitikz}}
\resposta{$V_1=\dfrac{56}{3}\,\si{\volt}$; $V_2=\dfrac{20}{3}\,\si{\volt}$; $i_s=\dfrac{5}{3}\,\si{\ampere}$ de 1 para 2}


\end{questao}
\begin{questao}[3]{}

Monte o sistema nodal $3\times3$ e determine $V_1$, $V_2$ e $V_3$.
\circuitoq{\begin{circuitikz}[american,scale=.82,transform shape]
\draw (0,0) to[isource,l_={$\SI{5}{\ampere}$}] (0,3)--(1.3,3) coordinate(n1);
\draw (n1) to[R,l={$\SI{4}{\ohm}$}] (1.3,0);
\draw (n1) to[R,l={$\SI{4}{\ohm}$}] (3.8,3) coordinate(n2);
\draw (n2) to[R,l={$\SI{8}{\ohm}$}] (3.8,0);
\draw (n2) to[R,l={$\SI{4}{\ohm}$}] (6.3,3) coordinate(n3);
\draw (n3) to[R,l={$\SI{4}{\ohm}$}] (6.3,0); \draw (6.3,0)--(0,0);
\node[ground] at(1.3,0){}; \node[above,Vcol] at(n1){$V_1$}; \node[above,Vcol] at(n2){$V_2$}; \node[above,Vcol] at(n3){$V_3$};
\end{circuitikz}}
\resposta{$V_1=\dfrac{40}{3}\,\si{\volt}$; $V_2=\dfrac{20}{3}\,\si{\volt}$; $V_3=\dfrac{10}{3}\,\si{\volt}$}


\end{questao}
\begin{questao}[3]{}

Determine a tensão $V$ pelo método nodal.
\circuitoq{\begin{circuitikz}[american,scale=.86,transform shape]
\draw (0,0) to[battery1,l_={$\SI{18}{\volt}$}] (0,2.8) to[R,l={$\SI{6}{\ohm}$}] (2.2,2.8) coordinate(n);
\draw (n) to[R,l={$\SI{6}{\ohm}$}] (2.2,0)--(0,0);
\draw (n) to[R,l={$\SI{3}{\ohm}$}] (4.7,2.8) to[battery1,l={$\SI{12}{\volt}$},invert] (4.7,0)--(2.2,0);
\draw (n)--(6.3,2.8) to[isource,l={$\SI{1}{\ampere}$}] (6.3,0)--(4.7,0);
\node[above,Vcol] at(n){$V$}; \node[ground] at(2.2,0){};
\end{circuitikz}}
\resposta{$V=\SI{9}{\volt}$}


\end{questao}
\begin{questao}[3]{}

O circuito possui uma fonte flutuante e três nós desconhecidos. Determine $V_1$, $V_2$ e $V_3$.
\circuitoq{\begin{circuitikz}[american,scale=.84,transform shape]
\draw (0,0) to[isource,l_={$\SI{2}{\ampere}$}] (0,3)--(1.3,3) coordinate(n1);
\draw (n1) to[R,l={$\SI{10}{\ohm}$}] (1.3,0);
\draw (n1) to[battery1,l={$\SI{8}{\volt}$}] (4.0,3) coordinate(n2);
\draw (n2) to[R,l={$\SI{5}{\ohm}$}] (6.5,3) coordinate(n3);
\draw (n3) to[R,l={$\SI{5}{\ohm}$}] (6.5,0)--(0,0);
\node[ground] at(1.3,0){}; \node[above,Vcol] at(n1){$V_1$}; \node[above,Vcol] at(n2){$V_2$}; \node[above,Vcol] at(n3){$V_3$};
\end{circuitikz}}
\resposta{$V_1=\SI{14}{\volt}$; $V_2=\SI{6}{\volt}$; $V_3=\SI{3}{\volt}$}


\end{questao}
\begin{questao}[3]{}

A corrente de controle é $i_x=V_1/4$. A fonte dependente injeta $0{,}5i_x$ no nó 2. Determine $V_1$ e $V_2$.
\circuitoq{\begin{circuitikz}[american,scale=.9,transform shape]
\draw (0,0) to[isource,l_={$\SI{4}{\ampere}$}] (0,3)--(1.5,3) coordinate(n1);
\draw (n1) to[R,l={$\SI{4}{\ohm}$},i>^={$i_x$}] (1.5,0);
\draw (n1) to[R,l={$\SI{8}{\ohm}$}] (4.5,3) coordinate(n2);
\draw (n2) to[R,l={$\SI{8}{\ohm}$}] (4.5,0);
\draw (6.3,0) to[cisource,l_={$0{,}5i_x$}] (6.3,3)--(n2); \draw (6.3,0)--(0,0);
\node[ground] at(1.5,0){}; \node[above,Vcol] at(n1){$V_1$}; \node[above,Vcol] at(n2){$V_2$};
\end{circuitikz}}
\resposta{$V_1=\SI{16}{\volt}$; $V_2=\SI{16}{\volt}$}


\end{questao}
\begin{questao}[3]{}

A fonte de tensão controlada impõe $V_2=2V_1$. Determine os potenciais nodais e a corrente no resistor de $\SI{8}{\ohm}$, positiva de 1 para 2.
\circuitoq{\begin{circuitikz}[american,scale=.9,transform shape]
\draw (0,0) to[isource,l_={$\SI{3}{\ampere}$}] (0,3)--(1.6,3) coordinate(n1);
\draw (n1) to[R,l={$\SI{4}{\ohm}$}] (1.6,0);
\draw (n1) to[R,l={$\SI{8}{\ohm}$}] (4.6,3) coordinate(n2);
\draw (n2) to[cvsource,l={$2V_1$}] (4.6,0)--(0,0);
\node[ground] at(1.6,0){}; \node[above,Vcol] at(n1){$V_1$}; \node[above,Vcol] at(n2){$V_2$};
\end{circuitikz}}
\resposta{$V_1=\SI{24}{\volt}$; $V_2=\SI{48}{\volt}$; $i_{12}=\SI{-3}{\ampere}$}


\end{questao}
\begin{questao}[3]{}

No circuito em ponte, determine os três potenciais e a corrente no resistor de $\SI{10}{\ohm}$ entre os nós 2 e 3.
\circuitoq{\begin{circuitikz}[american,scale=.80,transform shape]
\coordinate (g) at (0,-1.6);
\draw (g) to[isource,l_={$\SI{6}{\ampere}$}] (0,1.6)--(1.3,1.6) coordinate(n1);
\draw (n1) to[R,l={$\SI{4}{\ohm}$}] (4.2,3.0) coordinate(n2);
\draw (n1) to[R,l_={$\SI{6}{\ohm}$}] (4.2,0.2) coordinate(n3);
\draw (n2) to[R,l={$\SI{10}{\ohm}$}] (n3);
\draw (n2) to[R,l={$\SI{8}{\ohm}$}] (4.2,-1.6);
\draw (n3) to[R,l_={$\SI{12}{\ohm}$}] (4.2,-1.6); \draw (4.2,-1.6)--(g);
\node[ground] at(g){}; \node[above,Vcol] at(n1){$V_1$}; \node[above,Vcol] at(n2){$V_2$}; \node[right,Vcol] at(n3){$V_3$};
\end{circuitikz}}
\resposta{$V_1=\SI{43.2}{\volt}$; $V_2=V_3=\SI{28.8}{\volt}$; $i_{23}=0$}


\end{questao}
\begin{questao}[3]{}

Determine $V_1$ e $V_2$. Em seguida, transforme o ramo da fonte de $\SI{24}{\volt}$ e $\SI{6}{\ohm}$ em Norton e confirme o resultado.
\circuitoq{\begin{circuitikz}[american,scale=.88,transform shape]
\draw (0,0) to[battery1,l_={$\SI{24}{\volt}$}] (0,3) to[R,l={$\SI{6}{\ohm}$}] (2.4,3) coordinate(n1);
\draw (n1) to[R,l={$\SI{12}{\ohm}$}] (2.4,0)--(0,0);
\draw (n1) to[R,l={$\SI{4}{\ohm}$}] (5.2,3) coordinate(n2);
\draw (n2) to[R,l={$\SI{8}{\ohm}$}] (5.2,0)--(2.4,0);
\node[ground] at(2.4,0){}; \node[above,Vcol] at(n1){$V_1$}; \node[above,Vcol] at(n2){$V_2$};
\end{circuitikz}}
\resposta{$V_1=\SI{12}{\volt}$; $V_2=\SI{8}{\volt}$}


\end{questao}
\begin{questao}[3]{}

Determine $V_1$, $V_2$ e a corrente na fonte de $\SI{10}{\volt}$.
\circuitoq{\begin{circuitikz}[american,scale=.9,transform shape]
\draw (0,0) to[isource,l_={$\SI{2}{\ampere}$}] (0,3)--(1.5,3) coordinate(n1);
\draw (n1) to[R,l={$\SI{5}{\ohm}$}] (1.5,0);
\draw (n1) to[battery1,l={$\SI{10}{\volt}$}] (4.7,3) coordinate(n2);
\draw (n2) to[R,l={$\SI{10}{\ohm}$}] (4.7,0)--(0,0);
\node[ground] at(1.5,0){}; \node[above,Vcol] at(n1){$V_1$}; \node[above,Vcol] at(n2){$V_2$};
\end{circuitikz}}
\resposta{$V_1=\SI{10}{\volt}$; $V_2=\SI{0}{\volt}$; $i_s=0$}


\end{questao}
\begin{questao}[3]{}

A fonte dependente transfere corrente do nó 1 para o nó 2 no valor $0{,}05V_2$. Determine $V_1$, $V_2$ e a potência fornecida pela fonte de $\SI{5}{\ampere}$.
\circuitoq{\begin{circuitikz}[american,scale=.9,transform shape]
\draw (0,0) to[isource,l_={$\SI{5}{\ampere}$}] (0,3)--(1.5,3) coordinate(n1);
\draw (n1) to[R,l={$\SI{5}{\ohm}$}] (1.5,0);
\draw (n1) to[R,l={$\SI{10}{\ohm}$}] (4.5,3) coordinate(n2);
\draw (n2) to[R,l={$\SI{10}{\ohm}$}] (4.5,0);
\draw (n1) to[cisource,l={$0{,}05V_2$}] (n2); \draw (4.5,0)--(0,0);
\node[ground] at(1.5,0){}; \node[above,Vcol] at(n1){$V_1$}; \node[above,Vcol] at(n2){$V_2$};
\end{circuitikz}}
\resposta{$V_1=\SI{18.75}{\volt}$; $V_2=\SI{12.5}{\volt}$; $P=\SI{93.75}{\watt}$ fornecidos}


\end{questao}

\blocodesafio
\begin{questao}[4]{}

Resolva o circuito por \textbf{análise nodal modificada} (MNA), usando como incógnitas $V_1$, $V_2$ e a corrente $i_s$ da fonte de $\SI{10}{\volt}$.
\circuitoq{\begin{circuitikz}[american,scale=.9,transform shape]
\draw (0,0) to[isource,l_={$\SI{3}{\ampere}$}] (0,3)--(1.5,3) coordinate(n1);
\draw (n1) to[R,l={$\SI{10}{\ohm}$}] (1.5,0);
\draw (n1) to[battery1,l={$\SI{10}{\volt}$}] (4.7,3) coordinate(n2);
\draw (n2) to[R,l={$\SI{20}{\ohm}$}] (4.7,0)--(0,0);
\node[ground] at(1.5,0){}; \node[above,Vcol] at(n1){$V_1$}; \node[above,Vcol] at(n2){$V_2$};
\end{circuitikz}}
\resposta{$V_1=\dfrac{70}{3}\,\si{\volt}$; $V_2=\dfrac{40}{3}\,\si{\volt}$; $i_s=\dfrac{2}{3}\,\si{\ampere}$ de 1 para 2}


\end{questao}
\begin{questao}[4]{}

Determine o equivalente de Thévenin visto na porta $a$--terra usando análise nodal e uma fonte de teste para $R_{Th}$.
\circuitoq{\begin{circuitikz}[american,scale=.9,transform shape]
\draw (0,0) to[isource,l_={$\SI{4}{\ampere}$}] (0,3)--(1.5,3) coordinate(n1);
\draw (n1) to[R,l={$\SI{10}{\ohm}$}] (1.5,0);
\draw (n1) to[R,l={$\SI{20}{\ohm}$}] (4.7,3) coordinate(a);
\draw (a)--(5.8,3); \draw (5.8,2.75)--(5.8,3.25); \node[right] at(5.9,3){$a$};
\draw (1.5,0)--(0,0); \node[ground] at(1.5,0){};
\node[above,Vcol] at(n1){$V_1$};
\end{circuitikz}}
\resposta{$V_{Th}=\SI{40}{\volt}$; $R_{Th}=\SI{30}{\ohm}$}


\end{questao}
\begin{questao}[4]{}

Escolha $R_L$ para que a tensão nodal não ultrapasse $\SI{24}{\volt}$. Determine o maior valor admissível de $R_L$ e a potência nele dissipada no limite.
\circuitoq{\begin{circuitikz}[american,scale=.9,transform shape]
\draw (0,0) to[isource,l_={$\SI{2}{\ampere}$}] (0,3)--(1.7,3) coordinate(n);
\draw (n) to[R,l={$\SI{30}{\ohm}$}] (1.7,0)--(0,0);
\draw (n)--(4.2,3) to[R,l={$R_L$}] (4.2,0)--(1.7,0);
\node[ground] at(1.7,0){}; \node[above,Vcol] at(n){$V$};
\end{circuitikz}}
\resposta{$R_L\le\SI{20}{\ohm}$; no limite, $P_L=\SI{28.8}{\watt}$}


\end{questao}
\begin{questao}[4]{}

O circuito deve operar com $V=\SI{8}{\volt}$. Determine $R_L$ e a sensibilidade $\left.\dfrac{dV}{dR_L}\right|_{R_L}$.
\circuitoq{\begin{circuitikz}[american,scale=.9,transform shape]
\draw (0,0) to[battery1,l_={$\SI{12}{\volt}$}] (0,3) to[R,l={$\SI{4}{\ohm}$}] (2.6,3) coordinate(n);
\draw (n) to[R,l={$R_L$}] (2.6,0)--(0,0);
\node[ground] at(2.6,0){}; \node[above,Vcol] at(n){$V$};
\end{circuitikz}}
\resposta{$R_L=\SI{8}{\ohm}$; $\left.dV/dR_L\right|_{8\Omega}=\SI{0.333}{\volt\per\ohm}$}


\end{questao}
\begin{questao}[4]{}

Monte a matriz nodal da escada de quatro nós, resolva-a e determine a resistência de entrada vista pela fonte de corrente.
\circuitoq{\begin{circuitikz}[american,scale=.72,transform shape]
\draw (0,0) to[isource,l_={$\SI{4}{\ampere}$}] (0,3)--(1.1,3) coordinate(n1);
\draw (n1) to[R,l={$\SI{10}{\ohm}$}] (1.1,0);
\draw (n1) to[R,l={$\SI{10}{\ohm}$}] (3.3,3) coordinate(n2);
\draw (n2) to[R,l={$\SI{10}{\ohm}$}] (3.3,0);
\draw (n2) to[R,l={$\SI{10}{\ohm}$}] (5.5,3) coordinate(n3);
\draw (n3) to[R,l={$\SI{10}{\ohm}$}] (5.5,0);
\draw (n3) to[R,l={$\SI{10}{\ohm}$}] (7.7,3) coordinate(n4);
\draw (n4) to[R,l={$\SI{10}{\ohm}$}] (7.7,0); \draw (7.7,0)--(0,0);
\node[ground] at(1.1,0){}; \node[above,Vcol] at(n1){$V_1$}; \node[above,Vcol] at(n2){$V_2$}; \node[above,Vcol] at(n3){$V_3$}; \node[above,Vcol] at(n4){$V_4$};
\end{circuitikz}}
\resposta{$V_1=\dfrac{520}{21}\,\si{\volt}$; $V_2=\dfrac{200}{21}\,\si{\volt}$; $V_3=\dfrac{80}{21}\,\si{\volt}$; $V_4=\dfrac{40}{21}\,\si{\volt}$; $R_{in}=\dfrac{130}{21}\,\ohm\approx\SI{6.19}{\ohm}$}


\end{questao}
\begin{questao}[4]{}

A fonte de tensão controlada satisfaz $V_1-V_2=5i_x$, com $i_x=V_2/5$. Determine $V_1$, $V_2$, a corrente na fonte dependente e sua potência.
\circuitoq{\begin{circuitikz}[american,scale=.9,transform shape]
\draw (0,0) to[isource,l_={$\SI{5}{\ampere}$}] (0,3)--(1.5,3) coordinate(n1);
\draw (n1) to[R,l={$\SI{10}{\ohm}$}] (1.5,0);
\draw (n1) to[cvsource,l={$5i_x$}] (4.7,3) coordinate(n2);
\draw (n2) to[R,l={$\SI{5}{\ohm}$},i>^={$i_x$}] (4.7,0)--(0,0);
\node[ground] at(1.5,0){}; \node[above,Vcol] at(n1){$V_1$}; \node[above,Vcol] at(n2){$V_2$};
\end{circuitikz}}
\resposta{$V_1=\SI{25}{\volt}$; $V_2=\SI{12.5}{\volt}$; $i_s=\SI{2.5}{\ampere}$ de 1 para 2; a fonte dependente absorve $\SI{31.25}{\watt}$}


\end{questao}
\begin{questao}[4]{}

A fonte dependente injeta $gV$ no nó. Obtenha $V(g)$, determine o valor crítico de $g$ e calcule $V$ para $g=\SI{0.04}{\ensuremath{\mathrm{S}}}$.
\circuitoq{\begin{circuitikz}[american,scale=.95,transform shape]
\draw (0,0) to[isource,l_={$\SI{1}{\ampere}$}] (0,3)--(2,3) coordinate(n);
\draw (n) to[R,l={$\SI{20}{\ohm}$}] (2,0)--(0,0);
\draw (4.2,0) to[cisource,l_={$gV$}] (4.2,3)--(n); \draw (4.2,0)--(2,0);
\node[ground] at(2,0){}; \node[above,Vcol] at(n){$V$};
\end{circuitikz}}
\resposta{$V(g)=1/(0{,}05-g)$; $g_c=\SI{0.05}{\ensuremath{\mathrm{S}}}$; para $g=\SI{0.04}{\ensuremath{\mathrm{S}}}$, $V=\SI{100}{\volt}$}


\end{questao}
\begin{questao}[4]{}

Determine o equivalente de Thévenin visto na porta. A fonte dependente injeta $0{,}05V_1$ no nó 2. Para $R_{Th}$, mantenha a fonte dependente ativa e use uma fonte de teste.
\circuitoq{\begin{circuitikz}[american,scale=.9,transform shape]
\draw (0,0) to[isource,l_={$\SI{2}{\ampere}$}] (0,3)--(1.5,3) coordinate(n1);
\draw (n1) to[R,l={$\SI{10}{\ohm}$}] (1.5,0);
\draw (n1) to[R,l={$\SI{5}{\ohm}$}] (4.7,3) coordinate(n2);
\draw (6.5,0) to[cisource,l_={$0{,}05V_1$}] (6.5,3)--(n2); \draw (6.5,0)--(0,0);
\draw (n2)--(5.7,3); \draw (5.7,2.75)--(5.7,3.25); \node[right] at(5.8,3){porta};
\node[ground] at(1.5,0){}; \node[above,Vcol] at(n1){$V_1$}; \node[above,Vcol] at(n2){$V_2$};
\end{circuitikz}}
\resposta{$V_{Th}=\SI{50}{\volt}$; $R_{Th}=\SI{30}{\ohm}$}


\end{questao}
\begin{questao}[4]{}

A ponte é alimentada por uma fonte de corrente. Determine a condição algébrica para corrente nula no resistor central e, para os valores da figura, calcule $V_1$, $V_2$ e $V_3$.
\circuitoq{\begin{circuitikz}[american,scale=.80,transform shape]
\coordinate (g) at (0,-1.6);
\draw (g) to[isource,l_={$\SI{6}{\ampere}$}] (0,1.6)--(1.3,1.6) coordinate(n1);
\draw (n1) to[R,l={$R_1=\SI{4}{\ohm}$}] (4.2,3.0) coordinate(n2);
\draw (n1) to[R,l_={$R_3=\SI{6}{\ohm}$}] (4.2,0.2) coordinate(n3);
\draw (n2) to[R,l={$R_b=\SI{5}{\ohm}$}] (n3);
\draw (n2) to[R,l={$R_2=\SI{8}{\ohm}$}] (4.2,-1.6);
\draw (n3) to[R,l_={$R_4=\SI{12}{\ohm}$}] (4.2,-1.6); \draw (4.2,-1.6)--(g);
\node[ground] at(g){}; \node[above,Vcol] at(n1){$V_1$}; \node[above,Vcol] at(n2){$V_2$}; \node[right,Vcol] at(n3){$V_3$};
\end{circuitikz}}
\resposta{Equilíbrio quando $R_1/R_2=R_3/R_4$; para a figura, $V_1=\SI{43.2}{\volt}$ e $V_2=V_3=\SI{28.8}{\volt}$}


\end{questao}
\begin{questao}[4]{}

No circuito de síntese, deseja-se obter $V_1=\SI{18}{\volt}$ e $V_2=\SI{12}{\volt}$. Determine os valores e sentidos das fontes $I_1$ e $I_2$ e verifique o balanço de potência.
\circuitoq{\begin{circuitikz}[american,scale=.9,transform shape]
\draw (0,0) to[isource,l_={$I_1$}] (0,3)--(1.5,3) coordinate(n1);
\draw (n1) to[R,l={$\SI{6}{\ohm}$}] (1.5,0);
\draw (n1) to[R,l={$\SI{3}{\ohm}$}] (4.7,3) coordinate(n2);
\draw (n2) to[R,l={$\SI{4}{\ohm}$}] (4.7,0);
\draw (6.4,0) to[isource,l_={$I_2$}] (6.4,3)--(n2); \draw (6.4,0)--(0,0);
\node[ground] at(1.5,0){}; \node[above,Vcol] at(n1){$V_1=\SI{18}{\volt}$}; \node[above,Vcol] at(n2){$V_2=\SI{12}{\volt}$};
\end{circuitikz}}
\resposta{$I_1=\SI{5}{\ampere}$ e $I_2=\SI{1}{\ampere}$, ambas injetando corrente; potência total fornecida $=\SI{102}{\watt}$}


\end{questao}

\gabarito[Capítulo 7 --- Análise Nodal]

\section{Método de Maxwell (Análise de Malhas)}

\blocoaprofundamento

\blocoaprofundamento
\begin{questao}[2]{}

No circuito, a fonte de tensão dependente fornece $2{,}5I_1$ volts e atua na
malha 2. Determine $I_1$, $I_2$ e a corrente real no resistor central.

\circuitoq{
\begin{circuitikz}[american,scale=.92,transform shape,line width=.8pt]
\draw (0,0) to[\fonteV,l^={$\SI{20}{\volt}$}] (0,3.2)
      to[R,l={$\SI{5}{\ohm}$}] (3.2,3.2) coordinate(t);
\draw (t) to[R,l={$\SI{5}{\ohm}$}] (3.2,0) coordinate(b) -- (0,0);
\draw (t) to[R,l={$\SI{5}{\ohm}$}] (6.4,3.2)
      to[cvsource,l={$2{,}5I_1$}] (6.4,0) -- (b);
\draw[->,loopcol,line width=1.2pt] (1.75,1.00) arc[start angle=0,end angle=-310,radius=.46cm];
\node[loopcol,fill=white,inner sep=1pt] at (1.75,1.00){\footnotesize$I_1$};
\draw[->,loopcol,line width=1.2pt] (4.75,1.00) arc[start angle=0,end angle=-310,radius=.46cm];
\node[loopcol,fill=white,inner sep=1pt] at (4.75,1.00){\footnotesize$I_2$};
\end{circuitikz}}

\resposta{$I_1=\SI{3.2}{\ampere}$; $I_2=\SI{2.4}{\ampere}$; $I_c=\SI{0.8}{\ampere}$}


\end{questao}
\begin{questao}[2]{}

A fonte de corrente dependente está na periferia da malha 2 e impõe
$I_2=0{,}5I_1$. Determine $I_1$, $I_2$ e o módulo da tensão sobre essa fonte.

\circuitoq{
\begin{circuitikz}[american,scale=.92,transform shape,line width=.8pt]
\draw (0,0) to[\fonteV,l^={$\SI{30}{\volt}$}] (0,3.2)
      to[R,l={$\SI{5}{\ohm}$}] (3.2,3.2) coordinate(t);
\draw (t) to[R,l={$\SI{5}{\ohm}$}] (3.2,0) coordinate(b) -- (0,0);
\draw (t) to[R,l={$\SI{10}{\ohm}$}] (6.4,3.2)
      to[cisource,l={$0{,}5I_1$},invert] (6.4,0) -- (b);
\draw[->,loopcol,line width=1.2pt] (1.75,1.00) arc[start angle=0,end angle=-310,radius=.46cm];
\node[loopcol,fill=white,inner sep=1pt] at (1.75,1.00){\footnotesize$I_1$};
\draw[->,loopcol,line width=1.2pt] (4.75,1.00) arc[start angle=0,end angle=-310,radius=.46cm];
\node[loopcol,fill=white,inner sep=1pt] at (4.75,1.00){\footnotesize$I_2$};
\end{circuitikz}}

\resposta{$I_1=\SI{4}{\ampere}$; $I_2=\SI{2}{\ampere}$; $|v_s|=\SI{10}{\volt}$}


\end{questao}

\blocodesafio

\blocodesafio
\begin{questao}[3]{}

A fonte dependente vale $4i_x$ volts, com $i_x=I_1-I_2$ no resistor central.
Determine as correntes de malha.

\circuitoq{
\begin{circuitikz}[american,scale=.92,transform shape,line width=.8pt]
\draw (0,0) to[\fonteV,l^={$\SI{24}{\volt}$}] (0,3.2)
      to[R,l={$\SI{4}{\ohm}$}] (3.2,3.2) coordinate(t);
\draw (t) to[R,l={$\SI{4}{\ohm}$},i>^={$i_x$}] (3.2,0) coordinate(b) -- (0,0);
\draw (t) to[R,l={$\SI{8}{\ohm}$}] (6.4,3.2)
      to[cvsource,l={$4i_x$}] (6.4,0) -- (b);
\draw[->,loopcol,line width=1.2pt] (1.75,1.00) arc[start angle=0,end angle=-310,radius=.46cm];
\node[loopcol,fill=white,inner sep=1pt] at (1.75,1.00){\footnotesize$I_1$};
\draw[->,loopcol,line width=1.2pt] (4.75,1.00) arc[start angle=0,end angle=-310,radius=.46cm];
\node[loopcol,fill=white,inner sep=1pt] at (4.75,1.00){\footnotesize$I_2$};
\end{circuitikz}}

\resposta{$I_1=\SI{4}{\ampere}$; $I_2=\SI{2}{\ampere}$}


\end{questao}
\begin{questao}[3]{}

A tensão de controle $v_x$ é medida no resistor de $\SI{4}{\ohm}$ exclusivo da
malha 1. A fonte dependente vale $0{,}5v_x$. Determine $I_1$ e $I_2$.

\circuitoq{
\begin{circuitikz}[american,scale=.92,transform shape,line width=.8pt]
\draw (0,0) to[\fonteV,l^={$\SI{24}{\volt}$}] (0,3.2)
      to[R,l_={$\SI{4}{\ohm}$},v^>={$v_x$}] (3.2,3.2) coordinate(t);
\draw (t) to[R,l={$\SI{4}{\ohm}$}] (3.2,0) coordinate(b) -- (0,0);
\draw (t) to[R,l={$\SI{4}{\ohm}$}] (6.4,3.2)
      to[cvsource,l={$0{,}5v_x$}] (6.4,0) -- (b);
\draw[->,loopcol,line width=1.2pt] (1.75,1.00) arc[start angle=0,end angle=-310,radius=.46cm];
\node[loopcol,fill=white,inner sep=1pt] at (1.75,1.00){\footnotesize$I_1$};
\draw[->,loopcol,line width=1.2pt] (4.75,1.00) arc[start angle=0,end angle=-310,radius=.46cm];
\node[loopcol,fill=white,inner sep=1pt] at (4.75,1.00){\footnotesize$I_2$};
\end{circuitikz}}

\resposta{$I_1=\SI{4.8}{\ampere}$; $I_2=\SI{3.6}{\ampere}$}


\end{questao}
\begin{questao}[3]{}

A fonte de corrente dependente está no ramo comum e fornece corrente de baixo
para cima igual a $0{,}1v_x$. A tensão $v_x$ está indicada no resistor da malha
1. Determine $I_1$ e $I_2$ por supermalha.

\circuitoq{
\begin{circuitikz}[american,scale=.92,transform shape,line width=.8pt]
\draw (0,0) to[\fonteV,l^={$\SI{40}{\volt}$}] (0,3.2)
      to[R,l_={$\SI{5}{\ohm}$},v^>={$v_x$}] (3.2,3.2) coordinate(t);
\draw (t) to[cisource,l={$0{,}1v_x$},invert] (3.2,0) coordinate(b) -- (0,0);
\draw (t) to[R,l={$\SI{5}{\ohm}$}] (6.4,3.2) -- (6.4,0) -- (b);
\draw[->,loopcol,line width=1.2pt] (1.75,1.00) arc[start angle=0,end angle=-310,radius=.46cm];
\node[loopcol,fill=white,inner sep=1pt] at (1.75,1.00){\footnotesize$I_1$};
\draw[->,loopcol,line width=1.2pt] (4.75,1.00) arc[start angle=0,end angle=-310,radius=.46cm];
\node[loopcol,fill=white,inner sep=1pt] at (4.75,1.00){\footnotesize$I_2$};
\draw[dashed,laranja,line width=.9pt,rounded corners=7pt]
      (-.55,-.5) rectangle (6.95,4.10);
\node[laranja,font=\footnotesize\sffamily,fill=white,inner sep=1.5pt]
      at (3.2,-.80) {supermalha};
\end{circuitikz}}

\resposta{$I_1=\SI{3.2}{\ampere}$; $I_2=\SI{4.8}{\ampere}$}


\end{questao}
\begin{questao}[3]{}

A fonte de corrente dependente do ramo comum fornece, de baixo para cima,
$2i_x$, sendo $i_x=I_1$ no resistor esquerdo. Determine as correntes de malha
por supermalha.

\circuitoq{
\begin{circuitikz}[american,scale=.92,transform shape,line width=.8pt]
\draw (0,0) to[\fonteV,l^={$\SI{40}{\volt}$}] (0,3.2)
      to[R,l_={$\SI{5}{\ohm}$},i>^={$i_x$}] (3.2,3.2) coordinate(t);
\draw (t) to[cisource,l={$2i_x$},invert] (3.2,0) coordinate(b) -- (0,0);
\draw (t) to[R,l={$\SI{5}{\ohm}$}] (6.4,3.2) -- (6.4,0) -- (b);
\draw[->,loopcol,line width=1.2pt] (1.75,1.00) arc[start angle=0,end angle=-310,radius=.46cm];
\node[loopcol,fill=white,inner sep=1pt] at (1.75,1.00){\footnotesize$I_1$};
\draw[->,loopcol,line width=1.2pt] (4.75,1.00) arc[start angle=0,end angle=-310,radius=.46cm];
\node[loopcol,fill=white,inner sep=1pt] at (4.75,1.00){\footnotesize$I_2$};
\draw[dashed,laranja,line width=.9pt,rounded corners=7pt]
      (-.55,-.5) rectangle (6.95,4.10);
\node[laranja,font=\footnotesize\sffamily,fill=white,inner sep=1.5pt]
      at (3.2,-.80) {supermalha};
\end{circuitikz}}

\resposta{$I_1=\SI{2}{\ampere}$; $I_2=\SI{6}{\ampere}$}


\end{questao}
\begin{questao}[4]{}

No circuito de três malhas, a fonte dependente da terceira malha vale
$2{,}5I_1$ volts e impulsiona $I_3$. Determine $I_1$, $I_2$ e $I_3$.

\circuitoq{
\begin{circuitikz}[american,scale=.84,transform shape,line width=.8pt]
\draw (0,0) to[\fonteV,l^={$\SI{20}{\volt}$}] (0,3.2)
      to[R,l={$\SI{5}{\ohm}$}] (2.8,3.2) coordinate(t1);
\draw (t1) to[R,l_={$\SI{5}{\ohm}$}] (2.8,0) coordinate(b1) -- (0,0);
\draw (t1) to[R,l={$\SI{5}{\ohm}$}] (5.6,3.2) coordinate(t2);
\draw (t2) to[R,l_={$\SI{5}{\ohm}$}] (5.6,0) coordinate(b2) -- (b1);
\draw (t2) to[R,l={$\SI{5}{\ohm}$}] (8.4,3.2)
      to[cvsource,l={$2{,}5I_1$}] (8.4,0) -- (b2);
\foreach \x/\n in {1.25/1,4.15/2,7.05/3}{
  \draw[->,loopcol,line width=1.05pt] (\x,1.00) arc[start angle=0,end angle=-310,radius=.40cm];
  \node[loopcol,fill=white,inner sep=1pt] at (\x,1.00){\scriptsize$I_{\n}$};}
\end{circuitikz}}

\resposta{$I_1=\dfrac{8}{3}\,\si{\ampere}$; $I_2=I_3=\dfrac{4}{3}\,\si{\ampere}$}


\end{questao}
\begin{questao}[4]{}

A fonte de corrente dependente está entre as malhas 2 e 3 e impõe
$I_3-I_2=0{,}5I_1$. Determine as três correntes usando uma supermalha envolvendo
as malhas 2 e 3.

\circuitoq{
\begin{circuitikz}[american,scale=.84,transform shape,line width=.8pt]
\draw (0,0) to[\fonteV,l^={$\SI{30}{\volt}$}] (0,3.2)
      to[R,l={$\SI{5}{\ohm}$}] (2.8,3.2) coordinate(t1);
\draw (t1) to[R,l_={$\SI{5}{\ohm}$}] (2.8,0) coordinate(b1) -- (0,0);
\draw (t1) to[R,l={$\SI{5}{\ohm}$}] (5.6,3.2) coordinate(t2);
\draw (t2) to[cisource,l={$0{,}5I_1$},invert] (5.6,0) coordinate(b2) -- (b1);
\draw (t2) to[R,l={$\SI{10}{\ohm}$}] (8.4,3.2) -- (8.4,0) -- (b2);
\foreach \x/\n in {1.25/1,4.15/2,7.05/3}{
  \draw[->,loopcol,line width=1.05pt] (\x,1.00) arc[start angle=0,end angle=-310,radius=.40cm];
  \node[loopcol,fill=white,inner sep=1pt] at (\x,1.00){\scriptsize$I_{\n}$};}
\draw[dashed,laranja,line width=.9pt,rounded corners=7pt]
      (2.25,-.55) rectangle (8.95,4.05);
\node[laranja,font=\footnotesize\sffamily,fill=white,inner sep=1.5pt]
      at (5.6,-.85) {supermalha 2--3};
\end{circuitikz}}

\resposta{$I_1=\SI{3}{\ampere}$; $I_2=0$; $I_3=\SI{1.5}{\ampere}$}


\end{questao}
\begin{questao}[4]{}

A fonte dependente vale $kI_1$ volts. Determine $k$ para que
$I_2=\SI{2}{\ampere}$ e calcule $I_1$.

\circuitoq{
\begin{circuitikz}[american,scale=.92,transform shape,line width=.8pt]
\draw (0,0) to[\fonteV,l^={$\SI{20}{\volt}$}] (0,3.2)
      to[R,l={$\SI{5}{\ohm}$}] (3.2,3.2) coordinate(t);
\draw (t) to[R,l={$\SI{5}{\ohm}$}] (3.2,0) coordinate(b) -- (0,0);
\draw (t) to[R,l={$\SI{5}{\ohm}$}] (6.4,3.2)
      to[cvsource,l={$kI_1$}] (6.4,0) -- (b);
\draw[->,loopcol,line width=1.2pt] (1.75,1.00) arc[start angle=0,end angle=-310,radius=.46cm];
\node[loopcol,fill=white,inner sep=1pt] at (1.75,1.00){\footnotesize$I_1$};
\draw[->,loopcol,line width=1.2pt] (4.75,1.00) arc[start angle=0,end angle=-310,radius=.46cm];
\node[loopcol,fill=white,inner sep=1pt] at (4.75,1.00){\footnotesize$I_2$};
\end{circuitikz}}

\resposta{$I_1=\SI{3}{\ampere}$; $k=\dfrac{5}{3}\,\ohm\approx\SI{1.67}{\ohm}$}


\end{questao}
\begin{questao}[4]{}

A fonte de tensão dependente está no ramo comum e possui tensão, de cima para
baixo, $v_d=2I_1$. Determine $I_1$, $I_2$, a corrente real na fonte dependente e
sua potência, indicando se ela absorve ou fornece energia.

\circuitoq{
\begin{circuitikz}[american,scale=.92,transform shape,line width=.8pt]
\draw (0,0) to[\fonteV,l^={$\SI{20}{\volt}$}] (0,3.2)
      to[R,l={$\SI{4}{\ohm}$}] (3.2,3.2) coordinate(t);
\draw (t) to[cvsource,l_={$v_d$}] (3.2,0) coordinate(b) -- (0,0);
\draw (t) to[R,l={$\SI{4}{\ohm}$}] (6.4,3.2) -- (6.4,0) -- (b);
\draw[->,loopcol,line width=1.2pt] (1.75,1.00) arc[start angle=0,end angle=-310,radius=.46cm];
\node[loopcol,fill=white,inner sep=1pt] at (1.75,1.00){\footnotesize$I_1$};
\draw[->,loopcol,line width=1.2pt] (4.75,1.00) arc[start angle=0,end angle=-310,radius=.46cm];
\node[loopcol,fill=white,inner sep=1pt] at (4.75,1.00){\footnotesize$I_2$};
\end{circuitikz}}

\resposta{$I_1=\dfrac{10}{3}\,\si{\ampere}$; $I_2=\dfrac{5}{3}\,\si{\ampere}$;
$i_d=\dfrac{5}{3}\,\si{\ampere}$ para baixo; $P_d=\dfrac{100}{9}\,\si{\watt}$ absorvidos}


\end{questao}

\blocofixacao
\begin{questao}[1]{}

No desenho abaixo, a corrente de malha segue a convenção usada neste capítulo.
Identifique o sentido de $I_1$.
\circuitoq{
\begin{circuitikz}[american,scale=.9,transform shape]
\draw (0,0) to[\fonteV,l^={$E$}] (0,2.8) to[R,l={$R$}] (4,2.8) -- (4,0)--(0,0);
\fvMeshCW{2.0,1.25}{.42}{$I_1$}
\end{circuitikz}}
\resposta{Horário}


\end{questao}
\begin{questao}[1]{}

Duas malhas horárias compartilham um resistor $R_c$. Qual é a corrente real no
resistor, tomando como positivo o sentido de $I_1$ no ramo comum?
\resposta{$i_c=I_1-I_2$}


\end{questao}
\begin{questao}[1]{}

Um circuito planar conectado possui $b=11$ ramos e $n=7$ nós. Quantas correntes
de malha independentes são necessárias?
\resposta{$m=5$}


\end{questao}
\begin{questao}[1]{}

Uma fonte ideal de corrente está na periferia de uma única malha e sua seta
coincide com a corrente horária $I_2$. O que pode ser escrito imediatamente?
\resposta{$I_2=I_s$}


\end{questao}
\begin{questao}[1]{}

Uma fonte ideal de corrente está no ramo comum a duas malhas. Qual técnica deve
ser usada?
\resposta{Supermalha, mais a equação de restrição da fonte}


\end{questao}
\begin{questao}[1]{}

Na matriz de uma rede resistiva com correntes de malha todas horárias, o que
representam $R_{kk}$ e $R_{jk}$, $j\ne k$?
\resposta{$R_{kk}$ é a soma das resistências da malha $k$; $R_{jk}$ é o negativo
da resistência compartilhada}


\end{questao}
\begin{questao}[1]{}

Ao resolver uma malha adotada no sentido horário, obtém-se
$I_3=-\SI{0.8}{\ampere}$. Interprete fisicamente.
\resposta{A corrente real tem módulo $\SI{0.8}{\ampere}$ e circula no sentido anti-horário}


\end{questao}

\blocoaplicacao
\begin{questao}[2]{}

Duas malhas triangulares compartilham o resistor diagonal de
$\SI{6}{\ohm}$. Determine $I_1$, $I_2$ e a corrente no ramo compartilhado.
\circuitoq{
\begin{circuitikz}[american,scale=.88,transform shape,line width=.8pt]
\coordinate (A) at (0,0); \coordinate (B) at (3.4,4.2); \coordinate (C) at (6.8,0);
\draw (A) to[\fonteV,l^={$\SI{18}{\volt}$}] (0,2.0)
      to[R,l={$\SI{4}{\ohm}$}] (B);
\draw (B) to[R,l={$\SI{3}{\ohm}$}] (5.2,2.1)
      to[R,l={$\SI{3}{\ohm}$}] (C);
\draw (C)--(A);
\draw (B) to[R,l={$\SI{6}{\ohm}$}] (3.4,0)--(A);
\fvMeshCW{1.75,1.45}{.38}{$I_1$}
\fvMeshCW{5.05,1.45}{.38}{$I_2$}
\end{circuitikz}}
\resposta{$I_1=\dfrac{18}{7}\,\si{\ampere}$; $I_2=\dfrac{9}{7}\,\si{\ampere}$;
$i_c=\dfrac{9}{7}\,\si{\ampere}$}


\end{questao}
\begin{questao}[2]{}

A fonte de corrente de $\SI{1.5}{\ampere}$ está na periferia da malha 2 e sua
seta coincide com $I_2$. Determine $I_1$, $I_2$ e a corrente no resistor
central de $\SI{5}{\ohm}$.
\circuitoq{
\begin{circuitikz}[american,scale=.92,transform shape,line width=.8pt]
\draw (0,0) to[\fonteV,l^={$\SI{30}{\volt}$}] (0,3.2)
      to[R,l={$\SI{5}{\ohm}$}] (3.2,3.2) coordinate(t);
\draw (t) to[R,l_={$\SI{5}{\ohm}$}] (3.2,0) coordinate(b)--(0,0);
\draw (t) to[R,l={$\SI{10}{\ohm}$}] (6.4,3.2)
      to[isource,l={$\SI{1.5}{\ampere}$}] (6.4,0)--(b);
\fvMeshCW{1.55,1.05}{.40}{$I_1$}
\fvMeshCW{4.85,1.05}{.40}{$I_2$}
\end{circuitikz}}
\resposta{$I_1=\SI{3.75}{\ampere}$; $I_2=\SI{1.5}{\ampere}$;
$i_c=\SI{2.25}{\ampere}$}


\end{questao}
\begin{questao}[2]{}

O ramo comum contém um resistor de $\SI{4}{\ohm}$ em série com uma fonte de
$\SI{6}{\volt}$, com terminal positivo no topo. Determine $I_1$ e $I_2$.
\circuitoq{
\begin{circuitikz}[american,scale=.9,transform shape,line width=.8pt]
\draw (0,0) to[\fonteV,l^={$\SI{24}{\volt}$}] (0,3.6)
      to[R,l={$\SI{4}{\ohm}$}] (3.4,3.6) coordinate(t);
\draw (t) to[R,l_={$\SI{4}{\ohm}$}] (3.4,2.0)
      to[\fonteV,l_={$\SI{6}{\volt}$}] (3.4,0) coordinate(b)--(0,0);
\draw (t) to[R,l={$\SI{8}{\ohm}$}] (6.8,3.6)--(6.8,0)--(b);
\fvMeshCW{1.65,1.15}{.40}{$I_1$}
\fvMeshCW{5.15,1.15}{.40}{$I_2$}
\end{circuitikz}}
\resposta{$I_1=\SI{3}{\ampere}$; $I_2=\SI{1.5}{\ampere}$}


\end{questao}
\begin{questao}[2]{}

A rede possui três malhas triangulares em torno do nó central $O$. Monte a
matriz por inspeção e determine as três correntes.
\circuitoq{
\begin{circuitikz}[american,scale=.78,transform shape,line width=.8pt]
\coordinate (A) at (0,0); \coordinate (B) at (4,6); \coordinate (C) at (8,0); \coordinate (O) at (4,2.1);

\draw (A) to[\fonteV,l^={$\SI{16}{\volt}$}] (1.15,1.75)
      to[R,l={$\SI{4}{\ohm}$}] (B);

\draw (B) to[R,l={$\SI{5}{\ohm}$}] (6.85,1.75)
      to[\fonteV,l={$\SI{11}{\volt}$}] (C);

\draw (C) to[R,l_={$\SI{6}{\ohm}$}] (4.7,0)
      to[\fonteV,l_={$\SI{1}{\volt}$}] (A);

\draw (B) to[R,l={$\SI{2}{\ohm}$}] (O);
\draw (C) to[R,l={$\SI{3}{\ohm}$}] (O);
\draw (A) to[R,l_={$\SI{1}{\ohm}$}] (O);
\fvMeshCW{2.8,3.5}{.34}{$I_1$}
\fvMeshCW{5.4,3.5}{.34}{$I_2$}
\fvMeshCW{4.0,.95}{.34}{$I_3$}
\node[fill=white,inner sep=1pt] at(O){$O$};
\end{circuitikz}}
\resposta{$I_1=\SI{3}{\ampere}$; $I_2=\SI{2}{\ampere}$; $I_3=\SI{1}{\ampere}$}


\end{questao}
\begin{questao}[2]{}

As duas fontes do circuito atuam em oposição. Determine as correntes horárias e
interprete o sinal de $I_2$.
\circuitoq{
\begin{circuitikz}[american,scale=.92,transform shape,line width=.8pt]
\draw (0,0) to[\fonteV,l^={$\SI{20}{\volt}$}] (0,3.2)
      to[R,l={$\SI{6}{\ohm}$}] (3.2,3.2) coordinate(t);
\draw (t) to[R,l_={$\SI{4}{\ohm}$}] (3.2,0) coordinate(b)--(0,0);
\draw (t) to[R,l={$\SI{6}{\ohm}$}] (6.4,3.2)
      to[\fonteV,l_={$\SI{12}{\volt}$}] (6.4,0)--(b);
\fvMeshCW{1.55,1.05}{.40}{$I_1$}
\fvMeshCW{4.85,1.05}{.40}{$I_2$}
\end{circuitikz}}
\resposta{$I_1=\dfrac{38}{21}\,\si{\ampere}$;
$I_2=-\dfrac{10}{21}\,\si{\ampere}$}


\end{questao}
\begin{questao}[2]{}

Escolha o valor e a polaridade da fonte $E$ para que a corrente no resistor
compartilhado de $\SI{4}{\ohm}$ seja nula.
\circuitoq{
\begin{circuitikz}[american,scale=.92,transform shape,line width=.8pt]
\draw (0,0) to[\fonteV,l^={$\SI{24}{\volt}$}] (0,3.2)
      to[R,l={$\SI{4}{\ohm}$}] (3.2,3.2) coordinate(t);
\draw (t) to[R,l_={$\SI{4}{\ohm}$}] (3.2,0) coordinate(b)--(0,0);
\draw (t) to[R,l={$\SI{8}{\ohm}$}] (6.4,3.2)
      to[\fonteV,l_={$E$}] (6.4,0)--(b);
\fvMeshCW{1.55,1.05}{.40}{$I_1$}
\fvMeshCW{4.85,1.05}{.40}{$I_2$}
\end{circuitikz}}
\resposta{$E=\SI{48}{\volt}$, impulsionando $I_2$ no sentido horário}


\end{questao}
\begin{questao}[2]{}

Determine as correntes de malha e faça o balanço de potência.
\circuitoq{
\begin{circuitikz}[american,scale=.92,transform shape,line width=.8pt]
\draw (0,0) to[\fonteV,l^={$\SI{36}{\volt}$}] (0,3.2)
      to[R,l={$\SI{6}{\ohm}$}] (3.2,3.2) coordinate(t);
\draw (t) to[R,l_={$\SI{3}{\ohm}$}] (3.2,0) coordinate(b)--(0,0);
\draw (t) to[R,l={$\SI{6}{\ohm}$}] (6.4,3.2)
      to[\fonteV,l_={$\SI{12}{\volt}$},invert] (6.4,0)--(b);
\fvMeshCW{1.55,1.05}{.40}{$I_1$}
\fvMeshCW{4.85,1.05}{.40}{$I_2$}
\end{circuitikz}}
\resposta{$I_1=\SI{5}{\ampere}$; $I_2=\SI{3}{\ampere}$;
$P_{fontes}=P_R=\SI{216}{\watt}$}


\end{questao}
\begin{questao}[2]{}

A fonte de corrente da malha 2 fixa $I_2=\SI{2}{\ampere}$. Determine $I_1$ e a
tensão $v_s$ da fonte de corrente, definida positiva do terminal superior para
o inferior.
\circuitoq{
\begin{circuitikz}[american,scale=.92,transform shape,line width=.8pt]
\draw (0,0) to[\fonteV,l^={$\SI{20}{\volt}$}] (0,3.2)
      to[R,l={$\SI{5}{\ohm}$}] (3.2,3.2) coordinate(t);
\draw (t) to[R,l_={$\SI{5}{\ohm}$}] (3.2,0) coordinate(b)--(0,0);
\draw (t) to[R,l={$\SI{10}{\ohm}$}] (6.4,3.2)
      to[isource,l={$\SI{2}{\ampere}$}] (6.4,0)--(b);
\node[right] at(6.75,1.6){$v_s$};
\fvMeshCW{1.55,1.05}{.40}{$I_1$}
\fvMeshCW{4.85,1.05}{.40}{$I_2$}
\end{circuitikz}}
\resposta{$I_1=\SI{3}{\ampere}$; $v_s=\SI{-15}{\volt}$}


\end{questao}

\blocoaprofundamento
\begin{questao}[3]{}

A fonte de corrente de $\SI{2}{\ampere}$ está entre as malhas 1 e 2. A fonte de
$\SI{8}{\volt}$ da malha 3 se opõe ao sentido horário indicado. Resolva usando
uma supermalha 1--2.
\circuitoq{
\begin{circuitikz}[american,scale=.82,transform shape,line width=.8pt]
\draw (0,0) to[\fonteV,l^={$\SI{24}{\volt}$}] (0,3.4)
      to[R,l={$\SI{4}{\ohm}$}] (2.8,3.4) coordinate(t1);
\draw (t1) to[isource,l={$\SI{2}{\ampere}$},invert] (2.8,0) coordinate(b1)--(0,0);
\draw (t1) to[R,l={$\SI{2}{\ohm}$}] (5.6,3.4) coordinate(t2);
\draw (t2) to[R,l_={$\SI{4}{\ohm}$}] (5.6,0) coordinate(b2)--(b1);
\draw (t2) to[R,l={$\SI{4}{\ohm}$}] (8.4,3.4)
      to[\fonteV,l_={$\SI{8}{\volt}$}] (8.4,0)--(b2);
\fvMeshCW{1.25,1.05}{.34}{$I_1$}
\fvMeshCW{4.15,1.05}{.34}{$I_2$}
\fvMeshCW{7.05,1.05}{.34}{$I_3$}
\draw[dashed,laranja,line width=.9pt,rounded corners=7pt] (-.45,-.45) rectangle (6.05,3.85);
\node[laranja,fill=white,inner sep=1pt] at(2.8,-.72){\footnotesize supermalha 1--2};
\end{circuitikz}}
\resposta{$I_1=\SI{1.5}{\ampere}$; $I_2=\SI{3.5}{\ampere}$;
$I_3=\SI{0.75}{\ampere}$}


\end{questao}
\begin{questao}[3]{}

O ramo comum é uma fonte ideal de $\SI{6}{\volt}$, positiva no topo. A fonte de
$\SI{18}{\volt}$ impulsiona $I_1$, enquanto a de $\SI{12}{\volt}$ se opõe a
$I_2$. Determine as correntes e a potência da fonte compartilhada.
\circuitoq{
\begin{circuitikz}[american,scale=.92,transform shape,line width=.8pt]
\draw (0,0) to[\fonteV,l^={$\SI{18}{\volt}$}] (0,3.2)
      to[R,l={$\SI{6}{\ohm}$}] (3.2,3.2) coordinate(t);
\draw (t) to[\fonteV,l_={$\SI{6}{\volt}$}] (3.2,0) coordinate(b)--(0,0);
\draw (t) to[R,l={$\SI{4}{\ohm}$}] (6.4,3.2)
      to[\fonteV,l_={$\SI{12}{\volt}$}] (6.4,0)--(b);
\fvMeshCW{1.55,1.05}{.40}{$I_1$}
\fvMeshCW{4.85,1.05}{.40}{$I_2$}
\end{circuitikz}}
\resposta{$I_1=\SI{2}{\ampere}$; $I_2=\SI{-1.5}{\ampere}$;
$P_{6V}=\SI{21}{\watt}$ absorvidos}


\end{questao}
\begin{questao}[3]{}

As fontes de corrente periféricas fixam $I_1=\SI{2}{\ampere}$ e
$I_3=\SI{1}{\ampere}$. Determine $I_2$ e identifique qualquer ramo
compartilhado com corrente nula.
\circuitoq{
\begin{circuitikz}[american,scale=.82,transform shape,line width=.8pt]
\draw (0,0) to[isource,l^={$\SI{2}{\ampere}$}] (0,3.4)
      --(2.8,3.4) coordinate(t1);
\draw (t1) to[R,l_={$\SI{2}{\ohm}$}] (2.8,0) coordinate(b1)--(0,0);
\draw (t1) to[R,l={$\SI{4}{\ohm}$}] (4.5,3.4)
      to[\fonteV,l={$\SI{11}{\volt}$}] (5.8,3.4) coordinate(t2);
\draw (t2) to[R,l_={$\SI{3}{\ohm}$}] (5.8,0) coordinate(b2)--(b1);
\draw (t2)--(8.6,3.4) to[isource,l={$\SI{1}{\ampere}$}] (8.6,0)--(b2);
\fvMeshCW{1.25,1.05}{.34}{$I_1$}
\fvMeshCW{4.30,1.05}{.34}{$I_2$}
\fvMeshCW{7.25,1.05}{.34}{$I_3$}
\end{circuitikz}}
\resposta{$I_2=\SI{2}{\ampere}$; o resistor de $\SI{2}{\ohm}$ conduz corrente nula}


\end{questao}
\begin{questao}[3]{}

Analise a ponte resistiva pelo método das malhas e determine a corrente no
resistor de ponte de $\SI{5}{\ohm}$.
\circuitoq{
\begin{circuitikz}[american,scale=.78,transform shape,line width=.8pt]
\coordinate (A) at (5,5); \coordinate (D) at (5,0); \coordinate (B) at (2.6,2.5); \coordinate (C) at (7.4,2.5);
\draw (A)--(0.5,5) to[\fonteV,l^={$\SI{18}{\volt}$}] (0.5,0)--(D);
\draw (A) to[R,l_={$\SI{2}{\ohm}$}] (B) to[R,l_={$\SI{4}{\ohm}$}] (D);
\draw (A) to[R,l={$\SI{3}{\ohm}$}] (C) to[R,l={$\SI{6}{\ohm}$}] (D);
\draw (B) to[R,l={$\SI{5}{\ohm}$}] (C);
\fvMeshCW{2.1,2.55}{.34}{$I_1$}
\fvMeshCW{5.0,3.55}{.32}{$I_2$}
\fvMeshCW{5.0,1.45}{.32}{$I_3$}
\end{circuitikz}}
\resposta{$I_1=\SI{5}{\ampere}$; $I_2=I_3=\SI{2}{\ampere}$;
$i_{ponte}=0$}


\end{questao}
\begin{questao}[3]{}

Use uma fonte de teste de $\SI{12}{\volt}$ e análise de malhas para determinar
a resistência equivalente vista pela fonte.
\circuitoq{
\begin{circuitikz}[american,scale=.92,transform shape,line width=.8pt]
\draw (0,0) to[\fonteV,l^={$\SI{12}{\volt}$}] (0,3.2)
      to[R,l={$\SI{3}{\ohm}$}] (3.2,3.2) coordinate(t);
\draw (t) to[R,l_={$\SI{6}{\ohm}$}] (3.2,0) coordinate(b)--(0,0);
\draw (t) to[R,l={$\SI{6}{\ohm}$}] (6.4,3.2)--(6.4,0)--(b);
\fvMeshCW{1.55,1.05}{.40}{$I_1$}
\fvMeshCW{4.85,1.05}{.40}{$I_2$}
\end{circuitikz}}
\resposta{$I_1=\SI{2}{\ampere}$; $I_2=\SI{1}{\ampere}$;
$R_{eq}=\SI{6}{\ohm}$}


\end{questao}

\blocodesafio
\begin{questao}[4]{}

A rede em ``roda'' possui quatro malhas triangulares em torno do nó $O$.
Monte a matriz e determine as quatro correntes.
\circuitoq{
\begin{circuitikz}[american,scale=.72,transform shape,line width=.8pt]
\coordinate (A) at (-4,4); \coordinate (B) at (4,4); \coordinate (C) at (4,-4); \coordinate (D) at (-4,-4); \coordinate (O) at (0,0);

\draw (A) to[\fonteV,l^={$\SI{24}{\volt}$}] (-1.8,4) to[R,l={$\SI{4}{\ohm}$}] (B);
\draw (B) to[\fonteV,l={$\SI{12}{\volt}$}] (4,1.8) to[R,l={$\SI{4}{\ohm}$}] (C);
\draw (C) to[\fonteV,l_={$\SI{8}{\volt}$}] (1.8,-4) to[R,l_={$\SI{4}{\ohm}$}] (D);
\draw (D) to[\fonteV,l_={$\SI{4}{\volt}$},invert] (-4,1.8) to[R,l_={$\SI{4}{\ohm}$}] (A);

\draw (A) to[R,l={$\SI{2}{\ohm}$}] (O);
\draw (B) to[R,l_={$\SI{2}{\ohm}$}] (O);
\draw (C) to[R,l={$\SI{2}{\ohm}$}] (O);
\draw (D) to[R,l_={$\SI{2}{\ohm}$}] (O);
\fvMeshCW{0,2.45}{.34}{$I_1$}
\fvMeshCW{2.45,0}{.34}{$I_2$}
\fvMeshCW{0,-2.45}{.34}{$I_3$}
\fvMeshCW{-2.45,0}{.34}{$I_4$}
\node[fill=white,inner sep=1pt] at(O){$O$};
\end{circuitikz}}
\resposta{$I_1=\SI{4}{\ampere}$; $I_2=\SI{3}{\ampere}$;
$I_3=\SI{2}{\ampere}$; $I_4=\SI{1}{\ampere}$}


\end{questao}
\begin{questao}[4]{}

Determine o resistor compartilhado $R$ para que $I_2=I_1/2$. Calcule também
as correntes resultantes.
\circuitoq{
\begin{circuitikz}[american,scale=.92,transform shape,line width=.8pt]
\draw (0,0) to[\fonteV,l^={$\SI{20}{\volt}$}] (0,3.2)
      to[R,l={$\SI{4}{\ohm}$}] (3.2,3.2) coordinate(t);
\draw (t) to[R,l_={$R$}] (3.2,0) coordinate(b)--(0,0);
\draw (t) to[R,l={$\SI{6}{\ohm}$}] (6.4,3.2)--(6.4,0)--(b);
\fvMeshCW{1.55,1.05}{.40}{$I_1$}
\fvMeshCW{4.85,1.05}{.40}{$I_2$}
\end{circuitikz}}
\resposta{$R=\SI{6}{\ohm}$; $I_1=\dfrac{20}{7}\,\si{\ampere}$;
$I_2=\dfrac{10}{7}\,\si{\ampere}$}


\end{questao}
\begin{questao}[4]{}

O resistor compartilhado é a carga $R_L$. Determine $R_L$ para que sua potência
seja máxima e calcule $P_{L,\max}$.
\circuitoq{
\begin{circuitikz}[american,scale=.92,transform shape,line width=.8pt]
\draw (0,0) to[\fonteV,l^={$\SI{24}{\volt}$}] (0,3.2)
      to[R,l={$\SI{4}{\ohm}$}] (3.2,3.2) coordinate(t);
\draw (t) to[R,l_={$R_L$}] (3.2,0) coordinate(b)--(0,0);
\draw (t) to[R,l={$\SI{8}{\ohm}$}] (6.4,3.2)
      to[\fonteV,l_={$\SI{12}{\volt}$},invert] (6.4,0)--(b);
\fvMeshCW{1.55,1.05}{.40}{$I_1$}
\fvMeshCW{4.85,1.05}{.40}{$I_2$}
\end{circuitikz}}
\resposta{$R_L=\dfrac{8}{3}\,\ohm\approx\SI{2.67}{\ohm}$;
$P_{L,\max}=\SI{13.5}{\watt}$}


\end{questao}
\begin{questao}[4]{}

A fonte $E$ deve ser escolhida para impor $I_2=0$. Determine $E$, sua polaridade
relativa ao sentido horário e as correntes $I_1$ e $I_3$.
\circuitoq{
\begin{circuitikz}[american,scale=.82,transform shape,line width=.8pt]
\draw (0,0) to[\fonteV,l^={$\SI{24}{\volt}$}] (0,3.4)
      to[R,l={$\SI{6}{\ohm}$}] (2.8,3.4) coordinate(t1);
\draw (t1) to[R,l_={$\SI{2}{\ohm}$}] (2.8,0) coordinate(b1)--(0,0);
\draw (t1) to[R,l={$\SI{4}{\ohm}$}] (5.8,3.4) coordinate(t2);
\draw (t2) to[R,l_={$\SI{1}{\ohm}$}] (5.8,0) coordinate(b2)--(b1);
\draw (t2) to[R,l={$\SI{4}{\ohm}$}] (8.8,3.4)
      to[\fonteV,l_={$E$}] (8.8,0)--(b2);
\fvMeshCW{1.25,1.05}{.34}{$I_1$}
\fvMeshCW{4.30,1.05}{.34}{$I_2$}
\fvMeshCW{7.35,1.05}{.34}{$I_3$}
\end{circuitikz}}
\resposta{$I_1=\SI{3}{\ampere}$; $I_2=0$; $I_3=\SI{-6}{\ampere}$;
$E=\SI{30}{\volt}$ orientada contra $I_3$ horário}


\end{questao}
\begin{questao}[4]{}

Em ensaio de uma rede de duas malhas, mediram-se $I_1=\SI{3}{\ampere}$ e
$I_2=\SI{2}{\ampere}$. Determine a resistência desconhecida $R$ do ramo comum
e a potência dissipada nela.
\circuitoq{
\begin{circuitikz}[american,scale=.92,transform shape,line width=.8pt]
\draw (0,0) to[\fonteV,l^={$\SI{20}{\volt}$}] (0,3.2)
      to[R,l={$\SI{4}{\ohm}$}] (3.2,3.2) coordinate(t);
\draw (t) to[R,l_={$R$}] (3.2,0) coordinate(b)--(0,0);
\draw (t) to[R,l={$\SI{6}{\ohm}$}] (6.4,3.2)
      to[\fonteV,l_={$\SI{4}{\volt}$},invert] (6.4,0)--(b);
\fvMeshCW{1.55,1.05}{.40}{$I_1$}
\fvMeshCW{4.85,1.05}{.40}{$I_2$}
\end{circuitikz}}
\resposta{$R=\SI{8}{\ohm}$; $P_R=\SI{8}{\watt}$}


\end{questao}

\gabarito[Capítulo 9 --- Método de Maxwell]

\section{Capacitores e Indutores}

\subsection{Capacitores}

\blocofixacao
\begin{questao}[1]{}

A capacitância de um capacitor de placas paralelas:
\begin{alternativas}
\item aumenta com a distância entre as placas.
\item aumenta com a área das placas.
\item independe do meio entre as placas.
\item diminui ao inserir um dielétrico.
\item depende da tensão aplicada.
\end{alternativas}
\resposta{(b)}


\end{questao}
\begin{questao}[1]{}

Calcule a capacitância de um capacitor plano de área $A = \SI{100}{\centi\meter\squared}$
e distância entre placas $d = \SI{1}{\milli\meter}$, no vácuo.
\dados{$\varepsilon_0 = \SI{8.85e-12}{\farad\per\meter}$.}
\resposta{$C \approx \SI{88.5}{\pico\farad}$}


\end{questao}
\begin{questao}[1]{}

Um capacitor de $\SI{10}{\micro\farad}$ é submetido a uma tensão de
$\SI{12}{\volt}$. Determine a carga armazenada e a energia acumulada.
\resposta{$Q = \SI{120}{\micro\coulomb}$; $E = \SI{720}{\micro\joule}$}


\end{questao}
\begin{questao}[1]{}

Dois capacitores de $\SI{6}{\micro\farad}$ e $\SI{3}{\micro\farad}$ são ligados
em \textbf{série}. A capacitância equivalente é:
\begin{alternativas}[3]
\item \SI{2}{\micro\farad}
\item \SI{4.5}{\micro\farad}
\item \SI{9}{\micro\farad}
\item \SI{18}{\micro\farad}
\item \SI{0.5}{\micro\farad}
\end{alternativas}
\resposta{(a)}


\end{questao}
\begin{questao}[1]{}

Dois capacitores de $\SI{4}{\micro\farad}$ e $\SI{12}{\micro\farad}$ são ligados
em \textbf{paralelo}. A capacitância equivalente é:
\begin{alternativas}[3]
\item \SI{3}{\micro\farad}
\item \SI{8}{\micro\farad}
\item \SI{16}{\micro\farad}
\item \SI{48}{\micro\farad}
\item \SI{6}{\micro\farad}
\end{alternativas}
\resposta{(c)}


\end{questao}
\begin{questao}[1]{}

Um capacitor de $\SI{100}{\micro\farad}$ é carregado a $\SI{50}{\volt}$. A energia
armazenada é:
\begin{alternativas}[3]
\item \SI{2.5}{\milli\joule}
\item \SI{5}{\milli\joule}
\item \SI{125}{\milli\joule}
\item \SI{250}{\milli\joule}
\item \SI{5}{\joule}
\end{alternativas}
\resposta{(c)}


\end{questao}

\blocoaplicacao
\begin{questao}[2]{}

Uma associação mista tem $C_1 = \SI{6}{\micro\farad}$ e
$C_2 = \SI{3}{\micro\farad}$ em \emph{paralelo}, e esse conjunto em \emph{série}
com $C_3 = \SI{2}{\micro\farad}$. Calcule a capacitância total.

\circuitoq{
\begin{circuitikz}[scale=1,transform shape,line width=0.8pt]
 \draw (0,0.75) to[short,o-] (1,0.75) coordinate(a);
 \draw (a) -- (1,1.5) to[C,l={$C_1=\SI{6}{\micro\farad}$}] (3,1.5) -- (3,0.75) coordinate(b);
 \draw (a) -- (1,0) to[C,l_={$C_2=\SI{3}{\micro\farad}$}] (3,0) -- (3,0.75);
 \draw (b) to[C,l={$C_3=\SI{2}{\micro\farad}$}] (5,0.75) to[short,-o] (6,0.75);
\end{circuitikz}}

\resposta{$\Ceq \approx \SI{1.64}{\micro\farad}$}


\end{questao}
\begin{questao}[2]{}

Uma associação é feita com dois capacitores de $\SI{8}{\micro\farad}$ em
\emph{série}, e esse conjunto em \emph{paralelo} com um capacitor de
$\SI{4}{\micro\farad}$. Calcule a capacitância equivalente.
\resposta{$\Ceq = \SI{8}{\micro\farad}$}


\end{questao}
\begin{questao}[2]{}

Um circuito RC série tem $R = \SI{1}{\kilo\ohm}$ e $C = \SI{1}{\micro\farad}$,
ligado a uma fonte de $\SI{10}{\volt}$ em $t = 0$.
\begin{itensq}
\item Qual é a constante de tempo?
\item Qual é a tensão no capacitor após um intervalo igual a $\tau$?
\item Qual é a tensão em regime permanente?
\end{itensq}
\resposta{(a) $\tau = \SI{1}{\milli\second}$; (b) $\approx\SI{6.32}{\volt}$;
(c) \SI{10}{\volt}}


\end{questao}
\begin{questao}[2]{}

No circuito da figura, determine a tensão no capacitor em regime permanente.

\circuitoq{
\begin{circuitikz}[scale=1,transform shape,line width=0.8pt]
 \draw (0,0) to[battery1,l_={$\SI{12}{\volt}$}] (0,3)
       to[R,l={$R_1=\SI{4}{\ohm}$}] (3,3) coordinate(a);
 \draw (a) to[R,l={$R_2=\SI{2}{\ohm}$}] (3,0) coordinate(b) -- (0,0);
 \draw (a) -- (5,3) to[C,l={$C$}] (5,0) -- (b);
 \node[right,Vcol,font=\footnotesize] at (5.1,1.5) {$v_C$};
\end{circuitikz}}

\resposta{$v_C = \SI{4}{\volt}$}


\end{questao}
\begin{questao}[2]{}

Três capacitores de $\SI{6}{\micro\farad}$ cada são associados. Determine a
capacitância equivalente se estiverem ligados (a) todos em série; (b) todos em
paralelo; (c) dois em paralelo, e o conjunto em série com o terceiro.
\resposta{(a) \SI{2}{\micro\farad}; (b) \SI{18}{\micro\farad}; (c) \SI{4}{\micro\farad}}


\end{questao}

\blocoaprofundamento
\begin{questao}[3]{}

Um capacitor de $\SI{2}{\micro\farad}$ está carregado a $\SI{6}{\volt}$ e é
descarregado através de um resistor de $\SI{500}{\ohm}$.
\begin{itensq}
\item Qual é a constante de tempo da descarga?
\item Qual é a corrente inicial de descarga?
\item Qual é a carga restante no capacitor após um intervalo igual a $\tau$?
\end{itensq}
\resposta{(a) $\tau = \SI{1}{\milli\second}$; (b) $I_0 = \SI{12}{\milli\ampere}$;
(c) $\approx\SI{4.4}{\micro\coulomb}$}


\end{questao}
\begin{questao}[3]{}

O gráfico mostra a tensão $v_C(t)$ sobre um capacitor durante a carga em um
circuito RC alimentado por $\SI{10}{\volt}$.

\circuitoqq{
\begin{tikzpicture}
 \begin{axis}[width=9.5cm,height=5cm,xlabel={$t$ (ms)},ylabel={$v_C$ (V)},
   xmin=0,xmax=5,ymin=0,ymax=11,xtick={0,1,2,3,4,5},ytick={0,6.32,10},
   yticklabels={0,6.32,10},grid=both,axis lines=left,domain=0:5,samples=100]
   \addplot[Ccol,very thick] {10*(1-exp(-x))};
   \addplot[dashed,cinza] coordinates {(1,0)(1,6.32)};
   \addplot[dashed,cinza] coordinates {(0,6.32)(1,6.32)};
   \node[Ccol,anchor=west,font=\footnotesize] at (axis cs:1.1,6.0) {$(\tau,\,0{,}63\,U_0)$};
 \end{axis}
\end{tikzpicture}}{Carga de um capacitor em circuito RC.}

Sabendo que a tensão atinge $\SI{6.32}{\volt}$ em $t = \SI{1}{\milli\second}$,
determine a constante de tempo e a tensão final.
\resposta{$\tau = \SI{1}{\milli\second}$; $U_0 = \SI{10}{\volt}$}


\end{questao}
\begin{questao}[3]{}

Um capacitor $C_1 = \SI{10}{\micro\farad}$ é carregado a $\SI{100}{\volt}$ e
depois conectado, por meio de uma chave, a um capacitor descarregado
$C_2 = \SI{5}{\micro\farad}$.

\circuitoq{
\begin{circuitikz}[scale=1,transform shape,line width=0.8pt]
 \draw (0,0) to[C,l_={$C_1$}] (0,2.6) -- (1.6,2.6) coordinate(a);
 \draw (a) to[cspst,l={S}] (3.4,2.6) -- (4.6,2.6);
 \draw (4.6,2.6) to[C,l={$C_2$}] (4.6,0) -- (0,0);
 \node[font=\footnotesize] at (0,-0.45) {carregado a \SI{100}{\volt}};
 \node[font=\footnotesize] at (4.6,-0.45) {descarregado};
\end{circuitikz}}

Após fechar a chave:
\begin{itensq}
\item Determine a tensão final comum aos dois capacitores.
\item Calcule a energia armazenada antes e depois.
\item Explique o "paradoxo": para onde foi a energia que desapareceu?
\end{itensq}
\resposta{(a) $V_f \approx \SI{66.7}{\volt}$; (b) \SI{50}{\milli\joule} $\to$
\SI{33.3}{\milli\joule}; (c) ver comentário}


\end{questao}
\begin{questao}[3]{}

Um capacitor de placas paralelas com $C_0 = \SI{10}{\micro\farad}$ é carregado a
$U_0 = \SI{100}{\volt}$ e então \textbf{desconectado} da fonte. Em seguida,
introduz-se entre as placas um dielétrico de constante $\kappa = 4$.
\begin{itensq}
\item O que acontece com a carga, a capacitância, a tensão e a energia?
\item Calcule os valores finais de cada grandeza.
\item Para onde foi a energia "perdida"? Ela violou a conservação de energia?
\item E se o dielétrico fosse inserido com o capacitor \emph{ainda ligado} à
      fonte? Compare os dois casos.
\end{itensq}
\resposta{(b) $Q=\SI{1}{\milli\coulomb}$ (constante), $C=\SI{40}{\micro\farad}$,
$U=\SI{25}{\volt}$, $E$ cai de \SI{50}{\milli\joule} para \SI{12.5}{\milli\joule};
(c) e (d) ver comentário}


\end{questao}

\blocofixacao
\begin{questao}[1]{}

Um capacitor armazena $\SI{300}{\micro\coulomb}$ sob $\SI{20}{\volt}$.
Determine sua capacitância.
\resposta{$C=\SI{15}{\micro\farad}$}


\end{questao}
\begin{questao}[1]{}

Um capacitor de $\SI{100}{\micro\farad}$ armazena $\SI{8}{\milli\joule}$.
Determine a tensão em seus terminais.
\resposta{$U\approx\SI{12.6}{\volt}$}


\end{questao}
\begin{questao}[1]{}

Sobre a tensão em um capacitor ideal, é correto afirmar que ela:
\begin{alternativas}[2]
\item pode variar instantaneamente
\item não pode variar instantaneamente
\item é sempre igual à da fonte
\item é sempre nula em regime permanente CC
\end{alternativas}
\resposta{(b)}


\end{questao}
\begin{questao}[1]{}

Determine a constante de tempo de um circuito RC com
$R=\SI{4.7}{\kilo\ohm}$ e $C=\SI{10}{\micro\farad}$.
\resposta{$\tau=\SI{47}{\milli\second}$}


\end{questao}

\blocoaplicacao
\begin{questao}[2]{}

Dois capacitores de $\SI{10}{\micro\farad}$ em série são associados em paralelo
com um de $\SI{15}{\micro\farad}$. Determine $C_{eq}$.
\resposta{$C_{eq}=\SI{20}{\micro\farad}$}


\end{questao}
\begin{questao}[2]{}

Capacitores de $\SI{1}{\micro\farad}$ e $\SI{2}{\micro\farad}$ estão em série
sob $\SI{30}{\volt}$. Determine $C_{eq}$, a carga e a tensão em cada um.
\resposta{$C_{eq}=\SI{0.667}{\micro\farad}$; $Q=\SI{20}{\micro\coulomb}$;
$U_1=\SI{20}{\volt}$, $U_2=\SI{10}{\volt}$}


\end{questao}
\begin{questao}[2]{}

A tensão sobre um capacitor de $\SI{10}{\micro\farad}$ cresce linearmente à
taxa de $\SI{500}{\volt\per\second}$. Determine a corrente.
\resposta{$i=\SI{5}{\milli\ampere}$}


\end{questao}
\begin{questao}[2]{}

Aplica-se a um capacitor de $\SI{2.2}{\micro\farad}$ uma tensão
\textbf{triangular} que sobe $\SI{10}{\volt}$ em $\SI{1}{\milli\second}$ e
desce igualmente no milissegundo seguinte. Descreva e quantifique a corrente.
\resposta{Onda \textbf{quadrada} de $\pm\SI{22}{\milli\ampere}$}


\end{questao}
\begin{questao}[2]{}

Um circuito RC com $\tau=\SI{2}{\milli\second}$ é ligado a
$\SI{20}{\volt}$ com o capacitor inicialmente descarregado. Determine
$u_C$ em $t=\tau$, $2\tau$ e $3\tau$.
\resposta{$\SI{12.64}{\volt}$, $\SI{17.29}{\volt}$ e $\SI{19.00}{\volt}$}


\end{questao}
\begin{questao}[2]{}

No mesmo circuito, determine o instante em que a tensão atinge
$\SI{15}{\volt}$.
\resposta{$t\approx\SI{2.77}{\milli\second}$}


\end{questao}
\begin{questao}[2]{}

Uma fonte de $\SI{24}{\volt}$ alimenta $R_1=\SI{8}{\kilo\ohm}$ em série com
$R_2=\SI{4}{\kilo\ohm}$. Um capacitor está ligado em paralelo com $R_2$.
Determine a tensão no capacitor em regime permanente e a corrente que o
atravessa.
\resposta{$U_C=\SI{8}{\volt}$; $i_C=0$}


\end{questao}
\begin{questao}[2]{}

Dois capacitores armazenam a mesma energia. O primeiro tem
$C_1=\SI{100}{\micro\farad}$ sob $\SI{10}{\volt}$. Se
$C_2=\SI{25}{\micro\farad}$, determine sua tensão.
\resposta{$U_2=\SI{20}{\volt}$}


\end{questao}
\begin{questao}[2]{}

Um capacitor plano tem capacitância $C$. Determine a nova capacitância se:
\begin{itensq}
\item a área das placas for dobrada;
\item a distância entre placas for reduzida à metade;
\item ambas as alterações forem feitas.
\end{itensq}
\resposta{(a) $2C$; (b) $2C$; (c) $4C$}


\end{questao}
\begin{questao}[2]{}

Um capacitor carregado a $\SI{50}{\volt}$ é descarregado através de
$R=\SI{2.2}{\kilo\ohm}$. Determine a corrente inicial e explique sua evolução.
\resposta{$i_0=\SI{22.7}{\milli\ampere}$, decaindo exponencialmente}


\end{questao}

\blocoaprofundamento
\begin{questao}[3]{}

$C_1=\SI{4}{\micro\farad}$ e $C_2=\SI{12}{\micro\farad}$, ambos com tensão
máxima de trabalho de $\SI{50}{\volt}$, são ligados em série.
\begin{itensq}
\item Determine $C_{eq}$ e a distribuição de tensão sob $\SI{100}{\volt}$.
\item Determine a máxima tensão que a associação suporta.
\item Explique por que o resultado é contraintuitivo.
\end{itensq}
\resposta{(a) $\SI{3}{\micro\farad}$; $\SI{75}{\volt}$ e $\SI{25}{\volt}$;
(b) $\SI{66.7}{\volt}$}


\end{questao}
\begin{questao}[3]{}

Capacitores de $\SI{2}{\micro\farad}$, $\SI{3}{\micro\farad}$ e
$\SI{5}{\micro\farad}$ estão em paralelo sob $\SI{20}{\volt}$. Determine a
carga de cada um, a carga total e $C_{eq}$, e compare com o caso série.
\resposta{$40$, $60$ e $\SI{100}{\micro\coulomb}$; total
$\SI{200}{\micro\coulomb}$; $C_{eq}=\SI{10}{\micro\farad}$}


\end{questao}
\begin{questao}[3]{}

Uma fonte de $\SI{12}{\volt}$ alimenta $R_1=\SI{6}{\kilo\ohm}$ em série com
$R_2=\SI{3}{\kilo\ohm}$, e um capacitor de $\SI{100}{\nano\farad}$ está em
paralelo com $R_2$.
\begin{itensq}
\item Determine a tensão final no capacitor.
\item Determine $\tau$ pelo equivalente de Thévenin.
\item Escreva $u_C(t)$.
\end{itensq}
\resposta{(a) $\SI{4}{\volt}$; (b) $\tau=\SI{0.2}{\milli\second}$;
(c) $u_C=4\left(1-e^{-t/0{,}2\,\mathrm{ms}}\right)$}


\end{questao}
\begin{questao}[3]{}

Um capacitor já está a $\SI{2}{\volt}$ quando é ligado a um circuito que o
levaria a $\SI{10}{\volt}$, com $\tau=\SI{5}{\milli\second}$.
\begin{itensq}
\item Escreva $u_C(t)$.
\item Determine a tensão em $t=\SI{5}{\milli\second}$.
\item Determine a corrente inicial, sabendo que $R=\SI{5}{\kilo\ohm}$.
\end{itensq}
\resposta{(a) $u_C=10-8e^{-t/\tau}$; (b) $\SI{7.06}{\volt}$;
(c) $\SI{1.6}{\milli\ampere}$}


\end{questao}
\begin{questao}[3]{}

Um capacitor de $\SI{470}{\micro\farad}$ tem resistência de isolação de
$\SI{2}{\mega\ohm}$. Carregado a $\SI{100}{\volt}$ e deixado em circuito
aberto:
\begin{itensq}
\item Determine a constante de tempo de autodescarga.
\item Determine a tensão após $\SI{15}{\minute}$.
\item Comente a implicação de segurança.
\end{itensq}
\resposta{(a) $\tau=\SI{940}{\second}\approx\SI{15.7}{\minute}$;
(b) $\approx\SI{38}{\volt}$}


\end{questao}
\begin{questao}[3]{}

Um capacitor eletrolítico tem ESR (resistência equivalente em série) de
$\SI{0.1}{\ohm}$ e conduz corrente de ondulação de
$\SI{1.5}{\ampere}$ eficazes.
\begin{itensq}
\item Determine a potência dissipada internamente.
\item Explique por que essa dissipação limita a vida do componente.
\item Indique duas formas de reduzi-la.
\end{itensq}
\resposta{(a) $P=\SI{0.225}{\watt}$}


\end{questao}
\begin{questao}[3]{}

Um capacitor de $\SI{100}{\micro\farad}$, inicialmente descarregado, é
carregado a $\SI{20}{\volt}$ através de um resistor.
\begin{itensq}
\item Determine a energia fornecida pela fonte.
\item Determine a energia armazenada no capacitor.
\item Determine a energia dissipada no resistor e comente.
\end{itensq}
\resposta{(a) $\SI{40}{\milli\joule}$; (b) $\SI{20}{\milli\joule}$;
(c) $\SI{20}{\milli\joule}$ --- exatamente metade}


\end{questao}
\begin{questao}[3]{}

Um capacitor plano usa dielétrico com rigidez de
$\SI{20}{\kilo\volt\per\milli\meter}$ e espessura de
$\SI{0.1}{\milli\meter}$.
\begin{itensq}
\item Determine a tensão de ruptura.
\item Especifique a tensão nominal com margem de $50\%$.
\item Explique por que a rigidez limita o produto capacitância-tensão.
\end{itensq}
\resposta{(a) $\SI{2}{\kilo\volt}$; (b) $\SI{1}{\kilo\volt}$}


\end{questao}
\begin{questao}[3]{}

Em um circuito RC de carga, mediu-se tensão final de $\SI{10}{\volt}$,
corrente inicial de $\SI{2}{\milli\ampere}$ e tensão de $\SI{6.32}{\volt}$ em
$\SI{3}{\milli\second}$. Determine $R$ e $C$.
\resposta{$R=\SI{5}{\kilo\ohm}$; $C=\SI{0.6}{\micro\farad}$}


\end{questao}
\begin{questao}[3]{}

Compare o comportamento de capacitor e indutor preenchendo o raciocínio para os
instantes $t=0^{+}$ e $t\to\infty$, e justifique cada equivalência.
\resposta{Capacitor: curto em $t=0^{+}$ (se descarregado), aberto em regime;
indutor: aberto em $t=0^{+}$ (se sem corrente), curto em regime}


\end{questao}
\begin{questao}[3]{}

Um filtro RC passa-baixas tem $R=\SI{10}{\kilo\ohm}$ e
$C=\SI{47}{\nano\farad}$.
\begin{itensq}
\item Determine a frequência de corte.
\item Relacione-a com a constante de tempo.
\item Descreva o comportamento muito abaixo e muito acima de $f_c$.
\end{itensq}
\resposta{(a) $f_c\approx\SI{339}{\hertz}$; (b) $f_c=1/(2\pi\tau)$}


\end{questao}

\blocodesafio
\begin{questao}[4]{}

\textbf{(Dedução --- resposta do circuito RC.)} Um RC série com fonte $U$ tem a
chave fechada em $t=0$ e $u_C(0)=0$.
\begin{itensq}
\item Escreva a equação diferencial e resolva por separação de variáveis.
\item Obtenha $i(t)$ e verifique a LKT nos dois extremos.
\item Compare a estrutura da solução com a do circuito RL.
\end{itensq}
\resposta{$u_C=U\left(1-e^{-t/RC}\right)$;
$i=\dfrac{U}{R}e^{-t/RC}$}


\end{questao}
\begin{questao}[4]{}

\textbf{(Demonstração --- rendimento do carregamento.)} Um capacitor
descarregado é carregado por uma fonte $U$ através de $R$.
\begin{itensq}
\item Integre $p_R=Ri^{2}$ para obter a energia dissipada.
\item Mostre que ela vale $\frac12CU^{2}$, independentemente de $R$.
\item Explique como conversores chaveados contornam esse limite.
\end{itensq}
\resposta{$W_R=\frac12CU^{2}$; rendimento sempre $50\%$}


\end{questao}
\begin{questao}[4]{}

\textbf{(Projeto de banco.)} Deseja-se um banco de $\SI{100}{\micro\farad}$
capaz de operar em $\SI{100}{\volt}$, dispondo apenas de capacitores de
$\SI{100}{\micro\farad}$ e $\SI{50}{\volt}$.
\begin{itensq}
\item Determine a associação e o número de unidades.
\item Explique por que são necessários resistores de equalização.
\item Dimensione-os para uma corrente de fuga de até
      $\SI{10}{\micro\ampere}$.
\end{itensq}
\resposta{(a) duas séries de dois, em paralelo --- 4 unidades;
(c) $R\le\SI{50}{\kilo\ohm}$, dissipando $\SI{50}{\milli\watt}$ cada}


\end{questao}
\begin{questao}[4]{}

\textbf{(Projeto de filtro.)} Projete um passa-baixas RC de primeira ordem com
$f_c=\SI{100}{\hertz}$, usando $C=\SI{100}{\nano\farad}$.
\begin{itensq}
\item Determine $R$ e escolha o valor comercial E24.
\item Recalcule $f_c$ e avalie o erro.
\item Determine a atenuação em $\SI{1}{\kilo\hertz}$ e discuta a limitação da
      primeira ordem.
\end{itensq}
\resposta{(a) $R=\SI{15.9}{\kilo\ohm}\to\SI{16}{\kilo\ohm}$;
(b) $f_c=\SI{99.5}{\hertz}$, erro $0{,}5\%$;
(c) $\approx-\SI{20}{\decibel}$}


\end{questao}
\begin{questao}[4]{}

\textbf{(Propagação de incerteza.)} Analise o efeito das tolerâncias de $R$ e
$C$ sobre a constante de tempo.
\begin{itensq}
\item Obtenha a expressão da incerteza relativa de $\tau$.
\item Avalie para $R$ com $\pm1\%$ e $C$ eletrolítico com $\pm20\%$, nos
      critérios de pior caso e estatístico.
\item Conclua sobre onde investir em precisão.
\end{itensq}
\resposta{(a) $\dfrac{\Delta\tau}{\tau}=\dfrac{\Delta R}{R}+\dfrac{\Delta C}{C}$;
(b) $21\%$ (pior caso) e $20{,}0\%$ (estatístico)}


\end{questao}
\begin{questao}[4]{}

\textbf{(Dedução --- força entre as placas.)} Considere um capacitor plano de
área $A$ e separação $d$.
\begin{itensq}
\item Obtenha a força entre as placas pelo método da energia, com carga
      constante.
\item Expresse o resultado em função de $U$ e discuta o caso a tensão
      constante.
\item Avalie para $A=\SI{100}{\centi\meter\squared}$,
      $d=\SI{1}{\milli\meter}$ e $U=\SI{1}{\kilo\volt}$.
\end{itensq}
\dados{$\varepsilon_0=\pot{8.85}{-12}\,\si{\farad\per\meter}$}
\resposta{$F=\dfrac{\varepsilon_0AU^{2}}{2d^{2}}\approx\SI{44}{\milli\newton}$,
atrativa}


\end{questao}
\begin{questao}[4]{}

\textbf{(Projeto de temporizador.)} Um circuito deve atingir $70\%$ da tensão
final em exatamente $\SI{1}{\second}$.
\begin{itensq}
\item Determine $\tau$.
\item Escolha $C$ e determine $R$, com valores comerciais.
\item Avalie a repetibilidade real do temporizador.
\end{itensq}
\resposta{(a) $\tau=\SI{0.83}{\second}$; (b) $C=\SI{100}{\micro\farad}$ e
$R=\SI{8.2}{\kilo\ohm}$; (c) erro dominado pelo capacitor}


\end{questao}
\begin{questao}[4]{}

\textbf{(Modelagem --- autodescarga.)} Um capacitor real é modelado por $C$
ideal em paralelo com uma resistência de fuga $R_f$.
\begin{itensq}
\item Deduza a lei de decaimento da tensão em circuito aberto.
\item Defina a constante de tempo de autodescarga e relacione-a com a
      especificação de "corrente de fuga" de catálogo.
\item Avalie para $C=\SI{1000}{\micro\farad}$ e
      $R_f=\SI{100}{\kilo\ohm}$, com tensão inicial de $\SI{300}{\volt}$.
\end{itensq}
\resposta{$u=U_0e^{-t/R_fC}$; $\tau=\SI{100}{\second}$; após
$\SI{5}{\minute}$ restam $\approx\SI{15}{\volt}$}


\end{questao}
\begin{questao}[4]{}

\textbf{(ESR e ondulação.)} Um capacitor de filtro de fonte chaveada tem
$\mathrm{ESR}=\SI{50}{\milli\ohm}$ e deve conduzir ondulação de
$\SI{2}{\ampere}$ eficazes.
\begin{itensq}
\item Determine a dissipação e a ondulação de tensão devida \emph{apenas} à
      ESR.
\item Compare com a ondulação devida à capacitância, para
      $C=\SI{470}{\micro\farad}$ a $\SI{100}{\kilo\hertz}$.
\item Conclua sobre o critério de seleção.
\end{itensq}
\resposta{(a) $\SI{0.2}{\watt}$ e $\approx\SI{100}{\milli\volt}$ de pico;
(b) $\approx\SI{5}{\milli\volt}$ --- a ESR domina}


\end{questao}
\begin{questao}[4]{}

\textbf{(Comparação tecnológica.)} Um supercapacitor de $\SI{3000}{\farad}$ e
$\SI{2.7}{\volt}$ tem massa de $\SI{0.50}{\kilogram}$.
\begin{itensq}
\item Determine a energia armazenada e a densidade de energia.
\item Compare com uma bateria de íon-lítio ($\approx\SI{250}{\watt\hour\per\kilogram}$).
\item Aponte em que aspecto o supercapacitor é superior e por quê.
\end{itensq}
\resposta{(a) $\SI{10.9}{\kilo\joule}=\SI{3.04}{\watt\hour}$, ou
$\SI{6.1}{\watt\hour\per\kilogram}$; (b) cerca de $40$ vezes menos}


\end{questao}

\gabarito[10.1 Capacitores]

\subsection{Indutores}

\blocofixacao
\begin{questao}[1]{}

Um indutor ideal, em regime permanente CC (corrente constante), comporta-se como:
\begin{alternativas}
\item um circuito aberto.
\item um curto-circuito.
\item um resistor de valor $L$.
\item uma fonte de tensão.
\item um capacitor.
\end{alternativas}
\resposta{(b)}


\end{questao}
\begin{questao}[1]{}

Calcule a indutância de um solenoide com $N = 500$ espiras, comprimento
$\ell = \SI{0.3}{\meter}$ e área $A = \SI{5}{\centi\meter\squared}$.
\dados{$\mu_0 = 4\pi\times10^{-7}\,\si{\henry\per\meter}$.}
\resposta{$L \approx \SI{0.52}{\milli\henry}$}


\end{questao}
\begin{questao}[1]{}

Um indutor de $\SI{0.1}{\henry}$ tem a corrente variando de $0$ a
$\SI{2}{\ampere}$ em $\SI{0.2}{\second}$. Calcule a tensão média induzida.
\resposta{$v_L = \SI{1}{\volt}$}


\end{questao}
\begin{questao}[1]{}

Dois indutores de $\SI{3}{\henry}$ e $\SI{9}{\henry}$ são ligados em paralelo. A
indutância equivalente é:
\begin{alternativas}[3]
\item \SI{12}{\henry}
\item \SI{6}{\henry}
\item \SI{2.25}{\henry}
\item \SI{3}{\henry}
\item \SI{27}{\henry}
\end{alternativas}
\resposta{(c)}


\end{questao}

\blocoaplicacao
\begin{questao}[2]{}

Um indutor de $L = \SI{1}{\henry}$ é ligado em série com $R = \SI{100}{\ohm}$ a
uma fonte de $\SI{12}{\volt}$.
\begin{itensq}
\item Qual é a constante de tempo?
\item Qual é a corrente final (regime permanente)?
\item Qual é a corrente após um intervalo igual a $\tau$?
\end{itensq}
\resposta{(a) $\tau = \SI{10}{\milli\second}$; (b) \SI{0.12}{\ampere};
(c) $\approx\SI{75.8}{\milli\ampere}$}


\end{questao}
\begin{questao}[2]{}

Calcule a indutância equivalente para:
\begin{itensq}
\item $L_1 = \SI{2}{\milli\henry}$ e $L_2 = \SI{3}{\milli\henry}$ em série;
\item os mesmos dois indutores em paralelo.
\end{itensq}
\resposta{(a) \SI{5}{\milli\henry}; (b) \SI{1.2}{\milli\henry}}


\end{questao}
\begin{questao}[2]{Apostila original revisada}

Dois indutores ideais de $\SI{5}{\henry}$ são ligados em paralelo. Essa associação
é conectada em série com um indutor de $\SI{2}{\henry}$. Determine a indutância
equivalente.
\resposta{$L_{\mathrm{eq}}=\SI{4.5}{\henry}$}


\end{questao}

\blocoaprofundamento
\begin{questao}[3]{}

Compare, lado a lado, o comportamento de um capacitor e de um indutor em um
circuito de corrente contínua, preenchendo mentalmente a tabela abaixo e
justificando cada linha.

\begin{table}[H]
\centering
\small\sffamily
\setlength{\arraystretch}{1.3}
\begin{tabular}{@{}l c c@{}}
\toprule
 & \textbf{Capacitor} & \textbf{Indutor}\\
\midrule
Grandeza que não muda bruscamente & tensão $v_C$ & corrente $i_L$\\
Em regime permanente CC & circuito aberto & curto-circuito\\
Energia armazenada & $\tfrac12 C v^2$ (campo elétrico) & $\tfrac12 L i^2$ (campo magnético)\\
Constante de tempo & $\tau = RC$ & $\tau = L/R$\\
\bottomrule
\end{tabular}
\caption{Dualidade entre capacitor e indutor.}
\end{table}

Explique por que a tensão no capacitor e a corrente no indutor não podem variar
instantaneamente.
\resposta{Ver comentário}


\end{questao}
\begin{questao}[3]{}

Um indutor de $L = \SI{0.5}{\henry}$ conduz corrente estabilizada de
$\SI{4}{\ampere}$, alimentado por uma fonte através de um resistor. Em
$t=0$, a chave que liga a fonte é \textbf{aberta bruscamente}.
\begin{itensq}
\item Qual é a energia armazenada no indutor imediatamente antes da abertura?
\item Por que surge um arco elétrico (faísca) nos contatos da chave?
\item Estime a tensão induzida se a corrente for interrompida em
      $\SI{1}{\milli\second}$.
\item Qual solução prática se usa para proteger o circuito? Explique seu
      funcionamento.
\end{itensq}
\resposta{(a) $E=\SI{4}{\joule}$; (c) $\approx\SI{2000}{\volt}$; (b) e (d) ver comentário}


\end{questao}
\begin{questao}[3]{Apostila original revisada}

Um indutor de $L=\SI{2}{\milli\henry}$ permanece por muito tempo em série com um
resistor de $\SI{6}{\kilo\ohm}$ e uma fonte de $\SI{12}{\volt}$. Em $t=0$, a fonte
e o resistor de $\SI{6}{\kilo\ohm}$ são desconectados, e o indutor passa a descarregar
por um resistor de $\SI{2}{\kilo\ohm}$.
\begin{itensq}
\item Determine $i_L(0^-)$ e $i_L(0^+)$.
\item Calcule a constante de tempo da descarga.
\item Escreva $i_L(t)$ para $t\ge 0$ e determine o módulo da tensão inicial no resistor.
\end{itensq}
\resposta{$i_L(0^-)=i_L(0^+)=\SI{2}{\milli\ampere}$; $\tau=\SI{1}{\micro\second}$;
$i_L(t)=\SI{2}{\milli\ampere}e^{-t/(\SI{1}{\micro\second})}$; $|v_R(0^+)|=\SI{4}{\volt}$}


\end{questao}

\blocofixacao
\begin{questao}[1]{}

A tensão sobre um indutor vale $\SI{50}{\volt}$ quando a corrente varia à taxa
de $\SI{250}{\ampere\per\second}$. Determine a indutância.
\resposta{$L=\SI{0.2}{\henry}$}


\end{questao}
\begin{questao}[1]{}

Um indutor de $\SI{0.4}{\henry}$ conduz $\SI{5}{\ampere}$ constantes. Determine
a energia armazenada no campo magnético.
\resposta{$W=\SI{5}{\joule}$}


\end{questao}
\begin{questao}[1]{}

Sobre a corrente em um indutor ideal, é correto afirmar que ela:
\begin{alternativas}[2]
\item pode variar instantaneamente
\item não pode variar instantaneamente
\item é sempre constante
\item é sempre nula em corrente contínua
\end{alternativas}
\resposta{(b)}


\end{questao}
\begin{questao}[1]{}

Um indutor de $\SI{200}{\milli\henry}$ está em série com $\SI{50}{\ohm}$.
Determine a constante de tempo do circuito.
\resposta{$\tau=\SI{4}{\milli\second}$}


\end{questao}
\begin{questao}[1]{}

Determine o fluxo concatenado em um indutor de $\SI{0.25}{\henry}$ percorrido
por $\SI{8}{\ampere}$.
\resposta{$\lambda=\SI{2}{\weber}$}


\end{questao}
\begin{questao}[1]{}

Aplicam-se $\SI{30}{\volt}$ a um indutor de $\SI{0.15}{\henry}$ inicialmente
sem corrente. Determine a taxa inicial de variação da corrente.
\resposta{$\mathrm{d}i/\mathrm{d}t=\SI{200}{\ampere\per\second}$}


\end{questao}

\blocoaplicacao
\begin{questao}[2]{}

Um indutor de $\SI{0.5}{\henry}$ conduz $\SI{3}{\ampere}$.
\begin{itensq}
\item Determine a energia armazenada.
\item Determine a tensão necessária para levar a corrente a zero
      \emph{linearmente} em $\SI{50}{\milli\second}$.
\end{itensq}
\resposta{(a) $\SI{2.25}{\joule}$; (b) $\SI{30}{\volt}$}


\end{questao}
\begin{questao}[2]{}

A corrente em um indutor de $\SI{100}{\milli\henry}$ cai linearmente de
$\SI{4}{\ampere}$ a zero em $\SI{20}{\milli\second}$. Determine a tensão
induzida e sua polaridade.
\resposta{$\SI{20}{\volt}$, com polaridade \emph{invertida} em relação à fase
de crescimento}


\end{questao}
\begin{questao}[2]{}

Três indutores não acoplados de $\SI{6}{\milli\henry}$,
$\SI{12}{\milli\henry}$ e $\SI{12}{\milli\henry}$ são associados. Determine
$L_{eq}$ em série e em paralelo.
\resposta{Série: $\SI{30}{\milli\henry}$; paralelo: $\SI{3}{\milli\henry}$}


\end{questao}
\begin{questao}[2]{}

Um circuito RL série tem $L=\SI{0.6}{\henry}$, $R=\SI{30}{\ohm}$ e fonte de
$\SI{24}{\volt}$, com a chave fechada em $t=0$.
\begin{itensq}
\item Determine $\tau$ e a corrente final.
\item Determine a tensão sobre o indutor em $t=0^{+}$.
\end{itensq}
\resposta{(a) $\tau=\SI{20}{\milli\second}$, $I_f=\SI{0.8}{\ampere}$;
(b) $\SI{24}{\volt}$}


\end{questao}
\begin{questao}[2]{}

No circuito anterior, determine a corrente em $t=\tau$, $t=2\tau$ e
$t=3\tau$.
\resposta{$\SI{0.506}{\ampere}$, $\SI{0.692}{\ampere}$ e
$\SI{0.760}{\ampere}$}


\end{questao}
\begin{questao}[2]{}

Aplica-se a um indutor de $\SI{5}{\milli\henry}$ uma onda quadrada de tensão
que alterna entre $+\SI{10}{\volt}$ e $-\SI{10}{\volt}$, com meio período de
$\SI{1}{\milli\second}$. Determine a variação de corrente pico a pico em
regime.
\resposta{$\Delta i=\SI{2}{\ampere}$}


\end{questao}
\begin{questao}[2]{}

Um indutor \emph{real} de $\SI{0.5}{\henry}$ possui resistência de enrolamento
$r=\SI{8}{\ohm}$ e é ligado a uma fonte CC de $\SI{24}{\volt}$.
\begin{itensq}
\item Determine a corrente em regime permanente.
\item Determine a energia armazenada e a potência dissipada.
\end{itensq}
\resposta{(a) $\SI{3}{\ampere}$; (b) $W=\SI{2.25}{\joule}$ e
$P=\SI{72}{\watt}$}


\end{questao}
\begin{questao}[2]{}

Dois indutores armazenam a mesma energia: o primeiro tem
$L_1=\SI{2}{\henry}$ com $I_1=\SI{1}{\ampere}$. Se
$L_2=\SI{0.5}{\henry}$, determine $I_2$.
\resposta{$I_2=\SI{2}{\ampere}$}


\end{questao}
\begin{questao}[2]{}

Um solenoide de $N=200$ espiras tem indutância $L$. Determine a nova
indutância se o número de espiras for dobrado, mantidos o comprimento e a
seção.
\resposta{$4L$}


\end{questao}
\begin{questao}[2]{}

Um indutor com núcleo de ar tem $L_0=\SI{2}{\milli\henry}$. Introduz-se um
núcleo ferromagnético de permeabilidade relativa $\mu_r=500$. Determine a nova
indutância e comente as limitações.
\resposta{$L=\SI{1}{\henry}$}


\end{questao}
\begin{questao}[2]{}

No circuito RL com $\tau=\SI{20}{\milli\second}$ e $I_f=\SI{0.8}{\ampere}$,
determine o instante em que a corrente atinge $\SI{0.5}{\ampere}$.
\resposta{$t\approx\SI{19.6}{\milli\second}$}


\end{questao}
\begin{questao}[2]{}

Um indutor de $\SI{0.4}{\henry}$ conduz $\SI{2}{\ampere}$ quando a fonte é
substituída por um curto através de $R=\SI{20}{\ohm}$. Determine $\tau$ e a
corrente após $\SI{40}{\milli\second}$.
\resposta{$\tau=\SI{20}{\milli\second}$; $i=\SI{0.27}{\ampere}$}


\end{questao}

\blocoaprofundamento
\begin{questao}[3]{}

Um indutor de $\SI{0.2}{\henry}$ conduz $\SI{5}{\ampere}$ e é descarregado
através de um resistor.
\begin{itensq}
\item Determine a energia total dissipada no resistor.
\item Explique por que ela não depende do valor de $R$.
\item Determine como $R$ afeta o processo.
\end{itensq}
\resposta{(a) $\SI{2.5}{\joule}$; (b) conservação de energia; (c) $R$ altera a
\emph{duração} e a \emph{potência}, não a energia total}


\end{questao}
\begin{questao}[3]{}

Um indutor de $\SI{0.1}{\henry}$ conduz $\SI{2}{\ampere}$ quando uma chave
mecânica é aberta, interrompendo a corrente em cerca de
$\SI{1}{\micro\second}$.
\begin{itensq}
\item Estime a tensão induzida.
\item Explique por que esse valor não se observa na prática.
\item Indique a consequência real.
\end{itensq}
\resposta{(a) $\SI{200}{\kilo\volt}$ (estimativa idealizada);
(b) o ar rompe muito antes; (c) formação de arco elétrico nos contatos}


\end{questao}
\begin{questao}[3]{}

Um diodo de roda livre é colocado em antiparalelo com uma bobina de
$\SI{0.1}{\henry}$ e $r=\SI{5}{\ohm}$, que conduz $\SI{2}{\ampere}$.
\begin{itensq}
\item Explique o funcionamento no instante da abertura.
\item Determine a tensão máxima sobre a chave.
\item Determine a constante de tempo de extinção.
\end{itensq}
\resposta{(b) $\approx\SI{0.7}{\volt}$ acima da alimentação;
(c) $\tau=\SI{20}{\milli\second}$}


\end{questao}
\begin{questao}[3]{}

Dois indutores acoplados, $L_1=\SI{4}{\milli\henry}$ e
$L_2=\SI{9}{\milli\henry}$, têm indutância mútua $M=\SI{3}{\milli\henry}$.
\begin{itensq}
\item Determine $L_{eq}$ em série nas duas orientações possíveis.
\item Determine o coeficiente de acoplamento $k$.
\item Explique como identificar experimentalmente a orientação.
\end{itensq}
\resposta{(a) $\SI{19}{\milli\henry}$ e $\SI{7}{\milli\henry}$; (b) $k=0{,}5$}


\end{questao}
\begin{questao}[3]{}

Uma fonte de $\SI{12}{\volt}$ alimenta, através de $R_1=\SI{20}{\ohm}$, um nó
do qual saem $R_3=\SI{60}{\ohm}$ ao terra e um ramo com
$R_2=\SI{30}{\ohm}$ em série com $L=\SI{0.9}{\henry}$.
\begin{itensq}
\item Determine a resistência de Thévenin vista pelo indutor.
\item Determine $\tau$ e a corrente final no indutor.
\end{itensq}
\resposta{(a) $R_{Th}=\SI{45}{\ohm}$; (b) $\tau=\SI{20}{\milli\second}$ e
$I_f=\SI{0.2}{\ampere}$}


\end{questao}
\begin{questao}[3]{}

Um indutor de $\SI{0.5}{\henry}$ já conduz $\SI{1}{\ampere}$ quando é ligado a
uma fonte que estabeleceria $\SI{3}{\ampere}$ em regime, com
$R=\SI{25}{\ohm}$.
\begin{itensq}
\item Escreva $i(t)$.
\item Determine a corrente após $\SI{20}{\milli\second}$.
\item Determine a tensão inicial sobre o indutor.
\end{itensq}
\resposta{(a) $i(t)=3-2e^{-t/\tau}$ com $\tau=\SI{20}{\milli\second}$;
(b) $\SI{2.26}{\ampere}$; (c) $\SI{50}{\volt}$}


\end{questao}
\begin{questao}[3]{}

Em um circuito RL, mediu-se corrente final de $\SI{2}{\ampere}$ com fonte de
$\SI{50}{\volt}$, e verificou-se que a corrente atingiu $\SI{1.264}{\ampere}$
em $\SI{8}{\milli\second}$. Determine $R$ e $L$.
\resposta{$R=\SI{25}{\ohm}$; $L=\SI{0.2}{\henry}$}


\end{questao}
\begin{questao}[3]{}

Analise o sinal da potência instantânea $p=vi$ em um indutor durante um ciclo
completo de carga e descarga, e relacione-o com a energia armazenada.
\resposta{$p>0$ na carga (absorve), $p<0$ na descarga (devolve); a integral em
um ciclo completo é nula}


\end{questao}
\begin{questao}[3]{}

Determine a densidade de energia magnética em um núcleo onde
$B=\SI{1.2}{\tesla}$, supondo $\mu\approx\mu_0$ no entreferro.
\begin{itensq}
\item Calcule $u$.
\item Compare com a densidade de energia elétrica em um campo de
      $\SI{3}{\mega\volt\per\meter}$.
\end{itensq}
\dados{$\mu_0=4\pi\times10^{-7}\,\si{\henry\per\meter}$;
$\varepsilon_0=\pot{8.85}{-12}\,\si{\farad\per\meter}$}
\resposta{(a) $u\approx\SI{573}{\kilo\joule\per\meter\cubed}$;
(b) $\approx\SI{40}{\joule\per\meter\cubed}$}


\end{questao}
\begin{questao}[3]{}

Duas bobinas têm $L_1=\SI{50}{\milli\henry}$ e $L_2=\SI{200}{\milli\henry}$.
\begin{itensq}
\item Determine o valor máximo possível de $M$.
\item Se $M=\SI{80}{\milli\henry}$, determine $k$.
\item Interprete fisicamente $k=1$ e $k=0$.
\end{itensq}
\resposta{(a) $M_{\max}=\SI{100}{\milli\henry}$; (b) $k=0{,}8$}


\end{questao}
\begin{questao}[3]{}

Um indutor com núcleo apresenta $L=\SI{100}{\micro\henry}$ até
$\SI{3}{\ampere}$; acima disso, a saturação reduz $L$ a
$\SI{20}{\micro\henry}$. Ele opera com $\SI{12}{\volt}$ aplicados durante
$\SI{10}{\micro\second}$ a partir de $\SI{2}{\ampere}$.
\begin{itensq}
\item Determine a corrente ao final do pulso, ignorando a saturação.
\item Refaça considerando a saturação.
\item Comente a consequência prática.
\end{itensq}
\resposta{(a) $\SI{3.2}{\ampere}$; (b) $\SI{5}{\ampere}$; (c) risco de
destruição do interruptor}


\end{questao}
\begin{questao}[3]{}

Uma fonte de \textbf{corrente} de $\SI{2}{\ampere}$ alimenta um indutor de
$\SI{0.5}{\henry}$ em paralelo com $R=\SI{10}{\ohm}$, em regime permanente.
\begin{itensq}
\item Determine a corrente no indutor e no resistor.
\item A fonte é subitamente desligada (aberta). Determine a tensão inicial
      sobre o resistor.
\item Determine o tempo de extinção.
\end{itensq}
\resposta{(a) $\SI{2}{\ampere}$ no indutor, $\SI{0}{\ampere}$ no resistor;
(b) $\SI{-20}{\volt}$; (c) $\tau=\SI{50}{\milli\second}$}


\end{questao}

\blocodesafio
\begin{questao}[4]{}

\textbf{(Dedução --- resposta do circuito RL.)} Um circuito RL série com fonte
$V$ tem a chave fechada em $t=0$, com $i(0)=0$.
\begin{itensq}
\item Escreva a equação diferencial e resolva-a por separação de variáveis.
\item Identifique $\tau$ e o valor final.
\item Obtenha $v_L(t)$ e verifique a LKT em $t=0^{+}$ e $t\to\infty$.
\end{itensq}
\resposta{$i(t)=\dfrac{V}{R}\left(1-e^{-Rt/L}\right)$;
$\tau=L/R$; $v_L=Ve^{-Rt/L}$}


\end{questao}
\begin{questao}[4]{}

\textbf{(Integração --- energia na descarga.)} Um indutor com corrente inicial
$I_0$ é descarregado através de $R$.
\begin{itensq}
\item Escreva $i(t)$ e a potência instantânea no resistor.
\item Integre para obter a energia total dissipada.
\item Mostre que o resultado independe de $R$ e interprete.
\end{itensq}
\resposta{$W=\frac12LI_0^{2}$, independente de $R$}


\end{questao}
\begin{questao}[4]{}

\textbf{(Projeto de proteção.)} Uma bobina de relé com $L=\SI{0.2}{\henry}$ e
$r=\SI{100}{\ohm}$ opera em $\SI{24}{\volt}$. O transistor de comando suporta
no máximo $\SI{60}{\volt}$.
\begin{itensq}
\item Determine a corrente de operação e a energia armazenada.
\item Dimensione um grampeador resistivo (diodo em série com $R_g$) que
      respeite o limite do transistor.
\item Compare o tempo de extinção com o do diodo simples.
\end{itensq}
\resposta{(a) $\SI{0.24}{\ampere}$, $\SI{5.76}{\milli\joule}$;
(b) $R_g\le\SI{150}{\ohm}$; (c) $\SI{0.8}{\milli\second}$ contra
$\SI{2}{\milli\second}$}


\end{questao}
\begin{questao}[4]{}

\textbf{(Projeto de armazenamento.)} Deseja-se um indutor que armazene
$\SI{5}{\joule}$ conduzindo $\SI{10}{\ampere}$.
\begin{itensq}
\item Determine $L$.
\item Para um núcleo toroidal com $A=\SI{10}{\centi\meter\squared}$,
      $\ell=\SI{0.25}{\meter}$ e $\mu_r=2000$, determine o número de espiras.
\item Verifique a densidade de fluxo e avalie a viabilidade.
\end{itensq}
\resposta{(a) $L=\SI{0.1}{\henry}$; (b) $N\approx35$ espiras;
(c) $B\approx\SI{3.5}{\tesla}$ --- \textbf{inviável}, o núcleo satura}


\end{questao}
\begin{questao}[4]{}

\textbf{(Projeto de constante de tempo.)} Deseja-se $\tau=\SI{20}{\milli\second}$
com $R=\SI{50}{\ohm}$.
\begin{itensq}
\item Determine $L$.
\item Avalie a viabilidade física dessa indutância.
\item Proponha duas alternativas para obter o mesmo $\tau$ de forma prática.
\end{itensq}
\resposta{(a) $L=\SI{1}{\henry}$; (c) aumentar $R$, ou usar RC equivalente}


\end{questao}
\begin{questao}[4]{}

\textbf{(Demonstração --- acoplamento em série.)} Duas bobinas acopladas em
série conduzem a mesma corrente $i$.
\begin{itensq}
\item Deduza $L_{eq}=L_1+L_2\pm2M$ a partir das tensões induzidas.
\item Demonstre que $M\le\sqrt{L_1L_2}$, usando o fato de a energia ser não
      negativa.
\item Explique o significado da convenção dos pontos.
\end{itensq}
\resposta{(b) da condição $W\ge0$ para toda corrente segue
$M^{2}\le L_1L_2$}


\end{questao}
\begin{questao}[4]{}

\textbf{(Indutor real em regime.)} Um indutor com $L=\SI{0.5}{\henry}$ e
$r=\SI{2}{\ohm}$ é alimentado por $\SI{10}{\volt}$ CC.
\begin{itensq}
\item Determine a corrente, a energia armazenada e a potência dissipada em
      regime.
\item Determine o tempo necessário para atingir $99\%$ da corrente final.
\item Discuta o "custo de manutenção" da energia armazenada.
\end{itensq}
\resposta{(a) $\SI{5}{\ampere}$, $\SI{6.25}{\joule}$, $\SI{50}{\watt}$;
(b) $\approx\SI{1.15}{\second}$}


\end{questao}
\begin{questao}[4]{}

\textbf{(Conversor chaveado.)} Um conversor opera a $\SI{50}{\kilo\hertz}$ com
razão cíclica $D=0{,}5$, aplicando $\SI{12}{\volt}$ sobre um indutor de
$\SI{100}{\micro\henry}$ durante a condução.
\begin{itensq}
\item Determine a ondulação de corrente pico a pico.
\item Determine a corrente de pico se a componente contínua for
      $\SI{3}{\ampere}$.
\item O indutor satura em $\SI{3.5}{\ampere}$. Proponha e quantifique duas
      correções.
\end{itensq}
\resposta{(a) $\Delta i=\SI{1.2}{\ampere}$; (b) $I_{pico}=\SI{3.6}{\ampere}$;
(c) satura --- elevar $L$ ou a frequência}


\end{questao}
\begin{questao}[4]{}

\textbf{(Demonstração --- continuidade da corrente.)} Prove que a corrente em
um indutor ideal não pode sofrer descontinuidade, usando um argumento
energético.
\begin{itensq}
\item Suponha um salto de $I_1$ para $I_2$ em intervalo $\Delta t\to0$ e
      analise a tensão.
\item Analise a potência e a energia envolvidas.
\item Explique por que capacitores em paralelo apresentam o problema dual.
\end{itensq}
\resposta{Um salto exigiria tensão e potência infinitas; a energia é finita,
mas a taxa de transferência não pode ser}


\end{questao}
\begin{questao}[4]{}

\textbf{(Modelagem --- sensor indutivo.)} Um sensor de proximidade usa uma
bobina cuja indutância varia com a distância $x$ a um alvo ferromagnético,
segundo $L(x)=L_0\left(1+\dfrac{a}{x}\right)$, com
$L_0=\SI{10}{\milli\henry}$ e $a=\SI{1}{\milli\meter}$.
\begin{itensq}
\item Determine $L$ para $x=\SI{1}{\milli\meter}$ e $x=\SI{5}{\milli\meter}$.
\item Obtenha a sensibilidade $\mathrm{d}L/\mathrm{d}x$ e avalie nos dois
      pontos.
\item Discuta a consequência para o projeto do sensor.
\end{itensq}
\resposta{(a) $\SI{20}{\milli\henry}$ e $\SI{12}{\milli\henry}$;
(b) $-\SI{10}{\milli\henry\per\milli\meter}$ e
$-\SI{0.4}{\milli\henry\per\milli\meter}$}


\end{questao}

\gabarito[10.2 Indutores]

\section{Circuitos de Primeira e Segunda Ordem}

\blocofixacao
\begin{questao}[1]{}

Um capacitor ideal possui $v_C(0^-)=\SI{7}{\volt}$. No instante $t=0$ ocorre uma
comutação no circuito. Qual é o valor de $v_C(0^+)$?
\resposta{$v_C(0^+)=\SI{7}{\volt}$}


\end{questao}
\begin{questao}[1]{}

Um indutor ideal conduz $i_L(0^-)=\SI{350}{\milli\ampere}$. Uma chave muda de
posição em $t=0$. Determine $i_L(0^+)$.
\resposta{$i_L(0^+)=\SI{350}{\milli\ampere}$}


\end{questao}
\begin{questao}[1]{}

Um circuito RC tem $R=\SI{2}{\kilo\ohm}$ e $C=\SI{50}{\micro\farad}$. Determine
a constante de tempo.
\resposta{$\tau=\SI{100}{\milli\second}$}


\end{questao}
\begin{questao}[1]{}

Um circuito RL possui $L=\SI{80}{\milli\henry}$ e resistência equivalente
$R_{\mathrm{eq}}=\SI{20}{\ohm}$. Calcule $\tau$.
\resposta{$\tau=\SI{4}{\milli\second}$}


\end{questao}
\begin{questao}[1]{}

Um capacitor inicialmente descarregado é carregado por uma fonte CC através de
um resistor. Após um intervalo igual a uma constante de tempo, aproximadamente
qual fração da variação total de tensão já ocorreu?
\begin{alternativas}[3]
\item \SI{36.8}{\percent}
\item \SI{50}{\percent}
\item \SI{63.2}{\percent}
\item \SI{86.5}{\percent}
\item \SI{99.3}{\percent}
\end{alternativas}
\resposta{(c)}


\end{questao}
\begin{questao}[1]{}

Uma resposta natural de primeira ordem é proporcional a $e^{-t/\tau}$. Qual
percentual do valor inicial resta após $3\tau$?
\resposta{$e^{-3}\times100\%\approx\SI{4.98}{\percent}$}


\end{questao}
\begin{questao}[1]{}

Ao calcular a resistência equivalente vista por um capacitor para obter a
constante de tempo, como devem ser tratadas as fontes independentes ideais?
\begin{alternativas}
\item fontes de tensão abertas e fontes de corrente em curto.
\item todas as fontes abertas.
\item fontes de tensão em curto e fontes de corrente abertas.
\item todas as fontes em curto.
\item as fontes não devem ser modificadas.
\end{alternativas}
\resposta{(c)}


\end{questao}
\begin{questao}[1]{}

Um circuito ideal possui um capacitor e um indutor independentes entre si como
únicos elementos armazenadores de energia. Qual é, em geral, a ordem máxima do
circuito?
\resposta{Segunda ordem}


\end{questao}
\begin{questao}[1]{}

Para $L=\SI{0.2}{\henry}$ e $C=\SI{50}{\micro\farad}$, determine a frequência
natural não amortecida $\omega_0$.
\resposta{$\omega_0\approx\SI{316.2}{\radian\per\second}$}


\end{questao}
\begin{questao}[1]{}

Em um circuito RLC em série, $R=\SI{40}{\ohm}$, $L=\SI{0.1}{\henry}$ e
$C=\SI{100}{\micro\farad}$. Classifique o amortecimento.
\resposta{Subamortecido}


\end{questao}
\begin{questao}[1]{}

Determine a resistência crítica de um circuito RLC em série com $L=\SI{50}{\milli\henry}$ e
$C=\SI{200}{\micro\farad}$.
\resposta{$R_{\mathrm{crit}}\approx\SI{31.6}{\ohm}$}


\end{questao}

\gabarito[Nível 1]

\blocoaplicacao
\begin{questao}[2]{}

A tensão de um capacitor muda de $\SI{8}{\volt}$ para um valor final de
$\SI{2}{\volt}$ através de $R=\SI{4}{\kilo\ohm}$ e
$C=\SI{25}{\micro\farad}$. Determine $v_C(\SI{0.2}{\second})$.
\resposta{$v_C(\SI{0.2}{\second})\approx\SI{2.81}{\volt}$}


\end{questao}
\begin{questao}[2]{}

A corrente de um indutor passa de $\SI{1}{\ampere}$ para um valor final de
$\SI{4}{\ampere}$. Sabendo que $L=\SI{0.6}{\henry}$ e
$R_{\mathrm{eq}}=\SI{30}{\ohm}$, calcule $i_L(\SI{40}{\milli\second})$.
\resposta{$i_L\approx\SI{3.59}{\ampere}$}


\end{questao}
\begin{questao}[2]{}

Um circuito RC inicialmente descarregado possui $R=\SI{10}{\kilo\ohm}$ e
$C=\SI{10}{\micro\farad}$. Quanto tempo é necessário para atingir
\SI{90}{\percent} do valor final?
\resposta{$t\approx\SI{230}{\milli\second}$}


\end{questao}
\begin{questao}[2]{}

Um capacitor descarrega de $\SI{12}{\volt}$ para $\SI{3}{\volt}$ com
$\tau=\SI{50}{\milli\second}$. Determine o tempo decorrido.
\resposta{$t\approx\SI{69.3}{\milli\second}$}


\end{questao}
\begin{questao}[2]{}

A corrente natural de um circuito RL cai de $\SI{2}{\ampere}$ para
$\SI{0.2}{\ampere}$. Se $\tau=\SI{25}{\milli\second}$, determine o intervalo de
tempo.
\resposta{$t\approx\SI{57.6}{\milli\second}$}


\end{questao}
\begin{questao}[2]{}

No circuito abaixo, a chave conecta a fonte em $t=0$ e o capacitor está
inicialmente descarregado. Determine $v_C(\tau)$ e a constante de tempo.

\circuitoq{
\begin{circuitikz}[scale=1,transform shape,line width=0.8pt]
\draw (0,0) to[battery1,l_={$\SI{20}{\volt}$}] (0,3)
      to[R,l={$R_1=\SI{6}{\kilo\ohm}$}] (3,3) coordinate(a);
\draw (a) to[R,l={$R_2=\SI{3}{\kilo\ohm}$}] (3,0) -- (0,0);
\draw (a) -- (5,3) to[C,l={$C=\SI{50}{\micro\farad}$}] (5,0) -- (3,0);
\end{circuitikz}}

\resposta{$\tau=\SI{100}{\milli\second}$; $v_C(\tau)\approx\SI{4.21}{\volt}$}


\end{questao}
\begin{questao}[2]{}

Após uma comutação, um capacitor possui $v_C(0^+)=\SI{5}{\volt}$ e tende a
$-\SI{10}{\volt}$, com $\tau=\SI{20}{\milli\second}$. Determine
$v_C(\SI{10}{\milli\second})$.
\resposta{$v_C\approx-\SI{0.90}{\volt}$}


\end{questao}
\begin{questao}[2]{}

Um indutor de $\SI{40}{\milli\henry}$ está conectado, após uma
comutação, a uma rede cuja resistência equivalente vista em seus terminais é
$R_{\mathrm{eq}}=\SI{4}{\ohm}$. Sua corrente inicial é $\SI{0.5}{\ampere}$ e o valor final é
$\SI{2}{\ampere}$. Calcule a corrente após $\SI{20}{\milli\second}$.
\resposta{$i_L\approx\SI{1.80}{\ampere}$}


\end{questao}
\begin{questao}[2]{}

Um circuito RLC em série tem $R=\SI{10}{\ohm}$, $L=\SI{50}{\milli\henry}$ e
$C=\SI{200}{\micro\farad}$. Determine $\alpha$, $\omega_0$, $\omega_d$ e o tipo
de resposta.
\resposta{$\alpha=\SI{100}{\per\second}$; $\omega_0\approx\SI{316.2}{\radian\per\second}$;
$\omega_d=\SI{300}{\radian\per\second}$; subamortecida}


\end{questao}
\begin{questao}[2]{}

Para um circuito RLC em série com $R=\SI{40}{\ohm}$, $L=\SI{50}{\milli\henry}$ e
$C=\SI{200}{\micro\farad}$, classifique a resposta natural.
\resposta{Superamortecida}


\end{questao}
\begin{questao}[2]{}

Calcule $R_{\mathrm{crit}}$ para um circuito RLC em série com $L=\SI{0.1}{\henry}$ e
$C=\SI{100}{\micro\farad}$.
\resposta{$R_{\mathrm{crit}}\approx\SI{63.25}{\ohm}$}


\end{questao}
\begin{questao}[2]{}

Uma resposta subamortecida de tensão tem a forma
\[
v(t)=e^{-100t}\big(10\cos300t+B\sin300t\big).
\]
Sabendo que $\mathrm{d}v/\mathrm{d}t|_{0^+}=0$, determine $B$.
\resposta{$B\approx\SI{3.33}{\volt}$}


\end{questao}
\begin{questao}[2]{}

Em um capacitor de $C=\SI{100}{\micro\farad}$, a corrente em $0^+$ é
$\SI{20}{\milli\ampere}$ entrando no terminal positivo. Determine
$\mathrm{d}v_C/\mathrm{d}t|_{0^+}$.
\resposta{$\mathrm{d}v_C/\mathrm{d}t|_{0^+}=\SI{200}{\volt\per\second}$}


\end{questao}
\begin{questao}[2]{}

Um circuito RLC em paralelo possui $R=\SI{200}{\ohm}$, $C=\SI{100}{\micro\farad}$ e
$L=\SI{0.1}{\henry}$. Classifique a resposta natural da tensão.
\resposta{Subamortecida}


\end{questao}
\begin{questao}[2]{}

Em uma resposta de primeira ordem, qual erro percentual em relação ao valor final
permanece após $5\tau$?
\resposta{$\approx\SI{0.674}{\percent}$}


\end{questao}
\begin{questao}[2]{}

Um circuito LC ideal possui um capacitor de $C=\SI{100}{\micro\farad}$
com tensão inicial de $\SI{10}{\volt}$ e um indutor de
$L=\SI{50}{\milli\henry}$ com corrente inicial nula.
Determine a energia inicial e a corrente máxima ideal.
\resposta{$E=\SI{5}{\milli\joule}$; $I_{\max}\approx\SI{0.447}{\ampere}$}


\end{questao}
\begin{questao}[2]{}

Para $L=\SI{10}{\milli\henry}$ e $C=\SI{100}{\micro\farad}$, determine o
período natural ideal $T_0$.
\resposta{$T_0\approx\SI{6.28}{\milli\second}$}


\end{questao}

\gabarito[Nível 2]

\blocoaplicacao
\begin{questao}[2]{}

A chave do circuito permaneceu fechada por muito tempo e \textbf{abre} em
$t=0$. Determine $v_C(0^+)$, a constante de tempo para $t>0$ e
$v_C(\SI{0.30}{\second})$.

\circuitoqq{
\begin{circuitikz}[scale=1,transform shape,line width=0.8pt]
\draw (0,0) to[american voltage source,l_={$\SI{18}{\volt}$}] (0,3)
      to[switch,l={abre em $t=0$}] (1.7,3)
      to[R,l={$R_1=\SI{3}{\kilo\ohm}$}] (4,3) coordinate(a);
\draw (a) to[R,l={$R_2=\SI{6}{\kilo\ohm}$}] (4,0) -- (0,0);
\draw (a) -- (6,3)
      to[C,l={$C=\SI{100}{\micro\farad}$},v^>={$v_C$}] (6,0) -- (4,0);
\end{circuitikz}}{Circuito RC com mudança de topologia em $t=0$.}

\resposta{$v_C(0^+)=\SI{12}{\volt}$; $\tau=\SI{0.60}{\second}$;
$v_C(0{,}30\,\mathrm{s})\approx\SI{7.28}{\volt}$}


\end{questao}
\begin{questao}[2]{}

No circuito abaixo, a chave permaneceu fechada por muito tempo e abre em
$t=0$. Determine $i_L(0^+)$, a constante de tempo da descarga e
$i_L(\SI{10}{\milli\second})$.

\circuitoqq{
\begin{circuitikz}[scale=1,transform shape,line width=0.8pt]
\draw (0,0) to[american voltage source,l_={$\SI{24}{\volt}$}] (0,3)
      to[switch,l={abre em $t=0$}] (1.7,3)
      to[R,l={$R_1=\SI{6}{\ohm}$}] (4,3) coordinate(a);
\draw (a) to[L,l={$L=\SI{120}{\milli\henry}$},i>^={$i_L$}] (4,0) -- (0,0);
\draw (a) -- (6,3) to[R,l={$R_2=\SI{18}{\ohm}$}] (6,0) -- (4,0);
\end{circuitikz}}{Circuito RL com caminho de descarga resistivo.}

\resposta{$i_L(0^+)=\SI{4}{\ampere}$; $\tau\approx\SI{6.67}{\milli\second}$;
$i_L(10\,\mathrm{ms})\approx\SI{0.893}{\ampere}$}


\end{questao}
\begin{questao}[2]{}

A fonte de corrente já está ligada há muito tempo. Em $t=0$ a chave fecha,
inserindo $R_2$ em paralelo com o capacitor. Determine $v_C(0^+)$,
$v_C(\infty)$, $\tau$ e escreva $v_C(t)$ para $t>0$.

\circuitoqq{
\begin{circuitikz}[scale=1,transform shape,line width=0.8pt]
\draw (0,0) to[american current source,l_={$\SI{2}{\milli\ampere}$}] (0,3) -- (2,3) coordinate(a);
\draw (a) to[R,l={$R_1=\SI{4}{\kilo\ohm}$}] (2,0) -- (0,0);
\draw (a) -- (4,3) to[C,l={$C=\SI{50}{\micro\farad}$},v^>={$v_C$}] (4,0) -- (2,0);
\draw (a) -- (5.5,3) to[switch,l={fecha em $t=0$}] (6.8,3)
      to[R,l={$R_2=\SI{12}{\kilo\ohm}$}] (6.8,0) -- (4,0);
\end{circuitikz}}{Mudança da resistência vista pelo capacitor.}

\resposta{$v_C(0^+)=\SI{8}{\volt}$; $v_C(\infty)=\SI{6}{\volt}$;
$\tau=\SI{150}{\milli\second}$;
$v_C(t)=6+2e^{-t/0{,}15}\ \si{\volt}$}


\end{questao}
\begin{questao}[2]{}

A chave fecha em $t=0$, aplicando um degrau de $\SI{20}{\volt}$ ao circuito RLC
em série. Sem resolver ainda a resposta completa, determine $\alpha$, $\omega_0$,
$\omega_d$ e classifique o amortecimento.

\circuitoqq{
\begin{circuitikz}[scale=1,transform shape,line width=0.8pt]
\draw (0,0) to[american voltage source,l_={$\SI{20}{\volt}$}] (0,3)
      to[switch,l={fecha em $t=0$}] (1.6,3)
      to[R,l={$R=\SI{15}{\ohm}$}] (3.6,3)
      to[L,l={$L=\SI{50}{\milli\henry}$}] (5.8,3)
      to[C,l={$C=\SI{200}{\micro\farad}$}] (7.8,3) -- (7.8,0) -- (0,0);
\end{circuitikz}}{Circuito RLC em série excitado por degrau.}

\resposta{$\alpha=\SI{150}{\per\second}$;
$\omega_0\approx\SI{316.2}{\radian\per\second}$;
$\omega_d\approx\SI{278.4}{\radian\per\second}$; resposta subamortecida}


\end{questao}
\begin{questao}[2]{}

O circuito RLC em paralelo abaixo parte sem energia armazenada e recebe a fonte
$i_s(t)=3u(t)\,\si{\ampere}$. Determine $v(0^+)$, $i_L(0^+)$,
$\mathrm{d}v/\mathrm{d}t|_{0^+}$, $v(\infty)$ e $i_L(\infty)$.

\circuitoqq{
\begin{circuitikz}[scale=1,transform shape,line width=0.8pt]
\draw (0,0) to[american current source,l_={$3u(t)\,\si{\ampere}$}] (0,3) -- (6,3);
\draw (2,3) to[R,l={$R=\SI{6}{\ohm}$}] (2,0);
\draw (4,3) to[C,l={$C=\SI{0.5}{\farad}$},v^>={$v$}] (4,0);
\draw (6,3) to[L,l={$L=\SI{2}{\henry}$},i>^={$i_L$}] (6,0) -- (0,0);
\end{circuitikz}}{RLC em paralelo: condições iniciais, derivada inicial e regime final.}

\resposta{$v(0^+)=0$; $i_L(0^+)=0$;
$\dot v(0^+)=\SI{6}{\volt\per\second}$;
$v(\infty)=0$; $i_L(\infty)=\SI{3}{\ampere}$}


\end{questao}

\gabarito[Nível 2 --- circuitos e chaveamento]

\blocoaplicacao
\begin{questao}[2]{}

O circuito abaixo parte com o capacitor descarregado. Em $t=0$, a fonte de corrente
independente aplica um degrau de $\SI{3}{\milli\ampere}$. A fonte dependente é uma
VCCS de valor $g v_C$, com $g=\SI{0.25}{\milli\ensuremath{\mathrm{S}}}$. Determine a constante de
tempo, o valor final de $v_C$ e $v_C(\SI{0.10}{\second})$.

\circuitoqq{
\begin{circuitikz}[scale=0.96,transform shape,line width=0.8pt,every node/.append style={font=\scriptsize}]
\draw (0,0) to[american current source] (0,3) -- (7.0,3);
\node[anchor=east] at (-0.25,1.5) {$3u(t)\,\mathrm{mA}$};
\draw (2.0,3) to[R] (2.0,0);
\node[anchor=east] at (1.75,1.5) {$R=\SI{4}{\kilo\ohm}$};
\draw (4.5,3) to[C,v^>={$v_C$}] (4.5,0);
\node[anchor=east] at (4.25,1.5) {$C=\SI{50}{\micro\farad}$};
\draw (7.0,3) to[american controlled current source] (7.0,0) -- (0,0);
\node[anchor=west] at (7.25,1.5) {$g\,v_C$};
\end{circuitikz}}{Circuito RC com fonte de corrente controlada por tensão (VCCS).}

\resposta{$R_{\mathrm{eq}}=\SI{2}{\kilo\ohm}$; $\tau=\SI{0.10}{\second}$;
$v_C(\infty)=\SI{6}{\volt}$; $v_C(0{,}10\,\mathrm{s})\approx\SI{3.79}{\volt}$}


\end{questao}
\begin{questao}[2]{}

O capacitor do circuito natural abaixo está inicialmente a $\SI{10}{\volt}$.
A fonte dependente fornece a corrente $g v_C$, com
$g=\SI{0.5}{\milli\ensuremath{\mathrm{S}}}$. Determine $R_{\mathrm{eq}}$, $\tau$ e $v_C(t)$.

\circuitoqq{
\begin{circuitikz}[scale=0.98,transform shape,line width=0.8pt,every node/.append style={font=\scriptsize}]
\draw (0,0) -- (7.0,0);
\draw (0,3) -- (7.0,3);
\draw (1.8,3) to[R] (1.8,0);
\node[anchor=east] at (1.55,1.5) {$R=\SI{2}{\kilo\ohm}$};
\draw (4.2,3) to[C,v^>={$v_C$}] (4.2,0);
\node[anchor=east] at (3.95,1.5) {$C=\SI{100}{\micro\farad}$};
\draw (7.0,3) to[american controlled current source] (7.0,0);
\node[anchor=west] at (7.25,1.5) {$g\,v_C$};
\end{circuitikz}}{Resposta natural RC com uma VCCS em paralelo.}

\resposta{$R_{\mathrm{eq}}=\SI{1}{\kilo\ohm}$; $\tau=\SI{0.10}{\second}$;
$v_C(t)=10e^{-10t}\ \si{\volt}$}


\end{questao}
\begin{questao}[2]{}

No circuito abaixo, o indutor possui $i_L(0^+)=\SI{2}{\ampere}$. A fonte dependente
é uma CCCS que injeta para cima a corrente $\beta i_x$, em que
$i_x=v/R$ e $\beta=0{,}5$. Determine a resistência equivalente vista pelo indutor,
a constante de tempo e $i_L(t)$.

\circuitoqq{
\begin{circuitikz}[scale=0.98,transform shape,line width=0.8pt,every node/.append style={font=\scriptsize}]
\draw (0,0) -- (7.2,0);
\draw (0,3) -- (7.2,3);
\draw (1.8,3) to[R,i>^={$i_x$}] (1.8,0);
\node[anchor=east] at (1.55,1.5) {$R=\SI{20}{\ohm}$};
\draw (4.2,3) to[L,i>^={$i_L$}] (4.2,0);
\node[anchor=east] at (3.95,1.5) {$L=\SI{200}{\milli\henry}$};
\draw (7.2,0) to[american controlled current source] (7.2,3);
\node[anchor=west] at (7.45,1.5) {$\beta\,i_x$};
\end{circuitikz}}{Resposta natural RL com fonte de corrente controlada por corrente (CCCS).}

\resposta{$R_{\mathrm{eq}}=\SI{40}{\ohm}$; $\tau=\SI{5}{\milli\second}$;
$i_L(t)=2e^{-200t}\ \si{\ampere}$}


\end{questao}

\gabarito[Nível 2 --- fontes dependentes]

\blocoaprofundamento
\begin{questao}[3]{}

Um capacitor de $\SI{10}{\micro\farad}$, inicialmente descarregado, é ligado a
$\SI{12}{\volt}$ por $R_1=\SI{2}{\kilo\ohm}$ durante
$\SI{50}{\milli\second}$. Em seguida, a fonte é removida e ele passa a
descarregar por $R_2=\SI{5}{\kilo\ohm}$. Determine sua tensão
$\SI{40}{\milli\second}$ após a segunda comutação.
\resposta{$v_C\approx\SI{4.95}{\volt}$}


\end{questao}
\begin{questao}[3]{}

Um indutor de $\SI{100}{\milli\henry}$ é energizado por uma fonte de
$\SI{20}{\volt}$ através de $R_1=\SI{10}{\ohm}$ durante
$\SI{15}{\milli\second}$. Depois, a fonte é retirada e o indutor descarrega por
$R_2=\SI{25}{\ohm}$. Qual é a corrente após mais $\SI{8}{\milli\second}$?
\resposta{$i_L\approx\SI{0.210}{\ampere}$}


\end{questao}
\begin{questao}[3]{}

Após a supressão das fontes independentes, a resistência vista por um
capacitor de $C=\SI{20}{\micro\farad}$ é composta por
$R_3=\SI{1}{\kilo\ohm}$ em série com o paralelo de
$R_1=\SI{6}{\kilo\ohm}$ e $R_2=\SI{3}{\kilo\ohm}$. Determine a
constante de tempo.

\circuitoq{
\begin{circuitikz}[scale=1,transform shape,line width=0.8pt]
\draw (0,1.5) to[C,l={$C$},o-] (2,1.5) to[R,l={$R_3$}] (4,1.5) coordinate(a);
\draw (a) -- (4,2.6) to[R,l={$R_1$}] (6,2.6) -- (6,0.4);
\draw (a) -- (4,0.4) to[R,l_={$R_2$}] (6,0.4) -- (6,-0.2) node[ground]{};
\end{circuitikz}}

\resposta{$R_{\mathrm{eq}}=\SI{3}{\kilo\ohm}$; $\tau=\SI{60}{\milli\second}$}


\end{questao}
\begin{questao}[3]{}

Após uma comutação, $v_C(0^+)=-\SI{4}{\volt}$, $v_C(\infty)=\SI{8}{\volt}$,
$R_{\mathrm{eq}}=\SI{2.5}{\kilo\ohm}$ e $C=\SI{40}{\micro\farad}$. Em que
instante a tensão cruza zero?
\resposta{$t\approx\SI{40.5}{\milli\second}$}


\end{questao}
\begin{questao}[3]{}

Um indutor possui $i_L(0^+)=-\SI{1}{\ampere}$, $i_L(\infty)=\SI{3}{\ampere}$ e
$\tau=\SI{25}{\milli\second}$. Determine quando a corrente muda de sinal.
\resposta{$t\approx\SI{7.19}{\milli\second}$}


\end{questao}
\begin{questao}[3]{}

Um capacitor com $v_C(0)=\SI{2}{\volt}$ é ligado a uma fonte de
$\SI{10}{\volt}$ através de $R=\SI{2}{\kilo\ohm}$ e
$C=\SI{50}{\micro\farad}$. Determine a corrente inicial e a corrente após
$\SI{150}{\milli\second}$.
\resposta{$i(0^+)=\SI{4}{\milli\ampere}$; $i(150\,\mathrm{ms})\approx\SI{0.893}{\milli\ampere}$}


\end{questao}
\begin{questao}[3]{}

Um capacitor de $\SI{100}{\micro\farad}$ inicialmente a $\SI{20}{\volt}$ é
descarregado por $R=\SI{1}{\kilo\ohm}$. Mostre, por integração da potência no
resistor, que toda a energia inicialmente armazenada é dissipada e calcule-a.
\resposta{$E=\SI{20}{\milli\joule}$}


\end{questao}
\begin{questao}[3]{}

Projete um atraso RC para que uma tensão de $\SI{5}{\volt}$ alcance
$\SI{3.3}{\volt}$ em $\SI{0.5}{\second}$. Use $R=\SI{100}{\kilo\ohm}$ e
considere o capacitor inicialmente descarregado. Determine $C$ e indique um
valor comercial próximo.
\resposta{$C\approx\SI{4.63}{\micro\farad}$; valor comercial: \SI{4.7}{\micro\farad}}


\end{questao}
\begin{questao}[3]{}

Deseja-se que a corrente de um circuito RL alcance \SI{95}{\percent} de seu valor
final de $\SI{1}{\ampere}$ em $\SI{20}{\milli\second}$. Se
$L=\SI{50}{\milli\henry}$, determine a resistência e a tensão da fonte.
\resposta{$R\approx\SI{7.49}{\ohm}$; $V\approx\SI{7.49}{\volt}$}


\end{questao}
\begin{questao}[3]{}

Um circuito RLC em série tem $R=\SI{20}{\ohm}$, $L=\SI{0.1}{\henry}$ e
$C=\SI{100}{\micro\farad}$. Determine os polos naturais.
\resposta{$s_{1,2}=-100\pm j300\ \si{\per\second}$}


\end{questao}
\begin{questao}[3]{}

Uma resposta criticamente amortecida é
$v(t)=(A+Bt)e^{-200t}$. Se $v(0)=\SI{10}{\volt}$ e
$\dot v(0)=-\SI{500}{\volt\per\second}$, determine $A$ e $B$.
\resposta{$A=\SI{10}{\volt}$; $B=\SI{1500}{\volt\per\second}$}


\end{questao}
\begin{questao}[3]{}

Para um circuito RLC em série com $R=\SI{100}{\ohm}$, $L=\SI{0.1}{\henry}$ e
$C=\SI{100}{\micro\farad}$, determine os dois polos reais.
\resposta{$s_1\approx-\SI{112.7}{\per\second}$; $s_2\approx-\SI{887.3}{\per\second}$}


\end{questao}
\begin{questao}[3]{}

Determine a resistência crítica de um RLC \textbf{paralelo} com
$L=\SI{50}{\milli\henry}$ e $C=\SI{200}{\micro\farad}$.
\resposta{$R_{\mathrm{crit}}\approx\SI{7.91}{\ohm}$}


\end{questao}
\begin{questao}[3]{}

Para um sistema subamortecido de segunda ordem com $\zeta=0{,}5$, estime o
sobressinal percentual de pico usando
$M_p=e^{-\pi\zeta/\sqrt{1-\zeta^2}}$.
\resposta{$M_p\approx\SI{16.3}{\percent}$}


\end{questao}
\begin{questao}[3]{}

Um circuito RLC em série possui $R=\SI{20}{\ohm}$, $L=\SI{100}{\milli\henry}$ e
$C=\SI{100}{\micro\farad}$. Calcule o fator de amortecimento $\zeta$.
\resposta{$\zeta\approx0{,}316$}


\end{questao}
\begin{questao}[3]{}

Classifique os quatro casos da tabela.

\begin{table}[H]
\centering\small\sffamily
\begin{tabular}{@{}c c c c c@{}}
\toprule
\textbf{Caso} & \textbf{Topologia} & $R$ & $L$ & $C$\\
\midrule
A & série & \SI{20}{\ohm} & \SI{0.1}{\henry} & \SI{100}{\micro\farad}\\
B & série & \SI{63.25}{\ohm} & \SI{0.1}{\henry} & \SI{100}{\micro\farad}\\
C & série & \SI{100}{\ohm} & \SI{0.1}{\henry} & \SI{100}{\micro\farad}\\
D & paralelo & \SI{20}{\ohm} & \SI{50}{\milli\henry} & \SI{200}{\micro\farad}\\
\bottomrule
\end{tabular}
\end{table}
\resposta{A subamortecido; B aproximadamente crítico; C superamortecido; D subamortecido}


\end{questao}
\begin{questao}[3]{}

Em certo instante, o capacitor de $C=\SI{100}{\micro\farad}$ apresenta
tensão de $\SI{10}{\volt}$, enquanto o indutor de
$L=\SI{50}{\milli\henry}$ conduz corrente de $\SI{0.2}{\ampere}$. Qual é a energia total armazenada naquele instante?
\resposta{$E=\SI{6}{\milli\joule}$}


\end{questao}
\begin{questao}[3]{}

Antes de uma comutação, um capacitor possui $v_C=\SI{6}{\volt}$ e um indutor
possui $i_L=\SI{0.3}{\ampere}$. Logo após a comutação, ambos passam a integrar
uma nova rede resistiva. Quais condições iniciais devem ser usadas na nova rede?
\resposta{$v_C(0^+)=\SI{6}{\volt}$ e $i_L(0^+)=\SI{0.3}{\ampere}$}


\end{questao}
\begin{questao}[3]{}

Para um circuito RLC em série alimentado por um degrau de $\SI{10}{\volt}$, use
$R=\SI{20}{\ohm}$, $L=\SI{0.1}{\henry}$ e $C=\SI{100}{\micro\farad}$. Escreva
a equação diferencial para $v_C(t)$ após a comutação.
\resposta{$\ddot v_C+200\dot v_C+100000v_C=1000000$}


\end{questao}
\begin{questao}[3]{}

No circuito da questão anterior, admita energia inicial nula. Determine
$v_C(0^+)$, $\dot v_C(0^+)$ e $\ddot v_C(0^+)$.
\resposta{$v_C(0^+)=0$; $\dot v_C(0^+)=0$; $\ddot v_C(0^+)=\SI{1e6}{\volt\per\second\squared}$}


\end{questao}

\gabarito[Nível 3]

\blocoaprofundamento
\begin{questao}[3]{}

Um capacitor de $\SI{40}{\micro\farad}$ permaneceu ligado por muito tempo ao
circuito (a). Em $t=0$, a chave transfere instantaneamente o capacitor para o
circuito (b). Determine $v_C(t)$ para $t>0$, o instante em que a tensão cruza
zero e a corrente no resistor $R_2$ imediatamente após a comutação, tomando como
positivo o sentido do capacitor para a fonte de $-\SI{5}{\volt}$.

\circuitoqq{
\begin{circuitikz}[scale=0.92,transform shape,line width=0.8pt]
\draw (0,0) to[american voltage source,l_={$\SI{15}{\volt}$}] (0,3)
      to[R,l={$R_1=\SI{5}{\kilo\ohm}$}] (3,3)
      to[C,l={$C=\SI{40}{\micro\farad}$},v^>={$v_C$}] (3,0) -- (0,0);
\node[font=\sffamily\bfseries] at (1.5,-0.6) {(a) $t<0$};
\begin{scope}[xshift=7cm]
\draw (0,0) to[american voltage source,l_={$-\SI{5}{\volt}$}] (0,3)
      to[R,l={$R_2=\SI{10}{\kilo\ohm}$}] (3,3)
      to[C,l={$C=\SI{40}{\micro\farad}$},v^>={$v_C$}] (3,0) -- (0,0);
\node[font=\sffamily\bfseries] at (1.5,-0.6) {(b) $t>0$};
\end{scope}
\end{circuitikz}}{Circuitos equivalentes antes e depois da comutação.}

\resposta{$v_C(t)=-5+20e^{-t/0{,}4}\ \si{\volt}$;
$t_{v=0}\approx\SI{0.555}{\second}$; $i_{R_2}(0^+)=\SI{2}{\milli\ampere}$}


\end{questao}
\begin{questao}[3]{}

A chave permaneceu fechada por muito tempo e abre em $t=0$. O indutor passa a
liberar sua energia pela associação $R_2\parallel R_3$. Determine a corrente
$i_L(t)$ para $t>0$ e a energia ainda armazenada no indutor em
$t=\SI{50}{\milli\second}$.

\circuitoqq{
\begin{circuitikz}[scale=0.95,transform shape,line width=0.8pt]
\draw (0,0) to[american voltage source,l_={$\SI{30}{\volt}$}] (0,3)
      to[switch,l={abre em $t=0$}] (1.6,3)
      to[R,l={$R_1=\SI{5}{\ohm}$}] (3.6,3) coordinate(a);
\draw (a) to[L,l={$L=\SI{200}{\milli\henry}$},i>^={$i_L$}] (3.6,0) -- (0,0);
\draw (a) -- (5.6,3) to[R,l={$R_2=\SI{8}{\ohm}$}] (5.6,0) -- (3.6,0);
\draw (a) -- (7.6,3) to[R,l={$R_3=\SI{24}{\ohm}$}] (7.6,0) -- (5.6,0);
\end{circuitikz}}{Descarga de um indutor por uma rede resistiva equivalente.}

\resposta{$i_L(t)=6e^{-t/0{,}0333}\ \si{\ampere}$;
$W_L(50\,\mathrm{ms})\approx\SI{0.179}{\joule}$}


\end{questao}
\begin{questao}[3]{}

O capacitor possui $v_C(0^-)=\SI{10}{\volt}$. Em $t=0$ a chave fecha e o conecta
à rede mostrada. Determine $v_C(t)$, a resistência equivalente vista pelo
capacitor e o tempo necessário para que o erro em relação ao valor final seja
menor que \SI{1}{\percent} da variação inicial.

\circuitoqq{
\begin{circuitikz}[scale=0.95,transform shape,line width=0.8pt]
\draw (0,0) to[american voltage source,l_={$\SI{18}{\volt}$}] (0,3)
      to[R,l={$R_1=\SI{6}{\kilo\ohm}$}] (3,3) coordinate(b);
\draw (b) to[R,l={$R_2=\SI{3}{\kilo\ohm}$}] (3,0) -- (0,0);
\draw (b) to[R,l={$R_3=\SI{2}{\kilo\ohm}$}] (5.3,3)
      to[switch,l={fecha em $t=0$}] (6.7,3)
      to[C,l={$C=\SI{25}{\micro\farad}$},v^>={$v_C$}] (6.7,0) -- (3,0);
\end{circuitikz}}{Capacitor conectado a uma rede de Thévenin não trivial.}

\resposta{$R_{\mathrm{eq}}=\SI{4}{\kilo\ohm}$; $\tau=\SI{0.10}{\second}$;
$v_C(t)=6+4e^{-10t}\ \si{\volt}$; $t_{1\%}\approx\SI{0.461}{\second}$}


\end{questao}
\begin{questao}[3]{}

Para $t<0$, a fonte de corrente de $\SI{2}{\ampere}$ permaneceu conectada por
muito tempo ao RLC em paralelo. Em $t=0$ a chave abre e a fonte é removida. Determine
o tipo de resposta natural e obtenha $v(t)$ para $t>0$, considerando a polaridade
indicada.

\circuitoqq{
\begin{circuitikz}[scale=1,transform shape,line width=0.8pt]
\draw (0,0) to[american current source,l_={$\SI{2}{\ampere}$}] (0,3)
      to[switch,l={abre em $t=0$}] (1.7,3) -- (6.5,3);
\draw (2.5,3) to[R,l={$R=\SI{10}{\ohm}$}] (2.5,0);
\draw (4.5,3) to[C,l={$C=\SI{200}{\micro\farad}$},v^>={$v$}] (4.5,0);
\draw (6.5,3) to[L,l={$L=\SI{50}{\milli\henry}$},i>^={$i_L$}] (6.5,0) -- (0,0);
\end{circuitikz}}{Resposta natural de um circuito RLC em paralelo com energia inicialmente no indutor.}

\resposta{Subamortecida;
$v(t)\approx-51{,}64e^{-250t}\sin(193{,}65t)\ \si{\volt}$}


\end{questao}
\begin{questao}[3]{}

O circuito RLC em série parte sem energia armazenada e a chave fecha em $t=0$.
Determine a resposta completa $v_C(t)$ e o valor do primeiro pico da tensão do
capacitor.

\circuitoqq{
\begin{circuitikz}[scale=1,transform shape,line width=0.8pt]
\draw (0,0) to[american voltage source,l_={$\SI{10}{\volt}$}] (0,3)
      to[switch,l={fecha em $t=0$}] (1.5,3)
      to[R,l={$R=\SI{20}{\ohm}$}] (3.4,3)
      to[L,l={$L=\SI{100}{\milli\henry}$}] (5.5,3)
      to[C,l={$C=\SI{100}{\micro\farad}$},v^>={$v_C$}] (7.5,3)
      -- (7.5,0) -- (0,0);
\end{circuitikz}}{Resposta ao degrau de um circuito RLC em série subamortecido.}

\resposta{$v_C(t)=10\left[1-e^{-100t}\left(\cos300t+\frac13\sin300t\right)\right]\,\si{\volt}$;
$v_{C,\mathrm{pico}}\approx\SI{13.51}{\volt}$ em $t_p\approx\SI{10.47}{\milli\second}$}


\end{questao}
\begin{questao}[3]{}

Em laboratório, mediu-se a corrente natural do RLC em série mostrado. A distância
entre picos consecutivos é aproximadamente $T_d=\SI{8}{\milli\second}$ e a
envoltória cai à metade a cada $\SI{16}{\milli\second}$. Sabendo que
$R=\SI{8}{\ohm}$, estime $\alpha$, $\omega_d$, $L$ e $C$.

\circuitoqq{
\begin{circuitikz}[scale=1,transform shape,line width=0.8pt]
\draw (0,0) to[R,l={$R=\SI{8}{\ohm}$}] (2,0)
      to[L,l={$L$}] (4,0)
      to[C,l={$C$}] (6,0) -- (6,-2) -- (0,-2) -- (0,0);
\draw[->] (1.2,0.45) -- (2.3,0.45) node[midway,above] {$i(t)$};
\end{circuitikz}}{RLC em série identificado a partir de uma oscilação medida.}

\begin{figure}[H]
\centering
\begin{tikzpicture}
\begin{axis}[
 width=.88\linewidth,height=5.2cm,
 xlabel={$t$ (ms)},ylabel={$i(t)$ (u.a.)},xmin=0,xmax=32,ymin=-2.2,ymax=2.2,
 grid=both,samples=300,
 tick label style={font=\scriptsize},label style={font=\small\sffamily}]
\addplot[thick,azulprof,domain=0:32] {2*exp(-0.0433217*x)*sin(deg(0.785398*x))};
\addplot[dashed,laranja,domain=0:32] {2*exp(-0.0433217*x)};
\addplot[dashed,laranja,domain=0:32] {-2*exp(-0.0433217*x)};
\end{axis}
\end{tikzpicture}
\caption{Oscilação amortecida medida e envoltória exponencial.}
\end{figure}

\resposta{$\alpha\approx\SI{43.3}{\per\second}$;
$\omega_d\approx\SI{785.4}{\radian\per\second}$;
$L\approx\SI{92.3}{\milli\henry}$; $C\approx\SI{17.5}{\micro\farad}$}


\end{questao}

\gabarito[Nível 3 --- análise integrada e formas de onda]

\blocoaprofundamento
\begin{questao}[3]{}

O capacitor abaixo possui $v_C(0^+)=\SI{8}{\volt}$. A fonte dependente é uma VCVS
que impõe $v_b=\mu v_C$, com $\mu=0{,}5$. Determine a resistência equivalente vista
pelo capacitor, a constante de tempo e $v_C(t)$.

\circuitoqq{
\begin{circuitikz}[scale=0.96,transform shape,line width=0.8pt,every node/.append style={font=\scriptsize}]
\draw (0,0) -- (7.0,0);
\draw (0,3) to[C,v^>={$v_C$}] (0,0);
\node[anchor=west] at (0.25,1.5) {$C=\SI{25}{\micro\farad}$};
\draw (0,3) to[R] (3.8,3) coordinate(b);
\node[anchor=south] at (1.9,3.25) {$R=\SI{4}{\kilo\ohm}$};
\draw (b) to[american controlled voltage source] (3.8,0);
\node[anchor=west] at (4.05,1.5) {$\mu\,v_C$};
\node[anchor=south west] at ($(b)+(0.15,0.15)$) {$v_b$};
\end{circuitikz}}{Circuito RC com fonte de tensão controlada por tensão (VCVS).}

\resposta{$R_{\mathrm{eq}}=\SI{8}{\kilo\ohm}$; $\tau=\SI{0.20}{\second}$;
$v_C(t)=8e^{-5t}\ \si{\volt}$}


\end{questao}
\begin{questao}[3]{}

O indutor do circuito abaixo conduz inicialmente $\SI{3}{\ampere}$. A fonte
dependente é uma CCVS de valor $r_m i_x$, com $r_m=\SI{9}{\ohm}$, e sua polaridade
é indicada no desenho. Determine a resistência equivalente vista pelo indutor,
$\tau$ e $i_L(t)$.

\circuitoqq{
\begin{circuitikz}[scale=0.96,transform shape,line width=0.8pt,every node/.append style={font=\scriptsize}]
\draw (0,0) -- (7.0,0);
\draw (0,3) to[L,i>^={$i_L$}] (0,0);
\node[anchor=west] at (0.25,1.5) {$L=\SI{75}{\milli\henry}$};
\draw (0,3) to[R,i>^={$i_x$}] (3.7,3)
      to[american controlled voltage source] (3.7,0);
\node[anchor=south] at (1.85,3.25) {$R=\SI{6}{\ohm}$};
\node[anchor=west] at (3.95,1.5) {$r_m\,i_x$};
\end{circuitikz}}{Resposta natural RL com fonte de tensão controlada por corrente (CCVS).}

\resposta{$R_{\mathrm{eq}}=\SI{15}{\ohm}$; $\tau=\SI{5}{\milli\second}$;
$i_L(t)=3e^{-200t}\ \si{\ampere}$}


\end{questao}
\begin{questao}[3]{}

O RLC em paralelo abaixo não possui fonte independente para $t>0$. A VCCS fornece
uma corrente $g v$, com $g=\SI{50}{\milli\ensuremath{\mathrm{S}}}$. Considere
$v(0^+)=\SI{5}{\volt}$ e $i_L(0^+)=0$. Determine a equação diferencial da resposta
natural, classifique o amortecimento e obtenha $v(t)$.

\circuitoqq{
\begin{circuitikz}[scale=0.94,transform shape,line width=0.8pt,every node/.append style={font=\scriptsize}]
\draw (0,0) -- (8.6,0);
\draw (0,3) -- (8.6,3);
\draw (1.5,3) to[R] (1.5,0);
\node[anchor=east] at (1.25,1.5) {$R=\SI{20}{\ohm}$};
\draw (3.9,3) to[C,v^>={$v$}] (3.9,0);
\node[anchor=east] at (3.65,1.5) {$C=\SI{10}{\milli\farad}$};
\draw (6.2,3) to[L,i>^={$i_L$}] (6.2,0);
\node[anchor=east] at (5.95,1.5) {$L=\SI{1}{\henry}$};
\draw (8.6,3) to[american controlled current source] (8.6,0);
\node[anchor=west] at (8.85,1.5) {$g\,v$};
\end{circuitikz}}{RLC em paralelo com VCCS modificando o amortecimento.}

\resposta{$\ddot v+10\dot v+100v=0$; subamortecida;
$v(t)=e^{-5t}[5\cos(8{,}66t)-2{,}89\sin(8{,}66t)]\ \si{\volt}$}


\end{questao}
\begin{questao}[3]{}

No circuito RLC em série abaixo, a CCVS produz uma tensão $r_m i$ que se soma à
queda resistiva, com $r_m=\SI{30}{\ohm}$. Para
$R=\SI{10}{\ohm}$, $L=\SI{0.1}{\henry}$, $C=\SI{100}{\micro\farad}$,
$v_C(0^+)=\SI{10}{\volt}$ e $i(0^+)=0$, determine $v_C(t)$.

\circuitoqq{
\begin{circuitikz}[scale=0.90,transform shape,line width=0.8pt,every node/.append style={font=\scriptsize}]
\draw (0,0) to[R] (2.2,0)
      to[american controlled voltage source] (4.6,0)
      to[L,i>^={$i$}] (7.0,0)
      to[C,v^>={$v_C$}] (9.5,0)
      -- (9.5,-2.2) -- (0,-2.2) -- (0,0);
\node[anchor=south] at (1.1,0.25) {$R=\SI{10}{\ohm}$};
\node[anchor=south] at (3.4,0.25) {$r_m\,i$};
\node[anchor=south] at (5.8,0.25) {$L=\SI{0.1}{\henry}$};
\node[anchor=south] at (8.25,0.25) {$C=\SI{100}{\micro\farad}$};
\end{circuitikz}}{RLC em série com CCVS equivalente a amortecimento adicional.}

\resposta{$v_C(t)=e^{-200t}[10\cos(244{,}95t)+8{,}17\sin(244{,}95t)]\ \si{\volt}$}


\end{questao}

\gabarito[Nível 3 --- fontes dependentes]

\blocodesafio
\begin{questao}[4]{}

Em laboratório, um circuito RC com $R=\SI{5}{\kilo\ohm}$ e
$C=\SI{10}{\micro\farad}$ deveria apresentar $\tau=\SI{50}{\milli\second}$.
Entretanto, o ajuste exponencial do osciloscópio fornece
$\tau_{\mathrm{med}}=\SI{80}{\milli\second}$. A capacitância medida está
próxima do valor nominal. Estime a resistência efetiva vista pelo capacitor e proponha
uma causa elétrica plausível para a discrepância.
\resposta{$R_{\mathrm{eq}}\approx\SI{8}{\kilo\ohm}$; há cerca de \SI{3}{\kilo\ohm}
adicionais na trajetória equivalente}


\end{questao}
\begin{questao}[4]{}

Um capacitor de $\SI{470}{\micro\farad}$ será carregado por uma fonte de
$\SI{24}{\volt}$. Deseja-se limitar a corrente inicial a
$\SI{100}{\milli\ampere}$ e atingir pelo menos \SI{95}{\percent} da tensão final
em menos de $\SI{0.5}{\second}$. Projete o resistor mínimo em série, verifique o
tempo de carga e estime a energia total dissipada nele durante uma carga completa.
\resposta{$R=\SI{240}{\ohm}$; $t_{95}\approx\SI{0.338}{\second}$; $E_R\approx\SI{0.135}{\joule}$}


\end{questao}
\begin{questao}[4]{}

Uma bobina de relé tem $L=\SI{120}{\milli\henry}$ e resistência de cobre
$\SI{60}{\ohm}$, alimentada em $\SI{24}{\volt}$ até o regime permanente.
Compare duas estratégias de desligamento: (A) diodo de roda livre ideal; (B)
resistor de clamp de $\SI{180}{\ohm}$ no caminho de descarga. Determine a energia
inicial, as constantes de tempo aproximadas e a tensão inicial no resistor de
clamp. Explique o compromisso de projeto.
\resposta{$I_0=\SI{0.4}{\ampere}$; $E_L=\SI{9.6}{\milli\joule}$; $\tau_A\approx\SI{2}{\milli\second}$;
$\tau_B\approx\SI{0.5}{\milli\second}$; $v_{R,\mathrm{clamp}}(0^+)\approx\SI{72}{\volt}$}


\end{questao}
\begin{questao}[4]{}

Uma oscilação amortecida medida em um circuito RLC apresenta período entre picos
de aproximadamente $\SI{4}{\milli\second}$ e a envoltória cai à metade a cada
$\SI{10}{\milli\second}$. O capacitor é $C=\SI{10}{\micro\farad}$. Estime
$\alpha$, $\omega_d$, $\omega_0$, $\zeta$, $L$ e a resistência equivalente em série.

\begin{figure}[H]
\centering
\begin{tikzpicture}
\begin{axis}[
 width=0.88\linewidth,height=5.2cm,
 xlabel={$t$ (ms)},ylabel={$v(t)$},xmin=0,xmax=30,
 axis lines=middle,grid=both,samples=300,
 tick label style={font=\scriptsize},label style={font=\small\sffamily}]
\addplot[thick,azulprof,domain=0:30] {10*exp(-0.0693147*x)*cos(deg(1.5707963*x))};
\addplot[dashed,laranja,domain=0:30] {10*exp(-0.0693147*x)};
\addplot[dashed,laranja,domain=0:30] {-10*exp(-0.0693147*x)};
\end{axis}
\end{tikzpicture}
\caption{Exemplo de oscilação amortecida e sua envoltória exponencial.}
\end{figure}

\resposta{$\alpha\approx\SI{69.3}{\per\second}$; $\omega_d\approx\SI{1571}{\radian\per\second}$;
$\omega_0\approx\SI{1572}{\radian\per\second}$; $\zeta\approx0{,}044$;
$L\approx\SI{40.45}{\milli\henry}$; $R\approx\SI{5.61}{\ohm}$}


\end{questao}
\begin{questao}[4]{}

Um estágio RLC em série usa $L=\SI{20}{\milli\henry}$ e
$C=\SI{10}{\micro\farad}$. O requisito é obter a resposta mais rápida possível
sem oscilação. Determine o resistor ideal e compare qualitativamente o uso de
$\SI{45}{\ohm}$ e $\SI{180}{\ohm}$.
\resposta{$R_{\mathrm{crit}}\approx\SI{89.4}{\ohm}$; \SI{45}{\ohm} é subamortecido;
\SI{180}{\ohm} é superamortecido}


\end{questao}
\begin{questao}[4]{}

A tabela mostra a resposta de um RC inicialmente descarregado cuja tensão final é
$\SI{10}{\volt}$.

\begin{table}[H]
\centering\small\sffamily
\begin{tabular}{@{}c c@{}}
\toprule
$t$ & $v_C(t)$\\
\midrule
\SI{20}{\milli\second} & \SI{6.32}{\volt}\\
\SI{40}{\milli\second} & \SI{8.65}{\volt}\\
\bottomrule
\end{tabular}
\end{table}

Verifique se os dados são compatíveis com uma única constante de tempo e estime
$R$ se $C=\SI{4.7}{\micro\farad}$.
\resposta{$\tau\approx\SI{20}{\milli\second}$; $R\approx\SI{4.26}{\kilo\ohm}$}


\end{questao}
\begin{questao}[4]{}

Após uma comutação, um circuito RLC em série fica sem fonte independente e possui
$R=\SI{20}{\ohm}$, $L=\SI{50}{\milli\henry}$ e
$C=\SI{100}{\micro\farad}$. Adote $v_C(0^+)=\SI{5}{\volt}$ e
$i(0^+)=\SI{0.1}{\ampere}$ entrando no terminal positivo do capacitor. Determine
a expressão de $v_C(t)$.
\resposta{$v_C(t)=e^{-200t}[5\cos(400t)+5\sin(400t)]\ \si{\volt}$}


\end{questao}
\begin{questao}[4]{}

Projete uma rede de temporização para uma entrada de $\SI{3.3}{\volt}$ e
$C=\SI{1}{\micro\farad}$. Na carga, a tensão deve alcançar
$\SI{2}{\volt}$ em aproximadamente $\SI{8}{\milli\second}$. No desligamento,
deve cair de $\SI{3.3}{\volt}$ para menos de $\SI{0.5}{\volt}$ em
$\SI{5}{\milli\second}$. Mostre que um único resistor não atende bem às duas
metas e proponha resistores independentes de carga e descarga usando um diodo.
\resposta{$R_{\mathrm{carga}}\approx\SI{8.59}{\kilo\ohm}$; $R_{\mathrm{descarga}}\lesssim\SI{2.65}{\kilo\ohm}$;
valores práticos: \SI{8.2}{\kilo\ohm} e \SI{2.4}{\kilo\ohm}}


\end{questao}
\begin{questao}[4]{}

Para $L=\SI{20}{\milli\henry}$ e $C=\SI{10}{\micro\farad}$, compare os três
resistores da tabela e escolha o mais adequado quando se deseja resposta rápida
sem sobressinal.

\begin{table}[H]
\centering\small\sffamily
\begin{tabular}{@{}c c c@{}}
\toprule
\textbf{Opção} & $R$ & \textbf{Observação a determinar}\\
\midrule
A & \SI{40}{\ohm} & ?\\
B & \SI{89.4}{\ohm} & ?\\
C & \SI{180}{\ohm} & ?\\
\bottomrule
\end{tabular}
\end{table}

\resposta{A: $\zeta\approx0{,}447$, subamortecido; B: crítico; C: $\zeta\approx2{,}01$, superamortecido; escolher B}


\end{questao}
\begin{questao}[4]{}

Um projeto prevê $R=\SI{20}{\ohm}$, $L=\SI{0.1}{\henry}$ e
$C=\SI{100}{\micro\farad}$, portanto deveria ser subamortecido. Contudo, o ajuste
da resposta medida indica dois polos reais em aproximadamente
$-\SI{112.7}{\per\second}$ e $-\SI{887.3}{\per\second}$. Admitindo $L$ e $C$
corretos, estime a resistência efetiva em série e diagnostique a discrepância.
\resposta{$R_{\mathrm{efetivo}}\approx\SI{100}{\ohm}$; há cerca de \SI{80}{\ohm}
de resistência em série não contabilizada}


\end{questao}
\begin{questao}[4]{}

Um barramento CC de $\SI{400}{\volt}$ possui capacitor de
$\SI{470}{\micro\farad}$. Por segurança, após desligar a fonte deseja-se reduzir
a tensão para no máximo $\SI{50}{\volt}$ em $\SI{60}{\second}$. Dimensione o
maior resistor de descarga admissível, escolha um valor comercial abaixo dele,
estime a potência inicial e a energia que será dissipada.
\resposta{$R_{\max}\approx\SI{61.4}{\kilo\ohm}$; escolha \SI{56}{\kilo\ohm};
$t\approx\SI{54.7}{\second}$; $P_0\approx\SI{2.86}{\watt}$; $E\approx\SI{37.6}{\joule}$}


\end{questao}

\gabarito[Nível 4]

\blocodesafio
\begin{questao}[4]{}

Um sensor pode ser modelado por uma fonte de Thévenin de $\SI{5}{\volt}$ com
resistência interna $R_s=\SI{1}{\kilo\ohm}$. Ao fechar a chave em $t=0$, ele
alimenta uma entrada com $R_L=\SI{9}{\kilo\ohm}$ em paralelo com um capacitor
$C$, inicialmente descarregado. O sistema digital reconhece nível alto quando
$v_C\ge\SI{3}{\volt}$ e exige que isso ocorra em no máximo
$\SI{20}{\milli\second}$.

\circuitoqq{
\begin{circuitikz}[scale=1,transform shape,line width=0.8pt]
\draw (0,0) to[american voltage source,l_={$\SI{5}{\volt}$}] (0,3)
      to[R,l={$R_s=\SI{1}{\kilo\ohm}$}] (2.3,3)
      to[switch,l={fecha em $t=0$}] (3.8,3) -- (6.5,3) coordinate(a);
\draw (a) to[R,l={$R_L=\SI{9}{\kilo\ohm}$}] (6.5,0) -- (0,0);
\draw (a) -- (8.5,3) to[C,l={$C$},v^>={$v_C$}] (8.5,0) -- (6.5,0);
\end{circuitikz}}{Entrada RC carregada por uma fonte com resistência interna e carga finita.}

\begin{itensq}
\item Determine o valor final de $v_C$ e a resistência de Thévenin vista pelo capacitor.
\item Calcule o maior valor de $C$ que atende ao tempo especificado.
\item Escolha entre $\SI{18}{\micro\farad}$ e $\SI{22}{\micro\farad}$ e verifique o tempo de cruzamento de \SI{3}{\volt}.
\item Determine a corrente fornecida pela fonte em $0^+$ e em regime permanente.
\end{itensq}

\resposta{$V_f=\SI{4.5}{\volt}$; $R_{\mathrm{Th}}=\SI{900}{\ohm}$;
$C_{\max}\approx\SI{20.2}{\micro\farad}$; escolha: \SI{18}{\micro\farad};
$t_{3V}\approx\SI{17.8}{\milli\second}$;
$i_s(0^+)=\SI{5}{\milli\ampere}$; $i_s(\infty)=\SI{0.5}{\milli\ampere}$}


\end{questao}
\begin{questao}[4]{}

Um barramento CC possui um capacitor de $\SI{470}{\micro\farad}$ carregado a
$\SI{325}{\volt}$. Quando a alimentação é desconectada, um resistor de descarga
$R_b$ deve reduzir a tensão para no máximo $\SI{50}{\volt}$ em
$\SI{60}{\second}$. Despreze outras cargas.

\circuitoqq{
\begin{circuitikz}[scale=1,transform shape,line width=0.8pt]
\draw (0,0) to[american voltage source,l_={$\SI{325}{\volt}$}] (0,3)
      to[switch,l={desliga em $t=0$}] (1.8,3) -- (5,3) coordinate(a);
\draw (a) to[C,l={$C=\SI{470}{\micro\farad}$},v^>={$v_C$}] (5,0) -- (0,0);
\draw (a) -- (7.2,3) to[R,l={$R_b$}] (7.2,0) -- (5,0);
\end{circuitikz}}{Resistor de descarga de um barramento capacitivo.}

\begin{itensq}
\item Determine o maior valor permitido de $R_b$.
\item Escolha um valor comercial de $\SI{68}{\kilo\ohm}$ e verifique o tempo até \SI{50}{\volt}.
\item Calcule a potência instantânea inicial no resistor e a energia total que ele dissipará.
\end{itensq}

\resposta{$R_{b,\max}\approx\SI{68.2}{\kilo\ohm}$; com \SI{68}{\kilo\ohm},
$t_{50V}\approx\SI{59.8}{\second}$; $p_R(0^+)\approx\SI{1.55}{\watt}$;
$E\approx\SI{24.8}{\joule}$}


\end{questao}
\begin{questao}[4]{}

Projete o RLC em série abaixo para que a resposta do capacitor a um degrau seja
subamortecida com $\zeta=0{,}5$ e tempo de acomodação de \SI{2}{\percent} não
superior a $\SI{40}{\milli\second}$. Adote
$C=\SI{100}{\micro\farad}$ e use a aproximação
$t_s\approx4/(\zeta\omega_0)$. Escolha o menor $\omega_0$ que atende à
especificação. Determine $L$, $R$, o sobressinal percentual e o instante do
primeiro pico.

\circuitoqq{
\begin{circuitikz}[scale=1,transform shape,line width=0.8pt]
\draw (0,0) to[american voltage source,l_={$V_su(t)$}] (0,3)
      to[switch,l={fecha em $t=0$}] (1.5,3)
      to[R,l={$R$}] (3.3,3)
      to[L,l={$L$}] (5.3,3)
      to[C,l={$C=\SI{100}{\micro\farad}$},v^>={$v_C$}] (7.5,3)
      -- (7.5,0) -- (0,0);
\end{circuitikz}}{Projeto de um circuito RLC em série a partir de especificações transitórias.}

\resposta{$\omega_0=\SI{200}{\radian\per\second}$;
$L=\SI{0.25}{\henry}$; $R=\SI{50}{\ohm}$;
$M_p\approx\SI{16.3}{\percent}$; $t_p\approx\SI{18.1}{\milli\second}$}


\end{questao}
\begin{questao}[4]{}

O circuito RC em escada abaixo possui \textbf{dois capacitores independentes} e,
portanto, é de segunda ordem. Para $v_s(t)=10u(t)\,\si{\volt}$ e condições
iniciais nulas, obtenha as equações de estado em $v_1$ e $v_2$, determine os polos
naturais e encontre $v_2(t)$.

\circuitoqq{
\begin{circuitikz}[scale=1,transform shape,line width=0.8pt]
\draw (0,0) to[american voltage source,l_={$10u(t)\,\si{\volt}$}] (0,3)
      to[R,l={$R_1=\SI{1}{\kilo\ohm}$}] (3,3) coordinate(a)
      to[R,l={$R_2=\SI{2}{\kilo\ohm}$}] (6,3) coordinate(b);
\draw (a) to[C,l={$C_1=\SI{100}{\micro\farad}$},v^>={$v_1$}] (3,0);
\draw (b) to[C,l={$C_2=\SI{50}{\micro\farad}$},v^>={$v_2$}] (6,0) -- (0,0);
\end{circuitikz}}{Rede RC de segunda ordem com dois elementos armazenadores do mesmo tipo.}

\resposta{$\dot v_1=100u(t)-15v_1+5v_2$;
$\dot v_2=10v_1-10v_2$; polos $-5$ e $-20\ \si{\per\second}$;
$v_2(t)=10-\frac{40}{3}e^{-5t}+\frac{10}{3}e^{-20t}\ \si{\volt}$}


\end{questao}
\begin{questao}[4]{}

No circuito RLC em paralelo, a fonte de corrente vale $\SI{2}{\ampere}$ por muito
tempo e muda instantaneamente para $\SI{1}{\ampere}$ em $t=0$. Determine a
resposta completa $v(t)$, classifique o amortecimento e encontre o primeiro pico
negativo da tensão.

\circuitoqq{
\begin{circuitikz}[scale=1,transform shape,line width=0.8pt]
\draw (0,0) to[american current source,l_={$I_s:\ 2\,\mathrm{A}\to1\,\mathrm{A}$}] (0,3) -- (6.5,3);
\draw (2.2,3) to[R,l={$R=\SI{20}{\ohm}$}] (2.2,0);
\draw (4.3,3) to[C,l={$C=\SI{500}{\micro\farad}$},v^>={$v$}] (4.3,0);
\draw (6.5,3) to[L,l={$L=\SI{0.2}{\henry}$},i>^={$i_L$}] (6.5,0) -- (0,0);
\end{circuitikz}}{RLC em paralelo excitado por uma mudança de nível da fonte de corrente.}

\resposta{Subamortecida;
$v(t)\approx-23{,}09e^{-50t}\sin(86{,}60t)\ \si{\volt}$;
$t_p\approx\SI{12.1}{\milli\second}$;
$v_{\min}\approx-\SI{10.93}{\volt}$}


\end{questao}

\gabarito[Nível 4 --- projeto, diagnóstico e síntese]

\blocodesafio
\begin{questao}[4]{}

O circuito RC natural abaixo contém uma VCCS orientada para cima, de valor $g v_C$.
O capacitor está inicialmente carregado. Determine o valor crítico de $g$ que
separa decaimento, resposta constante e crescimento exponencial. Em seguida,
classifique os casos $g=\SI{0.4}{\milli\ensuremath{\mathrm{S}}}$,
$\SI{0.5}{\milli\ensuremath{\mathrm{S}}}$ e $\SI{0.6}{\milli\ensuremath{\mathrm{S}}}$ para
$R=\SI{2}{\kilo\ohm}$ e $C=\SI{100}{\micro\farad}$.

\circuitoqq{
\begin{circuitikz}[scale=1,transform shape,line width=0.8pt]
\draw (0,0) -- (6,0);
\draw (0,3) -- (6,3);
\draw (1.8,3) to[R,l={$R=\SI{2}{\kilo\ohm}$}] (1.8,0);
\draw (3.8,3) to[C,l={$C=\SI{100}{\micro\farad}$},v^>={$v_C$}] (3.8,0);
\draw (6,0) to[american controlled current source,l_={$g v_C$}] (6,3);
\end{circuitikz}}{RC com VCCS capaz de reduzir, anular ou inverter o amortecimento.}

\resposta{$g_{\mathrm{crit}}=1/R=\SI{0.5}{\milli\ensuremath{\mathrm{S}}}$; para
\SI{0.4}{\milli\ensuremath{\mathrm{S}}}: estável, $\tau=\SI{1}{\second}$; para
\SI{0.5}{\milli\ensuremath{\mathrm{S}}}: polo em zero; para \SI{0.6}{\milli\ensuremath{\mathrm{S}}}:
instável, crescimento com constante de tempo de magnitude \SI{1}{\second}}


\end{questao}
\begin{questao}[4]{}

Deseja-se usar uma VCCS $g v$ para ajustar o amortecimento do RLC em paralelo
abaixo. Determine o valor de $g$ para obter amortecimento crítico. Depois indique
se a resposta será subamortecida ou superamortecida para valores de $g$ menores ou
maiores que o valor calculado, supondo $g\ge0$.

\circuitoqq{
\begin{circuitikz}[scale=1,transform shape,line width=0.8pt]
\draw (0,0) -- (8,0);
\draw (0,3) -- (8,3);
\draw (1.5,3) to[R,l={$R=\SI{100}{\ohm}$}] (1.5,0);
\draw (3.7,3) to[C,l={$C=\SI{100}{\micro\farad}$},v^>={$v$}] (3.7,0);
\draw (5.8,3) to[L,l={$L=\SI{100}{\milli\henry}$}] (5.8,0);
\draw (8,3) to[american controlled current source,l_={$g v$}] (8,0);
\end{circuitikz}}{Projeto do amortecimento de um circuito RLC em paralelo usando uma VCCS.}

\resposta{$g_{\mathrm{crit}}\approx\SI{53.25}{\milli\ensuremath{\mathrm{S}}}$; abaixo desse
valor a resposta é subamortecida; acima, superamortecida}


\end{questao}
\begin{questao}[4]{}

No RLC em série abaixo, a VCVS tem valor $\mu v_C$ e polaridade tal que sua tensão
\textbf{auxilia} a tensão do capacitor na equação de malha, resultando no termo
restaurador $(1-\mu)v_C$. Para
$R=\SI{20}{\ohm}$, $L=\SI{0.1}{\henry}$ e
$C=\SI{100}{\micro\farad}$, derive a equação característica e determine os
intervalos de $\mu$ que produzem resposta subamortecida, criticamente amortecida,
superamortecida estável, polo em zero e instabilidade.

\circuitoqq{
\begin{circuitikz}[scale=1,transform shape,line width=0.8pt]
\draw (0,0) to[R,l={$R=\SI{20}{\ohm}$}] (2,0)
      to[L,l={$L=\SI{0.1}{\henry}$},i>^={$i$}] (4.3,0)
      to[C,l={$C=\SI{100}{\micro\farad}$},v^>={$v_C$}] (6.5,0)
      to[american controlled voltage source,l_={$\mu v_C$}] (8.5,0)
      -- (8.5,-2) -- (0,-2) -- (0,0);
\end{circuitikz}}{RLC com VCVS modificando o termo restaurador e a estabilidade.}

\resposta{$s^2+200s+100000(1-\mu)=0$; subamortecida para $\mu<0{,}9$;
crítica em $\mu=0{,}9$; superamortecida estável para $0{,}9<\mu<1$;
polo em zero para $\mu=1$; instável para $\mu>1$}


\end{questao}
\begin{questao}[4]{}

A rede abaixo possui dois capacitores e uma CCCS que injeta no nó $v_2$ a corrente
$\beta i_x$, em que $i_x=v_1/R_1$. Considere
$R_1=R_2=\SI{1}{\kilo\ohm}$ e $C_1=C_2=\SI{1}{\milli\farad}$.

\circuitoqq{
\begin{circuitikz}[scale=1,transform shape,line width=0.8pt]
\draw (0,0) -- (7.5,0);
\draw (0,3) coordinate(a) to[R,l={$R_2=\SI{1}{\kilo\ohm}$}] (4,3) coordinate(b);
\draw (a) to[R,l={$R_1=\SI{1}{\kilo\ohm}$},i>^={$i_x$}] (0,0);
\draw (a) -- (2,3) to[C,l={$C_1=\SI{1}{\milli\farad}$},v^>={$v_1$}] (2,0);
\draw (b) to[C,l={$C_2=\SI{1}{\milli\farad}$},v^>={$v_2$}] (4,0);
\draw (7.5,0) to[american controlled current source,l_={$\beta i_x$}] (7.5,3) -- (b);
\end{circuitikz}}{Rede RC de segunda ordem com CCCS e possibilidade de perda de estabilidade.}

\begin{itensq}
\item Obtenha as equações de estado em $v_1$ e $v_2$.
\item Determine a equação característica em função de $\beta$.
\item Encontre a faixa de $\beta$ que garante estabilidade assintótica.
\item Para $\beta=0{,}5$, $v_1(0)=\SI{1}{\volt}$ e $v_2(0)=0$, determine $v_2(t)$.
\end{itensq}

\resposta{$\dot v_1=-2v_1+v_2$;
$\dot v_2=(1+\beta)v_1-v_2$;
$s^2+3s+(1-\beta)=0$; estável para $\beta<1$;
para $\beta=0{,}5$,
$v_2(t)=\frac{1.5}{\sqrt7}\left(e^{(-3+\sqrt7)t/2}-e^{(-3-\sqrt7)t/2}\right)\ \si{\volt}$}


\end{questao}

\gabarito[Nível 4 --- fontes dependentes]

\section{Eletromagnetismo}

\subsection{Conceitos iniciais e materiais magnéticos}

\blocofixacao
\begin{questao}[1]{}

O princípio da inseparabilidade dos polos magnéticos garante que:
\begin{alternativas}
\item existem monopolos magnéticos isolados.
\item ao dividir um ímã, obtêm-se dois novos ímãs, cada um com dois polos.
\item o polo norte pode ser isolado do polo sul.
\item as linhas de campo magnético têm início e fim.
\item um ímã pode ter apenas um polo.
\end{alternativas}
\resposta{(b)}


\end{questao}
\begin{questao}[1]{}

Quanto à permeabilidade magnética relativa $\mu_r$, um material
\textbf{ferromagnético} caracteriza-se por:
\begin{alternativas}
\item $\mu_r \approx 1$.
\item $\mu_r$ ligeiramente menor que 1.
\item $\mu_r$ muito maior que 1 (centenas a milhares).
\item $\mu_r = 0$.
\item $\mu_r$ negativa.
\end{alternativas}
\resposta{(c)}


\end{questao}
\begin{questao}[1]{}

Dois materiais são submetidos ao mesmo campo externo $H$. Um apresenta
$B_1 = \SI{0.002}{\tesla}$ e o outro, $B_2 = \SI{1.8}{\tesla}$. Qual é o
ferromagnético?
\resposta{O de $B_2 = \SI{1.8}{\tesla}$}


\end{questao}
\begin{questao}[1]{Apostila original revisada}

Dois materiais apresentam permeabilidades relativas
$\mu_{rA}=1{,}0005$ e $\mu_{rB}=0{,}99998$. Classifique cada material como
paramagnético ou diamagnético e descreva, qualitativamente, como reage a um campo
magnético externo.
\resposta{$A$: paramagnético; $B$: diamagnético}


\end{questao}

\blocoaplicacao
\begin{questao}[2]{}

Explique por que os materiais ferromagnéticos são preferidos na construção de
núcleos de máquinas elétricas (transformadores, motores, geradores).
\resposta{Ver comentário}


\end{questao}

\blocoaprofundamento
\begin{questao}[3]{}

Um mesmo enrolamento é usado com dois núcleos diferentes: um de ar
($\mu_r \approx 1$) e outro de ferro-silício ($\mu_r = 4000$). Mantendo a mesma
corrente:
\begin{itensq}
\item Qual é a razão entre as induções magnéticas $B$ obtidas nos dois casos?
\item Se o núcleo de ferro \emph{saturar} a partir de certo ponto, o que acontece
      com essa razão ao se aumentar muito a corrente?
\item Por que a saturação impõe um limite de projeto às máquinas elétricas?
\end{itensq}
\resposta{(a) $B_{\text{Fe}}/B_{\text{ar}} = 4000$;
(b) a razão \emph{cai}; (c) ver comentário}


\end{questao}
\begin{questao}[3]{}

Um material ferromagnético submetido a um campo alternado percorre um
\textbf{ciclo de histerese}: a curva $B \times H$ não é uma reta e não retorna
pelo mesmo caminho.

\circuitoqq{
\begin{tikzpicture}
 \begin{axis}[width=8.2cm,height=5.6cm,xlabel={$H$ (A/m)},ylabel={$B$ (T)},
   xmin=-3,xmax=3,ymin=-1.6,ymax=1.6,xtick=\empty,ytick=\empty,
   axis lines=middle,axis line style={-Stealth}]
   \addplot[Icol,very thick,smooth] coordinates
     {(-2.6,-1.35)(-1.6,-1.25)(-0.6,-1.0)(0,-0.75)(0.5,-0.2)(1.0,0.75)(1.7,1.2)(2.6,1.35)};
   \addplot[Vcol,very thick,smooth] coordinates
     {(2.6,1.35)(1.6,1.25)(0.6,1.0)(0,0.75)(-0.5,0.2)(-1.0,-0.75)(-1.7,-1.2)(-2.6,-1.35)};
   \node[font=\footnotesize,Vcol] at (axis cs:0.42,0.95) {$B_r$};
   \node[font=\footnotesize,Icol] at (axis cs:-1.0,-0.35) {$H_c$};
 \end{axis}
\end{tikzpicture}}{Ciclo de histerese de um material ferromagnético.}

\begin{itensq}
\item Identifique o que representam $B_r$ (remanência) e $H_c$ (coercividade).
\item A \emph{área} interna do ciclo tem significado físico. Qual?
\item Compare os requisitos de um material para (i) núcleo de transformador e
      (ii) ímã permanente. Qual deve ter ciclo "largo" e qual "estreito"?
\end{itensq}
\resposta{Ver análise comentada}


\end{questao}

\blocofixacao
\begin{questao}[1]{}

Um material linear tem $\mu_r=500$ e é submetido a
$H=\SI{100}{\ampere\per\meter}$. Determine $B$.
\dados{$\mu_0=4\pi\times10^{-7}\,\si{\tesla\meter\per\ampere}$}
\resposta{$B=\SI{62.8}{\milli\tesla}$}


\end{questao}
\begin{questao}[1]{}

Determine a magnetização $M$ do material da questão anterior.
\resposta{$M=\SI{49.9}{\kilo\ampere\per\meter}$}


\end{questao}
\begin{questao}[1]{}

Define-se a suscetibilidade magnética como $\chi=\mu_r-1$. Classifique os
materiais com $\chi=+\pot{5}{-4}$ e $\chi=-\pot{2}{-5}$.
\resposta{Paramagnético e diamagnético, respectivamente}


\end{questao}
\begin{questao}[1]{}

Um enrolamento de $500$ espiras percorrido por $\SI{0.4}{\ampere}$ envolve um
núcleo de caminho médio $\SI{25}{\centi\meter}$. Determine $H$.
\resposta{$H=\SI{800}{\ampere\per\meter}$}


\end{questao}

\blocoaplicacao
\begin{questao}[2]{}

Um núcleo de $\SI{10}{\centi\meter}$ de caminho médio, seção
$\SI{1}{\centi\meter\squared}$ e $\mu_r=1000$ recebe $\SI{200}{}$
ampère-espiras.
\begin{itensq}
\item Determine a relutância.
\item Determine o fluxo e a densidade de fluxo.
\item Avalie o resultado.
\end{itensq}
\resposta{$\mathcal{R}=\pot{7.96}{5}\,\si{\per\henry}$;
$\Phi=\SI{0.251}{\milli\weber}$; $B=\SI{2.51}{\tesla}$ --- \textbf{saturado}}


\end{questao}
\begin{questao}[2]{}

Um mesmo enrolamento produz $H=\SI{500}{\ampere\per\meter}$. Determine $B$ para
núcleo de ar, de alumínio ($\mu_r=1{,}00002$) e de ferro-silício
($\mu_r=4000$).
\resposta{$\SI{628}{\micro\tesla}$; $\SI{628}{\micro\tesla}$;
$\SI{2.51}{\tesla}$}


\end{questao}
\begin{questao}[2]{}

Um material apresenta $B=\SI{1.2}{\tesla}$ para $H=\SI{400}{\ampere\per\meter}$
e $B=\SI{1.5}{\tesla}$ para $H=\SI{1200}{\ampere\per\meter}$.
\begin{itensq}
\item Determine $\mu_r$ em cada ponto.
\item Interprete a variação.
\end{itensq}
\resposta{$\mu_r=2390$ e $995$}


\end{questao}
\begin{questao}[2]{}

Explique a diferença entre a permeabilidade \textbf{normal}
($\mu=B/H$) e a permeabilidade \textbf{diferencial}
($\mu_d=\mathrm{d}B/\mathrm{d}H$), e indique quando cada uma se aplica.
\resposta{Normal para o ponto de operação; diferencial para pequenos sinais
sobrepostos}


\end{questao}
\begin{questao}[2]{}

Compare a temperatura de Curie do ferro ($\SI{770}{\celsius}$), do níquel
($\SI{358}{\celsius}$) e de uma ferrite típica
($\approx\SI{250}{\celsius}$), e explique o que ocorre acima dela.
\resposta{Acima de $T_C$ o material torna-se paramagnético e perde
praticamente toda a permeabilidade}


\end{questao}
\begin{questao}[2]{}

Um ímã permanente de ferrite tem coeficiente de temperatura de
$\SI{-0.2}{\percent\per\celsius}$ para a remanência. Determine a variação de
$B_r$ entre $\SI{20}{\celsius}$ e $\SI{120}{\celsius}$.
\resposta{Redução de $20\%$}


\end{questao}
\begin{questao}[2]{}

Um núcleo de aço-silício dissipa perdas por histerese proporcionais à área do
ciclo. Com área de $\SI{200}{\joule\per\meter\cubed}$ por ciclo e volume de
$\SI{1000}{\centi\meter\cubed}$, determine a potência a $\SI{60}{\hertz}$ e a
$\SI{400}{\hertz}$.
\resposta{$\SI{12}{\watt}$ e $\SI{80}{\watt}$}


\end{questao}
\begin{questao}[2]{}

Um material magnético \textbf{mole} tem $H_c=\SI{50}{\ampere\per\meter}$ e um
material \textbf{duro} tem $H_c=\SI{500}{\kilo\ampere\per\meter}$. Compare-os
e indique aplicações.
\resposta{Razão de $10^{4}$ na coercividade; moles para núcleos, duros para ímãs
permanentes}


\end{questao}
\begin{questao}[2]{}

Explique por que a curva de magnetização \textbf{inicial} de um material
virgem difere do ciclo de histerese subsequente.
\resposta{A curva inicial parte da origem; o ciclo nunca mais passa por ela}


\end{questao}
\begin{questao}[2]{}

Um transformador é desligado no pico da corrente de magnetização. Explique o
que ocorre na religação e por que a corrente de \emph{inrush} pode ser elevada.
\resposta{O fluxo remanente soma-se ao fluxo imposto, podendo levar o núcleo a
saturação profunda}


\end{questao}
\begin{questao}[2]{}

Um circuito magnético tem dois trechos em série: ferro
($\ell_1=\SI{30}{\centi\meter}$, $\mu_r=2000$) e ar
($\ell_2=\SI{2}{\milli\meter}$), ambos com seção
$\SI{5}{\centi\meter\squared}$. Determine a relutância total e a fração de cada
trecho.
\resposta{$\mathcal{R}_{fe}=\pot{2.39}{5}$, $\mathcal{R}_{ar}=\pot{3.18}{6}$;
o ar responde por $93\%$}


\end{questao}

\blocoaprofundamento
\begin{questao}[3]{}

Distinga rigorosamente $\vec H$, $\vec B$ e $\vec M$.
\begin{itensq}
\item Escreva a relação entre elas e as unidades.
\item Determine qual delas muda ao inserir um núcleo em um solenoide, mantida a
      corrente.
\item Explique qual é "mais fundamental" e por quê.
\end{itensq}
\resposta{$\vec B=\mu_0(\vec H+\vec M)$; ao inserir o núcleo, $H$ não muda e
$B$ aumenta}


\end{questao}
\begin{questao}[3]{}

A curva $B\times H$ de um material apresenta três regiões distintas.
\begin{itensq}
\item Descreva cada uma em termos de domínios.
\item Identifique onde $\mu_r$ é máxima.
\item Indique a região de operação recomendada.
\end{itensq}
\resposta{Deslocamento reversível, deslocamento irreversível e rotação;
$\mu_r$ máxima no trecho intermediário}


\end{questao}
\begin{questao}[3]{}

Um núcleo tem curva $B\times H$ não linear, e é preciso determinar o ponto de
operação com entreferro.
\begin{itensq}
\item Formule o problema.
\item Descreva o método gráfico da reta de entreferro.
\item Explique por que o entreferro estabiliza o ponto de operação.
\end{itensq}
\resposta{Interseção da curva $B\times H$ do material com a reta imposta pelo
entreferro}


\end{questao}
\begin{questao}[3]{}

Compare aço-silício e ferrite para aplicação em transformadores.
\begin{itensq}
\item Compare $B_{sat}$, resistividade e faixa de frequência.
\item Determine em qual faixa cada um é preferível.
\item Justifique fisicamente.
\end{itensq}
\resposta{Aço-silício: alto $B_{sat}$, baixa resistividade, até
$\si{\kilo\hertz}$; ferrite: baixo $B_{sat}$, altíssima resistividade,
até $\si{\mega\hertz}$}


\end{questao}
\begin{questao}[3]{}

Um ímã permanente é retirado de seu circuito magnético e deixado com os polos
livres ao ar.
\begin{itensq}
\item Explique o que ocorre com seu ponto de operação.
\item Determine o efeito de aproximar uma peça de ferro.
\item Justifique a prática de armazenar ímãs com "quebra-fluxo".
\end{itensq}
\resposta{Ao ar, o ponto de operação desce na curva de desmagnetização; o ferro
o eleva}


\end{questao}
\begin{questao}[3]{}

Um material tem $B_r=\SI{1.2}{\tesla}$ e $H_c=\SI{900}{\kilo\ampere\per\meter}$.
\begin{itensq}
\item Estime o produto de energia máximo.
\item Compare com ferrite ($B_r=\SI{0.4}{\tesla}$,
      $H_c=\SI{250}{\kilo\ampere\per\meter}$).
\item Explique o significado do produto de energia.
\end{itensq}
\resposta{$(BH)_{\max}\approx\SI{270}{\kilo\joule\per\meter\cubed}$ contra
$\SI{25}{\kilo\joule\per\meter\cubed}$}


\end{questao}
\begin{questao}[3]{}

Explique a diferença entre desmagnetização por \textbf{aquecimento} e por
\textbf{campo alternado decrescente}.
\begin{itensq}
\item Descreva cada processo.
\item Compare eficácia e aplicabilidade.
\end{itensq}
\resposta{Aquecimento acima de $T_C$ desordena os domínios; campo alternado os
distribui simetricamente}


\end{questao}
\begin{questao}[3]{}

Um núcleo opera com $B_{\max}=\SI{1.2}{\tesla}$ a $\SI{60}{\hertz}$. Analise o
efeito de elevar $B_{\max}$ para $\SI{1.5}{\tesla}$.
\begin{itensq}
\item Determine o efeito sobre o número de espiras.
\item Determine o efeito sobre as perdas.
\item Avalie o compromisso.
\end{itensq}
\resposta{$N$ cai $20\%$; perdas por histerese sobem cerca de $60\%$ e
parasitas $56\%$}


\end{questao}
\begin{questao}[3]{}

Explique por que a permeabilidade de um material ferromagnético não é uma
constante, e quais parâmetros a afetam.
\resposta{$\mu$ depende de $H$, da temperatura, da frequência e da história
magnética}


\end{questao}
\begin{questao}[3]{}

Um sensor de proximidade indutivo detecta metais pela variação de indutância.
\begin{itensq}
\item Explique o princípio para metais ferromagnéticos.
\item Explique para metais não ferromagnéticos (alumínio, cobre).
\item Determine por que o alcance difere entre eles.
\end{itensq}
\resposta{Ferromagnéticos aumentam $L$ por permeabilidade; não ferromagnéticos
reduzem $L$ e o fator de qualidade por correntes parasitas}


\end{questao}

\blocodesafio
\begin{questao}[4]{}

\textbf{(Energia do ciclo de histerese.)} Demonstre que a área do laço
$B\times H$ representa energia dissipada por unidade de volume e por ciclo.
\begin{itensq}
\item Deduza a expressão.
\item Determine a potência dissipada em função da frequência.
\item Explique a origem microscópica da perda.
\end{itensq}
\resposta{$w=\oint H\,\mathrm{d}B$ por ciclo; $P=fVw$}


\end{questao}
\begin{questao}[4]{}

\textbf{(Circuito magnético com dois materiais.)} Um circuito tem trecho de
ferro ($\ell_1$, $\mu_1$) em série com trecho de ferrite ($\ell_2$,
$\mu_2$), mesma seção.
\begin{itensq}
\item Formule as equações.
\item Determine $B$ e $H$ em cada trecho.
\item Explique qual grandeza é contínua na interface e por quê.
\end{itensq}
\resposta{$B$ é contínuo; $H$ é descontínuo, na razão inversa das
permeabilidades}


\end{questao}
\begin{questao}[4]{}

\textbf{(Ensaio do ciclo de histerese.)} Projete um experimento para levantar o
laço $B\times H$ de um material em osciloscópio.
\begin{itensq}
\item Descreva o circuito e o que cada canal mede.
\item Determine as constantes de conversão.
\item Aponte os cuidados experimentais.
\end{itensq}
\resposta{Canal horizontal $\propto H$ (corrente por shunt); vertical
$\propto B$ (f.e.m. integrada)}


\end{questao}
\begin{questao}[4]{}

\textbf{(Domínios magnéticos.)} Explique a estrutura de domínios e sua relação
com as propriedades macroscópicas.
\begin{itensq}
\item Explique por que os domínios existem.
\item Relacione a estrutura com $B_r$, $H_c$ e $\mu_r$.
\item Explique a diferença entre materiais moles e duros nesse nível.
\end{itensq}
\resposta{Domínios minimizam a energia magnetostática; a mobilidade das paredes
determina $\mu_r$ e $H_c$}


\end{questao}
\begin{questao}[4]{}

\textbf{(Projeto sob restrição de perdas.)} Um transformador de
$\SI{60}{\hertz}$ deve ter perdas no núcleo inferiores a $\SI{2}{\watt}$, com
material cuja perda específica é
$\SI{1.5}{\watt\per\kilogram}$ a $\SI{1.5}{\tesla}$ e
$\SI{0.9}{\watt\per\kilogram}$ a $\SI{1.2}{\tesla}$.
\begin{itensq}
\item Determine a massa admissível em cada condição.
\item Determine o efeito sobre o número de espiras.
\item Estabeleça o critério de escolha.
\end{itensq}
\resposta{$\SI{1.33}{\kilogram}$ a $\SI{1.5}{\tesla}$;
$\SI{2.22}{\kilogram}$ a $\SI{1.2}{\tesla}$}


\end{questao}
\begin{questao}[4]{}

\textbf{(Anisotropia e grão orientado.)} Explique o que é o aço-silício de grão
orientado e por que ele é usado em transformadores mas não em motores.
\begin{itensq}
\item Explique o processo e o efeito.
\item Compare as perdas nas duas direções.
\item Justifique a diferença de aplicação.
\end{itensq}
\resposta{Grãos alinhados numa direção preferencial; excelente ao longo dela,
ruim transversalmente}


\end{questao}
\begin{questao}[4]{}

\textbf{(Blindagem magnética com material de alta permeabilidade.)} Analise
quantitativamente a blindagem de campo estático por invólucro de mu-metal.
\begin{itensq}
\item Explique o mecanismo por relutância.
\item Estime o fator de blindagem para invólucro cilíndrico.
\item Aponte as limitações práticas.
\end{itensq}
\resposta{Fator $\approx\mu_r t/D$; o mecanismo é desvio de fluxo, não
cancelamento}


\end{questao}
\begin{questao}[4]{}

\textbf{(Ponto de operação de ímã permanente.)} Determine o ponto de operação
de um ímã em um circuito magnético com entreferro.
\begin{itensq}
\item Formule a equação da reta de carga.
\item Determine a condição de máximo aproveitamento.
\item Explique o efeito de aumentar o entreferro.
\end{itensq}
\resposta{A reta de carga tem inclinação
$-\dfrac{\ell_m A_g}{g A_m}$; o aproveitamento é máximo em
$(BH)_{\max}$}


\end{questao}

\gabarito[12.1 Conceitos iniciais e materiais magnéticos]

\subsection{Campo magnético}

\blocofixacao
\begin{questao}[1]{}

O campo magnético a uma distância $r$ de um fio retilíneo longo percorrido por
corrente $I$:
\begin{alternativas}
\item é proporcional a $r$.
\item é inversamente proporcional a $r$.
\item é inversamente proporcional a $r^2$.
\item independe de $r$.
\item é proporcional a $r^2$.
\end{alternativas}
\resposta{(b)}


\end{questao}
\begin{questao}[1]{}

Um fio reto longo conduz $I = \SI{10}{\ampere}$. Determine o módulo do campo
magnético a $\SI{5}{\centi\meter}$ do fio.
\dados{$\mu_0 = 4\pi\times10^{-7}\,\si{\tesla\meter\per\ampere}$.}
\resposta{$B = \pot{4}{-5}\,\si{\tesla} = \SI{40}{\micro\tesla}$}


\end{questao}
\begin{questao}[1]{}

Um fio retilíneo é percorrido por corrente para cima (saindo do chão em direção ao
teto). Usando a regra da mão direita, as linhas de campo magnético ao redor do fio:
\begin{alternativas}
\item apontam radialmente para fora do fio.
\item apontam radialmente para dentro do fio.
\item formam círculos concêntricos em torno do fio.
\item são paralelas ao fio.
\item não existem (fio reto não gera campo).
\end{alternativas}
\resposta{(c)}


\end{questao}

\blocoaplicacao
\begin{questao}[2]{}

Um solenoide de $\SI{1000}{}$ espiras por metro é percorrido por
$I = \SI{2}{\ampere}$. Determine o campo magnético no seu interior.
\dados{$\mu_0 = 4\pi\times10^{-7}\,\si{\tesla\meter\per\ampere}$.}
\resposta{$B \approx \SI{2.51}{\milli\tesla}$}


\end{questao}
\begin{questao}[2]{}

Uma espira circular de raio $R = \SI{10}{\centi\meter}$ conduz
$I = \SI{5}{\ampere}$. Determine o campo magnético no seu centro.
\dados{$\mu_0 = 4\pi\times10^{-7}\,\si{\tesla\meter\per\ampere}$.}
\resposta{$B \approx \SI{31.4}{\micro\tesla}$}


\end{questao}
\begin{questao}[2]{}

Um fio reto conduz $I = \SI{4}{\ampere}$. Determine o campo magnético a
$\SI{2}{\centi\meter}$ do fio.
\dados{$\mu_0 = 4\pi\times10^{-7}\,\si{\tesla\meter\per\ampere}$.}
\resposta{$B = \SI{40}{\micro\tesla}$}


\end{questao}

\blocoaprofundamento
\begin{questao}[3]{}

Dois fios paralelos, distantes $\SI{10}{\centi\meter}$, conduzem correntes de
$\SI{3}{\ampere}$ e $\SI{6}{\ampere}$ no \textbf{mesmo sentido}. Determine o campo
magnético resultante no ponto médio entre eles e diga se os fios se atraem ou se
repelem.
\dados{$\mu_0 = 4\pi\times10^{-7}\,\si{\tesla\meter\per\ampere}$.}
\resposta{$B \approx \SI{12}{\micro\tesla}$; os fios se atraem}


\end{questao}
\begin{questao}[3]{}

Dois fios paralelos e longos, separados por $\SI{5}{\centi\meter}$, conduzem
$\SI{10}{\ampere}$ cada, em sentidos \textbf{opostos}. Determine a força por
unidade de comprimento entre eles e diga se é de atração ou repulsão.
\dados{$\mu_0 = 4\pi\times10^{-7}\,\si{\tesla\meter\per\ampere}$.}
\resposta{$F/\ell = \SI{4e-4}{\newton\per\meter}$, repulsão}


\end{questao}
\begin{questao}[3]{}

Um \textbf{toroide} de raio médio $r = \SI{10}{\centi\meter}$ possui $N = 200$
espiras percorridas por $I = \SI{3}{\ampere}$.
\begin{itensq}
\item Aplicando a lei de Ampère a uma circunferência de raio $r$ interna ao
      toroide, deduza a expressão de $B$ e calcule seu valor.
\item Mostre que o campo \emph{fora} do toroide é nulo.
\item Compare com um solenoide reto de mesmo $N$ e comprimento igual ao
      perímetro médio do toroide. O que se conclui?
\end{itensq}
\dados{$\mu_0 = 4\pi\times10^{-7}\,\si{\tesla\meter\per\ampere}$.}
\resposta{(a) $B = \dfrac{\mu_0 N I}{2\pi r} = \SI{1.2}{\milli\tesla}$;
(b) e (c) ver comentário}


\end{questao}
\begin{questao}[3]{Apostila original revisada}

Duas espiras circulares, concêntricas e coplanares, têm raios
$R_1=10\pi\,\si{\centi\metre}$ e $R_2=2{,}5\pi\,\si{\centi\metre}$. A primeira é
percorrida por $I_1=\SI{2}{\ampere}$. Determine o valor e o sentido de $I_2$ para
que o campo magnético resultante no centro seja nulo.
\resposta{$I_2=\SI{0.5}{\ampere}$, em sentido oposto a $I_1$}


\end{questao}

\blocofixacao
\begin{questao}[1]{}

Um fio retilíneo longo conduz $\SI{2}{\ampere}$. Determine o campo magnético a
$\SI{1}{\centi\meter}$ dele.
\dados{$\mu_0=4\pi\times10^{-7}\,\si{\tesla\meter\per\ampere}$}
\resposta{$B=\SI{40}{\micro\tesla}$}


\end{questao}
\begin{questao}[1]{}

Um solenoide tem $\SI{2000}{}$ espiras por metro e conduz
$\SI{3}{\ampere}$. Determine o campo em seu interior.
\resposta{$B=\SI{7.54}{\milli\tesla}$}


\end{questao}
\begin{questao}[1]{}

Determine a corrente necessária para que um fio retilíneo produza
$B=\SI{100}{\micro\tesla}$ a $\SI{2}{\centi\meter}$.
\resposta{$I=\SI{10}{\ampere}$}


\end{questao}
\begin{questao}[1]{}

Sobre o campo magnético de um solenoide longo, é correto afirmar:
\begin{alternativas}[2]
\item é uniforme no interior e praticamente nulo no exterior
\item é maior no exterior
\item aumenta com a distância ao eixo, no interior
\item é nulo no interior
\end{alternativas}
\resposta{(a)}


\end{questao}
\begin{questao}[1]{}

O campo magnético da Terra vale aproximadamente $\SI{50}{\micro\tesla}$.
Expresse-o em gauss.
\dados{$\SI{1}{\tesla}=10^{4}\,\si{\gauss}$}
\resposta{$\SI{0.5}{\gauss}$}


\end{questao}

\blocoaplicacao
\begin{questao}[2]{}

Um solenoide de $100$ espiras e $\SI{10}{\centi\meter}$ de comprimento conduz
$\SI{1}{\ampere}$, com núcleo de ar.
\begin{itensq}
\item Determine a densidade de espiras.
\item Determine o campo interno.
\end{itensq}
\resposta{$n=\SI{1000}{\per\meter}$; $B=\SI{1.26}{\milli\tesla}$}


\end{questao}
\begin{questao}[2]{}

Um fio conduz corrente constante. Compare o campo a $\SI{4}{\centi\meter}$ com
o campo a $\SI{12}{\centi\meter}$.
\resposta{$B_1/B_2=3$}


\end{questao}
\begin{questao}[2]{}

Uma espira circular de raio $\SI{5}{\centi\meter}$ deve produzir
$\SI{100}{\micro\tesla}$ em seu centro. Determine a corrente.
\resposta{$I\approx\SI{7.96}{\ampere}$}


\end{questao}
\begin{questao}[2]{}

Um solenoide tem $N$ espiras e comprimento $L$. Determine o que ocorre com o
campo interno se $N$ e $L$ forem \textbf{ambos dobrados}.
\resposta{O campo não se altera}


\end{questao}
\begin{questao}[2]{}

Dois fios longos e paralelos, distantes $\SI{8}{\centi\meter}$, conduzem
$\SI{4}{\ampere}$ cada, no \textbf{mesmo} sentido. Determine o campo no ponto
médio entre eles.
\resposta{$B=0$}


\end{questao}
\begin{questao}[2]{}

Um fio reto é imerso em um meio de permeabilidade relativa $\mu_r=200$.
Determine o efeito sobre o campo, comparado ao ar.
\resposta{O campo fica $200$ vezes maior}


\end{questao}
\begin{questao}[2]{}

Determine a distância a que um fio de $\SI{50}{\ampere}$ produz campo igual ao
terrestre ($\SI{50}{\micro\tesla}$).
\resposta{$r=\SI{20}{\centi\meter}$}


\end{questao}
\begin{questao}[2]{}

Um toroide e um solenoide têm a mesma densidade de espiras e a mesma corrente.
Compare seus campos internos e externos.
\resposta{Campo interno praticamente igual; o toroide tem campo externo
\emph{rigorosamente} nulo}


\end{questao}
\begin{questao}[2]{}

Dois fios perpendiculares entre si, ambos no mesmo plano da folha, conduzem
$\SI{6}{\ampere}$ e $\SI{8}{\ampere}$. Determine o campo resultante em um
ponto a $\SI{3}{\centi\meter}$ do primeiro e $\SI{4}{\centi\meter}$ do
segundo.
\resposta{$B=\SI{56.6}{\micro\tesla}$}


\end{questao}

\blocoaprofundamento
\begin{questao}[3]{}

Um condutor cilíndrico maciço de raio $a=\SI{2}{\milli\meter}$ conduz
$\SI{100}{\ampere}$ uniformemente distribuídos na seção.
\begin{itensq}
\item Determine o campo a $\SI{1}{\milli\meter}$, na superfície e a
      $\SI{5}{\milli\meter}$ do eixo.
\item Identifique onde o campo é máximo.
\end{itensq}
\resposta{$5$, $10$ e $\SI{4}{\milli\tesla}$; máximo na superfície}


\end{questao}
\begin{questao}[3]{}

Um cabo coaxial tem condutor central de raio $a$ conduzindo $I$ e blindagem
externa conduzindo $I$ em sentido \textbf{oposto}, com $I=\SI{5}{\ampere}$.
\begin{itensq}
\item Determine o campo entre os condutores, a $\SI{2}{\milli\meter}$ do eixo.
\item Determine o campo fora do cabo.
\item Explique a consequência prática.
\end{itensq}
\resposta{$B=\SI{0.5}{\milli\tesla}$ entre os condutores; $B=0$ fora}


\end{questao}
\begin{questao}[3]{}

Uma espira de raio $R=\SI{10}{\centi\meter}$ conduz $\SI{5}{\ampere}$.
\begin{itensq}
\item Determine o campo no centro.
\item Determine o campo no eixo, a $\SI{10}{\centi\meter}$ e a
      $\SI{20}{\centi\meter}$ do centro.
\item Comente a taxa de queda.
\end{itensq}
\resposta{$31{,}4$; $11{,}1$ e $\SI{2.81}{\micro\tesla}$}


\end{questao}
\begin{questao}[3]{}

Compare o campo no centro de um solenoide \emph{finito} com o do solenoide
ideal, e determine o campo em sua extremidade.
\begin{itensq}
\item Indique a condição para a aproximação ideal.
\item Determine o campo na extremidade.
\item Discuta a consequência para projeto.
\end{itensq}
\resposta{Na extremidade, $B\approx\mu_0nI/2$ --- metade do valor central}


\end{questao}
\begin{questao}[3]{}

Três fios longos e paralelos, coplanares e igualmente espaçados de
$\SI{5}{\centi\meter}$, conduzem $\SI{10}{\ampere}$ cada, todos no mesmo
sentido. Determine o campo no fio central.
\resposta{$B=0$ no fio central}


\end{questao}
\begin{questao}[3]{}

Um eletroímã deve produzir $\SI{0.1}{\tesla}$ com núcleo de ar, usando
solenoide de $\SI{20}{\centi\meter}$.
\begin{itensq}
\item Determine o produto $NI$ necessário.
\item Avalie a viabilidade com $N=1000$ espiras.
\item Proponha alternativa.
\end{itensq}
\resposta{$NI=\SI{15.9}{\kilo\ampere}$-espira; exigiria $\SI{15.9}{\ampere}$
com $1000$ espiras}


\end{questao}
\begin{questao}[3]{}

Um cabo de par paralelo tem condutores separados por $\SI{2}{\milli\meter}$,
conduzindo $\SI{1}{\ampere}$ em sentidos opostos.
\begin{itensq}
\item Determine o campo no ponto médio entre eles.
\item Determine o campo a $\SI{10}{\centi\meter}$ do par.
\item Compare com um fio único.
\end{itensq}
\resposta{$\SI{400}{\micro\tesla}$ no meio; $\approx\SI{0.04}{\micro\tesla}$ a
$\SI{10}{\centi\meter}$}


\end{questao}
\begin{questao}[3]{}

Um sensor Hall indica $\SI{2.5}{\milli\volt}$ quando exposto a
$\SI{100}{\milli\tesla}$, com resposta linear e tensão de repouso nula.
\begin{itensq}
\item Determine a sensibilidade.
\item Determine o campo correspondente a uma leitura de
      $\SI{0.6}{\milli\volt}$.
\item Discuta as fontes de erro.
\end{itensq}
\resposta{$\SI{25}{\micro\volt\per\milli\tesla}$; $B=\SI{24}{\milli\tesla}$}


\end{questao}

\blocodesafio
\begin{questao}[4]{}

\textbf{(Dedução por Biot--Savart.)} Deduza o campo magnético no centro de uma
espira circular de raio $R$ percorrida por corrente $I$.
\begin{itensq}
\item Aplique a lei de Biot--Savart a um elemento de corrente.
\item Integre e obtenha o resultado.
\item Explique por que a integração é simples neste caso.
\end{itensq}
\resposta{$B=\dfrac{\mu_0 I}{2R}$}


\end{questao}
\begin{questao}[4]{}

\textbf{(Ampère contra Biot--Savart.)} Compare os dois métodos de cálculo de
campo magnético.
\begin{itensq}
\item Enuncie cada lei e sua condição de uso.
\item Determine quando a lei de Ampère é praticável.
\item Classifique quatro configurações quanto ao método adequado.
\end{itensq}
\resposta{Ampère exige simetria suficiente para tirar $B$ da integral;
Biot--Savart é geral mas trabalhoso}


\end{questao}
\begin{questao}[4]{}

\textbf{(Projeto de solenoide.)} Projete um solenoide de $500$ espiras e
$\SI{25}{\centi\meter}$ para produzir $\SI{1}{\milli\tesla}$.
\begin{itensq}
\item Determine a corrente necessária.
\item Determine a potência dissipada, para fio de resistência total
      $\SI{2}{\ohm}$.
\item Avalie a uniformidade do campo.
\end{itensq}
\resposta{$I=\SI{0.398}{\ampere}$; $P=\SI{0.32}{\watt}$}


\end{questao}
\begin{questao}[4]{}

\textbf{(Dipolo magnético.)} Uma espira de raio $R$ e corrente $I$ é observada
no eixo, a distância $x\gg R$.
\begin{itensq}
\item Obtenha a forma assintótica do campo.
\item Defina o momento de dipolo magnético e reescreva o resultado.
\item Compare com o dipolo elétrico e comente a ausência de monopolos.
\end{itensq}
\resposta{$B\to\dfrac{\mu_0 m}{2\pi x^{3}}$, com $m=I\pi R^{2}$}


\end{questao}
\begin{questao}[4]{}

\textbf{(Blindagem magnética.)} Compare a blindagem de campos magnéticos com a
de campos elétricos.
\begin{itensq}
\item Explique por que uma gaiola metálica não blinda campo magnético estático.
\item Descreva os dois mecanismos disponíveis.
\item Discuta a dependência com a frequência.
\end{itensq}
\resposta{Não há "cargas magnéticas" a redistribuir; usa-se desvio por material
de alto $\mu_r$ ou correntes induzidas}


\end{questao}
\begin{questao}[4]{}

\textbf{(Mapeamento experimental.)} Elabore um procedimento para mapear o campo
magnético de uma bobina usando sensor Hall.
\begin{itensq}
\item Descreva o procedimento e o tratamento do \emph{offset}.
\item Indique como obter as três componentes do campo.
\item Estabeleça o critério de validação dos dados.
\end{itensq}
\resposta{Medir com inversão de orientação para cancelar \emph{offset};
três orientações ortogonais para as componentes}


\end{questao}
\begin{questao}[4]{}

\textbf{(Limites de eletroímãs.)} Analise o que limita, na prática, o campo
magnético alcançável.
\begin{itensq}
\item Determine o limite imposto pela saturação do núcleo.
\item Determine o limite térmico de bobinas de cobre.
\item Descreva as soluções empregadas para superá-los.
\end{itensq}
\resposta{Núcleos saturam entre $1{,}5$ e $\SI{2}{\tesla}$; acima disso é
preciso bobina sem núcleo, com corrente muito alta}


\end{questao}
\begin{questao}[4]{}

\textbf{(Superposição e simetria.)} Um condutor cilíndrico de raio $a$ possui
uma cavidade cilíndrica de raio $b$, paralela ao eixo e deslocada de $d$.
A corrente $I$ distribui-se uniformemente na seção restante.
\begin{itensq}
\item Formule a estratégia de solução por superposição.
\item Determine o campo dentro da cavidade.
\item Interprete o resultado.
\end{itensq}
\resposta{O campo na cavidade é \textbf{uniforme}, de módulo
$\mu_0Jd/2$}


\end{questao}

\gabarito[12.2 Campo magnético]

\subsection{Fluxo magnético}

\blocofixacao
\begin{questao}[1]{}

O fluxo magnético através de uma espira é máximo quando o plano da espira está:
\begin{alternativas}
\item paralelo ao campo.
\item perpendicular ao campo (campo atravessando a espira frontalmente).
\item inclinado a $\SI{45}{\degree}$ com o campo.
\item em qualquer orientação (o fluxo não depende do ângulo).
\item girando rapidamente.
\end{alternativas}
\resposta{(b)}


\end{questao}
\begin{questao}[1]{}

Uma espira de área $\SI{20}{\centi\meter\squared}$ é atravessada
perpendicularmente por um campo uniforme de $\SI{0.5}{\tesla}$. Calcule o fluxo
magnético.
\resposta{$\Phi = \SI{1}{\milli\weber}$}


\end{questao}

\blocoaplicacao
\begin{questao}[2]{}

Uma espira circular de raio $\SI{8}{\centi\meter}$ está imersa em um campo
uniforme de $\SI{0.3}{\tesla}$, perpendicular ao seu plano.
\begin{itensq}
\item Calcule o fluxo magnético através da espira.
\item Se o campo variar linearmente de $\SI{0.3}{}$ para $\SI{0.1}{\tesla}$ em
      $\SI{0.5}{\second}$, qual é a fem induzida?
\end{itensq}
\resposta{(a) $\Phi \approx \SI{6.03}{\milli\weber}$; (b) $\varepsilon \approx \SI{8.04}{\milli\volt}$}


\end{questao}
\begin{questao}[2]{Apostila original revisada}

Um campo magnético uniforme de módulo $B=\SI{1e-4}{\tesla}$ atravessa uma espira de
área $A=\SI{1e-5}{\metre\squared}$. O campo forma $\SI{45}{\degree}$ com a normal
ao plano da espira. Calcule o fluxo magnético.
\resposta{$\Phi=\SI{7.07e-10}{\weber}$}


\end{questao}

\blocoaprofundamento
\begin{questao}[3]{}

O gráfico mostra a variação do fluxo magnético $\Phi(t)$ através de uma espira.
Determine a fem induzida em cada trecho e esboce mentalmente a forma de onda da
fem.

\circuitoqq{
\begin{tikzpicture}
 \begin{axis}[width=10cm,height=5cm,xlabel={$t$ (s)},ylabel={$\Phi$ (mWb)},
   xmin=0,xmax=6,ymin=0,ymax=5,xtick={0,1,2,3,4,5,6},ytick={0,2,4},
   grid=both,axis lines=left]
   \addplot[Icol,very thick] coordinates {(0,0)(2,4)(4,4)(6,0)};
 \end{axis}
\end{tikzpicture}}{Fluxo magnético em função do tempo.}

\resposta{$0$--$2\,$s: $\varepsilon=\SI{2}{\milli\volt}$; $2$--$4\,$s:
$\varepsilon=0$; $4$--$6\,$s: $\varepsilon=\SI{-2}{\milli\volt}$}


\end{questao}

\blocodesafio
\begin{questao}[4]{}

Uma espira retangular de área $A = \SI{200}{\centi\meter\squared}$ gira com
velocidade angular constante $\omega$ em um campo magnético uniforme
$B = \SI{0.4}{\tesla}$, completando $\SI{1800}{}$ rotações por minuto.
\begin{itensq}
\item Escreva a expressão do fluxo $\Phi(t)$ e calcule seu valor máximo.
\item Obtenha a fem induzida $\varepsilon(t)$ e seu valor de pico.
\item Determine a frequência da tensão gerada e explique por que a saída é
      \emph{senoidal}.
\item Se a espira tivesse $N = 100$ voltas, o que mudaria?
\end{itensq}
\resposta{(a) $\Phi_{\max}=\SI{8}{\milli\weber}$;
(b) $\varepsilon_{\text{pico}}\approx\SI{1.51}{\volt}$; (c) \SI{30}{\hertz};
(d) tudo $\times100$}


\end{questao}

\blocofixacao
\begin{questao}[1]{}

Uma superfície plana de $\SI{0.01}{\meter\squared}$ está em campo uniforme de
$\SI{0.1}{\tesla}$, com a normal formando $\ang{60}$ com o campo. Determine o
fluxo.
\resposta{$\Phi=\SI{0.5}{\milli\weber}$}


\end{questao}
\begin{questao}[1]{}

Determine o fluxo quando o campo é \textbf{paralelo} ao plano da superfície.
\resposta{$\Phi=0$}


\end{questao}
\begin{questao}[1]{}

Um fluxo de $\SI{2}{\milli\weber}$ atravessa perpendicularmente uma superfície
em campo de $\SI{0.4}{\tesla}$. Determine a área.
\resposta{$A=\SI{50}{\centi\meter\squared}$}


\end{questao}
\begin{questao}[1]{}

Uma espira de $\SI{30}{\centi\meter\squared}$ é atravessada
perpendicularmente por fluxo de $\SI{1.5}{\milli\weber}$. Determine a densidade
de fluxo.
\resposta{$B=\SI{0.5}{\tesla}$}


\end{questao}
\begin{questao}[1]{}

O fluxo magnético líquido através de qualquer superfície \textbf{fechada} é:
\begin{alternativas}[2]
\item proporcional à corrente envolvida
\item sempre nulo
\item proporcional ao volume
\item máximo quando a superfície é esférica
\end{alternativas}
\resposta{(b)}


\end{questao}
\begin{questao}[1]{}

Converta $\SI{5}{\milli\weber}$ para maxwells.
\dados{$\SI{1}{\weber}=10^{8}\,\si{\maxwell}$}
\resposta{$\pot{5}{5}\,\si{\maxwell}$}


\end{questao}

\blocoaplicacao
\begin{questao}[2]{}

Uma bobina de $50$ espiras é atravessada por fluxo de
$\SI{1}{\milli\weber}$ por espira.
\begin{itensq}
\item Determine o fluxo concatenado.
\item Determine a indutância, se a corrente for $\SI{2}{\ampere}$.
\end{itensq}
\resposta{$\lambda=\SI{50}{\milli\weber}$-espira; $L=\SI{25}{\milli\henry}$}


\end{questao}
\begin{questao}[2]{}

Uma espira de $\SI{200}{\centi\meter\squared}$ está em campo de
$\SI{0.25}{\tesla}$, com a normal a $\ang{60}$ do campo. Determine o fluxo e o
valor que ele teria na posição de fluxo máximo.
\resposta{$\Phi=\SI{2.5}{\milli\weber}$; máximo de $\SI{5}{\milli\weber}$}


\end{questao}
\begin{questao}[2]{}

Um solenoide de $\SI{1000}{}$ espiras por metro conduz $\SI{2}{\ampere}$ e tem
seção de $\SI{10}{\centi\meter\squared}$. Determine o fluxo por espira.
\resposta{$\Phi=\SI{2.51}{\micro\weber}$}


\end{questao}
\begin{questao}[2]{}

O fluxo através de uma bobina de $100$ espiras varia linearmente de zero a
$\SI{4}{\milli\weber}$ em $\SI{20}{\milli\second}$. Determine a f.e.m.
induzida.
\resposta{$\varepsilon=\SI{20}{\volt}$}


\end{questao}
\begin{questao}[2]{}

Duas superfícies diferentes apoiam-se na \textbf{mesma} espira: uma plana e
outra em forma de calota. Compare o fluxo através delas.
\resposta{São iguais}


\end{questao}
\begin{questao}[2]{}

Uma espira quadrada de lado $\SI{20}{\centi\meter}$ é girada de
$\ang{0}$ a $\ang{90}$ (normal em relação ao campo) em campo de
$\SI{0.6}{\tesla}$. Determine a variação de fluxo.
\resposta{$\Delta\Phi=\SI{-24}{\milli\weber}$}


\end{questao}
\begin{questao}[2]{}

Um núcleo toroidal tem seção de $\SI{4}{\centi\meter\squared}$ e conduz fluxo
de $\SI{0.6}{\milli\weber}$. Determine a densidade de fluxo e avalie a
proximidade da saturação.
\resposta{$B=\SI{1.5}{\tesla}$ --- próximo da saturação típica}


\end{questao}
\begin{questao}[2]{}

Uma bobina de $200$ espiras e área $\SI{15}{\centi\meter\squared}$ está em
campo de $\SI{0.4}{\tesla}$ perpendicular. Determine o fluxo concatenado.
\resposta{$\lambda=\SI{120}{\milli\weber}$-espira}


\end{questao}
\begin{questao}[2]{}

Uma espira está totalmente fora de uma região de campo uniforme e é empurrada
até ficar totalmente dentro. Descreva a variação do fluxo.
\resposta{Cresce de zero a $BA$, com taxa constante enquanto a espira atravessa
a fronteira}


\end{questao}
\begin{questao}[2]{}

Explique a diferença entre fluxo magnético e densidade de fluxo, e indique
quando cada um é a grandeza adequada.
\resposta{$\Phi$ é integral (weber); $B$ é local (tesla). $\Phi$ governa a
indução; $B$ governa forças e saturação}


\end{questao}

\blocoaprofundamento
\begin{questao}[3]{}

Um campo não uniforme $B(x)=B_0\dfrac{x}{L}$, perpendicular ao plano, atravessa
uma superfície retangular de largura $L$ (na direção $x$) e altura $h$, com
$B_0=\SI{0.5}{\tesla}$, $L=\SI{20}{\centi\meter}$ e
$h=\SI{10}{\centi\meter}$.
\begin{itensq}
\item Determine o fluxo por integração.
\item Compare com as estimativas usando $B_0$ e usando o valor médio.
\end{itensq}
\resposta{$\Phi=\SI{5}{\milli\weber}$ --- igual à estimativa pelo valor médio}


\end{questao}
\begin{questao}[3]{}

Um fio retilíneo longo conduz $\SI{10}{\ampere}$. Uma espira retangular de
$\SI{20}{\centi\meter}$ de comprimento (paralelo ao fio) está no mesmo plano,
com os lados a $\SI{5}{\centi\meter}$ e $\SI{20}{\centi\meter}$ do fio.
\begin{itensq}
\item Deduza a expressão do fluxo.
\item Calcule seu valor.
\end{itensq}
\resposta{$\Phi=\dfrac{\mu_0Ia}{2\pi}\ln\dfrac{d+b}{d}=\SI{0.55}{\micro\weber}$}


\end{questao}
\begin{questao}[3]{}

Um núcleo magnético de ferro tem comprimento médio $\SI{50}{\centi\meter}$,
seção de $\SI{4}{\centi\meter\squared}$ e $\mu_r=2000$, com enrolamento de
$\SI{1000}{}$ ampère-espiras.
\begin{itensq}
\item Determine a relutância do núcleo.
\item Determine o fluxo e a densidade de fluxo.
\item Avalie o resultado.
\end{itensq}
\resposta{$\mathcal{R}=\pot{4.97}{5}\,\si{\per\henry}$;
$\Phi=\SI{2}{\milli\weber}$; $B=\SI{5}{\tesla}$ --- \textbf{impossível}}


\end{questao}
\begin{questao}[3]{}

O mesmo núcleo da questão anterior recebe um \textbf{entreferro} de
$\SI{1}{\milli\meter}$.
\begin{itensq}
\item Determine a relutância do entreferro e a total.
\item Determine o novo fluxo e $B$.
\item Explique por que o entreferro domina.
\end{itensq}
\resposta{$\mathcal{R}_g=\pot{1.99}{6}$, quatro vezes a do ferro;
$B=\SI{1.0}{\tesla}$}


\end{questao}
\begin{questao}[3]{}

Duas bobinas compartilham parte do fluxo. A bobina $1$ produz
$\SI{5}{\milli\weber}$, dos quais $\SI{4}{\milli\weber}$ atravessam a bobina
$2$.
\begin{itensq}
\item Determine o fluxo de dispersão.
\item Determine o coeficiente de acoplamento aproximado.
\item Discuta como aumentá-lo.
\end{itensq}
\resposta{$\SI{1}{\milli\weber}$ de dispersão; $k\approx0{,}8$}


\end{questao}
\begin{questao}[3]{}

Um núcleo tem seção variável: $\SI{6}{\centi\meter\squared}$ em um trecho e
$\SI{2}{\centi\meter\squared}$ em outro, sem dispersão.
\begin{itensq}
\item Compare o fluxo nos dois trechos.
\item Compare a densidade de fluxo.
\item Identifique onde ocorre a saturação primeiro.
\end{itensq}
\resposta{Mesmo fluxo; $B$ três vezes maior na seção estreita, que satura
primeiro}


\end{questao}
\begin{questao}[3]{}

Uma espira de $\SI{25}{\centi\meter\squared}$ está parcialmente dentro de uma
região de campo de $\SI{0.8}{\tesla}$, com $60\%$ de sua área na região.
\begin{itensq}
\item Determine o fluxo.
\item Determine a f.e.m. se a espira for retirada completamente em
      $\SI{40}{\milli\second}$.
\end{itensq}
\resposta{$\Phi=\SI{1.2}{\milli\weber}$; $\varepsilon=\SI{30}{\milli\volt}$}


\end{questao}
\begin{questao}[3]{}

Um transformador tem núcleo de seção $\SI{10}{\centi\meter\squared}$ e opera
com $B_{\max}=\SI{1.2}{\tesla}$.
\begin{itensq}
\item Determine o fluxo máximo.
\item Determine o número de espiras do primário para
      $\SI{220}{\volt}$ a $\SI{60}{\hertz}$.
\end{itensq}
\resposta{$\Phi_{\max}=\SI{1.2}{\milli\weber}$; $N\approx688$ espiras}


\end{questao}
\begin{questao}[3]{}

Uma superfície fechada envolve completamente um ímã permanente. Determine o
fluxo magnético líquido através dela e justifique.
\resposta{Nulo, independentemente da posição e da intensidade do ímã}


\end{questao}
\begin{questao}[3]{}

Uma bobina de $500$ espiras e área $\SI{8}{\centi\meter\squared}$ é usada como
sensor de campo. Ela é retirada rapidamente do campo, e um integrador registra
$\SI{6}{\milli\volt\second}$.
\begin{itensq}
\item Determine o fluxo concatenado inicial.
\item Determine o campo.
\item Explique a vantagem do método.
\end{itensq}
\resposta{$\lambda=\SI{6}{\milli\weber}$-espira; $B=\SI{15}{\milli\tesla}$}


\end{questao}
\begin{questao}[3]{}

Avalie o erro de assumir campo uniforme ao calcular o fluxo através de uma
espira de $\SI{10}{\centi\meter}$ de lado, situada entre $\SI{5}{}$ e
$\SI{15}{\centi\meter}$ de um fio retilíneo.
\begin{itensq}
\item Determine o fluxo exato e o estimado pelo campo no centro.
\item Determine o erro.
\end{itensq}
\resposta{Erro de cerca de $+2{,}5\%$ na estimativa pelo centro}


\end{questao}

\blocodesafio
\begin{questao}[4]{}

\textbf{(Demonstração --- fluxo por superfície fechada.)} Prove que
$\oint\vec B\cdot\mathrm{d}\vec A=0$ e explore suas consequências.
\begin{itensq}
\item Demonstre a partir da inexistência de monopolos.
\item Deduza a conservação do fluxo em circuitos magnéticos.
\item Compare com a lei de Gauss elétrica.
\end{itensq}
\resposta{Linhas fechadas $\Rightarrow$ fluxo líquido nulo $\Rightarrow$ o fluxo
se conserva ao longo de um circuito magnético}


\end{questao}
\begin{questao}[4]{}

\textbf{(Independência da superfície.)} Prove que o fluxo através de qualquer
superfície apoiada em uma dada curva fechada é o mesmo.
\begin{itensq}
\item Demonstre.
\item Explique a importância para a lei de Faraday.
\item Indique o que ocorre se as normais forem orientadas incoerentemente.
\end{itensq}
\resposta{Duas superfícies de mesma borda formam superfície fechada, cujo fluxo
é nulo}


\end{questao}
\begin{questao}[4]{}

\textbf{(Circuito magnético.)} Estabeleça a analogia formal entre circuitos
magnéticos e elétricos, e discuta seus limites.
\begin{itensq}
\item Construa a tabela de correspondências.
\item Deduza a expressão da relutância.
\item Aponte três diferenças que quebram a analogia.
\end{itensq}
\resposta{$\mathcal{F}=\Phi\mathcal{R}$, com
$\mathcal{R}=\ell/\mu A$; a analogia falha por não linearidade, dispersão e
ausência de dissipação}


\end{questao}
\begin{questao}[4]{}

\textbf{(Projeto de entreferro.)} Um indutor de $\SI{100}{\micro\henry}$ deve
conduzir $\SI{5}{\ampere}$ de pico, com núcleo de ferrite de seção
$\SI{1}{\centi\meter\squared}$, caminho médio $\SI{6}{\centi\meter}$,
$\mu_r=2000$ e $B_{sat}=\SI{0.3}{\tesla}$.
\begin{itensq}
\item Determine a energia armazenada e o número mínimo de espiras.
\item Determine o entreferro necessário.
\item Determine a fração de energia armazenada no entreferro.
\end{itensq}
\resposta{$W=\SI{1.25}{\milli\joule}$; $N=17$ espiras;
$g\approx\SI{0.33}{\milli\meter}$; $92\%$ da energia no entreferro}


\end{questao}
\begin{questao}[4]{}

\textbf{(Dispersão e acoplamento.)} Analise o efeito da dispersão de fluxo no
desempenho de um transformador.
\begin{itensq}
\item Relacione $k$ com as indutâncias própria e mútua.
\item Determine o efeito sobre a tensão de saída.
\item Discuta quando a dispersão é desejável.
\end{itensq}
\resposta{$k=M/\sqrt{L_1L_2}$; a dispersão aparece como indutância série e
causa queda de tensão com a carga}


\end{questao}
\begin{questao}[4]{}

\textbf{(Fluxo e energia.)} Deduza a energia armazenada em um circuito
magnético e mostre onde ela reside.
\begin{itensq}
\item Deduza $W=\frac12\mathcal{R}\Phi^{2}$.
\item Determine a densidade de energia em função de $B$ e $\mu$.
\item Determine a fração de energia no entreferro do exemplo anterior.
\end{itensq}
\resposta{$u=B^{2}/2\mu$; praticamente toda a energia fica no entreferro}


\end{questao}
\begin{questao}[4]{}

\textbf{(Núcleo de seção variável.)} Formule o cálculo do fluxo em um núcleo
cuja seção varia ao longo do caminho, sem dispersão.
\begin{itensq}
\item Escreva a relutância total.
\item Determine o critério de dimensionamento.
\item Analise o caso de um núcleo com estreitamento localizado.
\end{itensq}
\resposta{$\mathcal{R}=\displaystyle\int\frac{\mathrm{d}\ell}{\mu A(\ell)}$;
dimensiona-se pela \emph{menor} seção}


\end{questao}

\gabarito[12.3 Fluxo magnético]

\subsection{Força magnética}

\blocofixacao
\begin{questao}[1]{}

A força magnética sobre uma carga em movimento é \textbf{nula} quando:
\begin{alternativas}
\item a velocidade é perpendicular ao campo.
\item a velocidade é paralela (ou antiparalela) ao campo.
\item a carga é positiva.
\item o campo é muito intenso.
\item a velocidade é máxima.
\end{alternativas}
\resposta{(b)}


\end{questao}
\begin{questao}[1]{}

Um condutor retilíneo de $\SI{30}{\centi\meter}$ conduz $\SI{2}{\ampere}$
perpendicularmente a um campo de $\SI{0.5}{\tesla}$. Determine a força sobre ele.
\resposta{$F = \SI{0.3}{\newton}$}


\end{questao}
\begin{questao}[1]{}

Uma carga positiva move-se horizontalmente para a direita, em uma região onde o
campo magnético aponta para dentro do plano da página. A força magnética sobre a
carga aponta:
\begin{alternativas}[2]
\item para cima.
\item para baixo.
\item para a direita.
\item para a esquerda.
\item para dentro da página.
\end{alternativas}
\resposta{(a)}


\end{questao}

\blocoaplicacao
\begin{questao}[2]{}

Uma carga de $\SI{1.6e-19}{\coulomb}$ move-se a
$\SI{e6}{\meter\per\second}$, perpendicularmente a um campo de $\SI{0.2}{\tesla}$.
Determine o módulo da força magnética e descreva a trajetória resultante.
\resposta{$F = \SI{3.2e-14}{\newton}$; trajetória circular}


\end{questao}
\begin{questao}[2]{}

Uma espira retangular percorrida por corrente é colocada em um campo magnético
uniforme. Explique por que ela tende a \textbf{girar}, e não a se deslocar, e
relacione isso com o funcionamento de um motor elétrico.
\resposta{Ver comentário}


\end{questao}

\blocoaprofundamento
\begin{questao}[3]{}

Um próton ($q = \SI{1.6e-19}{\coulomb}$, $m = \SI{1.67e-27}{\kilogram}$) entra
perpendicularmente em um campo uniforme de $\SI{0.5}{\tesla}$ com velocidade
$\SI{2e6}{\meter\per\second}$. Determine o raio da trajetória circular descrita.
\resposta{$r \approx \SI{4.2}{\centi\meter}$}


\end{questao}
\begin{questao}[3]{}

Um seletor de velocidades é formado por campos elétrico e magnético
\textbf{cruzados} (perpendiculares entre si e à velocidade). Partículas
carregadas entram com velocidades variadas e apenas as de uma velocidade
específica atravessam sem desviar.

\circuitoq{
\begin{tikzpicture}[scale=1,line width=0.8pt,>=Stealth]
 \draw[cinza!50,fill=cinza!8] (0,-1.2) rectangle (5.4,1.2);
 
 \foreach \x in {0.7,1.6,2.5,3.4,4.3}
   \draw[->,Vcol!70] (\x,1.05)--(\x,-1.05);
 \node[Vcol,font=\footnotesize] at (5.05,0.9) {$\vec E$};
 
 \foreach \x in {1.15,2.05,2.95,3.85}
   \foreach \y in {-0.6,0,0.6}
     {\node[Icol,font=\scriptsize] at (\x,\y) {$\otimes$};}
 \node[Icol,font=\footnotesize] at (0.35,-0.9) {$\vec B$};
 
 \draw[->,laranja,line width=1.2pt] (-1.3,0)--(0,0)
       node[midway,above,font=\footnotesize]{$v$};
 \draw[->,laranja,line width=1.2pt] (5.4,0)--(6.5,0);
 \fill[laranja] (-1.3,0) circle (0.1);
\end{tikzpicture}}{Seletor de velocidades com campos cruzados.}

\begin{itensq}
\item Deduza a condição para que a partícula atravesse sem desvio.
\item Mostre que essa condição \emph{não depende} da carga nem da massa da
      partícula.
\item Com $E = \SI{3e5}{\volt\per\meter}$ e $B = \SI{0.15}{\tesla}$, qual a
      velocidade selecionada?
\item Após o seletor, coloca-se um campo magnético $B'$ que curva a trajetória.
      Explique como essa combinação permite medir a razão $q/m$.
\end{itensq}
\resposta{(a) $qE = qvB \Rightarrow v = E/B$; (c) $v = \SI{2e6}{\meter\per\second}$;
(b) e (d) ver comentário}


\end{questao}

\blocofixacao
\begin{questao}[1]{}

Uma carga $q=+\SI{2}{\micro\coulomb}$ move-se a $\SI{500}{\meter\per\second}$
perpendicularmente a um campo de $\SI{0.2}{\tesla}$. Determine a força.
\resposta{$F=\SI{0.2}{\milli\newton}$}


\end{questao}
\begin{questao}[1]{}

Um condutor retilíneo de $\SI{40}{\centi\meter}$ conduz $\SI{3}{\ampere}$
perpendicularmente a $B=\SI{0.25}{\tesla}$. Determine a força.
\resposta{$F=\SI{0.3}{\newton}$}


\end{questao}
\begin{questao}[1]{}

A mesma carga da primeira questão move-se agora formando $\ang{30}$ com o
campo. Determine a força.
\resposta{$F=\SI{0.1}{\milli\newton}$}


\end{questao}
\begin{questao}[1]{}

Uma carga positiva move-se para a \textbf{direita} em uma região onde
$\vec B$ aponta para \textbf{cima}. Determine o sentido da força.
\begin{alternativas}[2]
\item para dentro do plano
\item para fora do plano
\item para cima
\item para a direita
\end{alternativas}
\resposta{(a)}


\end{questao}
\begin{questao}[1]{}

A partir de $F=qvB$, mostre que $\SI{1}{\tesla}$ equivale a
$\SI{1}{\newton\second\per\coulomb\per\meter}$.
\resposta{$\si{\tesla}=\dfrac{\si{\newton}}{\si{\coulomb}\cdot
\si{\meter\per\second}}$}


\end{questao}

\blocoaplicacao
\begin{questao}[2]{}

Um fio de $\SI{10}{\centi\meter}$ conduz $\SI{2}{\ampere}$ em um campo de
$\SI{0.2}{\tesla}$, formando $\ang{30}$ com ele. Determine a força.
\resposta{$F=\SI{20}{\milli\newton}$}


\end{questao}
\begin{questao}[2]{}

Um próton entra perpendicularmente em um campo de $\SI{0.5}{\tesla}$ com
$v=\pot{2}{6}\,\si{\meter\per\second}$.
\begin{itensq}
\item Determine o raio da trajetória.
\item Determine o período.
\end{itensq}
\dados{$e=\pot{1.6}{-19}\,\si{\coulomb}$; $m_p=\pot{1.67}{-27}\,\si{\kilogram}$}
\resposta{$r=\SI{4.18}{\centi\meter}$; $T=\SI{131}{\nano\second}$}


\end{questao}
\begin{questao}[2]{}

Um elétron e um próton, com a mesma velocidade, entram no mesmo campo
magnético perpendicularmente. Compare seus raios e períodos.
\resposta{$r_p/r_e=1836$; $T_p/T_e=1836$}


\end{questao}
\begin{questao}[2]{}

Dois fios paralelos longos, separados por $\SI{5}{\centi\meter}$, conduzem
$\SI{10}{\ampere}$ cada. Determine a força por unidade de comprimento.
\dados{$\mu_0=4\pi\times10^{-7}\,\si{\tesla\meter\per\ampere}$}
\resposta{$F/L=\SI{0.4}{\milli\newton\per\meter}$}


\end{questao}
\begin{questao}[2]{}

Uma espira retangular de área $\SI{100}{\centi\meter\squared}$ com $50$ voltas
conduz $\SI{2}{\ampere}$ em um campo de $\SI{0.3}{\tesla}$.
\begin{itensq}
\item Determine o momento de dipolo magnético.
\item Determine o torque máximo.
\end{itensq}
\resposta{$m=\SI{1}{\ampere\meter\squared}$; $\tau=\SI{0.3}{\newton\meter}$}


\end{questao}
\begin{questao}[2]{}

Explique por que a força magnética não altera o módulo da velocidade de uma
partícula carregada.
\resposta{Ela é sempre perpendicular a $\vec v$, logo não realiza trabalho}


\end{questao}
\begin{questao}[2]{}

Um elétron é abandonado \textbf{em repouso} em uma região onde coexistem
$E=\SI{500}{\volt\per\meter}$ e $B=\SI{0.02}{\tesla}$, perpendiculares entre
si.
\begin{itensq}
\item Determine a força e a aceleração no instante inicial.
\item Descreva qualitativamente o movimento subsequente.
\end{itensq}
\dados{$e=\pot{1.6}{-19}\,\si{\coulomb}$; $m_e=\pot{9.11}{-31}\,\si{\kilogram}$}
\resposta{$F=\SI{8e-17}{\newton}$ (só elétrica);
$a=\pot{8.8}{13}\,\si{\meter\per\second\squared}$}


\end{questao}
\begin{questao}[2]{}

Um fio em forma de semicircunferência de raio $R$ conduz corrente $I$ em um
campo uniforme $\vec B$ perpendicular ao plano do semicírculo. Determine a
força resultante.
\resposta{$F=2RIB$, como se o fio fosse retilíneo de comprimento $2R$}


\end{questao}
\begin{questao}[2]{}

Uma balança de corrente tem um fio horizontal de $\SI{20}{\centi\meter}$ em um
campo vertical de $\SI{0.4}{\tesla}$. Determine a corrente necessária para
equilibrar uma massa de $\SI{2}{\gram}$.
\dados{$g=\SI{9.8}{\meter\per\second\squared}$}
\resposta{$I=\SI{0.245}{\ampere}$}


\end{questao}
\begin{questao}[2]{}

Determine a velocidade de uma partícula que atravessa sem desvio uma região com
$E=\SI{2e4}{\volt\per\meter}$ e $B=\SI{0.1}{\tesla}$ perpendiculares entre si
e à velocidade.
\resposta{$v=\pot{2}{5}\,\si{\meter\per\second}$}


\end{questao}
\begin{questao}[2]{}

Uma espira quadrada de lado $\SI{10}{\centi\meter}$ conduz $\SI{5}{\ampere}$
em um campo uniforme de $\SI{0.2}{\tesla}$ paralelo ao seu plano.
\begin{itensq}
\item Determine a força resultante.
\item Determine o torque.
\end{itensq}
\resposta{$\vec F=\vec 0$; $\tau=\SI{10}{\milli\newton\meter}$}


\end{questao}

\blocoaprofundamento
\begin{questao}[3]{}

Uma partícula de carga $q$ e massa $m$ entra perpendicularmente em um campo
uniforme $B$ com velocidade $v$.
\begin{itensq}
\item Deduza o raio e o período da trajetória.
\item Determine a frequência angular (frequência de cíclotron).
\item Explique por que o período independe de $v$.
\end{itensq}
\resposta{$r=\dfrac{mv}{qB}$; $T=\dfrac{2\pi m}{qB}$;
$\omega=\dfrac{qB}{m}$}


\end{questao}
\begin{questao}[3]{}

Uma partícula entra em um campo uniforme formando $\ang{30}$ com $\vec B$, com
$v=\pot{1}{6}\,\si{\meter\per\second}$ e $B=\SI{0.1}{\tesla}$ (próton).
\begin{itensq}
\item Decomponha a velocidade e descreva a trajetória.
\item Determine o raio e o passo da hélice.
\end{itensq}
\resposta{Hélice com $r=\SI{5.2}{\centi\meter}$ e passo de
$\SI{0.57}{\meter}$}


\end{questao}
\begin{questao}[3]{}

Um espectrômetro de massa usa um seletor de velocidades ($v=\pot{1}{5}\,
\si{\meter\per\second}$) seguido de um campo de $\SI{0.1}{\tesla}$. Íons de
neônio de massas $20\,u$ e $22\,u$, todos com carga $+e$, entram no segundo
campo.
\begin{itensq}
\item Determine o raio de cada trajetória.
\item Determine a separação entre os pontos de chegada no detector.
\end{itensq}
\dados{$u=\pot{1.66}{-27}\,\si{\kilogram}$}
\resposta{$r_{20}=\SI{20.8}{\centi\meter}$, $r_{22}=\SI{22.8}{\centi\meter}$;
separação de $\SI{4.2}{\centi\meter}$}


\end{questao}
\begin{questao}[3]{}

A definição do ampère baseava-se na força entre condutores paralelos.
\begin{itensq}
\item Determine a força por metro entre dois fios distantes
      $\SI{1}{\meter}$, conduzindo $\SI{1}{\ampere}$ cada.
\item Explique o papel dessa relação na definição da unidade.
\item Comente a mudança recente.
\end{itensq}
\resposta{$F/L=\pot{2}{-7}\,\si{\newton\per\meter}$}


\end{questao}
\begin{questao}[3]{}

Um sensor Hall de espessura $\SI{0.1}{\milli\meter}$ conduz
$\SI{1}{\ampere}$ em campo de $\SI{0.5}{\tesla}$.
\begin{itensq}
\item Determine a tensão Hall para cobre ($n=\pot{8.5}{28}\,\si{\per\meter\cubed}$).
\item Repita para um semicondutor com $n=\pot{1}{22}\,\si{\per\meter\cubed}$.
\item Explique a diferença.
\end{itensq}
\resposta{$\SI{0.37}{\micro\volt}$ no cobre; $\SI{3.1}{\volt}$ no
semicondutor}


\end{questao}
\begin{questao}[3]{}

Um motor elementar tem espira de $200$ voltas, área
$\SI{50}{\centi\meter\squared}$, corrente de $\SI{3}{\ampere}$ e campo de
$\SI{0.4}{\tesla}$.
\begin{itensq}
\item Determine o torque máximo.
\item Determine a potência mecânica a $\SI{1200}{rpm}$, supondo torque médio
      igual a $2/\pi$ do máximo.
\end{itensq}
\resposta{$\tau_{\max}=\SI{1.2}{\newton\meter}$;
$P\approx\SI{96}{\watt}$}


\end{questao}
\begin{questao}[3]{}

Uma espira de corrente é colocada em um campo magnético \textbf{não uniforme},
mais intenso de um lado.
\begin{itensq}
\item Determine se há força resultante.
\item Determine o sentido dessa força.
\item Relacione com o comportamento de ímãs.
\end{itensq}
\resposta{Há força resultante, dirigida para a região de campo mais intenso (se
$\vec m$ estiver alinhado com $\vec B$)}


\end{questao}
\begin{questao}[3]{}

Compare a força entre dois fios paralelos conduzindo correntes iguais no mesmo
sentido e em sentidos opostos.
\begin{itensq}
\item Determine o sentido em cada caso.
\item Determine o módulo.
\item Explique a diferença em termos de campo.
\end{itensq}
\resposta{Mesmo sentido: atração; opostos: repulsão. Módulos iguais}


\end{questao}
\begin{questao}[3]{}

Um próton atravessa uma região com $E=\SI{2e4}{\volt\per\meter}$ e
$B=\SI{0.1}{\tesla}$, perpendiculares entre si e à velocidade, dispostos de
modo que as forças se oponham.
\begin{itensq}
\item Determine as duas parcelas da força de Lorentz para
      $v=\pot{1}{5}$, $\pot{2}{5}$ e $\pot{4}{5}\,\si{\meter\per\second}$.
\item Identifique qual parcela domina em cada caso.
\item Interprete o resultado.
\end{itensq}
\resposta{Domina a elétrica em $\pot{1}{5}$, equilíbrio em $\pot{2}{5}$ e a
magnética em $\pot{4}{5}\,\si{\meter\per\second}$}


\end{questao}
\begin{questao}[3]{}

Uma partícula carregada move-se em uma região com campos $\vec E$ e $\vec B$
perpendiculares entre si, mas com velocidade \emph{diferente} de $E/B$.
\begin{itensq}
\item Descreva qualitativamente o movimento.
\item Determine a velocidade média de deriva.
\item Indique uma aplicação.
\end{itensq}
\resposta{Movimento cicloidal, com deriva média
$\vec v_d=\dfrac{\vec E\times\vec B}{B^{2}}$}


\end{questao}
\begin{questao}[3]{}

Um condutor de $\SI{50}{\centi\meter}$ e $\SI{100}{\gram}$ está apoiado sobre
dois trilhos horizontais em campo vertical de $\SI{0.8}{\tesla}$. O
coeficiente de atrito estático é $0{,}25$.
\begin{itensq}
\item Determine a corrente mínima para iniciar o movimento.
\item Determine a aceleração inicial se a corrente for o dobro, com atrito
      cinético de $0{,}20$.
\end{itensq}
\resposta{$I=\SI{0.613}{\ampere}$; $a\approx\SI{2.9}{\meter\per\second\squared}$}


\end{questao}

\blocodesafio
\begin{questao}[4]{}

\textbf{(Demonstração --- trabalho nulo.)} Prove que a força magnética nunca
realiza trabalho sobre uma carga puntiforme.
\begin{itensq}
\item Demonstre formalmente.
\item Concilie com o fato de que motores elétricos realizam trabalho.
\item Explique o que ocorre com a energia em um motor.
\end{itensq}
\resposta{$\vec F\cdot\vec v=0$ sempre; o trabalho no motor vem da fonte, não do
campo}


\end{questao}
\begin{questao}[4]{}

\textbf{(Dedução do torque.)} Deduza a expressão do torque sobre uma espira
retangular em campo uniforme.
\begin{itensq}
\item Analise as forças em cada lado.
\item Obtenha $\tau=NIAB\sin\theta$.
\item Generalize para espira de forma qualquer.
\end{itensq}
\resposta{$\vec\tau=\vec m\times\vec B$, com $\vec m=NI\vec A$}


\end{questao}
\begin{questao}[4]{}

\textbf{(Projeto de espectrômetro.)} Projete um espectrômetro capaz de separar
isótopos com $\Delta m/m=1\%$.
\begin{itensq}
\item Determine a separação espacial em função dos parâmetros.
\item Estabeleça o requisito de colimação do feixe.
\item Discuta as limitações práticas.
\end{itensq}
\resposta{$\Delta r/r=\Delta m/m$; a resolução exige feixe com dispersão angular
e de velocidade menores que $1\%$}


\end{questao}
\begin{questao}[4]{}

\textbf{(Projeto de atuador linear.)} Projete um atuador que produza
$\SI{5}{\newton}$ com curso de $\SI{2}{\centi\meter}$.
\begin{itensq}
\item Determine as combinações de $B$, $I$ e $L$ possíveis.
\item Avalie o compromisso entre corrente e número de espiras.
\item Determine a potência dissipada e o trabalho realizado.
\end{itensq}
\resposta{Com $B=\SI{0.5}{\tesla}$ e $L=\SI{0.2}{\meter}$: $I=\SI{50}{\ampere}$
com uma espira, ou $\SI{2.5}{\ampere}$ com vinte}


\end{questao}
\begin{questao}[4]{}

\textbf{(Efeito Hall.)} Deduza a expressão da tensão Hall e analise suas
aplicações.
\begin{itensq}
\item Deduza $V_H=\dfrac{IB}{nqt}$.
\item Explique como o sinal revela o tipo de portador.
\item Cite três aplicações e o que cada uma mede.
\end{itensq}
\resposta{Equilíbrio entre força magnética e força elétrica transversal}


\end{questao}
\begin{questao}[4]{}

\textbf{(Estabilidade de espira.)} Analise o equilíbrio de uma espira de
corrente em campo magnético uniforme.
\begin{itensq}
\item Determine as posições de equilíbrio.
\item Classifique cada uma quanto à estabilidade.
\item Deduza o período de pequenas oscilações.
\end{itensq}
\resposta{$\theta=0$ é estável, $\theta=\pi$ é instável;
$T=2\pi\sqrt{J/(mB)}$}


\end{questao}
\begin{questao}[4]{}

\textbf{(Força de Lorentz --- estrutura e consequências.)} Analise a expressão
completa $\vec F=q\left(\vec E+\vec v\times\vec B\right)$.
\begin{itensq}
\item Compare as duas parcelas quanto à dependência da velocidade, à direção e
      ao trabalho realizado.
\item Avalie um próton com $v=\pot{3}{4}\,\si{\meter\per\second}$ em
      $E=\SI{500}{\volt\per\meter}$ e $B=\SI{0.02}{\tesla}$
      \emph{perpendiculares entre si}, com $\vec v$ paralelo a $\vec E$.
\item Explique por que aceleradores usam campo elétrico para acelerar e
      magnético para guiar.
\end{itensq}
\resposta{$F_E=\SI{8e-17}{\newton}$, $F_B=\SI{9.6e-17}{\newton}$,
perpendiculares; $|F|=\SI{1.25e-16}{\newton}$ a $\ang{50.2}$ de $\vec E$}


\end{questao}
\begin{questao}[4]{}

\textbf{(Projeto de cíclotron.)} Um cíclotron acelera prótons até
$\SI{10}{\mega\electronvolt}$ usando campo de $\SI{1.5}{\tesla}$.
\begin{itensq}
\item Determine a frequência de operação.
\item Determine o raio máximo e o número de voltas, com
      $\SI{50}{\kilo\volt}$ de ganho por passagem.
\item Determine se a correção relativística é necessária.
\end{itensq}
\dados{$e=\pot{1.6}{-19}\,\si{\coulomb}$; $m_p=\pot{1.67}{-27}\,\si{\kilogram}$;
$\SI{1}{\mega\electronvolt}=\pot{1.6}{-13}\,\si{\joule}$}
\resposta{$f=\SI{22.9}{\mega\hertz}$; $r=\SI{30.5}{\centi\meter}$;
$100$ voltas; correção de $1{,}1\%$}


\end{questao}
\begin{questao}[4]{}

\textbf{(Força entre correntes e o limite de barramentos.)} Um barramento de
subestação tem dois condutores paralelos separados por $\SI{15}{\centi\meter}$,
que devem suportar um curto-circuito de $\SI{40}{\kilo\ampere}$.
\begin{itensq}
\item Determine a força por metro durante o curto.
\item Determine o espaçamento máximo entre suportes, para deflexão admissível.
\item Discuta o caráter dessa solicitação.
\end{itensq}
\resposta{$F/L\approx\SI{2.1}{\kilo\newton\per\meter}$ --- comparável ao peso
de uma pessoa por metro}


\end{questao}

\gabarito[12.4 Força magnética]

\subsection{Lei de Lenz e indução eletromagnética}

\blocofixacao
\begin{questao}[1]{}

Segundo a lei de Lenz, a corrente induzida em uma espira tem sentido tal que:
\begin{alternativas}
\item reforça a variação do fluxo que a originou.
\item se opõe à variação do fluxo que a originou.
\item é sempre horária.
\item independe do sentido da variação do fluxo.
\item é proporcional ao fluxo, e não à sua variação.
\end{alternativas}
\resposta{(b)}


\end{questao}
\begin{questao}[1]{}

Um campo magnético variável atravessa perpendicularmente uma espira, com fluxo
$\Phi(t) = 0{,}02\,t$ (em Wb, com $t$ em segundos). Determine a fem induzida.
\resposta{$\varepsilon = \SI{20}{\milli\volt}$}


\end{questao}

\blocoaplicacao
\begin{questao}[2]{}

O campo magnético no interior de uma bobina aumenta linearmente de $0$ a
$\SI{0.5}{\tesla}$ em $\SI{2}{\second}$. A bobina tem $\SI{50}{}$ espiras e área de
$\SI{20}{\centi\meter\squared}$. Calcule a fem induzida.
\resposta{$\varepsilon = \SI{25}{\milli\volt}$}


\end{questao}
\begin{questao}[2]{}

Uma bobina de $\SI{100}{}$ espiras e área $\SI{30}{\centi\meter\squared}$ está em
um campo que varia segundo $B(t) = 0{,}5 - 0{,}1\,t$ (em T). Calcule a fem induzida
e diga se a corrente induzida cresce, decresce ou é constante.
\resposta{$\varepsilon = \SI{30}{\milli\volt}$, constante}


\end{questao}
\begin{questao}[2]{}

Um ímã é aproximado, com o polo norte à frente, de uma espira condutora, como na
figura. Determine o sentido da corrente induzida e a natureza da força entre o
ímã e a espira.

\circuitoq{
\begin{tikzpicture}[scale=0.9,line width=0.8pt,>=Stealth]
 
 \draw[fill=Icol!20,draw=Icol] (-4,-0.35) rectangle (-2.4,0.35);
 \node at (-3.5,0){\footnotesize S}; \node at (-2.7,0){\footnotesize N};
 \draw[->,Vcol,thick] (-2.2,0)--(-1.2,0); \node[Vcol,font=\footnotesize] at (-1.7,0.3){$v$};
 
 \draw[azulprof,very thick] (0,0) ellipse (0.35 and 1.1);
 \node[azulprof,font=\footnotesize] at (0,-1.4){espira};
\end{tikzpicture}}{Ímã aproximando-se de uma espira.}

\begin{alternativas}
\item A corrente induzida cria um polo norte voltado para o ímã, repelindo-o.
\item A corrente induzida cria um polo sul voltado para o ímã, atraindo-o.
\item Não há corrente induzida, pois o ímã não toca a espira.
\item A corrente induzida não gera força sobre o ímã.
\item A espira é atraída, acelerando o ímã.
\end{alternativas}
\resposta{(a)}


\end{questao}
\begin{questao}[2]{}

Um transformador ideal tem no primário uma bobina onde o fluxo magnético varia à
taxa de $\SI{0.01}{\weber\per\second}$. Se o enrolamento tem $\SI{200}{}$ espiras,
qual é a fem induzida?
\resposta{$\varepsilon = \SI{2}{\volt}$}


\end{questao}

\blocoaprofundamento
\begin{questao}[3]{}

O fluxo magnético em uma espira varia senoidalmente:
$\Phi(t) = \Phi_0\cos(\omega t)$, com $\Phi_0 = \SI{0.01}{\weber}$ e
$\omega = \SI{100}{\radian\per\second}$.
\begin{itensq}
\item Obtenha a expressão da fem induzida.
\item Qual é o valor máximo da fem?
\end{itensq}
\resposta{(a) $\varepsilon(t) = \Phi_0\omega\sin(\omega t)$;
(b) $\varepsilon_{\max} = \SI{1}{\volt}$}


\end{questao}
\begin{questao}[3]{Apostila original revisada}

A corrente induzida em um circuito puramente resistivo é
$i(t)=I_0\sin(\omega t)$. Determine a expressão da fem induzida e indique em quais
instantes a corrente muda de sentido durante um período.
\resposta{$\mathcal{E}(t)=RI_0\sin(\omega t)$; inversões em $t=0$, $\pi/\omega$ e $2\pi/\omega$}


\end{questao}

\blocodesafio
\begin{questao}[4]{}

Uma barra condutora de comprimento $\ell = \SI{40}{\centi\meter}$ desliza sem
atrito, com velocidade constante $v = \SI{3}{\meter\per\second}$, sobre dois
trilhos paralelos fechados por um resistor $R = \SI{2}{\ohm}$. O conjunto está
imerso em campo magnético uniforme $B = \SI{0.5}{\tesla}$, perpendicular ao plano
dos trilhos.

\circuitoq{
\begin{tikzpicture}[scale=1,line width=0.8pt,>=Stealth]
 
 \draw[cinza,line width=1.4pt] (0,0)--(6,0);
 \draw[cinza,line width=1.4pt] (0,2.2)--(6,2.2);
 
 \draw (0,0) to[R,l_={$R$}] (0,2.2);
 
 \draw[Vcol,line width=2.2pt] (3.6,-0.15)--(3.6,2.35);
 \node[Vcol,font=\footnotesize,above] at (3.6,2.4) {barra};
 \draw[->,laranja,line width=1.3pt] (3.9,1.1)--(5.1,1.1)
       node[midway,above,font=\footnotesize]{$v$};
 
 \foreach \x in {1.0,1.8,2.6}
   \foreach \y in {0.55,1.1,1.65}
     {\node[Icol,font=\scriptsize] at (\x,\y) {$\odot$};}
 \node[Icol,font=\footnotesize] at (1.8,0.15) {$\vec B$ (saindo)};
 \draw[<->,cinza] (6.3,0)--(6.3,2.2) node[midway,right,font=\footnotesize]{$\ell$};
\end{tikzpicture}}{Barra deslizante sobre trilhos --- fem de movimento.}

\begin{itensq}
\item Calcule a fem induzida e a corrente no circuito.
\item Determine a força magnética que atua sobre a barra e seu sentido.
\item Calcule a potência dissipada no resistor e a potência mecânica necessária
      para manter a barra em movimento uniforme. Compare-as.
\item Interprete o resultado à luz da lei de Lenz e da conservação de energia.
\end{itensq}
\resposta{(a) $\varepsilon=\SI{0.6}{\volt}$, $i=\SI{0.3}{\ampere}$;
(b) $F=\SI{0.06}{\newton}$, opondo-se ao movimento;
(c) ambas $=\SI{0.18}{\watt}$; (d) ver comentário}


\end{questao}

\blocofixacao
\begin{questao}[1]{}

Uma bobina de $200$ espiras sofre variação de fluxo de
$\SI{2}{\milli\weber}$ por espira em $\SI{5}{\milli\second}$. Determine a
f.e.m. média induzida.
\resposta{$\varepsilon=\SI{80}{\volt}$}


\end{questao}
\begin{questao}[1]{}

Uma barra de $\SI{30}{\centi\meter}$ move-se a $\SI{4}{\meter\per\second}$
perpendicularmente a um campo de $\SI{0.5}{\tesla}$. Determine a f.e.m.
induzida.
\resposta{$\varepsilon=\SI{0.6}{\volt}$}


\end{questao}
\begin{questao}[1]{}

Determine o número de espiras necessário para obter $\SI{12}{\volt}$ com
variação de fluxo de $\SI{0.04}{\weber\per\second}$.
\resposta{$N=300$}


\end{questao}
\begin{questao}[1]{}

Uma bobina de $100$ espiras deve gerar $\SI{50}{\volt}$ com variação de fluxo
de $\SI{5}{\milli\weber}$. Determine o tempo em que essa variação deve ocorrer.
\resposta{$\Delta t=\SI{10}{\milli\second}$}


\end{questao}
\begin{questao}[1]{}

O sinal negativo na lei de Faraday expressa:
\begin{alternativas}[2]
\item a lei de Lenz e a conservação da energia
\item que a f.e.m. é sempre negativa
\item um erro de convenção
\item que a corrente é sempre horária
\end{alternativas}
\resposta{(a)}


\end{questao}
\begin{questao}[1]{}

Uma espira está imersa em campo magnético uniforme e \textbf{constante}, e é
deslocada rapidamente dentro da região. Determine a f.e.m. induzida.
\resposta{$\varepsilon=0$}


\end{questao}
\begin{questao}[1]{}

Mostre que $\si{\weber\per\second}$ equivale a volt.
\resposta{$\si{\weber}=\si{\volt\second}$}


\end{questao}
\begin{questao}[1]{}

Segundo a lei de Lenz, quando o fluxo através de uma espira \textbf{aumenta},
a corrente induzida cria um campo próprio que:
\begin{alternativas}[2]
\item reforça o campo externo
\item opõe-se ao campo externo
\item é perpendicular ao campo externo
\item é nulo
\end{alternativas}
\resposta{(b)}


\end{questao}

\blocoaplicacao
\begin{questao}[2]{}

Uma barra de $\SI{50}{\centi\meter}$ desliza a $\SI{3}{\meter\per\second}$
sobre trilhos em campo de $\SI{0.4}{\tesla}$, com resistência total de
$\SI{2}{\ohm}$.
\begin{itensq}
\item Determine a f.e.m. e a corrente.
\item Determine a força necessária para manter a velocidade.
\item Verifique o balanço de potência.
\end{itensq}
\resposta{$\varepsilon=\SI{0.6}{\volt}$; $I=\SI{0.3}{\ampere}$;
$F=\SI{60}{\milli\newton}$; $P=\SI{0.18}{\watt}$ pelos dois caminhos}


\end{questao}
\begin{questao}[2]{}

Uma bobina de $100$ espiras e área $\SI{200}{\centi\meter\squared}$ gira a
$\SI{377}{\radian\per\second}$ em campo de $\SI{0.5}{\tesla}$.
\begin{itensq}
\item Determine a f.e.m. máxima.
\item Determine a f.e.m. eficaz.
\end{itensq}
\resposta{$\varepsilon_{\max}=\SI{377}{\volt}$;
$\varepsilon_{ef}=\SI{267}{\volt}$}


\end{questao}
\begin{questao}[2]{}

Duas bobinas têm indutância mútua $M=\SI{50}{\milli\henry}$. A corrente na
primeira varia à taxa de $\SI{200}{\ampere\per\second}$. Determine a f.e.m.
induzida na segunda.
\resposta{$\varepsilon_2=\SI{10}{\volt}$}


\end{questao}
\begin{questao}[2]{}

Um transformador ideal tem $1000$ espiras no primário e $100$ no secundário,
com $\SI{220}{\volt}$ na entrada.
\begin{itensq}
\item Determine a tensão de saída.
\item Determine a corrente de entrada para carga que exige
      $\SI{5}{\ampere}$.
\end{itensq}
\resposta{$V_2=\SI{22}{\volt}$; $I_1=\SI{0.5}{\ampere}$}


\end{questao}
\begin{questao}[2]{}

O fluxo em uma espira varia segundo $\Phi(t)=0{,}05\,t^{2}$ (weber, com $t$ em
segundos). Determine a f.e.m. em $t=\SI{2}{\second}$ e a f.e.m. média nos dois
primeiros segundos.
\resposta{$\varepsilon(2)=\SI{0.2}{\volt}$; média de $\SI{0.1}{\volt}$}


\end{questao}
\begin{questao}[2]{}

Uma espira de resistência $\SI{10}{\ohm}$ e $50$ espiras sofre variação de
fluxo de $\SI{4}{\milli\weber}$. Determine a \textbf{carga} total que circula.
\resposta{$q=\SI{20}{\milli\coulomb}$}


\end{questao}
\begin{questao}[2]{}

Uma espira quadrada de lado $\SI{10}{\centi\meter}$ entra com velocidade
constante de $\SI{2}{\meter\per\second}$ em uma região de campo de
$\SI{0.6}{\tesla}$.
\begin{itensq}
\item Determine a f.e.m. durante a entrada.
\item Determine a f.e.m. quando totalmente dentro.
\end{itensq}
\resposta{$\SI{0.12}{\volt}$ durante a entrada; zero depois}


\end{questao}
\begin{questao}[2]{}

Um disco metálico gira entre os polos de um ímã. Explique o efeito observado
sobre seu movimento.
\resposta{Ele freia --- correntes de Foucault opõem-se ao movimento}


\end{questao}
\begin{questao}[2]{}

Uma bobina com núcleo de ferro tem $\SI{500}{}$ espiras. O fluxo varia de
$\SI{0}{}$ a $\SI{1.2}{\milli\weber}$ em $\SI{8}{\milli\second}$. Determine a
f.e.m. média.
\resposta{$\varepsilon=\SI{75}{\volt}$}


\end{questao}
\begin{questao}[2]{}

Uma espira é conectada a um resistor e um ímã é afastado dela. Determine o
sentido da força que a espira exerce sobre o ímã.
\resposta{Atrativa --- ela tenta impedir o afastamento}


\end{questao}
\begin{questao}[2]{}

Uma bobina de $200$ espiras, área $\SI{100}{\centi\meter\squared}$ e
resistência $\SI{40}{\ohm}$ está em campo de $\SI{0.8}{\tesla}$
perpendicular. O campo é desligado em $\SI{20}{\milli\second}$.
\begin{itensq}
\item Determine a f.e.m. e a corrente induzidas.
\item Determine o sentido da corrente em relação ao campo original.
\item Determine a carga total.
\end{itensq}
\resposta{$\varepsilon=\SI{80}{\volt}$; $I=\SI{2}{\ampere}$;
$q=\SI{40}{\milli\coulomb}$}


\end{questao}

\blocoaprofundamento
\begin{questao}[3]{}

Uma barra de $\SI{100}{\gram}$ e $\SI{40}{\centi\meter}$ desliza sem atrito em
trilhos \textbf{verticais}, em campo horizontal de $\SI{0.5}{\tesla}$, com
resistência total de $\SI{0.5}{\ohm}$. Ela é solta do repouso.
\begin{itensq}
\item Escreva a equação do movimento.
\item Determine a velocidade-limite.
\item Determine a potência dissipada nesse regime.
\end{itensq}
\dados{$g=\SI{9.8}{\meter\per\second\squared}$}
\resposta{$v_{lim}=\SI{12.25}{\meter\per\second}$; $P=\SI{12}{\watt}$}


\end{questao}
\begin{questao}[3]{}

Uma bobina de $N$ espiras e resistência $R$ é retirada de um campo. Demonstre
que a carga induzida independe da velocidade do movimento.
\begin{itensq}
\item Deduza a expressão.
\item Explique fisicamente.
\item Indique a aplicação em instrumentação.
\end{itensq}
\resposta{$q=\dfrac{N\Delta\Phi}{R}$ --- independe de $\Delta t$}


\end{questao}
\begin{questao}[3]{}

Um núcleo de transformador maciço de $\SI{10}{\milli\meter}$ de espessura é
substituído por $20$ lâminas isoladas de $\SI{0.5}{\milli\meter}$.
\begin{itensq}
\item Determine a redução das perdas por correntes parasitas.
\item Explique o mecanismo.
\item Indique a limitação prática.
\end{itensq}
\resposta{Redução de $400$ vezes}


\end{questao}
\begin{questao}[3]{}

Um solenoide longo de $n_1$ espiras por metro e seção $A$ tem, em torno de si,
uma bobina de $N_2$ espiras.
\begin{itensq}
\item Deduza a indutância mútua.
\item Avalie para $n_1=\SI{2000}{\per\meter}$, $A=\SI{5}{\centi\meter\squared}$
      e $N_2=100$.
\item Explique por que $M$ não depende do raio da bobina externa.
\end{itensq}
\resposta{$M=\mu_0n_1N_2A$; $M=\SI{1.26}{\milli\henry}$}


\end{questao}
\begin{questao}[3]{}

Uma espira retangular está próxima a um fio retilíneo cuja corrente varia
segundo $i(t)=I_0\sin\omega t$.
\begin{itensq}
\item Escreva o fluxo através da espira.
\item Determine a f.e.m. induzida.
\item Explique por que essa interferência é difícil de eliminar.
\end{itensq}
\resposta{$\varepsilon=-M\omega I_0\cos\omega t$, com
$M=\dfrac{\mu_0a}{2\pi}\ln\dfrac{d+b}{d}$}


\end{questao}
\begin{questao}[3]{}

Uma bobina de $\SI{200}{\milli\henry}$ conduz corrente que varia de
$\SI{2}{\ampere}$ a zero em $\SI{5}{\milli\second}$.
\begin{itensq}
\item Determine a f.e.m. autoinduzida.
\item Determine a energia envolvida.
\item Explique o risco associado.
\end{itensq}
\resposta{$\varepsilon=\SI{80}{\volt}$; $W=\SI{0.4}{\joule}$}


\end{questao}
\begin{questao}[3]{}

Um gerador fornece $\SI{500}{\watt}$ a uma carga. Analise o torque necessário
para acioná-lo e sua relação com a lei de Lenz.
\begin{itensq}
\item Determine o torque a $\SI{1800}{rpm}$, supondo rendimento de
      $\SI{90}{\percent}$.
\item Explique o que acontece se a carga for desligada.
\end{itensq}
\resposta{$\tau=\SI{2.95}{\newton\meter}$; sem carga, o torque cai quase a
zero}


\end{questao}
\begin{questao}[3]{}

Um anel condutor é colocado sobre o núcleo de um solenoide alimentado por
corrente alternada. Observa-se que ele é \textbf{arremessado} para cima.
\begin{itensq}
\item Explique o fenômeno.
\item Determine o que muda se o anel for cortado.
\item Determine o que muda com corrente contínua em regime.
\end{itensq}
\resposta{Repulsão por corrente induzida; anel cortado não salta; em CC de
regime, nada acontece}


\end{questao}
\begin{questao}[3]{}

Uma espira de área $\SI{50}{\centi\meter\squared}$ está em campo que varia
segundo $B(t)=0{,}4\sin(100\pi t)$ tesla, e \emph{simultaneamente} sua área
diminui a $\SI{10}{\centi\meter\squared\per\second}$.
\begin{itensq}
\item Escreva a expressão geral da f.e.m.
\item Identifique as duas contribuições.
\end{itensq}
\resposta{$\varepsilon=-\left(A\dfrac{\mathrm{d}B}{\mathrm{d}t}
+B\dfrac{\mathrm{d}A}{\mathrm{d}t}\right)$}


\end{questao}
\begin{questao}[3]{}

Uma bobina de $500$ espiras, área $\SI{20}{\centi\meter\squared}$ e resistência
$\SI{25}{\ohm}$ é girada de $\ang{180}$ em campo de $\SI{0.3}{\tesla}$.
\begin{itensq}
\item Determine a variação de fluxo por espira.
\item Determine a carga total induzida.
\end{itensq}
\resposta{$\Delta\Phi=\SI{1.2}{\milli\weber}$; $q=\SI{24}{\milli\coulomb}$}


\end{questao}
\begin{questao}[3]{}

Um transformador com $k=0{,}95$ e $L_1=\SI{100}{\milli\henry}$,
$L_2=\SI{25}{\milli\henry}$.
\begin{itensq}
\item Determine a indutância mútua.
\item Determine a f.e.m. no secundário para
      $\mathrm{d}i_1/\mathrm{d}t=\SI{50}{\ampere\per\second}$.
\item Determine a fração de fluxo que não se acopla.
\end{itensq}
\resposta{$M=\SI{47.5}{\milli\henry}$; $\varepsilon_2=\SI{2.38}{\volt}$;
$5\%$ de dispersão}


\end{questao}
\begin{questao}[3]{}

Analise o perfil de f.e.m. produzido por um fluxo que cresce linearmente por
$\SI{2}{\second}$, permanece constante por $\SI{3}{\second}$ e decresce
linearmente por $\SI{1}{\second}$, com amplitude de $\SI{6}{\milli\weber}$ e
$N=100$.
\begin{itensq}
\item Determine a f.e.m. em cada trecho.
\item Verifique a área total sob a curva de f.e.m.
\end{itensq}
\resposta{$+0{,}3$; $0$; $\SI{-0.6}{\volt}$; área líquida nula}


\end{questao}
\begin{questao}[3]{}

Um freio magnético produz força proporcional à velocidade, $F=bv$, com
$b=\SI{20}{\newton\second\per\meter}$, atuando sobre um corpo de
$\SI{5}{\kilogram}$ inicialmente a $\SI{10}{\meter\per\second}$.
\begin{itensq}
\item Escreva a equação do movimento e resolva.
\item Determine a constante de tempo e a distância percorrida.
\item Explique por que o freio não para o corpo.
\end{itensq}
\resposta{$v=v_0e^{-t/\tau}$ com $\tau=\SI{0.25}{\second}$;
$d=\SI{2.5}{\meter}$}


\end{questao}

\blocodesafio
\begin{questao}[4]{}

\textbf{(Lenz e conservação da energia.)} Demonstre que a lei de Lenz é uma
consequência necessária da conservação da energia.
\begin{itensq}
\item Suponha o contrário e derive a contradição.
\item Quantifique com o exemplo da barra em trilhos.
\item Discuta o que a lei de Lenz \emph{não} determina.
\end{itensq}
\resposta{Sinal oposto levaria a crescimento ilimitado de energia}


\end{questao}
\begin{questao}[4]{}

\textbf{(Projeto de bobina sensora.)} Projete uma bobina que gere
$\SI{5}{\volt}$ médios quando o fluxo variar $\SI{0.2}{\milli\weber}$ em
$\SI{4}{\milli\second}$.
\begin{itensq}
\item Determine o número de espiras.
\item Determine a resistência admissível para corrente de saída inferior a
      $\SI{1}{\milli\ampere}$ em carga de $\SI{10}{\kilo\ohm}$.
\item Avalie a resposta em frequência.
\end{itensq}
\resposta{$N=100$ espiras}


\end{questao}
\begin{questao}[4]{}

\textbf{(Freio por correntes parasitas.)} Modele quantitativamente um freio
magnético de disco.
\begin{itensq}
\item Estabeleça a dependência da força com os parâmetros.
\item Determine por que a força é proporcional à velocidade.
\item Analise o comportamento em velocidades muito altas.
\end{itensq}
\resposta{$F\propto\sigma B^{2}tv$; proporcionalidade a $v$ decorre de
$\varepsilon\propto v$ e $i\propto\varepsilon$}


\end{questao}
\begin{questao}[4]{}

\textbf{(Convenção dos pontos na indução mútua.)} Explique o significado físico
da convenção dos pontos e sua determinação experimental.
\begin{itensq}
\item Enuncie a convenção.
\item Determine o sinal de $M$ nas equações de malha.
\item Descreva o procedimento experimental.
\end{itensq}
\resposta{Pontos marcam terminais magneticamente correspondentes; o sinal de
$M$ depende de as correntes entrarem ou não pelos pontos}


\end{questao}
\begin{questao}[4]{}

\textbf{(Verificação experimental da lei de Faraday.)} Projete um experimento
para verificar $\varepsilon\propto N$ e $\varepsilon\propto v$.
\begin{itensq}
\item Descreva o arranjo e o procedimento.
\item Estabeleça o critério de validação.
\item Aponte as principais fontes de erro.
\end{itensq}
\resposta{Ímã caindo por tubo com bobinas de $N$ variável; verificar
linearidade em $N$ e em $v$}


\end{questao}
\begin{questao}[4]{}

\textbf{(Transformador em corrente contínua.)} Explique por que um
transformador não transfere potência em regime de corrente contínua.
\begin{itensq}
\item Analise o regime permanente.
\item Analise o transitório de ligação.
\item Determine o que ocorre com o enrolamento primário.
\end{itensq}
\resposta{Fluxo constante $\Rightarrow$ $\mathrm{d}\Phi/\mathrm{d}t=0
\Rightarrow$ nenhuma f.e.m. no secundário}


\end{questao}
\begin{questao}[4]{}

\textbf{(Espira entrando em campo.)} Deduza a f.e.m. durante a entrada de uma
espira retangular em uma região de campo uniforme, com velocidade constante.
\begin{itensq}
\item Deduza pela variação de área e pela força sobre portadores.
\item Determine a força necessária para manter a velocidade.
\item Analise o balanço de energia.
\end{itensq}
\resposta{$\varepsilon=B\ell v$ pelos dois caminhos;
$F=\dfrac{B^{2}\ell^{2}v}{R}$}


\end{questao}
\begin{questao}[4]{}

\textbf{(Balanço volt-segundo.)} Demonstre que, em regime periódico, a tensão
média sobre um indutor é nula, e explore as consequências.
\begin{itensq}
\item Demonstre.
\item Aplique a um conversor com dois patamares de tensão.
\item Explique o que ocorre se o balanço não for respeitado.
\end{itensq}
\resposta{$\bar v_L=0$; daí $V_1t_1=V_2t_2$ em regime}


\end{questao}
\begin{questao}[4]{}

\textbf{(Correntes parasitas em condutores.)} Analise por que condutores
maciços de grande seção apresentam perdas superiores às previstas em corrente
alternada.
\begin{itensq}
\item Explique o efeito pelicular.
\item Determine a profundidade de penetração e avalie para cobre a
      $\SI{60}{\hertz}$.
\item Indique as soluções empregadas.
\end{itensq}
\resposta{$\delta=\sqrt{\dfrac{2\rho}{\omega\mu}}\approx\SI{8.5}{\milli\meter}$
para cobre a $\SI{60}{\hertz}$}


\end{questao}
\begin{questao}[4]{}

 Um feixe de prótons atravessa inicialmente uma região onde atuam simultaneamente um campo elétrico uniforme de intensidade
\[
E = 70\,\mathrm{kV/m}
\]
e um campo magnético uniforme de intensidade
\[
B = 50\,\mathrm{mT},
\]
perpendiculares entre si e também à direção de movimento das partículas. Os campos são ajustados de modo que apenas os prótons que atravessam essa região sem sofrer desvio alcancem uma segunda região, na qual existe apenas o campo magnético de mesma intensidade, conforme ilustrado na figura.

Desprezando os efeitos gravitacionais e considerando
\[
m_p = 1{,}67\times10^{-27}\,\mathrm{kg},
\qquad
q_p = 1{,}60\times10^{-19}\,\mathrm{C},
\]
determine:

\begin{itensq}
    \item a velocidade dos prótons que atravessam o seletor sem sofrer desvio;
    \item o raio \(r\) da trajetória circular descrita pelos prótons na região onde atua apenas o campo magnético;
    \item a distância \(2r\) entre o ponto de entrada na região magnética e o ponto em que os prótons atingem a placa detectora.
\end{itensq}



 \begin{tikzpicture}[
    >=Latex,
    line cap=round,
    line join=round,
    font=\sffamily
]




\def\r{2.35}
\def\xleft{0.0}
\def\xfilter{3.4}
\def\xslit{6.45}
\pgfmathsetmacro{\xfield}{\xslit-0.35}
\def\ybeam{0.0}
\def\ybottom{-4.7}


\coordinate (C) at (\xfield,\ybottom/2);
\coordinate (TopArc) at ($(C)+(90:\r)$);
\coordinate (BotArc) at ($(C)+(-90:\r)$);




\draw[thick] (-1.05,0) circle (0.23);
\node at (-1.05,0) {\large $+$};


\draw[thick,dashed] (-0.75,0) -- (\xfield,0);




\draw[line width=3.6pt] (0.12,0.32) -- (0.12,0.82);
\draw[line width=3.6pt] (0.12,-0.32) -- (0.12,-0.82);





\draw[line width=2.2pt] (2.15,0.78) -- (4.65,0.78);
\draw[line width=2.2pt] (2.15,-0.78) -- (4.65,-0.78);


\draw[line width=1.8pt] (3.40,0.78) -- (3.40,1.62);
\draw[line width=1.8pt] (3.40,-0.78) -- (3.40,-1.62);


\node[font=\Large] at (2.75,1.15) {$-$};
\node[font=\Large] at (2.75,-1.15) {$+$};


\foreach \x in {2.45,2.95,3.45,3.95,4.45}{
    \draw[myE,line width=1.5pt,-{Latex[length=3mm,width=2mm]}]
        (\x,-0.62) -- (\x,0.62);
}


\foreach \x in {2.55,3.25,3.95,4.35}{
    \foreach \y in {-0.42,0,0.42}{
        \fill[myB] (\x,\y) circle (2.3pt);
    }
}


\node[myE,above] at (4.95,0.25) {$\vec E$};
\node[myB,right] at (4.55,0.0) {$\vec B$};


\draw[thick,-{Latex[length=3mm,width=2mm]}] (1.05,0.28) -- (1.95,0.28);
\node[above] at (1.5,0.28) {$\vec v$};




\draw[line width=3.6pt] (\xslit-0.35,0.32) -- (\xslit-0.35,0.82);
\draw[line width=3.6pt] (\xslit-0.35,-0.32) -- (\xslit-0.35,-0.82);




\foreach \x in {6.7,7.4,8.1,8.8,9.5}{
    \foreach \y in {0.45,-0.15,-0.85,-1.55,-2.25,-2.95,-3.65,-4.35, -5}{
        \fill[myB] (\x,\y) circle (2.5pt);
    }
}
\foreach \x/\y in {
    7.05/-0.50,7.75/-1.20,8.45/-0.50,9.15/-1.20,
    7.05/-1.90,7.75/-2.60,8.45/-1.90,9.15/-2.60,
    7.05/-3.30,7.75/-4.00,8.45/-3.30,9.15/-4.00
}{
    \fill[myB] (\x,\y) circle (2.5pt);
}

\node[myB] at (9.75,-0.25) {$\vec B$};




\draw[thick,dashed,-{Latex[length=3mm,width=2mm]}]
    (TopArc) arc[start angle=90,end angle=-90,radius=\r];




\draw[thick,{Latex[length=2.8mm]}-{Latex[length=2.8mm]}]
    (5.65,-0.02) -- (5.65,\ybottom+0.02);

\draw[thin] (5.42,0) -- (5.90,0);
\draw[thin] (5.42,\ybottom) -- (5.90,\ybottom);

\node[fill=white,inner sep=1.5pt,font=\Large] at (5.65,-2.35) {$2r$};




\draw[line width=4pt] (\xslit-0.35,\ybottom+0.55) -- (\xslit-0.35,\ybottom-0.55);
\draw[line width=1.7pt] (\xslit-0.35,\ybottom) -- (\xslit+0.15,\ybottom);

\node[align=center,font=\Large] at (3.95,-4.55) {Detector\\plate};




\node[anchor=west,text=myE,font=\bfseries\Large] at (-0.55,-2.05)
    {$E = 70\ \mathrm{kV/m}$};

\node[anchor=west,text=myB,font=\bfseries\Large] at (-0.55,-2.85)
    {$B = 50\ \mathrm{mT}$};

\end{tikzpicture}

\resposta{$v=\dfrac{E}{B}=\SI{1.4e6}{\meter\per\second}$,
\quad
$r=\dfrac{m_pv}{q_pB}\approx\SI{0.292}{\meter}$,
\quad
$2r\approx\SI{0.585}{\meter}$}



\end{questao}
\begin{questao}[4]{}


Um elétron parte praticamente do repouso e é acelerado por uma diferença de
potencial de
\[
\Delta V=\SI{1.8}{\kilo\volt}.
\]
Após atravessar uma pequena abertura, ele penetra perpendicularmente em uma
região onde existe um campo magnético uniforme de intensidade
\[
B=\SI{25}{\milli\tesla},
\]
perpendicular ao plano da figura.

Considere
\[
m_e=\SI{9.11e-31}{\kilogram},
\qquad
|q_e|=\SI{1.60e-19}{\coulomb},
\]
e despreze efeitos relativísticos.

Determine:

\begin{itensq}
    \item a velocidade do elétron ao deixar a região de aceleração;
    \item o raio da trajetória circular no campo magnético;
    \item o período do movimento circular;
    \item o tempo necessário para o elétron percorrer um quarto de circunferência.
\end{itensq}

\begin{center}
\begin{tikzpicture}[
    >=Latex,
    line cap=round,
    line join=round,
    scale=0.95,
    font=\sffamily
]





\node at (-4.6,0) {\Large $e^-$};

\draw[thick,-{Latex[length=3mm]}]
    (-4.25,0) -- (-3.65,0)
    node[midway,above] {$\vec v$};


\draw[line width=2.5pt] (-3.3,-1.1) -- (-3.3,1.1);
\draw[line width=2.5pt] (-1.3,-1.1) -- (-1.3,1.1);

\node[font=\Large] at (-3.65,0.65) {$-$};
\node[font=\Large] at (-0.95,0.65) {$+$};


\foreach \y in {-0.7,-0.35,0,0.35,0.7}{
    \draw[myE,thick,-{Latex[length=2.5mm]}]
        (-1.55,\y) -- (-3.05,\y);
}

\node[myE] at (-2.3,1.45) {$\Delta V=\SI{1.8}{\kilo\volt}$};


\draw[line width=3pt] (-0.65,0.28) -- (-0.65,1.0);
\draw[line width=3pt] (-0.65,-0.28) -- (-0.65,-1.0);


\draw[thick,dashed,-{Latex[length=3mm]}]
    (-0.6,0) -- (0.6,0);





\foreach \x in {0.8,1.5,2.2,2.9,3.6}{
    \foreach \y in {-2.1,-1.4,-0.7,0,0.7,1.4,2.1}{
        \draw[myB,thick] (\x-0.08,\y-0.08) -- (\x+0.08,\y+0.08);
        \draw[myB,thick] (\x-0.08,\y+0.08) -- (\x+0.08,\y-0.08);
    }
}

\node[myB] at (3.65,2.55)
    {$\vec B$ entrando no plano};


\draw[very thick,dashed,-{Latex[length=3mm]}]
    (0.6,0)
    arc[start angle=90,end angle=0,radius=2.0];


\fill (2.6,-2.0) circle (1.5pt);

\draw[thin,dashed]
    (2.6,-2.0) -- (0.6,0);

\node at (1.55,-1.05) {$r$};

\end{tikzpicture}
\end{center}

\resposta{$v\approx\SI{2.52e7}{\meter\per\second}$,
\quad
$r\approx\SI{5.72}{\milli\meter}$,
\quad
$T\approx\SI{1.43}{\nano\second}$,
\quad
$t_{1/4}\approx\SI{0.357}{\nano\second}$}




\end{questao}
\begin{questao}[4]{}


Um espectrômetro de massas é utilizado para separar dois isótopos de neônio,
$^{20}\mathrm{Ne}^{+}$ e $^{22}\mathrm{Ne}^{+}$. Inicialmente, os íons
atravessam um seletor de velocidades constituído por campos elétrico e
magnético perpendiculares, de intensidades
\[
E=\SI{30}{\kilo\volt\per\meter},
\qquad
B_s=\SI{20}{\milli\tesla}.
\]

Os íons selecionados entram em seguida em uma região contendo apenas um campo
magnético uniforme de intensidade
\[
B_a=\SI{0.35}{\tesla}.
\]

Considere
\[
u=\SI{1.66e-27}{\kilogram},
\qquad
q=+e=\SI{1.60e-19}{\coulomb}.
\]

Determine:

\begin{itensq}
    \item a velocidade dos íons que atravessam o seletor sem desvio;
    \item os raios das trajetórias dos dois isótopos;
    \item a distância entre os pontos de impacto na placa detectora;
    \item qual isótopo atinge a região mais distante do ponto de entrada.
\end{itensq}

\begin{center}
\begin{tikzpicture}[
    >=Latex,
    line cap=round,
    line join=round,
    scale=0.88,
    font=\sffamily
]


\draw[thick] (-4.7,0) circle (0.3);
\node at (-4.7,0) {$+$};

\draw[thick,-{Latex[length=3mm]}]
    (-4.35,0) -- (-3.65,0);


\draw[line width=2pt] (-3.2,0.8) -- (-1.0,0.8);
\draw[line width=2pt] (-3.2,-0.8) -- (-1.0,-0.8);

\node at (-2.85,1.12) {$-$};
\node at (-2.85,-1.12) {$+$};

\foreach \x in {-2.9,-2.4,-1.9,-1.4}{
    \draw[myE,thick,-{Latex[length=2.5mm]}]
        (\x,-0.6) -- (\x,0.6);
}

\foreach \x in {-2.7,-2.0,-1.3}{
    \foreach \y in {-0.4,0,0.4}{
        \fill[myB] (\x,\y) circle (2pt);
    }
}

\node at (-2.1,1.55) {Seletor de velocidades};


\draw[line width=3pt] (-0.45,0.3) -- (-0.45,1);
\draw[line width=3pt] (-0.45,-0.3) -- (-0.45,-1);


\foreach \x in {0.2,0.9,1.6,2.3,3.0,3.7}{
    \foreach \y in {-4.1,-3.4,-2.7,-2.0,-1.3,-0.6,0.1}{
        \fill[myB] (\x,\y) circle (2.3pt);
    }
}

\node[myB] at (3.5,0.7) {$\vec B_a$};


\draw[very thick,dashed]
    (-0.35,0)
    arc[start angle=90,end angle=-90,radius=1.8];

\node at (1.25,-1.1) {$^{20}\mathrm{Ne}^{+}$};


\draw[very thick,dash dot]
    (-0.35,0)
    arc[start angle=90,end angle=-90,radius=2.15];

\node at (2.65,-2.2) {$^{22}\mathrm{Ne}^{+}$};


\draw[line width=4pt] (-0.45,-3.2) -- (-0.45,-4.8);

\node[rotate=90] at (-0.85,-4.0) {Detector};


\fill (-0.35,-3.6) circle (2.2pt);
\fill (-0.35,-4.3) circle (2.2pt);

\draw[<->,thick]
    (-1.25,-3.6) -- (-1.25,-4.3)
    node[midway,left] {$\Delta d$};

\end{tikzpicture}
\end{center}

\resposta{$v=\SI{1.50e6}{\meter\per\second}$,
\quad
$r_{20}\approx\SI{0.889}{\meter}$,
\quad
$r_{22}\approx\SI{0.977}{\meter}$,
\quad
$\Delta d\approx\SI{0.178}{\meter}$}




\end{questao}
\begin{questao}[4]{}


Uma partícula alfa entra em uma região onde existe um campo magnético uniforme
de intensidade
\[
B=\SI{0.18}{\tesla}.
\]
A velocidade inicial da partícula possui módulo
\[
v=\SI{7.0e5}{\meter\per\second}
\]
e forma um ângulo de
\[
\theta=40^\circ
\]
com a direção do campo magnético.

Considere
\[
m_\alpha=\SI{6.64e-27}{\kilogram},
\qquad
q_\alpha=\SI{3.20e-19}{\coulomb}.
\]

Determine:

\begin{itensq}
    \item as componentes $v_{\parallel}$ e $v_{\perp}$ da velocidade;
    \item o raio da trajetória helicoidal;
    \item o período de uma volta;
    \item o passo $p$ da hélice;
    \item o que aconteceria com a trajetória se $\theta=90^\circ$.
\end{itensq}

\begin{center}
\begin{tikzpicture}[
    >=Latex,
    line cap=round,
    line join=round,
    scale=1.0,
    font=\sffamily
]


\draw[myB,line width=1.3pt,-{Latex[length=3mm]}]
    (-4,0) -- (4.7,0)
    node[right] {$\vec B$};


\foreach \y in {-1.4,-0.7,0.7,1.4}{
    \draw[myB!70,thin,-{Latex[length=2mm]}]
        (-3.8,\y) -- (4.3,\y);
}


\draw[very thick]
plot[
    domain=0:720,
    samples=150,
    smooth
]
({-3.3+0.010*\x},{0.60*sin(\x)});


\fill (-3.3,0) circle (2.8pt);
\node[above left] at (-3.3,0) {$\alpha$};


\draw[thick,-{Latex[length=3mm]}]
    (-3.3,0) -- (-2.35,0.8)
    node[above] {$\vec v$};


\draw[dashed,thick,-{Latex[length=2.5mm]}]
    (-3.3,-0.15) -- (-2.0,-0.15)
    node[midway,below] {$v_{\parallel}$};


\draw[dashed,thick,-{Latex[length=2.5mm]}]
    (-3.3,0) -- (-3.3,0.85)
    node[left] {$v_{\perp}$};


\draw (-2.75,0)
    arc[start angle=0,end angle=40,radius=0.55];

\node at (-2.55,0.25) {$\theta$};


\draw[<->,thick]
    (0.3,-1.15) -- (3.9,-1.15)
    node[midway,fill=white,inner sep=2pt] {$p$};

\node at (0,-2.0)
    {Trajetória helicoidal};

\end{tikzpicture}
\end{center}

\resposta{$v_{\parallel}\approx\SI{5.36e5}{\meter\per\second}$,
\quad
$v_{\perp}\approx\SI{4.50e5}{\meter\per\second}$,
\quad
$r\approx\SI{5.18}{\centi\meter}$,
\quad
$T\approx\SI{0.724}{\micro\second}$,
\quad
$p\approx\SI{0.388}{\meter}$}




\end{questao}
\begin{questao}[4]{}


Um cíclotron é utilizado para acelerar prótons em uma região de campo magnético
uniforme de intensidade
\[
B=\SI{0.80}{\tesla}.
\]
Os prótons descrevem órbitas de raio progressivamente maior até atingirem o
raio máximo
\[
R=\SI{25}{\centi\meter}.
\]

Considere
\[
m_p=\SI{1.67e-27}{\kilogram},
\qquad
q_p=\SI{1.60e-19}{\coulomb},
\]
e despreze correções relativísticas.

Determine:

\begin{itensq}
    \item a frequência que deve ser aplicada entre os eletrodos do cíclotron;
    \item o período do movimento dos prótons;
    \item a velocidade do próton ao atingir o raio máximo;
    \item a energia cinética máxima do próton, em joules e em elétron-volts.
\end{itensq}

\begin{center}
\begin{tikzpicture}[
    >=Latex,
    line cap=round,
    line join=round,
    scale=0.95,
    font=\sffamily
]


\foreach \x in {-3,-2.25,-1.5,-0.75,0.75,1.5,2.25,3}{
    \foreach \y in {-2.2,-1.5,-0.8,0,0.8,1.5,2.2}{
        \fill[myB!75] (\x,\y) circle (1.8pt);
    }
}


\draw[line width=2pt]
    (0.18,2.55)
    arc[start angle=90,end angle=270,radius=2.55];

\draw[line width=2pt]
    (-0.18,-2.55)
    arc[start angle=-90,end angle=90,radius=2.55];


\draw[line width=2pt] (-0.18,-2.55) -- (-0.18,2.55);
\draw[line width=2pt] (0.18,-2.55) -- (0.18,2.55);

\node at (-1.45,0) {\Large D};
\node at (1.45,0) {\Large D};


\draw[myE,very thick,<->]
    (-0.13,0.75) -- (0.13,0.75);

\node[myE,right] at (0.25,0.75) {$\vec E(t)$};


\draw[very thick,-{Latex[length=3mm]}]
plot[
    domain=0:900,
    samples=250,
    smooth
]
({0.0027*\x*cos(\x)},
 {0.0027*\x*sin(\x)});


\fill (0,0) circle (2.5pt);


\draw[dashed,thick]
    (0,0) -- (2.43,0);

\node[below] at (1.25,0) {$R$};


\node[myB] at (3.7,2.0)
    {$\vec B$};


\draw[thick,-{Latex[length=3mm]}]
    (2.45,-0.35) -- (3.65,-0.7);

\node[right] at (3.65,-0.7)
    {feixe};

\end{tikzpicture}
\end{center}

\resposta{$f\approx\SI{12.2}{\mega\hertz}$,
\quad
$T\approx\SI{82.0}{\nano\second}$,
\quad
$v_{\max}\approx\SI{1.92e7}{\meter\per\second}$,
\quad
$K_{\max}\approx\SI{3.07e-13}{\joule}
\approx\SI{1.92}{\mega\electronvolt}$}




\end{questao}

\gabarito[12.5 Lei de Lenz e indução eletromagnética]
